% Response & coordination in organisms — Biology Upper Sixth — Cours signé OnBuch+
% Official MINESEC syllabus. Compiler avec : tectonic BIO02-response-coordination-in-organisms.tex
\newcommand{\DOCMATIERE}{Biology}
\newcommand{\DOCNIVEAU}{Upper Sixth}
\newcommand{\DOCMODULE}{Upper Sixth — Biology}
\newcommand{\DOCLECON}{2}
\newcommand{\DOCTITRE}{Response \& coordination in organisms}
% OnBuch+ — préambule « cours signés » ENRICHI (compilé avec tectonic / XeLaTeX)
% Système visuel « Pop » d'OnBuch+ : fond crème (jamais blanc), bordures encre
% épaisses, ombres « dures » décalées, titres Archivo Black en capitales.
% Version MATHÉMATIQUES : graphes (pgfplots/tikz), boîtes pédagogiques
% (définition, propriété, méthode, attention, le savais-tu, exercice résolu,
% exercice, TP, bilan), page de garde et signature OnBuch+ ancrée en bas.
%
% Macros attendues dans main.tex AVANT \input{preamble} :
%   \DOCMATIERE  \DOCNIVEAU  \DOCMODULE  \DOCLECON (numéro)  \DOCTITRE
\documentclass[11pt]{article}

\usepackage{fontspec}
\usepackage[english]{babel}
\usepackage[a4paper,margin=2cm,top=3.4cm,bottom=2.8cm,headheight=1.4cm,footskip=1.3cm]{geometry}
\usepackage{amsmath,amssymb,amsfonts}
\usepackage{mathtools} % \xmapsto, \coloneqq... souvent utilisés par les modèles
\usepackage{cancel} % simplifications barrées (bilans de forces)
% \abs{z} (module, valeur absolue) et \norm{u} (norme), utilisés par les modèles
\providecommand{\abs}{}\let\abs\relax\DeclarePairedDelimiter{\abs}{\lvert}{\rvert}
\providecommand{\norm}{}\let\norm\relax\DeclarePairedDelimiter{\norm}{\lVert}{\rVert}
\usepackage[table]{xcolor} % [table] : fournit aussi \rowcolors (alternance de lignes)
\usepackage{pagecolor}
\usepackage{tikz}
\usetikzlibrary{arrows.meta,positioning,calc,shapes.geometric,shapes.misc,shapes.arrows,decorations.pathmorphing,decorations.pathreplacing,decorations.markings,patterns,shadows,mindmap,trees,fit,backgrounds,matrix,intersections,angles,quotes}
\usepackage{pgfplots}
\usepackage[version=4]{mhchem} % équations nucléaires \ce{^{235}_{92}U}
\usepackage[european,siunitx]{circuitikz} % schémas électriques
\pgfplotsset{compat=1.18}
\usepgfplotslibrary{fillbetween,groupplots} % aires entre courbes (calcul intégral)
\usepackage{enumitem}
\usepackage{fancyhdr}
% [most] charge toutes les bibliothèques tcolorbox (listings, minted, raster,
% documentation...) et gonfle inutilement la mémoire TeX sur un document long
% (~300 boîtes) : on ne charge que ce qui est réellement utilisé.
\usepackage[skins,breakable]{tcolorbox}
\usepackage{graphicx}
\usepackage{colortbl}
\usepackage{booktabs}
\usepackage{tabularx}
\usepackage{diagbox} % case d'en-tête diagonale des tableaux à double entrée
% Colonnes extensibles centrée / gauche / droite (C, L, R), souvent utilisées par les modèles
\newcolumntype{C}{>{\centering\arraybackslash}X}
\newcolumntype{L}{>{\raggedright\arraybackslash}X}
\newcolumntype{R}{>{\raggedleft\arraybackslash}X}
\usepackage{array}
\usepackage{multicol}
\usepackage{siunitx}
% parse-numbers=false : accepte aussi \SI{2,5 \times 10^{-3}}{...}
\sisetup{locale=UK,output-decimal-marker={.},group-minimum-digits=4,parse-numbers=false}
\DeclareSIUnit\ohm{\ensuremath{\symup{\Omega}}} % U+2126 absent de la police
\DeclareSIUnit\division{div} % réglages d'oscilloscope (ms/div, V/div)
\DeclareSIUnit\tr{tr} % tours (vitesse de rotation, ex. \SI{1500}{\tr\per\minute})
\usepackage{qrcode}
% \degree : commande fréquemment utilisée par les modèles pour un angle ou une
% température en dehors de \SI{}{} (non fournie par les paquets chargés).
\providecommand{\degree}{^{\circ}}
% \qed (fin de démonstration), utilisé par les modèles sans amsthm
\providecommand{\qed}{\hfill$\square$}
% \barre : double barre des valeurs interdites dans un tableau de variations (tkz-tab)
\providecommand{\barre}{\ensuremath{\|}}
% environnement proof (amsthm non chargé) utilisé par les modèles
\ifdefined\proof\else\newenvironment{proof}[1][Proof]{\par\noindent\textit{#1.}\ }{\qed\par}\fi
\usepackage{lastpage}
\usepackage{hyperref}
\hypersetup{hidelinks}

% ============================================================
% Palette « Pop » OnBuch+ — mêmes hex que la classe P (plus_ui.dart)
% ============================================================
\definecolor{poporange}{HTML}{F59321}
\definecolor{popdark}{HTML}{E07A0C}
\definecolor{popcream}{HTML}{FFF6E8}
\definecolor{popink}{HTML}{1C1714}
\definecolor{poppurple}{HTML}{7A5AE0}
\definecolor{popgreen}{HTML}{2E9E6A}
\definecolor{poppink}{HTML}{E0457B}
\definecolor{popblue}{HTML}{2F7DE1}
\definecolor{popgold}{HTML}{C98A12}
\definecolor{popmuted}{HTML}{8A7F76}
\definecolor{popline}{HTML}{EADFD2}
\definecolor{popgray}{HTML}{8A7F76}
\colorlet{popgrey}{popgray}
% Alias tolérants (le modèle nomme parfois les couleurs différemment)
\colorlet{popwhite}{white}
\colorlet{popblack}{popink}
\colorlet{popred}{poppink}
\colorlet{popyellow}{popgold}
% Teintes claires pour fonds de tableaux / zones
\colorlet{popblueL}{popblue!10!white}
\colorlet{poporangeL}{poporange!14!white}
\colorlet{popgreenL}{popgreen!12!white}
\colorlet{poppurpleL}{poppurple!12!white}
\colorlet{poppinkL}{poppink!12!white}
\colorlet{popgoldL}{popgold!14!white}

\pagecolor{popcream}

% ============================================================
% Typographie
% ============================================================
\defaultfontfeatures{Ligatures=TeX}
\setmainfont{PlusJakartaSans-Medium.ttf}[Path=../../../fonts/,BoldFont=PlusJakartaSans-Bold.ttf]
% Maths en sans-serif (Fira Math) avec chiffres et lettres droites de Plus
% Jakarta, pour que les formules se fondent dans le texte.
\usepackage[math-style=ISO,bold-style=ISO]{unicode-math}
\setmathfont{FiraMath-Regular.otf}[Path=../../../fonts/]
\setmathfont{PlusJakartaSans-Medium.ttf}[Path=../../../fonts/,range={up/num,up/{Latin,latin},\mathup}]
% (icomma retiré : décimales avec point en anglais)
% Caractères mathématiques Unicode absents des polices de texte (Plus Jakarta,
% Archivo Black) : rendus par leur équivalent mathématique, en texte comme en
% formule — sinon ils s'affichent en carré vide (ex. « Arithmétique dans ℤ »).
\usepackage{newunicodechar}
\newunicodechar{ℝ}{\ensuremath{\mathbb{R}}}
\newunicodechar{ℤ}{\ensuremath{\mathbb{Z}}}
\newunicodechar{ℂ}{\ensuremath{\mathbb{C}}}
\newunicodechar{ℕ}{\ensuremath{\mathbb{N}}}
\newunicodechar{ℚ}{\ensuremath{\mathbb{Q}}}
\newunicodechar{𝔻}{\ensuremath{\mathbb{D}}}
\newunicodechar{α}{\ensuremath{\alpha}}
\newunicodechar{β}{\ensuremath{\beta}}
\newunicodechar{γ}{\ensuremath{\gamma}}
\newunicodechar{δ}{\ensuremath{\delta}}
\newunicodechar{ε}{\ensuremath{\varepsilon}}
\newunicodechar{θ}{\ensuremath{\theta}}
\newunicodechar{λ}{\ensuremath{\lambda}}
\newunicodechar{μ}{\ensuremath{\mu}}
\newunicodechar{σ}{\ensuremath{\sigma}}
\newunicodechar{τ}{\ensuremath{\tau}}
\newunicodechar{φ}{\ensuremath{\varphi}}
\newunicodechar{ψ}{\ensuremath{\psi}}
\newunicodechar{ω}{\ensuremath{\omega}}
\newunicodechar{Σ}{\ensuremath{\Sigma}}
% majuscules produites par \MakeUppercase dans les titres (α-aminés -> Α)
\newunicodechar{Α}{\ensuremath{\alpha}}
\newunicodechar{Β}{\ensuremath{\beta}}
\newunicodechar{Γ}{\ensuremath{\Gamma}}
\newunicodechar{Δ}{\ensuremath{\Delta}}
\newunicodechar{Ε}{\ensuremath{\varepsilon}}
\newunicodechar{Θ}{\ensuremath{\Theta}}
\newunicodechar{Λ}{\ensuremath{\Lambda}}
\newunicodechar{Μ}{\ensuremath{\mu}}
\newunicodechar{Τ}{\ensuremath{\tau}}
\newunicodechar{Φ}{\ensuremath{\Phi}}
\newunicodechar{Ψ}{\ensuremath{\Psi}}
\newunicodechar{Ω}{\ensuremath{\Omega}}
\newunicodechar{↦}{\ensuremath{\mapsto}}
\newunicodechar{∈}{\ensuremath{\in}}
\newunicodechar{∉}{\ensuremath{\notin}}
\newunicodechar{∧}{\ensuremath{\wedge}}
\newunicodechar{⇔}{\ensuremath{\Leftrightarrow}}
\newunicodechar{⟺}{\ensuremath{\Longleftrightarrow}}
\newunicodechar{⇒}{\ensuremath{\Rightarrow}}
\newunicodechar{⟹}{\ensuremath{\Longrightarrow}}
\newunicodechar{∘}{\ensuremath{\circ}}
\newunicodechar{∪}{\ensuremath{\cup}}
\newunicodechar{∩}{\ensuremath{\cap}}
\newunicodechar{≡}{\ensuremath{\equiv}}
\newunicodechar{⊂}{\ensuremath{\subset}}
\newunicodechar{⊥}{\ensuremath{\perp}}
\newunicodechar{∃}{\ensuremath{\exists}}
\newunicodechar{∀}{\ensuremath{\forall}}
\newunicodechar{∛}{\ensuremath{\sqrt[3]{\,}}}
\newunicodechar{∜}{\ensuremath{\sqrt[4]{\,}}}
\newunicodechar{⃗}{\ensuremath{{}^{\to}}}
\newunicodechar{ⁿ}{\ifmmode^{n}\else\textsuperscript{\ensuremath{n}}\fi}
\newunicodechar{ˣ}{\ifmmode^{x}\else\textsuperscript{\ensuremath{x}}\fi}
\newunicodechar{ᵇ}{\ifmmode^{b}\else\textsuperscript{\ensuremath{b}}\fi}
\newunicodechar{ʳ}{\ifmmode^{r}\else\textsuperscript{\ensuremath{r}}\fi}
\newunicodechar{ᵘ}{\ifmmode^{u}\else\textsuperscript{\ensuremath{u}}\fi}
\newunicodechar{ʸ}{\ifmmode^{y}\else\textsuperscript{\ensuremath{y}}\fi}
\newunicodechar{ᵃ}{\ifmmode^{a}\else\textsuperscript{\ensuremath{a}}\fi}
\newunicodechar{ᵉ}{\ifmmode^{e}\else\textsuperscript{\ensuremath{e}}\fi}
\newunicodechar{ᵏ}{\ifmmode^{k}\else\textsuperscript{\ensuremath{k}}\fi}
\newunicodechar{ᵖ}{\ifmmode^{p}\else\textsuperscript{\ensuremath{p}}\fi}
\newunicodechar{ᴺ}{\ifmmode^{N}\else\textsuperscript{\ensuremath{N}}\fi}
\newunicodechar{ᶜ}{\ifmmode^{c}\else\textsuperscript{\ensuremath{c}}\fi}
\newunicodechar{ʰ}{\ifmmode^{h}\else\textsuperscript{\ensuremath{h}}\fi}
\newunicodechar{ᵠ}{\ifmmode^{q}\else\textsuperscript{\ensuremath{q}}\fi}
\newunicodechar{ₙ}{\ifmmode_{n}\else\textsubscript{\ensuremath{n}}\fi}
\newunicodechar{ₐ}{\ifmmode_{a}\else\textsubscript{\ensuremath{a}}\fi}
\newunicodechar{ᵢ}{\ifmmode_{i}\else\textsubscript{\ensuremath{i}}\fi}
\newunicodechar{ᵣ}{\ifmmode_{r}\else\textsubscript{\ensuremath{r}}\fi}
\newunicodechar{ₖ}{\ifmmode_{k}\else\textsubscript{\ensuremath{k}}\fi}
\newunicodechar{ᵦ}{\ifmmode_{\beta}\else\textsubscript{\ensuremath{\beta}}\fi}
\newunicodechar{ₓ}{\ifmmode_{x}\else\textsubscript{\ensuremath{x}}\fi}
\newfontfamily\popdisplay{ArchivoBlack-Regular.ttf}[Path=../../../fonts/]
\newfontfamily\poplogo{SpaceGrotesk-Bold.ttf}[Path=../../../fonts/]
\newfontfamily\popsemi{PlusJakartaSans-SemiBold.ttf}[Path=../../../fonts/]
\newfontfamily\popxbold{PlusJakartaSans-ExtraBold.ttf}[Path=../../../fonts/]
\newfontfamily\popxboldls{PlusJakartaSans-ExtraBold.ttf}[Path=../../../fonts/,LetterSpace=3.0]

\setlist[enumerate]{leftmargin=1.7em,itemsep=5pt,topsep=4pt}
\setlist[itemize]{leftmargin=1.7em,itemsep=4pt,topsep=4pt,label={\color{poporange}\textbullet}}
\setlength{\parindent}{0pt}
\setlength{\parskip}{5pt}
\renewcommand{\arraystretch}{1.25}

% pgfplots : style Pop par défaut
\pgfplotsset{
  every axis/.append style={axis line style={popink,line width=1pt},tick style={popink},
    grid style={popline,line width=0.5pt},label style={font=\small},tick label style={font=\footnotesize}},
  popcurve/.style={poporange,line width=1.8pt,no markers,smooth},
}
% Même style disponible dans un \draw TikZ simple (hors pgfplots)
\tikzset{popcurve/.style={poporange,line width=1.8pt}}
\tikzset{popgrid/.style={popline,very thin}}
\colorlet{popcurve}{poporange} % le nom du style est parfois utilisé comme couleur

% ============================================================
% Pill / squiggle
% ============================================================
\newcommand{\poppill}[3]{%
  \tikz[baseline=-0.55ex]\node[draw=popink,line width=1.2pt,rounded corners=2.6pt,%
    fill=#2,inner xsep=6.5pt,inner ysep=3pt]{%
    \popxboldls\fontsize{7}{7}\selectfont\color{#3}\MakeUppercase{#1}};%
}
\newcommand{\popsquiggle}[1]{%
  \tikz[baseline=-0.4ex]{
    \draw[#1,line width=2.6pt,line cap=round]
      plot [smooth] coordinates {(0,0.09) (0.34,0.01) (0.68,0.21) (1.02,0.01) (1.36,0.10)};
  }%
}
% Mot-clé surligné (marqueur orange sous le texte)
\newcommand{\cle}[1]{{\popxbold\color{popdark}#1}}

% ============================================================
% En-tête / pied
% ============================================================
\pagestyle{fancy}
\fancyhf{}
\renewcommand{\headrulewidth}{0pt}
\renewcommand{\footrulewidth}{0pt}
\lhead{\raisebox{-0.15ex}{\poplogo\fontsize{13}{13}\selectfont\color{popink}OnBuch\color{poporange}+}}
\rhead{\poppill{\DOCMATIERE\ \textbullet\ \DOCNIVEAU\ \textbullet\ Chapter \DOCLECON}{white}{popink}}
\lfoot{\popxbold\fontsize{7.6}{7.6}\selectfont\color{popmuted}\MakeUppercase{Signed course OnBuch+}\ \textbullet\ {\popsemi\DOCTITRE}}
\rfoot{\tikz[baseline=-0.6ex]\node[rounded corners=8.5pt,fill=popink,minimum height=17pt,minimum width=17pt,inner xsep=5pt,inner sep=0]{\popxbold\fontsize{8}{8}\selectfont\color{popcream}\thepage\,/\,\pageref*{LastPage}};}

% ============================================================
% Titres
% ============================================================
\newcommand{\chaptitre}[2]{%
  \begin{tcolorbox}[enhanced,colback=poporange,colframe=popink,boxrule=2.6pt,arc=11pt,
      drop shadow=popink,left=15pt,right=15pt,top=12pt,bottom=11pt,before skip=3pt,after skip=18pt]
    \poppill{#2}{popink}{popcream}\par\vspace{9pt}
    {\raggedright\hyphenpenalty=10000\exhyphenpenalty=10000\popdisplay\fontsize{22}{24}\selectfont\color{popcream}\MakeUppercase{#1}\par}
  \end{tcolorbox}
}
% Section numérotée : pastille orange avec le numéro + titre Archivo
\newcounter{popsec}
\newcommand{\coursec}[1]{%
  \stepcounter{popsec}\setcounter{popsub}{0}%
  \par\vspace{16pt}\noindent
  \begin{minipage}{\linewidth}
  \tikz[baseline=-0.7ex]\node[draw=popink,line width=1.6pt,fill=poporange,rounded corners=4pt,minimum size=22pt,inner sep=0]{\popdisplay\fontsize{12}{12}\selectfont\color{popcream}\thepopsec};%
  \hspace{8pt}{\popdisplay\fontsize{15}{17}\selectfont\color{popink}\MakeUppercase{#1}}\par
  \vspace{4pt}\noindent{\color{poporange}\rule{36pt}{3.6pt}}
  \end{minipage}\par\nopagebreak\vspace{8pt}
}
\newcounter{popsub}
\newcommand{\courssub}[1]{%
  \stepcounter{popsub}%
  \par\vspace{9pt}\noindent{\popxbold\fontsize{11.5}{13}\selectfont\color{popink}{\color{poporange}\thepopsec.\thepopsub}\hspace{6pt}#1}\par\nopagebreak\vspace{4pt}
}

% ============================================================
% Boîtes pédagogiques — même recette : fond blanc, bordure épaisse colorée,
% ombre dure encre, badge plein. Titre optionnel [#1].
% ============================================================
\tcbset{popbase/.style={breakable,enhanced,colback=white,boxrule=2.3pt,arc=9pt,
  shadow={3pt}{-3pt}{0pt}{popink},left=14pt,right=14pt,top=11pt,bottom=11pt,before skip=11pt,after skip=15pt,
  fontupper=\popsemi\fontsize{10.6}{14.5}\selectfont\color{popink},
  fontlower=\popsemi\fontsize{10.6}{14.5}\selectfont\color{popink}}}
% Coche dessinée (le glyphe ✓ n'existe pas dans les polices du document)
\newcommand{\popcheck}[1]{\tikz[baseline=-0.5ex]\draw[#1,line width=1.6pt,line cap=round,line join=round](-0.13,0.0)--(-0.04,-0.09)--(0.13,0.1);}
\providecommand{\coche}{\popcheck{popgreen}}
\providecommand{\resultat}[1]{\par\smallskip\noindent\fcolorbox{popgreen}{popgreenL}{\parbox{\dimexpr\linewidth-2\fboxsep-2\fboxrule\relax}{\textbf{Result.} #1}}\par\smallskip}
\newcommand{\popboxhead}[3]{\poppill{#1}{#2}{white}\ifx\relax#3\relax\else\hspace{6pt}{\popxbold\fontsize{10.6}{12}\selectfont\color{popink}#3}\fi\par\vspace{7pt}}

\newtcolorbox{aretenir}[1][]{popbase,colframe=popgreen,before upper={\popboxhead{Key points}{popgreen}{#1}}}
\newtcolorbox{exemplebox}[1][]{popbase,colframe=poporange,before upper={\popboxhead{Exemple}{poporange}{#1}}}
\newtcolorbox{definition}[1][]{popbase,colframe=popblue,before upper={\popboxhead{Definition}{popblue}{#1}}}
\newtcolorbox{propriete}[1][]{popbase,colframe=poppurple,before upper={\popboxhead{Property}{poppurple}{#1}}}
\newtcolorbox{methode}[1][]{popbase,colframe=popgold,before upper={\popboxhead{Method}{popgold}{#1}}}
\newtcolorbox{attention}[1][]{popbase,colframe=poppink,before upper={\popboxhead{Warning}{poppink}{#1}}}
\newtcolorbox{experience}[1][]{popbase,colframe=popblue,colback=popblueL,before upper={\popboxhead{Experiment}{popblue}{#1}}}
% « Le savais-tu ? » : carte violette pleine (contexte, culture, Cameroun)
\newtcolorbox{savaistu}[1][]{popbase,colback=poppurple,colframe=popink,
  fontupper=\popsemi\fontsize{10.4}{14.2}\selectfont\color{white},
  before upper={\poppill{Did you know?}{white}{poppurple}\ifx\relax#1\relax\else\hspace{6pt}{\popxbold\color{white}#1}\fi\par\vspace{7pt}}}
% Exercice résolu : énoncé en haut, \tcblower puis la solution sur fond crème
\newcounter{popexo}\newcounter{popexr}
\newtcolorbox{exoresolu}[1][]{popbase,colframe=poporange,colbacklower=popcream,
  segmentation style={popink,line width=1.2pt,dashed},
  before upper={\stepcounter{popexr}\popboxhead{Worked example \thepopexr}{poporange}{#1}},
  before lower={\poppill{Solution}{popink}{popcream}\par\vspace{6pt}}}
% Exercice (non résolu en place) — la correction est dans la partie Corrigés
\newtcolorbox{exercice}[1][]{popbase,colframe=popink,
  before upper={\stepcounter{popexo}\popboxhead{Exercise \thepopexo}{popink}{#1}}}
% Corrigé
\newtcolorbox{corrige}[1][]{popbase,colframe=popgreen,colback=popgreenL,
  before upper={\popboxhead{Solution}{popgreen}{#1}}}
% Objectifs / prérequis / bilan
\newtcolorbox{objectifs}{popbase,colframe=popink,colback=poporangeL,
  before upper={\popboxhead{Chapter objectives}{popink}{}}}
\newtcolorbox{prerequis}{popbase,colframe=popink,colback=popblueL,
  before upper={\popboxhead{Prerequisites}{popblue}{}}}
\newtcolorbox{bilan}{popbase,colframe=popink,colback=white,boxrule=2.8pt,
  before upper={\popboxhead{Fiche bilan}{poporange}{L'essentiel à connaître pour l'examen}}}
% Figure encadrée avec légende
\newtcolorbox{popfigure}[1][]{enhanced,breakable=false,colback=white,colframe=popline,boxrule=1.4pt,arc=8pt,
  left=10pt,right=10pt,top=10pt,bottom=8pt,before skip=10pt,after skip=12pt,halign=center,
  lower separated=false,halign lower=center,
  fontlower=\popsemi\fontsize{9}{11.5}\selectfont\color{popmuted},
  #1}
\newcommand{\legende}[1]{\par\vspace{6pt}{\popsemi\fontsize{9}{11.5}\selectfont\color{popmuted}#1}}

% ============================================================
% Signature OnBuch+
% ============================================================
\newcommand{\poplogoinline}[1]{{\poplogo\fontsize{#1}{#1}\selectfont\color{popink}OnBuch\color{poporange}+}}
\newcommand{\signatureonbuch}{%
  \begin{tcolorbox}[enhanced,colback=popink,colframe=popink,boxrule=2.2pt,arc=10pt,
      drop shadow=poporange,left=14pt,right=14pt,top=10pt,bottom=10pt,before skip=0pt,after skip=0pt]
    \tikz[baseline=-0.7ex]\node[circle,fill=popgreen,draw=popcream,line width=1.3pt,minimum size=22pt,inner sep=0]{\popcheck{white}};%
    \hspace{9pt}\begin{minipage}[c]{0.62\linewidth}
      {\popxbold\fontsize{12}{13}\selectfont\color{popcream}Signed\ \textbullet\ {\poplogo OnBuch\color{poporange}+}}\par\vspace{2pt}
      {\raggedright\popsemi\fontsize{8}{10.5}\selectfont\color{popline}\MakeUppercase{OnBuch+ teaching team}\\\MakeUppercase{Follows the official MINESEC syllabus}\par}
    \end{minipage}\hfill
    {\poplogo\fontsize{22}{22}\selectfont\color{popcream}OnBuch\color{poporange}+}
  \end{tcolorbox}%
}

% ============================================================
% Page de garde
% #1 titre · #2 matière · #3 niveau · #4 module · #5 n° leçon · #6 durée ·
% #7 description / objectifs (texte ou itemize)
% ============================================================
\newcommand{\pagegarde}[7]{%
  \thispagestyle{empty}
  \newgeometry{a4paper,margin=1.7cm,top=1.6cm,bottom=1.4cm}%
  \begin{tikzpicture}[remember picture,overlay]
    \fill[poporange!18] ([xshift=-3.2cm,yshift=-3.6cm]current page.north east) circle (4.6cm);
    \fill[poppurple!14] ([xshift=2.4cm,yshift=7.6cm]current page.south west) circle (3.8cm);
  \end{tikzpicture}%
  {\poplogo\fontsize{30}{30}\selectfont\color{popink}OnBuch\color{poporange}+}\hfill
  \raisebox{0.6ex}{\poppill{Signed course \textbullet\ Premium}{popink}{popcream}}\par
  \vspace{1pt}\popsquiggle{poppurple}\par
  \vspace{8pt}
  \begin{tcolorbox}[enhanced,colback=poporange,colframe=popink,boxrule=2.8pt,arc=13pt,
      drop shadow=popink,left=18pt,right=18pt,top=12pt,bottom=13pt,after skip=10pt]
    \poppill{#2 \textbullet\ #3}{popink}{popcream}\hspace{5pt}\poppill{#4}{white}{popink}\par\vspace{8pt}
    {\popdisplay\fontsize{14}{14}\selectfont\color{popink}CHAPTER #5}\par\vspace{4pt}
    {\raggedright\hyphenpenalty=10000\exhyphenpenalty=10000\popdisplay\fontsize{25}{27}\selectfont\color{popcream}\MakeUppercase{#1}\par}
    \vspace{8pt}
    {\popxbold\fontsize{9.5}{9.5}\selectfont\color{popink}\MakeUppercase{Suggested time: #6}}
  \end{tcolorbox}
  \begin{tcolorbox}[enhanced,colback=white,colframe=popink,boxrule=2.2pt,arc=10pt,
      drop shadow=popink,left=16pt,right=16pt,top=10pt,bottom=10pt,after skip=8pt]
    \poppill{In this chapter}{poppurple}{white}\par\vspace{5pt}
    \popsemi\fontsize{9.3}{12.6}\selectfont\color{popink}#7
  \end{tcolorbox}
  \noindent\begin{minipage}[t]{0.487\linewidth}
    \begin{tcolorbox}[enhanced,colback=popgreen,colframe=popink,boxrule=2.2pt,arc=10pt,
        drop shadow=popink,left=11pt,right=11pt,top=10pt,bottom=10pt,height=3.1cm,valign=center]
      \tcbox[colback=white,colframe=popink,boxrule=1.3pt,arc=5pt,boxsep=0pt,nobeforeafter,
        left=3pt,right=3pt,top=3pt,bottom=3pt,valign=center]{\qrcode[height=1.9cm]{https://play.google.com/store/apps/details?id=com.luvvix.onbuch&hl=en}}%
      \hspace{8pt}\begin{minipage}[c]{0.5\linewidth}
        {\popdisplay\fontsize{11}{12.5}\selectfont\color{white}\MakeUppercase{Download OnBuch}}\par\vspace{4pt}
        {\raggedright\popsemi\fontsize{8.4}{11}\selectfont\color{white}Free \textbullet\ Google Play\\AI tutor, past papers, results\par}
      \end{minipage}
    \end{tcolorbox}
  \end{minipage}\hfill
  \begin{minipage}[t]{0.487\linewidth}
    \begin{tcolorbox}[enhanced,colback=poppurple,colframe=popink,boxrule=2.2pt,arc=10pt,
        drop shadow=popink,left=11pt,right=11pt,top=10pt,bottom=10pt,height=3.1cm,valign=center]
      \tikz[baseline=-0.7ex]\node[circle,fill=white,draw=popink,line width=1.3pt,minimum size=1.6cm,inner sep=0]{\popdisplay\fontsize{24}{24}\selectfont\color{poppurple}+};%
      \hspace{8pt}\begin{minipage}[c]{0.6\linewidth}
        {\popdisplay\fontsize{11}{12.5}\selectfont\color{white}\MakeUppercase{Go OnBuch+}}\par\vspace{4pt}
        {\raggedright\popsemi\fontsize{8.4}{11}\selectfont\color{white}All signed courses, unlimited AI tutor, PDF sheets — OnBuch+ tab in the app\par}
      \end{minipage}
    \end{tcolorbox}
  \end{minipage}
  \vspace{8pt}
  \popcontenu
  \vfill
  \signatureonbuch
  \restoregeometry
  \newpage
}

% Rangée « Ce cours contient » (page de garde)
\newcommand{\poptile}[3]{% #1 couleur #2 gros texte #3 légende
  \begin{tcolorbox}[enhanced,colback=#1,colframe=popink,boxrule=2pt,arc=9pt,shadow={2.5pt}{-2.5pt}{0pt}{popink},
      left=6pt,right=6pt,top=9pt,bottom=9pt,halign=center,valign=center,height=1.75cm,width=\linewidth,nobeforeafter]
    {\popdisplay\fontsize{12.5}{13}\selectfont\color{white}\MakeUppercase{#2}}\par\vspace{3pt}
    {\centering\popsemi\fontsize{7.8}{9.5}\selectfont\color{white}#3\par}
  \end{tcolorbox}}
\newcommand{\popcontenu}{%
  \par\noindent{\popxboldls\fontsize{7.5}{7.5}\selectfont\color{popmuted}THIS SIGNED COURSE CONTAINS}\par\vspace{7pt}
  \noindent\begin{minipage}[t]{0.235\linewidth}\poptile{popblue}{Course}{complete, illustrated, step by step}\end{minipage}\hfill
  \begin{minipage}[t]{0.235\linewidth}\poptile{poporange}{Methods}{and worked examples}\end{minipage}\hfill
  \begin{minipage}[t]{0.235\linewidth}\poptile{poppink}{Exercises}{exam style + solutions}\end{minipage}\hfill
  \begin{minipage}[t]{0.235\linewidth}\poptile{popgreen}{Summary sheet}{the essentials to remember}\end{minipage}\par}

% Sommaire
\newtcolorbox{sommaire}{enhanced,colback=white,colframe=popink,boxrule=2.2pt,arc=10pt,
  drop shadow=popink,left=18pt,right=18pt,top=13pt,bottom=13pt,before skip=6pt,after skip=20pt,
  before upper={\poppill{Contents}{poppurple}{white}\par\vspace{9pt}\popsemi\fontsize{10.6}{14.5}\selectfont\color{popink}}}

% Signature de fin de cours — poussée tout en bas de la dernière page
\newcommand{\findecours}{%
  \par\vfill\nopagebreak
  \begin{center}\popsquiggle{poporange}\end{center}\vspace{4pt}
  \signatureonbuch
}
\AtBeginDocument{\def\llbracket{\mathopen{[\mkern-3.5mu[}}\def\rrbracket{\mathclose{]\mkern-3.5mu]}}} % intervalles d'entiers (glyphe ⟦ absent de la police)
% Glyphes absents de Fira Math : redéfinis à partir de glyphes disponibles
\AtBeginDocument{%
  \def\setminus{\mathbin{\backslash}}\let\smallsetminus\setminus
  \def\vdots{\mathord{\vbox{\baselineskip3.2pt\lineskiplimit0pt\kern4pt\hbox{.}\hbox{.}\hbox{.}}}}%
  \def\ddots{\mathinner{\mkern1mu\raise6.5pt\hbox{.}\mkern2mu\raise3.6pt\hbox{.}\mkern2mu\raise0.7pt\hbox{.}\mkern1mu}}%
  \def\triangle{\mathord{\scalebox{0.9}{$\symup{\Delta}$}}}%
  \def\nabla{\mathord{\raisebox{\depth}{\scalebox{1}[-1]{$\symup{\Delta}$}}}}%
  \def\checkmark{\popcheck{popdark}}%
  \def\star{\mathbin{\ast}}\def\bigstar{\mathord{\ast}}%
  \def\diamond{\mathbin{\scalebox{0.6}{$\Diamond$}}}%
  \def\bigcup{\mathop{\mathchoice{\raisebox{-0.3em}{\scalebox{1.7}{$\cup$}}}{\scalebox{1.25}{$\cup$}}{\cup}{\cup}}\displaylimits}%
  \def\bigcap{\mathop{\mathchoice{\raisebox{-0.3em}{\scalebox{1.7}{$\cap$}}}{\scalebox{1.25}{$\cap$}}{\cap}{\cap}}\displaylimits}%
  \def\lesseqqgtr{\mathrel{\lessgtr}}\def\bigoplus{\mathop{\oplus}\displaylimits}%
}
\providecommand{\ssi}{if and only if} % abréviation fréquente
\AtBeginDocument{\providecommand{\wideparen}{\overparen}} % arc de cercle

%% FIGFIX-BEGIN v5 — correctifs globaux des figures (docs/onbuch-generation/tools/figfix)
\usepackage{adjustbox}\usepackage{etoolbox}\tcbuselibrary{raster}
% toute figure TikZ/pgfplots plus large que la ligne est réduite à la largeur disponible
\BeforeBeginEnvironment{tikzpicture}{\begin{adjustbox}{max width=\linewidth}}
\AfterEndEnvironment{tikzpicture}{\end{adjustbox}}
% légendes pgfplots lisibles par-dessus les courbes
\pgfplotsset{every axis/.append style={legend style={fill=white,fill opacity=0.92,text opacity=1,draw=black!15,inner sep=3pt}}}
% étiquettes d'arêtes (syntaxe edge["..."]) avec fond blanc
\tikzset{every edge quotes/.append style={fill=white,inner sep=1.2pt}}
%% FIGFIX-END
\begin{document}

\pagegarde{Response \& coordination in organisms}{Biology}{Upper Sixth}{Upper Sixth — Biology}{2}{25 periods}{This chapter explores how organisms detect stimuli and coordinate responses through sense organs, the nervous system and the endocrine system in animals, and through hormones and tropic movements in plants. It provides the structural, physiological and molecular foundations needed for GCE Advanced Level questions on reception, coordination and control.
\vspace{4pt}
\begin{multicols}{2}\small
\begin{itemize}
\item By the end of this chapter I can define stimulus, response, receptor, effector and coordination, and classify stimuli and responses.
\item By the end of this chapter I can explain the principles of stimulus reception, including sensory thresholds, adaptation and the generator potential.
\item By the end of this chapter I can describe the structure of the mammalian eye and relate each part to its function.
\item By the end of this chapter I can explain image formation, accommodation, control of light intensity, light/dark adaptation and the photochemistry of photoreception.
\item By the end of this chapter I can describe the structure of the ear and explain the mechanisms of hearing and balance.
\item By the end of this chapter I can describe the organisation of the nervous system, the structure of neurones and the generation and transmission of nerve impulses.
\end{itemize}\end{multicols}}

\begin{sommaire}
\begin{multicols}{2}
\begin{enumerate}
\item Stimuli, Responses and Principles of Reception
\item The Eye: Structure and Image Formation
\item Control of Light Intensity, Adaptation and Accommodation
\item Photoreception and the Ear: Hearing and Balance
\item The Nervous System and Neurones
\item Transmission of Impulses and the Synapse
\item Reflex Actions, Brain, Spinal Cord and Autonomic Nervous System
\item The Endocrine System and Mechanism of Hormone Action
\item Functions of the Endocrine Glands
\item Plant Movements, Plant Hormones, Light and Plant Growth
\item Integration activity
\item Methods and worked examples
\item Exercises
\item Solutions to the exercises
\item Summary sheet
\end{enumerate}
\end{multicols}
\end{sommaire}

\begin{objectifs}
\begin{itemize}
\item By the end of this chapter I can define stimulus, response, receptor, effector and coordination, and classify stimuli and responses.
\item By the end of this chapter I can explain the principles of stimulus reception, including sensory thresholds, adaptation and the generator potential.
\item By the end of this chapter I can describe the structure of the mammalian eye and relate each part to its function.
\item By the end of this chapter I can explain image formation, accommodation, control of light intensity, light/dark adaptation and the photochemistry of photoreception.
\item By the end of this chapter I can describe the structure of the ear and explain the mechanisms of hearing and balance.
\item By the end of this chapter I can describe the organisation of the nervous system, the structure of neurones and the generation and transmission of nerve impulses.
\item By the end of this chapter I can explain synaptic transmission, reflex actions, and the roles of the brain, spinal cord and autonomic nervous system.
\item By the end of this chapter I can describe the endocrine system, the mechanism of hormone action and the functions of the main endocrine glands.
\item By the end of this chapter I can compare nervous and hormonal coordination and explain plant movements, plant hormones and the effect of light on plant growth.
\end{itemize}
\end{objectifs}

\begin{prerequis}
\begin{itemize}
\item Cell structure and ultrastructure (Lower Sixth): membranes, mitochondria, organelles
\item Movement of substances across membranes: diffusion, osmosis, active transport
\item Protein structure and enzymes as biological catalysts
\item Basic chemistry of ions (Na+, K+) and concentration gradients
\item Cell division and growth in plants (Lower Sixth)
\end{itemize}
\end{prerequis}

\newpage

\coursec{Stimuli, Responses and Principles of Reception}

During an evening football match at the Stade Omnisports in Yaoundé, a striker hears the referee's whistle, sees the ball emerging from the shadows into the floodlights, and volleys it into the net — all in less than a second. How did his eyes adjust instantly to the changing light intensity? How did the sound of the whistle travel from his ear to his leg muscles? And why do the bean seedlings on his grandmother's farm in Bafoussam bend toward the morning sun? Behind each of these events lies the same fundamental biological machinery: \cle{receptors} that detect changes, \cle{nerves} and \cle{hormones} that carry and integrate information, and \cle{effectors} that produce a precise, coordinated response. In this first section we lay the foundations: what stimuli and responses are, how receptors work, and how the strength of a stimulus is coded in the language of the nervous system.

\courssub{Stimulus, Response and Coordination: the Basic Vocabulary}

Every living organism, from an \emph{Amoeba} in a pond near Lake Chad to the striker on the pitch, survives only because it can detect changes in its environment and react to them appropriately. A wilting hibiscus plant recovers after watering; a goat in Ngaoundéré runs when it hears a dog bark; you withdraw your hand from a hot cooking pot before you even feel the pain. All these are examples of the same sequence: a change is detected, information is processed, and a reaction follows.

\begin{definition}[Stimulus, response, receptor, effector, coordination]
\begin{itemize}
\item A \cle{stimulus} is any change in the internal or external environment of an organism that is capable of provoking a response.
\item A \cle{response} is the observable reaction of an organism (or of part of it) to a stimulus.
\item A \cle{receptor} (or sensor) is a cell or organ specialised to detect a particular type of stimulus and convert it into an electrical signal.
\item An \cle{effector} is an organ that carries out the response — in animals, usually a \textbf{muscle} (which contracts) or a \textbf{gland} (which secretes).
\item \cle{Coordination} is the linking together and control of the activities of receptors, conducting pathways and effectors so that the response is appropriate, correctly timed and correctly proportioned. In animals, coordination is achieved by the \textbf{nervous system} (fast, electrical) and the \textbf{endocrine system} (slower, chemical); in plants, by plant growth substances (hormones).
\end{itemize}
\end{definition}

Stimuli are classified by their origin. \textbf{External stimuli} come from outside the body: light, sound, touch, the smell of ripe mangoes, a change in air temperature. \textbf{Internal stimuli} arise inside the body: a fall in blood glucose, a rise in blood \ce{CO2}, stretching of the full bladder, thirst caused by concentrated blood plasma. Internal stimuli are just as important — they drive \emph{homeostasis}, the maintenance of a constant internal environment.

\begin{aretenir}[The coordination pathway]
Every coordinated response follows the same five-step plan:
\[
\text{stimulus} \;\longrightarrow\; \text{receptor} \;\longrightarrow\; \text{coordinator} \;\longrightarrow\; \text{effector} \;\longrightarrow\; \text{response}
\]
The coordinator (brain, spinal cord, or a simple nerve net in lower animals) receives the message, \emph{integrates} it, and dispatches instructions to the effector.
\end{aretenir}

\begin{popfigure}
\begin{center}
\begin{tikzpicture}[
box/.style={draw=popblue, thick, rounded corners=3pt, fill=popblueL, text width=2.4cm, align=center, minimum height=1cm, font=\small},
arr/.style={-{Stealth[length=3mm]}, thick, popdark}]
\node[box, align=center] (s) at (0,0){\textbf{Stimulus}\\ e.g. whistle};
\node[box, align=center] (r) at (3.4,0){\textbf{Receptor}\\ ear (cochlea)};
\node[box, align=center] (c) at (6.8,0){\textbf{Coordinator}\\ brain / spinal cord};
\node[box, align=center] (e) at (10.2,0){\textbf{Effector}\\ leg muscles};
\node[box, fill=popgreenL, draw=popgreen, align=center] (resp) at (13.6,0){\textbf{Response}\\ striker shoots};
\draw[arr] (s) -- (r);
\draw[arr] (r) -- (c) node[midway, above, font=\scriptsize, text=popmuted]{impulse};
\draw[arr] (c) -- (e) node[midway, above, font=\scriptsize, text=popmuted]{impulse};
\draw[arr] (e) -- (resp);
\end{tikzpicture}
\end{center}
\legende{The five stages of a coordinated response, illustrated by the striker reacting to the referee's whistle.}
\end{popfigure}

\courssub{Classification of Stimuli and of Receptors}

Stimuli exist in many physical forms, and evolution has produced a receptor specialised for each. The main forms of stimulus energy are:

\begin{center}
\begin{tabularx}{\linewidth}{|l|X|X|}
\hline
\rowcolor{popblueL} \textbf{Stimulus energy} & \textbf{Examples of stimuli} & \textbf{Receptor type (examples)} \\
\hline
Light (electromagnetic) & brightness, colour, day length & \textbf{Photoreceptors}: rods and cones of the retina; plant phytochrome \\
\hline
Sound / vibration (mechanical) & whistle, footsteps, predator approach & \textbf{Mechanoreceptors}: hair cells of the cochlea \\
\hline
Chemical & taste, smell, blood glucose, pheromones & \textbf{Chemoreceptors}: taste buds, olfactory cells, carotid bodies \\
\hline
Mechanical (pressure, touch, stretch) & touch, blood pressure, muscle stretch & \textbf{Mechanoreceptors}: touch corpuscles in skin, stretch receptors in muscle \\
\hline
Temperature & heat of a cooking pot, cold harmattan wind & \textbf{Thermoreceptors}: warm and cold receptors in skin and hypothalamus \\
\hline
Electrical & electric fields in water & \textbf{Electroreceptors}: organs of electric fish (e.g. mormyrids of the River Sanaga) \\
\hline
\end{tabularx}
\end{center}

Receptors are therefore named after the stimulus they detect: \cle{photoreceptors} (light), \cle{mechanoreceptors} (pressure, touch, sound, gravity), \cle{chemoreceptors} (chemicals), \cle{thermoreceptors} (temperature) and \cle{electroreceptors} (electric fields). Pain receptors, which respond to tissue damage by several kinds of stimulus, are called \textbf{nociceptors}.

\begin{attention}[Do not confuse response and impulse]
A \textbf{response} is the visible action of the whole organism or of an organ (the hand withdraws, the muscle contracts, the seedling bends). An \textbf{impulse} (action potential) is the invisible electrical signal travelling along a neurone. The impulse \emph{causes} the response but is not itself the response. In examinations, writing "the response travels along the nerve" is a classic error — it is the \emph{impulse} that travels.
\end{attention}

\courssub{Types of Responses in Organisms}

Responses vary enormously in speed, complexity and flexibility. At GCE Advanced Level you must recognise five broad categories:

\begin{itemize}
\item \textbf{Taxes (singular: taxis)} — a directed movement of a whole motile organism (or cell) toward or away from a stimulus. Movement \emph{toward} light is \emph{positive phototaxis} (\emph{Euglena} swimming toward light); movement \emph{away} from light is \emph{negative phototaxis} (blowfly larvae). Chemotaxis, geotaxis and rheotaxis (response to water current) follow the same logic.
\item \textbf{Kineses (singular: kinesis)} — a non-directional change in the \emph{rate} of movement in response to the intensity of a stimulus. Woodlice move faster and turn more often in dry, bright conditions, and slow down in humid, dark ones; by chance they accumulate where conditions favour survival. There is no "toward" or "away" — only a change of activity.
\item \textbf{Tropisms} — directional \emph{growth} movements of a fixed part of a plant in response to a directional stimulus: \emph{phototropism} (shoots bending toward light), \emph{geotropism/gravitropism} (roots growing downward), \emph{hydrotropism} and \emph{thigmotropism} (tendrils coiling around a support). These will be studied in detail in Section 10.
\item \textbf{Reflexes} — rapid, automatic, inborn (innate) responses to specific stimuli, usually protective: withdrawal from a hot object, blinking, the knee-jerk. They do not require learning and often do not initially involve the brain.
\item \textbf{Learned behaviour} — responses modified by experience: habituation, conditioning, trial-and-error learning and insight. A goalkeeper learns to anticipate a penalty taker's run-up; no reflex arc can explain that.
\end{itemize}

\courssub{Properties of Receptors}

However diverse they are, all receptors share four fundamental properties that you must be able to define and illustrate.

\textbf{1. Specificity.} Each receptor responds best to one particular form of stimulus energy — its \emph{adequate stimulus}. The retina responds to light, not to sound; the cochlea to vibration, not to taste. This is why each receptor type has a characteristic structure.

\textbf{2. Sensitivity and sensory threshold.} Receptors are remarkably sensitive: the human eye can detect a flash of only a few photons under dark-adapted conditions. But every receptor has a \cle{sensory threshold} — the minimum stimulus intensity that produces a detectable response (an impulse in the sensory neurone). Below threshold, nothing happens; above threshold, impulses are fired.

\textbf{3. Adaptation.} If a constant stimulus is maintained, many receptors gradually reduce their firing rate — they \emph{adapt}. You stop noticing your shirt on your skin or the smell of your own kitchen minutes after entering. Adaptation prevents the nervous system from being flooded with unchanging, unimportant information, leaving it free to detect \emph{changes}. Receptors such as touch receptors (Meissner's corpuscles) adapt rapidly; pain receptors (nociceptors) and muscle stretch receptors (muscle spindles) adapt slowly or not at all — a slow-adapting pain receptor keeps warning you as long as the damage persists.

\textbf{4. Transduction.} This is the master property, stated formally below.

\begin{propriete}[Transduction: receptors are biological energy converters]
A receptor converts (transduces) the energy of the stimulus — light, mechanical, chemical or thermal — into an electrical signal: a local depolarisation of its membrane called the \cle{generator potential} (or receptor potential). If the generator potential reaches the threshold of the sensory neurone, it triggers \textbf{action potentials} (nerve impulses) that travel to the central nervous system. (The mechanism of the action potential itself is examinable in Section 6.)
\end{propriete}

The generator potential has two crucial features: it is \emph{graded} (its size increases with stimulus intensity, unlike the all-or-nothing action potential) and it is \emph{local} (it does not propagate; it only triggers impulses if threshold is reached).

\begin{savaistu}[Receptor diversity in Cameroon's ecosystems]
The mormyrid electric fish of the River Sanaga and other Cameroonian rivers possess specialised electroreceptors (Knollenorgans and mormyromasts) that detect weak electric fields. They use these for navigation in murky water, communication, and locating prey — a remarkable example of a sensory modality absent in most terrestrial vertebrates.
\end{savaistu}

\courssub{Frequency Coding: How Stimulus Intensity is Represented}

A single action potential is always the same size — it obeys the \emph{all-or-nothing law}. So how does the brain know whether the light is dim or dazzling, the touch gentle or painful? The answer is \cle{frequency coding}: as stimulus intensity increases, the generator potential becomes larger, the threshold is re-crossed sooner after each impulse, and therefore the \textbf{frequency of action potentials increases}. A strong stimulus also recruits \emph{more} sensory neurones (spatial summation). The brain reads intensity from the number of impulses per second (and the number of active fibres), not from the size of each impulse.

\begin{popfigure}
\begin{center}
\begin{tikzpicture}
\begin{axis}[
width=11cm, height=6.5cm,
xlabel={Stimulus intensity (arbitrary units)},
ylabel={Impulse frequency (impulses s$^{-1}$)},
xmin=0, xmax=10, ymin=0, ymax=110,
xtick={0,2,4,6,8,10}, ytick={0,20,40,60,80,100},
grid=both, grid style={popline, dashed, very thin},
axis lines=left, thick,
legend style={at={(0.03,0.97)}, anchor=north west, font=\small}]
\addplot[popblue, very thick, smooth, domain=0:10, samples=60] {100*(1 - exp(-0.45*x))};
\addlegendentry{Impulse frequency}
\addplot[poporange, dashed, thick, domain=0:10, samples=2] {100};
\addlegendentry{Maximum (refractory limit)}
\draw[popdark, dashed] (axis cs:1.2,0) -- (axis cs:1.2,41.6);
\node[font=\scriptsize, text=popdark, anchor=west] at (axis cs:1.25,20) {threshold};
\end{axis}
\end{tikzpicture}
\end{center}
\legende{Frequency coding: the frequency of action potentials rises with stimulus intensity and levels off at a maximum set by the refractory period of the neurone. Below the threshold intensity, no impulses are fired.}
\end{popfigure}

\begin{exemplebox}[Reading a frequency-coded message]
A pressure receptor in the skin fires at 10 impulses per second when a feather touches it and at 80 impulses per second when a finger presses firmly. Each impulse is identical in size (about $+40\ \mathrm{mV}$ at its peak). The brain interprets the \emph{eightfold increase in frequency} — not any change in impulse size — as "firmer pressure". If the pressure becomes painful, additional nociceptor fibres are recruited, adding a second channel of information.
\end{exemplebox}

\courssub{Investigating Responses: the Choice Chamber Experiment}

The distinction between taxes and kineses, and the reality of stimulus--response behaviour, can be demonstrated in the school laboratory with small invertebrates such as woodlice or earthworms.

\begin{methode}[Using a choice chamber to study responses of woodlice to humidity and light]
\begin{enumerate}
\item \textbf{Set up the chamber.} A choice chamber is a shallow dish divided into two (or four) connected compartments. To test humidity, place damp filter paper (or moist cotton wool) under one half and a drying agent such as anhydrous calcium chloride or silica gel under the other. To test light, cover one half with opaque black card and leave the other half transparent under a lamp.
\item \textbf{Introduce the animals.} Place about 10 woodlice in the centre through a small hole, and leave them undisturbed for about 10 minutes to settle.
\item \textbf{Record.} Count the number of woodlice in each compartment at regular intervals (e.g. every minute for 10 minutes).
\item \textbf{Repeat and replicate.} Repeat the experiment several times, rotating the chamber 180\textdegree\ between runs to cancel any bias from the room (e.g. a window on one side), and pool results from several groups.
\item \textbf{Analyse.} Compare the numbers in each condition with the numbers expected by chance (a chi-squared test may be used). A consistent, significant accumulation in the dark, humid side demonstrates a behavioural response.
\end{enumerate}
\end{methode}

\begin{experience}[Expected results and interpretation]
Woodlice typically accumulate in the \textbf{dark, humid} compartment. Careful observation shows they do not crawl straight there: in the dry, bright side they move faster and turn more frequently (orthokinesis and klinokinesis); in the humid, dark side they slow down and settle. This is a \textbf{kinesis}, not a directed taxis. Its survival value is clear: woodlice lose water easily through their permeable cuticle, so behaviour that traps them in humid refuges — under stones and rotting logs — reduces desiccation. Earthworms tested similarly avoid bright light (negative phototactic behaviour), keeping them in moist soil where gas exchange through the skin is possible.
\end{experience}

\begin{attention}[Controls and fair testing]
Common mistakes in this practical: forgetting to rotate the chamber (so a room light gradient fakes a "preference"); testing too few animals or a single run (chance clumping is misread as a response); and describing woodlice behaviour as "moving toward humidity" — a kinesis is \emph{not} directional movement toward the stimulus. Always state the control, replicate, and use correct vocabulary.
\end{attention}

\begin{exoresolu}[Worked exercise: stimuli, receptors and coding]
\textbf{Statement.} A student touches a warm water bath at $35\ ^{\circ}\mathrm{C}$, then a hot bath at $55\ ^{\circ}\mathrm{C}$. (a) Identify the stimulus, the receptor type, the effector and the response when she withdraws her hand from the hot bath. (b) Explain, in terms of generator potentials and frequency coding, how her nervous system distinguishes the warm bath from the hot bath. (c) State the type of response shown by the withdrawal, and justify your answer.
\tcblower
\textbf{Solution.}
(a) The \emph{stimulus} is the high temperature of the water (an external, thermal stimulus). The \emph{receptors} are thermoreceptors (and, at $55\ ^{\circ}\mathrm{C}$, nociceptors) in the skin of the fingers. The \emph{effectors} are the flexor muscles of the arm, and the \emph{response} is contraction of these muscles, withdrawing the hand.
(b) Heat energy is transduced by the thermoreceptors into a generator potential. At $35\ ^{\circ}\mathrm{C}$ the generator potential is small and produces a low frequency of action potentials in the sensory neurones. At $55\ ^{\circ}\mathrm{C}$ the generator potential is much larger, so threshold is re-crossed more rapidly and the impulse frequency is much higher; in addition, more receptor fibres (including nociceptors) are recruited. The central nervous system interprets the higher frequency and greater number of active fibres as a more intense, potentially damaging stimulus. The size of each action potential does not change — intensity is coded by frequency.
(c) The withdrawal is a \emph{reflex} (the withdrawal reflex): it is rapid, automatic, innate and protective, and it is mediated by a reflex arc through the spinal cord, so it occurs before the sensation of pain is consciously perceived.
\end{exoresolu}

\begin{exercice}[Check your understanding]
\begin{enumerate}
\item Define the terms \emph{stimulus}, \emph{receptor} and \emph{coordination}.
\item Classify each of the following receptors by the stimulus energy detected: (i) rods of the retina; (ii) taste buds; (iii) hair cells of the cochlea; (iv) cold receptors of the skin.
\item Distinguish between a taxis and a kinesis, giving one example of each.
\item Explain why a constant touch on your arm eventually goes unnoticed, naming the receptor property involved.
\item Sketch and annotate the graph of impulse frequency against stimulus intensity, indicating the threshold and the maximum frequency.
\end{enumerate}
\end{exercice}

\begin{corrige}[Exercise 1]
\begin{enumerate}
\item Stimulus: any change in the internal or external environment capable of provoking a response. Receptor: a cell or organ specialised to detect a specific stimulus and convert it into an electrical signal. Coordination: the linking and control of receptors, conducting pathways and effectors so that responses are appropriate and well timed.
\item (i) Photoreceptor (light); (ii) chemoreceptor (chemicals); (iii) mechanoreceptor (sound vibration); (iv) thermoreceptor (temperature).
\item A taxis is a \emph{directed} movement of a whole organism toward or away from a stimulus (e.g. \emph{Euglena} showing positive phototaxis). A kinesis is a \emph{non-directional} change in the rate of movement with stimulus intensity (e.g. woodlice moving faster in dry conditions).
\item The touch receptors show \emph{adaptation}: with a constant, unchanging stimulus their generator potentials decline and the frequency of impulses falls, so the brain is no longer informed of the unchanged stimulus.
\item The curve starts at zero below the threshold intensity, rises steeply once threshold is crossed, then levels off at a maximum frequency fixed by the refractory period of the neurone (as in the figure above). Axes must be labelled: stimulus intensity (x) and impulse frequency in impulses s$^{-1}$ (y).
\end{enumerate}
\end{corrige}

\begin{aretenir}[Section summary]
\begin{itemize}
\item A stimulus is a detectable environmental change; the response is the organism's reaction, carried out by effectors (muscles, glands) under the control of a coordinator.
\item The universal pathway is: stimulus $\rightarrow$ receptor $\rightarrow$ coordinator $\rightarrow$ effector $\rightarrow$ response.
\item Receptors are classified by the energy they detect: photo-, mechano-, chemo-, thermo- and electroreceptors.
\item Responses include taxes, kineses, tropisms, reflexes and learned behaviour.
\item Receptors are specific, sensitive (with a threshold), show adaptation, and perform \textbf{transduction}: stimulus energy $\rightarrow$ generator potential $\rightarrow$ action potentials.
\item Stimulus intensity is coded by the \emph{frequency} of action potentials (and the number of fibres recruited), never by impulse size.
\item The choice chamber demonstrates kinesis in woodlice: they accumulate in dark, humid conditions by changing their rate of movement, not by directed travel.
\end{itemize}
\end{aretenir}

\coursec{The Eye: Structure and Image Formation}

In the previous section we established that a \cle{stimulus} is any change in the internal or external environment capable of eliciting a response, and that \cle{reception} is the process by which specialised cells (receptors) convert stimulus energy into nerve impulses. The mammalian eye is the supreme example of a photoreceptor organ: it collects, focuses and detects light, then transmits coded information to the brain for interpretation. We now examine its anatomy, the physics of image formation, and the common optical defects that GCE examiners frequently test.

\courssub{Gross Structure of the Mammalian Eye}

The eye is a hollow, roughly spherical organ housed in the bony orbit of the skull. Its wall consists of three concentric layers (tunics), and its interior is filled with transparent humours that maintain shape and transmit light. The table below summarises the major structures and their functions --- learn one function for every label; the GCE often asks for \emph{structure $\to$ function} matching.

\begin{center}
\begin{tabularx}{\linewidth}{|l|X|}\hline
\rowcolor{popblueL}
\textbf{Structure} & \textbf{Main Function(s)} \\ \hline
\cle{Sclera} & Tough, white, fibrous outer coat; protects the eye and maintains its shape; provides attachment for the extrinsic eye muscles. \\ \hline
\cle{Choroid} & Middle vascular layer, rich in blood vessels and dark pigment (melanin); supplies nutrients and oxygen to the retina; the pigment absorbs stray light, preventing internal reflection. \\ \hline
\cle{Retina} & Innermost light-sensitive (nervous) layer; contains the photoreceptors (rods and cones) and their neurones; site of conversion of light into nerve impulses. \\ \hline
\cle{Cornea} & Transparent anterior continuation of the sclera; provides \emph{most} of the eye's refractive (light-bending) power; avascular, nourished by tears and aqueous humour. \\ \hline
\cle{Conjunctiva} & Thin, transparent membrane covering the anterior sclera and lining the eyelids; protects and lubricates the front of the eye. \\ \hline
\cle{Iris} & Pigmented muscular diaphragm (containing circular and radial muscles) that controls the diameter of the pupil and hence the amount of light entering the eye. \\ \hline
\cle{Pupil} & Central aperture (hole) in the iris through which light enters; appears black because the entering light is absorbed inside the eye. \\ \hline
\cle{Lens} & Transparent, biconvex, elastic structure; fine-tunes focusing by changing shape (accommodation). \\ \hline
\cle{Ciliary body} & Ring of smooth muscle (ciliary muscle); secretes aqueous humour; alters lens shape via the suspensory ligaments. \\ \hline
\cle{Suspensory ligaments} & Fibres connecting the ciliary body to the lens; transmit tension that stretches and flattens the lens. \\ \hline
\cle{Aqueous humour} & Clear, watery fluid in the chambers in front of the lens; nourishes the cornea and lens; maintains the shape of the front of the eye. \\ \hline
\cle{Vitreous humour} & Clear, jelly-like substance filling the large chamber behind the lens; supports the retina and maintains the spherical shape of the eyeball. \\ \hline
\cle{Optic nerve} & Bundle of sensory neurones carrying visual impulses from the retina to the brain. \\ \hline
\cle{Fovea (yellow spot)} & Tiny region of the retina on the visual axis; contains \emph{only cones} at very high density; region of \emph{sharpest, most detailed colour vision}. \\ \hline
\cle{Blind spot (optic disc)} & Point where the optic nerve leaves the eye and blood vessels enter; contains \emph{no photoreceptors}, so no image is detected there. \\ \hline
\end{tabularx}
\end{center}

\begin{popfigure}
\begin{tikzpicture}[scale=1.05]
% Outer wall: sclera
\draw[popdark, line width=1.6pt] (0,0) ellipse (4 and 3);
% Choroid
\draw[poppurple, line width=1pt] (0,0) ellipse (3.8 and 2.8);
% Retina (inner, posterior two-thirds)
\draw[popblue, line width=1.2pt] (3.55,1.4) arc (21:339:3.65 and 2.65);
% Cornea (anterior bulge, right side)
\draw[popgreen, line width=1.6pt] (3.8,1.6) .. controls (4.9,0.9) and (4.9,-0.9) .. (3.8,-1.6);
% Iris
\draw[popgold, line width=1.6pt] (2.6,1.5) -- (3.3,0.55);
\draw[popgold, line width=1.6pt] (2.6,-1.5) -- (3.3,-0.55);
% Lens (biconvex ellipse)
\draw[poporange, line width=1.4pt] (2.2,0) ellipse (0.45 and 1.25);
% Ciliary bodies
\draw[poppink, line width=2.2pt] (2.55,1.35) -- (2.95,1.75);
\draw[poppink, line width=2.2pt] (2.55,-1.35) -- (2.95,-1.75);
% Suspensory ligaments
\draw[popmuted, thin] (2.2,1.25) -- (2.7,1.55);
\draw[popmuted, thin] (2.2,-1.25) -- (2.7,-1.55);
% Optic nerve (posterior, left)
\draw[popdark, line width=2.4pt] (-3.9,0.45) -- (-5.1,0.45);
\draw[popdark, line width=2.4pt] (-3.9,-0.45) -- (-5.1,-0.45);
% Fovea marker (slightly above optic axis on retina)
\fill[popgreen] (-3.55,0.95) circle (0.07);
% Blind spot marker
\fill[popink] (-3.62,0) circle (0.08);
% Labels with leader lines
\node[popdark, anchor=east] at (-5.3,2.4) {Sclera};
\draw[popdark, thin] (-5.2,2.3) -- (-2.6,2.75);
\node[poppurple, anchor=east] at (-5.3,1.7) {Choroid};
\draw[poppurple, thin] (-5.2,1.6) -- (-2.4,2.55);
\node[popblue, anchor=east] at (-5.3,1.0) {Retina};
\draw[popblue, thin] (-5.2,0.9) -- (-2.2,2.35);
\node[popdark, anchor=east] at (-5.3,-0.9) {Optic nerve};
\draw[popdark, thin] (-5.2,-0.8) -- (-4.6,-0.2);
\node[popink, anchor=east] at (-5.3,-1.6) {Blind spot};
\draw[popink, thin] (-5.2,-1.5) -- (-3.7,-0.1);
\node[popgreen, anchor=east] at (-5.3,-2.3) {Fovea (yellow spot)};
\draw[popgreen, thin] (-5.2,-2.2) -- (-3.5,0.85);
\node[popgreen, anchor=west] at (5.1,2.2) {Cornea};
\draw[popgreen, thin] (5.0,2.1) -- (4.5,0.9);
\node[popgold, anchor=west] at (5.1,1.4) {Iris};
\draw[popgold, thin] (5.0,1.3) -- (3.35,0.7);
\node[popdark, anchor=west] at (5.1,0.6) {Pupil};
\draw[popdark, thin] (5.0,0.5) -- (3.35,0.1);
\node[poporange, anchor=west] at (5.1,-0.3) {Lens};
\draw[poporange, thin] (5.0,-0.4) -- (2.65,-0.3);
\node[poppink, anchor=west] at (5.1,-1.1) {Ciliary body};
\draw[poppink, thin] (5.0,-1.2) -- (2.95,-1.6);
\node[popmuted, anchor=west] at (5.1,-1.9) {Suspensory ligament};
\draw[popmuted, thin] (5.0,-2.0) -- (2.5,-1.45);
\node[popblue, anchor=west] at (5.1,-2.7) {Aqueous humour};
\draw[popblue, thin] (5.0,-2.8) -- (3.1,-0.95);
\node[popblue, anchor=north] at (0,-3.25) {Vitreous humour};
\draw[popblue, thin] (0,-3.15) -- (0,-1.6);
% Light rays entering from the right
\draw[->, poporange, thick] (7.2,0.8) -- (4.7,0.45);
\draw[->, poporange, thick] (7.2,0) -- (4.8,0);
\draw[->, poporange, thick] (7.2,-0.8) -- (4.7,-0.45);
\node[poporange, anchor=south] at (7.2,0.9) {Light};
\end{tikzpicture}
\legende{Horizontal section through the mammalian eye. Note the three concentric layers (sclera, choroid, retina), the transparent cornea continuous with the sclera, the biconvex lens held by suspensory ligaments from the ciliary body, and the fovea lying on the visual axis near the blind spot where the optic nerve leaves.}
\end{popfigure}

\begin{definition}[Fovea (Yellow Spot)]
The \cle{fovea} is a small region of the retina lying on the visual axis, close to the blind spot. It contains a very high density of \emph{cones only} (no rods), and each cone is connected to its own bipolar neurone and its own ganglion cell (a 1:1:1 arrangement). This ``private line'' to the brain makes the fovea the region of \emph{maximum visual acuity} (ability to distinguish fine detail) and the centre of colour vision.
\end{definition}

\begin{methode}[Drawing and Labelling the Eye in the Examination]
\textbf{Step 1:} Orient yourself --- the \cle{optic nerve} marks the back (posterior) of the eye; the \cle{cornea} bulges at the front (anterior).\par
\textbf{Step 2:} Draw a large, clear outline (at least half a page) with a slight anterior bulge for the cornea.\par
\textbf{Step 3:} Draw the three layers at the back and sides: sclera (outermost), choroid, retina (innermost).\par
\textbf{Step 4:} Add the biconvex \cle{lens}, held by thin \cle{suspensory ligaments} attached to the thickened \cle{ciliary body}.\par
\textbf{Step 5:} Draw the \cle{iris} in front of the lens, leaving a central gap for the \cle{pupil}.\par
\textbf{Step 6:} Mark the \cle{blind spot} where the optic nerve leaves, and the \cle{fovea} on the visual axis.\par
\textbf{Step 7:} Label with a ruler: label lines must be straight, must not cross, and must touch the structure exactly. Give your drawing a title.
\end{methode}

\begin{attention}[Protective Accessory Structures --- Easy Marks]
Do not forget the \emph{accessory structures} of the eye; they are part of the syllabus:
\begin{itemize}
\item \textbf{Eyelids:} close reflexively (blink reflex) to protect the eye and spread tears over the cornea.
\item \textbf{Eyelashes:} trap dust and trigger the blink reflex when touched.
\item \textbf{Lacrimal (tear) glands:} secrete tears which moisten and nourish the cornea and contain the antibacterial enzyme lysozyme.
\item \textbf{Extrinsic eye muscles:} move the eyeball within the orbit so that both eyes can be directed at the same object.
\end{itemize}
\end{attention}

\courssub{Refraction and Image Formation}

Light from an object is bent (refracted) as it passes through the eye, mainly at two places:
\begin{enumerate}
\item the \cle{cornea}, where light passes from air into the denser corneal tissue --- this is where \emph{most} of the refraction occurs, because the difference in density between air and cornea is large and the corneal curvature is fixed;
\item the \cle{lens}, whose curvature is \emph{adjustable}, allowing fine focusing of objects at different distances (this is \cle{accommodation}, studied in Section 3).
\end{enumerate}
The aqueous and vitreous humours transmit the light and maintain the shape of the eyeball but contribute little to refraction.

\begin{propriete}[Nature of the Retinal Image]
The cornea and lens together act as a \emph{converging lens system}. For an object viewed by a normal (emmetropic) eye, the image formed on the retina is:
\begin{itemize}
\item \textbf{Real} --- formed by the actual convergence of light rays;
\item \textbf{Inverted} --- upside-down and reversed left-to-right;
\item \textbf{Diminished} --- smaller than the object.
\end{itemize}
The brain interprets the inverted image so that we perceive the world the right way up.
\end{propriete}

\begin{attention}[The Inverted Image Trap]
\textbf{Common mistake:} stating that the image on the retina is ``upright'' or ``virtual''. \textbf{Always remember:} the retinal image is \textbf{real, inverted and diminished}. It is the brain that re-interprets it. If a GCE question asks ``State the nature of the image formed on the retina'', the expected answer is exactly: \emph{real, inverted and diminished}.
\end{attention}

\courssub{The Retina: Rods and Cones}

The retina contains two types of photoreceptor cell, \cle{rods} and \cle{cones}, which convert light energy into nerve impulses. Their different properties explain why our vision in dim light differs from our vision in bright light.

\begin{center}
\begin{tabularx}{\linewidth}{|l|X|X|}\hline
\rowcolor{popblueL}
\textbf{Feature} & \textbf{Rods} & \textbf{Cones} \\ \hline
\cle{Photopigment} & Rhodopsin (visual purple) & Iodopsin (three forms, sensitive to blue, green or red light) \\ \hline
\cle{Sensitivity to light} & Very high --- function in dim light & Low --- function only in bright light \\ \hline
\cle{Colour vision} & None --- give black-and-white vision & Yes --- responsible for colour vision \\ \hline
\cle{Visual acuity} & Low --- many rods connect to one sensory neurone, so impulses are summed and detail is lost & High --- each cone (at the fovea) has its own connection to the brain \\ \hline
\cle{Distribution} & Absent from the fovea; concentrated in the peripheral retina & Concentrated at the \cle{fovea}; few in the periphery \\ \hline
\end{tabularx}
\end{center}

\begin{definition}[Blind Spot (Optic Disc)]
The \cle{blind spot} is the region of the retina where the axons of the ganglion cells leave the eye as the optic nerve. It contains \emph{no rods and no cones}, so light falling on it produces no sensation. We are normally unaware of it because the brain fills in the missing information and the two eyes cover each other's blind spot.
\end{definition}

\begin{savaistu}[Demonstrating Your Own Blind Spot]
Close your left eye. Stare at the \textbf{cross} below with your right eye only, and slowly bring the page closer to your face from about 30 cm. At a certain distance the \textbf{dot} disappears: its image has fallen on the blind spot of your right eye.
\begin{center}
\begin{tikzpicture}
\node[popdark, font=\Large] at (-2,0) {$\times$};
\node[popdark, font=\Large] at (2,0) {$\bullet$};
\end{tikzpicture}
\end{center}
\end{savaistu}

\courssub{Binocular Vision}

Humans have \cle{binocular vision}: the two eyes face forward and view the same object from slightly different positions. The brain fuses the two slightly different images into a single three-dimensional perception (\cle{stereoscopic vision}), which allows accurate judgement of distance and depth. This is a frequent GCE question: the \emph{advantages of binocular vision} are \textbf{accurate depth/distance judgement} and \textbf{overlapping fields of view}, so the blind spot of one eye is covered by the other. Predators (lion, human) have forward-facing eyes for good depth judgement; prey animals (antelope, grasscutter) have eyes on the sides of the head giving a wide field of view to detect danger, but poorer depth perception.

\courssub{Common Eye Defects and Their Correction}

Refractive errors occur when the eyeball length and the refractive power of the cornea and lens do not match, so light is not focused exactly on the retina. They are corrected with spectacle lenses that shift the point of focus back onto the retina.

\begin{propriete}[Myopia (Short-sightedness)]
\textbf{Cause:} the eyeball is too long, or the cornea/lens refracts too strongly. Parallel rays from a distant object are focused \emph{in front of} the retina, so distant objects appear blurred while near objects are seen clearly.\par
\textbf{Correction:} a \cle{diverging (concave) lens} is placed in front of the eye. It spreads the incoming rays slightly so that the eye's own optics bring them to a focus exactly \emph{on} the retina.
\end{propriete}

\begin{propriete}[Hypermetropia (Long-sightedness)]
\textbf{Cause:} the eyeball is too short, or the cornea/lens refracts too weakly. Rays from a near object would be focused \emph{behind} the retina, so near objects appear blurred while distant objects may be seen clearly.\par
\textbf{Correction:} a \cle{converging (convex) lens} is placed in front of the eye. It begins to converge the rays before they enter the eye, so that they are focused exactly \emph{on} the retina.
\end{propriete}

\begin{propriete}[Astigmatism]
\textbf{Cause:} the cornea (or lens) is unevenly curved --- more curved in one plane than in the perpendicular plane --- so rays in different planes are focused at different points and the image is distorted.\par
\textbf{Correction:} a \cle{cylindrical lens} whose curvature compensates for the uneven curvature of the cornea.
\end{propriete}

\begin{popfigure}
\begin{tikzpicture}[scale=1.15]
% ---- Normal eye ----
\begin{scope}[xshift=-6.5cm]
\draw[popdark, thick] (-1.5,-1) -- (-1.5,1);
\draw[popblue, thick] (-2.6,-0.6) .. controls (-2.2,0) .. (-2.6,0.6);
\draw[popblue, thick] (-2.6,-0.6) .. controls (-3.0,0) .. (-2.6,0.6);
\draw[->, poporange, thick] (-4.6,0.6) -- (-2.8,0.35);
\draw[->, poporange, thick] (-4.6,-0.6) -- (-2.8,-0.35);
\draw[poporange, thick] (-2.8,0.35) -- (-1.5,0);
\draw[poporange, thick] (-2.8,-0.35) -- (-1.5,0);
\node[popdark, align=center] at (-3,-1.5) {Normal eye\\focus on retina};
\end{scope}
% ---- Myopia ----
\begin{scope}[xshift=0cm]
\draw[popdark, thick] (-1.5,-1) -- (-1.5,1);
\draw[popblue, thick] (-2.6,-0.6) .. controls (-2.2,0) .. (-2.6,0.6);
\draw[popblue, thick] (-2.6,-0.6) .. controls (-3.0,0) .. (-2.6,0.6);
\draw[->, poporange, thick] (-4.6,0.6) -- (-2.8,0.35);
\draw[->, poporange, thick] (-4.6,-0.6) -- (-2.8,-0.35);
\draw[poporange, thick] (-2.8,0.35) -- (-1.9,0);
\draw[poporange, thick] (-2.8,-0.35) -- (-1.9,0);
\draw[poporange, dashed, thick] (-1.9,0) -- (-1.5,0.15);
\draw[poporange, dashed, thick] (-1.9,0) -- (-1.5,-0.15);
\fill[popink] (-1.9,0) circle (0.05);
\node[popink, align=center] at (-3,-1.5) {Myopia\\focus in front of retina};
\end{scope}
% ---- Hypermetropia ----
\begin{scope}[xshift=6.5cm]
\draw[popdark, thick] (-1.5,-1) -- (-1.5,1);
\draw[popblue, thick] (-2.6,-0.6) .. controls (-2.2,0) .. (-2.6,0.6);
\draw[popblue, thick] (-2.6,-0.6) .. controls (-3.0,0) .. (-2.6,0.6);
\draw[->, poporange, thick] (-4.6,0.6) -- (-2.8,0.35);
\draw[->, poporange, thick] (-4.6,-0.6) -- (-2.8,-0.35);
\draw[poporange, thick] (-2.8,0.35) -- (-1.5,0.12);
\draw[poporange, thick] (-2.8,-0.35) -- (-1.5,-0.12);
\draw[poporange, dashed, thick] (-1.5,0.12) -- (-1.0,0);
\draw[poporange, dashed, thick] (-1.5,-0.12) -- (-1.0,0);
\fill[popink] (-1.0,0) circle (0.05);
\node[popink, align=center] at (-3,-1.5) {Hypermetropia\\focus behind retina};
\end{scope}
\end{tikzpicture}
\legende{Focusing in the normal eye, in myopia (image forms in front of the retina) and in hypermetropia (image would form behind the retina; dashed lines show the rays continuing to their uncorrected focal point).}
\end{popfigure}

\begin{popfigure}
\begin{tikzpicture}[scale=1.15]
% ---- Myopia correction ----
\begin{scope}[xshift=-3.5cm]
\draw[popdark, thick] (-1.5,-1) -- (-1.5,1);
\draw[popblue, thick] (-2.6,-0.6) .. controls (-2.2,0) .. (-2.6,0.6);
\draw[popblue, thick] (-2.6,-0.6) .. controls (-3.0,0) .. (-2.6,0.6);
% concave lens
\draw[popgreen, line width=1.4pt] (-3.6,-0.8) .. controls (-3.75,0) .. (-3.6,0.8);
\draw[popgreen, line width=1.4pt] (-3.6,-0.8) .. controls (-3.45,0) .. (-3.6,0.8);
\draw[->, poporange, thick] (-5.2,0.6) -- (-3.7,0.5);
\draw[->, poporange, thick] (-5.2,-0.6) -- (-3.7,-0.5);
\draw[poporange, thick] (-3.5,0.5) -- (-2.8,0.35);
\draw[poporange, thick] (-3.5,-0.5) -- (-2.8,-0.35);
\draw[poporange, thick] (-2.8,0.35) -- (-1.5,0);
\draw[poporange, thick] (-2.8,-0.35) -- (-1.5,0);
\node[popgreen, align=center] at (-3.6,1.2) {Concave lens};
\node[popdark, align=center] at (-3.3,-1.5) {Myopia corrected};
\end{scope}
% ---- Hypermetropia correction ----
\begin{scope}[xshift=3.5cm]
\draw[popdark, thick] (-1.5,-1) -- (-1.5,1);
\draw[popblue, thick] (-2.6,-0.6) .. controls (-2.2,0) .. (-2.6,0.6);
\draw[popblue, thick] (-2.6,-0.6) .. controls (-3.0,0) .. (-2.6,0.6);
% convex lens
\draw[popgreen, line width=1.4pt] (-3.6,-0.8) .. controls (-3.35,0) .. (-3.6,0.8);
\draw[popgreen, line width=1.4pt] (-3.6,-0.8) .. controls (-3.85,0) .. (-3.6,0.8);
\draw[->, poporange, thick] (-5.2,0.7) -- (-3.7,0.55);
\draw[->, poporange, thick] (-5.2,-0.7) -- (-3.7,-0.55);
\draw[poporange, thick] (-3.5,0.55) -- (-2.8,0.35);
\draw[poporange, thick] (-3.5,-0.55) -- (-2.8,-0.35);
\draw[poporange, thick] (-2.8,0.35) -- (-1.5,0);
\draw[poporange, thick] (-2.8,-0.35) -- (-1.5,0);
\node[popgreen, align=center] at (-3.6,1.2) {Convex lens};
\node[popdark, align=center] at (-3.3,-1.5) {Hypermetropia corrected};
\end{scope}
\end{tikzpicture}
\legende{Correction of refractive errors: a diverging (concave) lens corrects myopia; a converging (convex) lens corrects hypermetropia. In each case the focal point is shifted back onto the retina.}
\end{popfigure}

\begin{exoresolu}[Identifying and Correcting a Defect]
A student can read a book held at 25 cm clearly, but the writing on the classroom board is blurred. (a) Name the defect. (b) Explain its cause. (c) State the type of lens needed for correction and justify your answer.
\tcblower
\textbf{(a)} The defect is \textbf{myopia} (short-sightedness): near objects are seen clearly but distant objects are blurred.\par
\textbf{(b)} The eyeball is too long (or the cornea and lens refract too strongly), so parallel rays from the distant board are brought to a focus \emph{in front of} the retina. The retina therefore receives a blurred image.\par
\textbf{(c)} A \textbf{diverging (concave) lens} is required. It diverges the parallel rays slightly before they enter the eye, so that the eye's own converging system now focuses them exactly \emph{on} the retina, restoring clear distance vision.
\end{exoresolu}

\begin{exoresolu}[Identifying and Correcting Long Sight]
An elderly farmer in Buea must hold his newspaper at arm's length (about 100 cm) to read it, although his distance vision is good. (a) Name the defect. (b) Explain its cause. (c) State the correcting lens.
\tcblower
\textbf{(a)} \textbf{Hypermetropia} (long-sightedness) --- in an older person this is specifically \emph{presbyopia}, long sight due to loss of lens elasticity with age.\par
\textbf{(b)} The eyeball is too short, or (in presbyopia) the lens can no longer become sufficiently rounded because it has lost elasticity and the ciliary muscles are weaker. Rays from near objects are therefore focused \emph{behind} the retina.\par
\textbf{(c)} A \textbf{converging (convex) lens}, which pre-converges the rays so that they focus on the retina.
\end{exoresolu}

\begin{aretenir}[Key Points for GCE Revision]
\begin{itemize}
\item The eye wall has three layers: \textbf{sclera} (protection and shape), \textbf{choroid} (nutrition and absorption of stray light), \textbf{retina} (contains the photoreceptors).
\item The \textbf{cornea} provides most of the refraction; the \textbf{lens} provides the adjustable refraction used in accommodation.
\item The retinal image is \textbf{real, inverted and diminished}.
\item \textbf{Rods:} very sensitive, dim-light (black-and-white) vision, low acuity, peripheral retina. \textbf{Cones:} bright-light and colour vision, high acuity, concentrated at the fovea.
\item \textbf{Fovea} = region of sharpest vision (cones only, one-to-one connections). \textbf{Blind spot} = exit point of the optic nerve (no photoreceptors).
\item \textbf{Myopia:} image in front of retina $\to$ concave (diverging) lens. \textbf{Hypermetropia:} image behind retina $\to$ convex (converging) lens. \textbf{Astigmatism:} uneven corneal curvature $\to$ cylindrical lens.
\item \textbf{Binocular vision} gives stereoscopic (3-D) vision and accurate judgement of distance.
\end{itemize}
\end{aretenir}

\begin{exercice}[GCE-Style Practice]
\begin{enumerate}
\item Draw a large, labelled diagram of a section through the mammalian eye. On your diagram, show the path of two light rays from a distant object to the retina. [8 marks]
\item State \emph{two} differences between rods and cones, and explain how each difference is related to the function of the cell. [6 marks]
\item A student cannot see objects clearly unless they are closer than 30 cm. Name the defect, explain its cause, and state the type of lens needed for correction, justifying your answer. [5 marks]
\item Explain why the image formed on the retina is inverted, and how we nevertheless perceive objects the right way up. [3 marks]
\item Describe the distribution of rods and cones in the retina and relate this distribution to visual acuity and colour vision. [4 marks]
\end{enumerate}
\end{exercice}

\begin{corrige}[Exercise 1]
\textbf{Diagram:} large outline with anterior corneal bulge; three layers at the back (sclera outermost, choroid, retina innermost); cornea continuous with the sclera; biconvex lens suspended by suspensory ligaments from the ciliary body; iris with central pupil; optic nerve leaving at the blind spot; fovea marked on the visual axis. Two parallel rays from a distant object enter through the cornea and pupil, are refracted by the cornea and lens, cross over, and converge to a point on the retina (inverted image). Labels: sclera, choroid, retina, cornea, iris, pupil, lens, ciliary body, suspensory ligaments, aqueous humour, vitreous humour, optic nerve, blind spot, fovea. Marks awarded for accuracy of proportions, correct relative positions, complete labels and correct light path.
\end{corrige}

\begin{corrige}[Exercise 2]
\begin{itemize}
\item \textbf{Photopigment:} rods contain \emph{rhodopsin}, which is very sensitive to light and is bleached in bright light $\to$ rods function in dim light. Cones contain \emph{iodopsin} (in three forms) which is less sensitive $\to$ cones function in bright light and the three forms make colour vision possible.
\item \textbf{Connections:} many rods connect to a single sensory neurone $\to$ their impulses are summed, giving high sensitivity but poor detail (low acuity). Each cone at the fovea connects to its own neurone $\to$ no summation, so fine detail is resolved (high acuity).
\end{itemize}
(Also acceptable: distribution --- rods peripheral, cones at fovea --- related to where acuity and colour vision are best.)
\end{corrige}

\begin{corrige}[Exercise 3]
\textbf{Defect:} myopia (short-sightedness). \textbf{Cause:} the eyeball is too long or the cornea/lens refracts too strongly, so light from distant objects is focused in front of the retina; only near objects (closer than 30 cm) focus on the retina. \textbf{Correction:} a diverging (concave) lens, which spreads the incoming rays so that the eye focuses them on the retina.
\end{corrige}

\begin{corrige}[Exercise 4]
The cornea and lens form a converging lens system. Rays from the top of an object are refracted to the bottom of the retina, and rays from the bottom of the object to the top of the retina; the rays cross over, so the image is inverted (and diminished). The brain receives the impulses from the retina and, through experience and coordination with the other senses, interprets the image so that the object is perceived the right way up.
\end{corrige}

\begin{corrige}[Exercise 5]
\textbf{Cones} are concentrated at the \textbf{fovea}, are few in the peripheral retina and absent from the blind spot. \textbf{Rods} are absent from the fovea and are concentrated in the \textbf{peripheral retina}. \textbf{Visual acuity} is therefore greatest at the fovea (high cone density and one-to-one neurone connections) and decreases towards the periphery (rod dominance and many-to-one connections). \textbf{Colour vision} occurs only where cones are present (fovea and central retina); peripheral vision, mediated by rods, is in black and white.
\end{corrige}

\coursec{Control of Light Intensity, Adaptation and Accommodation}

Having examined the gross structure of the eye and the physics of image formation on the retina, we now turn to the dynamic physiological mechanisms that ensure the image is not only formed, but formed \emph{clearly} and under \emph{optimal illumination}. The eye is not a static camera; it actively adjusts its aperture (the pupil), its sensitivity (retinal adaptation), and its focal length (accommodation) to cope with the enormous range of light intensities and object distances encountered in daily life — from the intense equatorial noon sun over the Adamawa plateau to the dim light of a forest understory on Mount Cameroon.

\courssub{The Iris Reflex: Controlling Light Entry}

The amount of light entering the eye is regulated by the diameter of the pupil, which is controlled by two sets of smooth muscles in the iris: the \cle{circular muscles} (sphincter pupillae) and the \cle{radial muscles} (dilator pupillae). These muscles are antagonistic effectors of the autonomic nervous system.

\begin{popfigure}
\centering
\begin{tikzpicture}[scale=1.2, every node/.style={font=\small}]
    % Bright light - left
    \begin{scope}[xshift=-4cm]
        % Sclera outline
        \draw[popdark, thick] (0,0) ellipse (1.5 and 1);
        % Iris - constricted pupil
        \fill[popblue!30] (0,0) ellipse (1.2 and 0.8);
        % Pupil - small
        \fill[black] (0,0) ellipse (0.25 and 0.25);
        % Circular muscles - contracted (thick ring)
        \draw[popink, very thick] (0,0) ellipse (0.9 and 0.6);
        % Radial muscles - relaxed (thin lines)
        \foreach \angle in {0,30,...,330} {
            \draw[poporange, thin] ({0.9*cos(\angle)},{0.6*sin(\angle)}) -- ({1.2*cos(\angle)},{0.8*sin(\angle)});
        }
        % Labels
        \node[popink] at (0, -1.4) {\textbf{Bright Light}};
        \node[popink, align=center] at (-1.8, 0.5) {Circular muscles\\ \textbf{contract}};
        \node[poporange, align=center] at (1.8, 0.5) {Radial muscles\\ relax};
        \node[black, align=center] at (0, -1.7){Pupil \textbf{constricts}\\(miosis)};
    \end{scope}

    % Dim light - right
    \begin{scope}[xshift=4cm]
        % Sclera outline
        \draw[popdark, thick] (0,0) ellipse (1.5 and 1);
        % Iris - dilated pupil
        \fill[popblue!30] (0,0) ellipse (1.2 and 0.8);
        % Pupil - large
        \fill[black] (0,0) ellipse (0.85 and 0.85);
        % Circular muscles - relaxed (thin ring)
        \draw[popink, thin] (0,0) ellipse (0.9 and 0.6);
        % Radial muscles - contracted (thick lines)
        \foreach \angle in {0,30,...,330} {
            \draw[poporange, very thick] ({0.9*cos(\angle)},{0.6*sin(\angle)}) -- ({1.2*cos(\angle)},{0.8*sin(\angle)});
        }
        % Labels
        \node[poporange] at (0, -1.4) {\textbf{Dim Light}};
        \node[popink, align=center] at (-1.8, 0.5) {Circular muscles\\ relax};
        \node[poporange, align=center] at (1.8, 0.5) {Radial muscles\\ \textbf{contract}};
        \node[black, align=center] at (0, -1.7){Pupil \textbf{dilates}\\(mydriasis)};
    \end{scope}
\end{tikzpicture}
\legende{Antagonistic action of iris muscles. In bright light (left), parasympathetic stimulation contracts circular muscles $\to$ pupil constricts. In dim light (right), sympathetic stimulation contracts radial muscles $\to$ pupil dilates.}
\end{popfigure}

\begin{attention}[The Bright Light Trap]
\textbf{Never invert the muscle actions!} In \textbf{bright light}, the \textbf{circular muscles contract} (parasympathetic) $\to$ pupil \textbf{constricts} (miosis). In \textbf{dim light}, the \textbf{radial muscles contract} (sympathetic) $\to$ pupil \textbf{dilates} (mydriasis). A common GCE error is stating that radial muscles contract in bright light. Remember: "Bright $\to$ Constrict $\to$ Circular".
\end{attention}

\paragraph{Nervous Pathway of the Pupil Reflex (Autonomic Control)}
The pupil reflex is a \cle{reflex arc} involving the autonomic nervous system. It is \emph{consensual}: shining light in one eye causes constriction in both eyes.

\begin{enumerate}
    \item \textbf{Receptor:} Photoreceptors (rods/cones) in the retina detect light intensity.
    \item \textbf{Sensory Neurone:} Optic nerve (CN II) carries impulses to the \textbf{pretectal nucleus} in the midbrain (not the visual cortex — this is a reflex, not conscious perception).
    \item \textbf{Interneurones:} Connect pretectal nucleus to the \textbf{Edinger-Westphal nuclei} (parasympathetic component of oculomotor nerve nuclei) on \emph{both} sides (decussation via posterior commissure).
    \item \textbf{Motor Neurone (Parasympathetic):} Preganglionic fibres run in the \textbf{oculomotor nerve (CN III)} $\to$ synapse in the \textbf{ciliary ganglion} $\to$ postganglionic \textbf{short ciliary nerves} $\to$ \textbf{circular muscles} contract $\to$ pupil constricts.
    \item \textbf{Sympathetic Pathway (Dim Light):} Hypothalamus $\to$ spinal cord (T1) $\to$ superior cervical ganglion $\to$ long ciliary nerves $\to$ \textbf{radial muscles} contract $\to$ pupil dilates.
\end{enumerate}

\begin{definition}[Pupil Reflex (Light Reflex)]
A protective autonomic reflex that adjusts pupil diameter in response to light intensity falling on the retina, preventing photobleaching of photoreceptors and improving depth of field. It is mediated by the parasympathetic division (constriction) and sympathetic division (dilation) of the autonomic nervous system.
\end{definition}

\courssub{Light and Dark Adaptation: The Chemistry of Vision}

While the iris reflex adjusts light entry rapidly (seconds), the retina itself adjusts its \emph{sensitivity} over a much wider range (orders of magnitude) through the biochemistry of photopigments. This is \cle{retinal adaptation}.

\begin{propriete}[Rhodopsin Cycle: Bleaching and Regeneration]
\textbf{Rhodopsin (Visual Purple)} is the photopigment in rod cells. It consists of the protein \cle{opsin} bound to the aldehyde of vitamin A (\cle{retinal}, specifically \emph{11-cis-retinal}).
\begin{itemize}
    \item \textbf{In Bright Light (Bleaching):} Photons isomerise 11-cis-retinal to \emph{all-trans-retinal}. This conformational change activates opsin $\to$ triggers transduction cascade $\to$ rhodopsin splits into opsin + all-trans-retinal (bleached, colourless). Sensitivity \textbf{decreases}.
    \item \textbf{In Darkness (Regeneration):} All-trans-retinal is isomerised back to 11-cis-retinal (requiring ATP) in the retinal pigment epithelium (RPE). It recombines with opsin to reform rhodopsin (purple). Sensitivity \textbf{increases}.
\end{itemize}
\textbf{Vitamin A deficiency} $\to$ insufficient retinal $\to$ impaired rhodopsin regeneration $\to$ \cle{night blindness} (nyctalopia).
\textbf{Examinable:} Descriptive knowledge of the rhodopsin cycle and the role of vitamin A is required; no formal proof of the transduction cascade is expected.
\end{propriete}

\begin{savaistu}[Vitamin A and Public Health in Cameroon]
Night blindness due to vitamin A deficiency remains a public health concern in parts of Cameroon, particularly in the Far North and among children with limited dietary diversity (low intake of palm oil, mangoes, liver, eggs). The National Vitamin A Supplementation Programme targets children 6--59 months and postpartum women. Understanding the rhodopsin cycle explains \emph{why} supplementation restores night vision: it provides the retinal precursor for rhodopsin regeneration.
\end{savaistu}

\begin{definition}[Light Adaptation \& Dark Adaptation]
\begin{itemize}
    \item \textbf{Light Adaptation:} Rapid process (seconds to minutes) where sensitivity decreases. Cones take over (photopic vision); rhodopsin is bleached faster than it regenerates. Pupil constricts.
    \item \textbf{Dark Adaptation:} Slow process (20--30 minutes for full rod sensitivity). Rods take over (scotopic vision); rhodopsin regenerates in the dark. Pupil dilates. The \emph{Purkinje shift} occurs: peak spectral sensitivity shifts from $\approx 555\,\mathrm{nm}$ (cones, green-yellow) to $\approx 505\,\mathrm{nm}$ (rods, blue-green).
\end{itemize}
\end{definition}

\begin{exemplebox}[Walking from Sunlight into a Dark Room]
Imagine leaving a brightly lit classroom in Buea at noon and entering a dark storeroom.
\begin{enumerate}
    \item \textbf{0--5 sec:} Pupils dilate rapidly (sympathetic iris reflex) — immediate but limited gain ($\times 16$ area).
    \item \textbf{5 sec -- 5 min:} Cone sensitivity recovers (fast). You start seeing large shapes.
    \item \textbf{5--30 min:} Rod sensitivity recovers (slow). Rhodopsin regenerates. Sensitivity increases up to $\times 10^5$. You see fine detail in dim light.
    \item \textbf{Re-entering sunlight:} Intense light bleaches rhodopsin almost instantly — glare, temporary blindness until cones dominate and pupil constricts.
\end{enumerate}
\end{exemplebox}

\courssub{Accommodation: Focusing for Near and Far}

The iris controls \emph{quantity} of light; accommodation controls \emph{quality} (sharpness) by changing the focal length of the eye. The cornea provides $\approx 2/3$ of the total refractive power (fixed), while the \cle{crystalline lens} provides the remaining $\approx 1/3$ (variable). Accommodation is the process by which the lens changes shape to focus objects at different distances onto the retina.

\begin{popfigure}
\centering
\begin{tikzpicture}[scale=1.1, every node/.style={font=\small}]
    % Distant Vision - Left
    \begin{scope}[xshift=-5cm]
        % Sclera
        \draw[popdark, thick] (-2, -1.5) -- (-2, 1.5) .. controls (-1, 2) and (1, 2) .. (2, 1.5) -- (2, -1.5) .. controls (1, -2) and (-1, -2) .. (-2, -1.5);
        % Lens - flat
        \draw[popblue, very thick] (-0.8, -0.8) .. controls (-0.5, -0.9) and (0.5, -0.9) .. (0.8, -0.8);
        \draw[popblue, very thick] (-0.8, 0.8) .. controls (-0.5, 0.9) and (0.5, 0.9) .. (0.8, 0.8);
        % Suspensory ligaments - TAUT
        \draw[poporange, thick] (-2, 0) -- (-0.8, 0.8);
        \draw[poporange, thick] (-2, 0) -- (-0.8, -0.8);
        \draw[poporange, thick] (2, 0) -- (0.8, 0.8);
        \draw[poporange, thick] (2, 0) -- (0.8, -0.8);
        % Ciliary muscle - RELAXED (wide ring)
        \draw[popgreen, thick] (-1.8, 0.9) .. controls (-1.5, 1.1) and (1.5, 1.1) .. (1.8, 0.9);
        \draw[popgreen, thick] (-1.8, -0.9) .. controls (-1.5, -1.1) and (1.5, -1.1) .. (1.8, -0.9);
        % Light rays - parallel
        \draw[->, popgray] (-3, 0.6) -- (-2, 0.6);
        \draw[->, popgray] (-3, -0.6) -- (-2, -0.6);
        \draw[popgray] (0.8, 0.6) .. controls (1.2, 0.4) and (1.8, 0) .. (2.5, 0);
        \draw[popgray] (0.8, -0.6) .. controls (1.2, -0.4) and (1.8, 0) .. (2.5, 0);
        % Labels
        \node at (0, -1.9) {\textbf{Distant Vision ($>6\,\mathrm{m}$)}};
        \node[poporange] at (-1.5, 1.3) {Ligaments \textbf{taut}};
        \node[popgreen] at (0, 1.3) {Ciliary muscle \textbf{relaxed}};
        \node[popblue] at (0, -1.3) {Lens \textbf{flat} (low power)};
    \end{scope}

    % Near Vision - Right
    \begin{scope}[xshift=5cm]
        % Sclera
        \draw[popdark, thick] (-2, -1.5) -- (-2, 1.5) .. controls (-1, 2) and (1, 2) .. (2, 1.5) -- (2, -1.5) .. controls (1, -2) and (-1, -2) .. (-2, -1.5);
        % Lens - fat/round
        \draw[popblue, very thick] (-0.6, -1) .. controls (-0.3, -1.2) and (0.3, -1.2) .. (0.6, -1);
        \draw[popblue, very thick] (-0.6, 1) .. controls (-0.3, 1.2) and (0.3, 1.2) .. (0.6, 1);
        % Suspensory ligaments - SLACK
        \draw[poporange, dashed, thick] (-2, 0) -- (-0.6, 1);
        \draw[poporange, dashed, thick] (-2, 0) -- (-0.6, -1);
        \draw[poporange, dashed, thick] (2, 0) -- (0.6, 1);
        \draw[poporange, dashed, thick] (2, 0) -- (0.6, -1);
        % Ciliary muscle - CONTRACTED (narrow ring)
        \draw[popgreen, thick] (-1.2, 0.6) .. controls (-0.8, 0.9) and (0.8, 0.9) .. (1.2, 0.6);
        \draw[popgreen, thick] (-1.2, -0.6) .. controls (-0.8, -0.9) and (0.8, -0.9) .. (1.2, -0.6);
        % Light rays - diverging
        \draw[->, popgray] (-3, 0.8) -- (-2, 0.5);
        \draw[->, popgray] (-3, -0.8) -- (-2, -0.5);
        \draw[popgray] (0.6, 0.5) .. controls (1.0, 0.2) and (1.6, 0) .. (2.5, 0);
        \draw[popgray] (0.6, -0.5) .. controls (1.0, -0.2) and (1.6, 0) .. (2.5, 0);
        % Labels
        \node at (0, -1.9) {\textbf{Near Vision ($<6\,\mathrm{m}$)}};
        \node[poporange] at (-1.5, 1.3) {Ligaments \textbf{slack}};
        \node[popgreen] at (0, 1.3) {Ciliary muscle \textbf{contracted}};
        \node[popblue] at (0, -1.3) {Lens \textbf{round} (high power)};
    \end{scope}
\end{tikzpicture}
\legende{Mechanism of accommodation. \textbf{Left (Distant):} Ciliary muscle relaxed $\to$ ligaments taut $\to$ lens flat. \textbf{Right (Near):} Ciliary muscle contracts $\to$ ligaments slack $\to$ lens elastic recoil makes it round.}
\end{popfigure}

\begin{definition}[Accommodation]
The process by which the eye changes the focal length of the crystalline lens to form a sharp image on the retina for objects at different distances. It is achieved by the contraction/relaxation of the \textbf{ciliary muscle} altering tension on the \textbf{suspensory ligaments} (zonular fibres), allowing the naturally elastic lens to change shape.
\end{definition}

\begin{methode}[Step-by-Step Reasoning: Focusing a Near Object]
To answer GCE questions on accommodation, follow this causal chain:
\begin{center}
\begin{tabularx}{\linewidth}{|l|X|}\hline
\textbf{Step} & \textbf{Physiological Event} \\ \hline
1 & Object is near $\to$ light rays diverge strongly entering the eye. \\ \hline
2 & Blurred image forms \textbf{behind} retina (hypermetropic blur detected). \\ \hline
3 & Visual cortex $\to$ parasympathetic signal via CN III (oculomotor) to \textbf{ciliary muscle}. \\ \hline
4 & \textbf{Ciliary muscle contracts} (circular fibres shorten). \\ \hline
5 & Ciliary body moves \textbf{forward and inward}. \\ \hline
6 & Tension on \textbf{suspensory ligaments} is \textbf{released} (slack). \\ \hline
7 & \textbf{Lens elasticity} causes it to become \textbf{more convex} (rounder). \\ \hline
8 & Refractive power of lens \textbf{increases}. \\ \hline
9 & Diverging rays bent more sharply $\to$ image focused \textbf{on retina}. \\ \hline
\end{tabularx}
\end{center}
\textbf{Reverse every step for distant vision.}
\end{methode}

\begin{attention}[Presbyopia -- The Ageing Lens]
With age, the lens loses elasticity (denaturation of crystallin proteins) and the ciliary muscle weakens. The near point recedes (from $\approx 10\,\mathrm{cm}$ at age 10 to $\approx 100\,\mathrm{cm}$ at age 60). This is \cle{presbyopia} ("old eye"). It is \emph{not} a refractive error like myopia/hypermetropia; it is a loss of \emph{accommodative amplitude}. Corrected by convex reading glasses (add power).
\end{attention}

\paragraph{The Near Response Triad (Convergence Reflex)}
Accommodation for near objects is never isolated; it is part of a synergistic \cle{near response triad}:
\begin{enumerate}
    \item \textbf{Accommodation:} Lens rounds up (parasympathetic, CN III).
    \item \textbf{Convergence:} Medial recti muscles contract $\to$ eyes rotate medially to fixate on the near object (somatic motor, CN III).
    \item \textbf{Pupillary Constriction (Miosis):} Circular iris muscles contract $\to$ pupil constricts $\to$ increases depth of field and reduces spherical aberration (parasympathetic, CN III).
\end{enumerate}
This is why your pupils constrict when you look at your phone screen close to your face, even if the room lighting hasn't changed.

\begin{exoresolu}[Analysing the Near Response]
\textbf{Statement:} A student reads a textbook held at $25\,\mathrm{cm}$. Describe the changes in the eye and the nervous control involved.
\tcblower
\textbf{Detailed Solution:}
\begin{enumerate}
    \item \textbf{Stimulus:} Diverging light rays from near object; blur detected on retina (image would form behind retina).
    \item \textbf{Afferent Pathway:} Optic nerve (CN II) $\to$ visual cortex $\to$ association areas.
    \item \textbf{Efferent Pathway (Parasympathetic):} Edinger-Westphal nucleus $\to$ Oculomotor nerve (CN III) $\to$ Ciliary ganglion.
    \item \textbf{Effectors \& Changes:}
        \begin{itemize}
            \item \textbf{Ciliary muscle:} Contracts $\to$ moves forward/inward.
            \item \textbf{Suspensory ligaments:} Slacken.
            \item \textbf{Lens:} Becomes more convex (increased curvature) $\to$ increased refractive power $\to$ image focused on retina.
            \item \textbf{Pupil:} Circular muscles contract $\to$ pupil constricts (miosis) $\to$ increased depth of field, reduced aberrations.
            \item \textbf{Medial rectus muscles:} Contract (somatic CN III) $\to$ eyes converge (medial rotation).
        \end{itemize}
    \item \textbf{Result:} Clear, single, binocular image of the text on both retinas.
\end{enumerate}
\end{exoresolu}

\courssub{Summary and Key Points}

\begin{aretenir}[Control of Light Intensity, Adaptation \& Accommodation]
\begin{itemize}
    \item \textbf{Iris Reflex:} Bright light $\to$ Circular muscles contract (Parasympathetic, CN III) $\to$ Pupil constricts (Miosis). Dim light $\to$ Radial muscles contract (Sympathetic) $\to$ Pupil dilates (Mydriasis). Consensual.
    \item \textbf{Light Adaptation:} Fast (secs/mins). Cones function. Rhodopsin bleached $>$ regenerated. Low sensitivity, high acuity, colour vision.
    \item \textbf{Dark Adaptation:} Slow (20--30 mins). Rods function. Rhodopsin regenerates (needs Vitamin A/Retinal). High sensitivity, low acuity, no colour (Purkinje shift).
    \item \textbf{Accommodation:} Near $\to$ Ciliary muscle contracts $\to$ Ligaments slack $\to$ Lens rounds up (elastic recoil) $\to$ Power increases. Distant $\to$ Ciliary muscle relaxes $\to$ Ligaments taut $\to$ Lens flattens $\to$ Power decreases.
    \item \textbf{Near Response Triad:} Accommodation + Convergence + Pupillary Constriction (all Parasympathetic CN III).
    \item \textbf{Presbyopia:} Age-related loss of lens elasticity $\to$ reduced accommodation $\to$ corrected by convex lenses.
\end{itemize}
\end{aretenir}

\begin{exercice}[GCE Style Practice]
\begin{enumerate}
    \item Explain why a person with Vitamin A deficiency experiences night blindness but can see normally in bright daylight.
    \item Describe the sequence of events from the moment light from a near object hits the retina to the point where a sharp image is formed. Include the role of the ciliary muscle, suspensory ligaments, and lens elasticity.
    \item A doctor shines a penlight into a patient's left eye and observes constriction of the right pupil. Name this reflex and explain the neural pathway responsible for the consensual response.
    \item Distinguish between \emph{accommodation} and \emph{adaptation}. State two differences.
\end{enumerate}
\end{exercice}

\begin{corrige}[Exercise 1]
Vitamin A (retinol) is the precursor for retinal (11-cis-retinal), the chromophore of rhodopsin. In deficiency, rhodopsin regeneration in darkness is impaired. Rods (responsible for scotopic/dim light vision) cannot recover sensitivity after bleaching. Cones (photopic/bright light vision) use different photopigments (photopsins) and function adequately in bright light where rhodopsin is mostly bleached anyway. Hence night blindness (rod failure) with normal day vision (cone function).
\end{corrige}

\begin{corrige}[Exercise 2]
1. Diverging rays from near object enter eye $\to$ image focused behind retina (blur). 2. Blur detected $\to$ visual cortex $\to$ parasympathetic outflow via CN III (oculomotor) to ciliary ganglion. 3. Postganglionic fibres stimulate \textbf{ciliary muscle contraction}. 4. Ciliary ring moves forward/inward $\to$ \textbf{suspensory ligaments slacken}. 5. \textbf{Lens elasticity} causes lens to become \textbf{more convex} (rounder). 6. Increased refractive power bends rays more sharply $\to$ image focused on retina.
\end{corrige}

\begin{corrige}[Exercise 3]
\textbf{Consensual Pupil Light Reflex.} Pathway: Light $\to$ Left Retina $\to$ Left Optic Nerve (CN II) $\to$ Optic Chiasm (partial decussation) $\to$ Left \& Right Optic Tracts $\to$ Left \& Right Pretectal Nuclei (Midbrain) $\to$ Interneurones cross via Posterior Commissure $\to$ Bilateral Edinger-Westphal Nuclei $\to$ Bilateral CN III (Parasympathetic) $\to$ Ciliary Ganglia $\to$ Short Ciliary Nerves $\to$ Circular Muscles of \textbf{both} irises contract $\to$ Bilateral Miosis.
\end{corrige}

\begin{corrige}[Exercise 4]
\begin{center}
\begin{tabularx}{\linewidth}{|l|X|X|}\hline
\textbf{Feature} & \textbf{Accommodation} & \textbf{Adaptation (Retinal)} \\ \hline
\textbf{Mechanism} & Physical change in lens shape (curvature) & Biochemical change in photopigment (rhodopsin) concentration/state \\ \hline
\textbf{Structures} & Ciliary muscle, suspensory ligaments, lens & Rods, cones, retinal pigment epithelium, Vitamin A \\ \hline
\textbf{Stimulus} & Object distance (blur/disparity) & Light intensity (photon flux) \\ \hline
\textbf{Time Scale} & Fast (milliseconds to seconds) & Light adapt: seconds; Dark adapt: 20--30 minutes \\ \hline
\textbf{Purpose} & Focus image on retina (sharpness) & Adjust retinal sensitivity (dynamic range) \\ \hline
\end{tabularx}
\end{center}
\end{corrige}

\coursec{Photoreception and the Ear: Hearing and Balance}

In the previous section we saw how the eye controls the amount of light entering it (the iris reflex), how it adapts to light and darkness, and how it focuses images by accommodation. But forming a sharp image on the retina is only half the story. The retina must now \emph{convert} light energy into electrical signals that the brain can interpret. This conversion is the work of the photoreceptor cells — the \cle{rods} and \cle{cones} — and it is the subject of \cle{photoreception}. We shall then turn to a second sense organ, the \cle{ear}, which performs two apparently different jobs with the same basic strategy: detecting sound (\cle{hearing}) and detecting head position and movement (\cle{balance}).

\begin{definition}[Photoreception]
\cle{Photoreception} is the process by which specialised receptor cells (photoreceptors) absorb light energy and convert it into electrical signals (receptor/generator potentials) that are transmitted to the central nervous system. In mammals the photoreceptors are the \cle{rod cells} and \cle{cone cells} of the retina. The conversion of one form of energy into an electrical signal by a receptor is called \cle{transduction}.
\end{definition}

\courssub{Structure of Rod and Cone Cells}

Rods and cones are modified neurones lying in the deepest layer of the retina (light must pass through the other retinal layers to reach them). Each cell has four recognisable regions:

\begin{itemize}
\item The \cle{outer segment}: a stack of flattened membranous discs whose membranes carry the \cle{visual pigments} (rhodopsin in rods, iodopsins in cones). This is where light is actually absorbed.
\item A narrow connecting \cle{cilium} linking the outer and inner segments.
\item The \cle{inner segment}: packed with mitochondria (to supply ATP for the active pumping of ions) and containing the nucleus.
\item The \cle{synaptic terminal} (synaptic body): forms synapses with \cle{bipolar cells}, which in turn connect to \cle{ganglion cells} whose axons form the optic nerve.
\end{itemize}

\begin{popfigure}
\begin{tikzpicture}[scale=0.95, font=\small]
% Rod cell
\begin{scope}
\draw[fill=popblueL, draw=popblue] (0,0) rectangle (1.4,2.6);
\foreach \y in {0.2,0.5,0.8,1.1,1.4,1.7,2.0,2.3}{\draw[popblue] (0.08,\y) -- (1.32,\y);}
\node[align=center] at (0.7,-0.45) {outer segment\\ (discs with rhodopsin)};
\draw[fill=popgreenL, draw=popgreen] (0.35,2.6) rectangle (1.05,4.1);
\draw[fill=popgold, draw=popdark] (0.7,3.35) circle (0.28);
\node[align=center] at (2.6,3.4) {inner segment\\ (mitochondria, nucleus)};
\draw[fill=poppink!40, draw=poppink] (0.2,4.1) -- (1.2,4.1) -- (1.0,4.9) -- (0.4,4.9) -- cycle;
\node[align=center] at (2.6,4.6) {synaptic terminal\\ (to bipolar cell)};
\node[align=center, text=popblue] at (-1.3,1.3) {light\\ $\rightarrow$};
\node[font=\bfseries] at (0.7,5.4) {Rod cell};
\end{scope}
% Rhodopsin cycle
\begin{scope}[xshift=8.2cm]
\node[font=\bfseries] at (1.6,5.4) {The rhodopsin cycle};
\node[draw=popblue, fill=popblueL, rounded corners, align=center, minimum width=3.4cm] (rh) at (1.6,4.2) {rhodopsin\\ (11-\emph{cis}-retinal + opsin)};
\node[draw=poporange, fill=poporange!20, rounded corners, align=center, minimum width=3.4cm] (br) at (1.6,1.6) {retinal (\emph{trans}) + opsin\\ (bleached products)};
\draw[->, very thick, popdark] (rh.east) .. controls (4.2,4.2) and (4.2,1.6) .. node[right, align=center]{light absorbed\\ (bleaching)} (br.east);
\draw[->, very thick, popgreen] (br.west) .. controls (-0.9,1.6) and (-0.9,4.2) .. node[left, align=center]{dark: enzyme + ATP\\ (resynthesis)} (rh.west);
\node[align=center, text=popmuted] at (1.6,0.3) {bleaching triggers a generator\\ potential in the rod};
\end{scope}
\end{tikzpicture}
\legende{A rod cell (left) and the photochemical cycle of rhodopsin (right).}
\end{popfigure}

\courssub{The Photochemistry of Vision}

The key event in photoreception is a \emph{chemical} one. The visual pigment of rods, \cle{rhodopsin} (visual purple), consists of a protein, \cle{opsin}, bound to a light-sensitive derivative of vitamin A, \cle{retinal} (retinene), in the 11-\emph{cis} form.

\begin{enumerate}
\item In the \textbf{dark}, rhodopsin is intact. The rod membrane maintains a \emph{dark current}: sodium ions leak in through cGMP-gated channels, while the Na$^+$/K$^+$ ATPase actively pumps them out. The inside of the cell is partially negative (around $-40$ mV) and the rod \emph{continuously releases} the neurotransmitter glutamate onto its bipolar cell.
\item When \textbf{light} strikes the outer segment, a photon is absorbed by retinal, which changes shape from the \emph{cis} to the \emph{trans} form and splits away from opsin — rhodopsin is said to be \cle{bleached}.
\item Activated opsin (metarhodopsin II) triggers a G-protein (transducin) cascade that activates phosphodiesterase (PDE). PDE hydrolyses cGMP to GMP, causing the cGMP-gated Na$^+$ channels to \emph{close}.
\item The Na$^+$/K$^+$ pump continues working, so the inside of the cell becomes \emph{more} negative (hyperpolarises to about $-70$ mV). This change in membrane potential is the \cle{generator potential} (receptor potential).
\item The generator potential reduces glutamate release, and this change is detected by the bipolar cell and relayed via ganglion cells along the optic nerve to the visual cortex.
\item In the dark, retinal is converted back to the \emph{cis} form by retinal isomerase and recombines with opsin (using ATP): rhodopsin is \cle{resynthesised}. This is the chemical basis of \cle{dark adaptation} studied earlier.
\end{enumerate}

\begin{attention}[A surprising sign]
Unlike most receptors, rods \emph{hyperpolarise} in response to their stimulus and release neurotransmitter in the \emph{dark}. Do not write that light "depolarises" the rod — this is a classic error that loses marks.
\end{attention}

\courssub{Cones and Colour Vision}

Cones contain pigments called \cle{iodopsins}, built on the same retinal + opsin plan but with three slightly different opsins. This gives \textbf{three types of cone}, each most sensitive to a different region of the spectrum:

\begin{itemize}
\item \cle{red-sensitive} cones (L-cones, peak absorption $\approx 560$ nm, long wavelengths),
\item \cle{green-sensitive} cones (M-cones, peak $\approx 530$ nm, medium wavelengths),
\item \cle{blue-sensitive} cones (S-cones, peak $\approx 420$ nm, short wavelengths).
\end{itemize}

According to the \cle{trichromatic theory} of colour vision, every colour we perceive results from the \emph{relative} stimulation of these three cone types. For example, yellow light ($\approx 580$ nm) stimulates red and green cones roughly equally and blue cones weakly; the brain interprets this particular ratio as "yellow". White light stimulates all three strongly and about equally. A person lacking one functional cone type is \cle{colour-blind} (e.g.\ red–green colour blindness, a sex-linked condition commoner in males).

\courssub{Rods versus Cones}

\begin{propriete}[Sensitivity and acuity]
\cle{Rods} give \textbf{high sensitivity} in dim light (they contain rhodopsin, which is bleached even by weak light, and many rods converge onto one bipolar cell, so their signals summate) but \textbf{low acuity} and no colour vision. \cle{Cones} need bright light but give \textbf{high acuity} (each cone in the fovea has its own bipolar cell, so signals are kept separate) and \textbf{colour vision}.
\end{propriete}

\begin{aretenir}[Comparing rods and cones — a classic GCE table]
\begin{center}
\begin{tabularx}{\linewidth}{|l|X|X|}
\hline
\rowcolor{popblueL} \textbf{Feature} & \textbf{Rods} & \textbf{Cones} \\
\hline
Pigment & Rhodopsin (one type) & Iodopsins (three types: red, green, blue) \\
\hline
Sensitivity to light & High — work in dim light & Low — need bright light \\
\hline
Visual acuity & Low (many rods per bipolar cell: convergence) & High (one cone per bipolar cell at fovea) \\
\hline
Colour vision & No (black, white, grey) & Yes (trichromatic) \\
\hline
Distribution & Peripheral retina; absent from fovea & Concentrated at the fovea \\
\hline
Function & Night vision, movement detection & Daylight vision, fine detail, colour \\
\hline
\end{tabularx}
\end{center}
\end{aretenir}

\begin{exemplebox}[Why stars "disappear" when you look straight at them]
Astronomers — and hunters on dark nights on the Adamawa plateau — know that a faint star is easier to see when you look \emph{slightly to one side} of it. Looking directly at the star places its image on the \cle{fovea}, which contains only cones, and cones are insensitive in dim light. Looking slightly aside moves the image onto the peripheral retina, rich in highly sensitive rods, so the star becomes visible.
\end{exemplebox}

\courssub{Structure of the Ear}

The ear is a double sense organ: the \cle{cochlea} detects sound, while the \cle{vestibular apparatus} (semicircular canals, utricle and saccule) detects head position and movement. Anatomically it has three parts.

\textbf{The outer ear.} The \cle{pinna} (visible ear flap) collects sound waves and funnels them into the \cle{auditory canal} (external auditory meatus), at the end of which lies the \cle{eardrum} (tympanic membrane), a thin membrane that vibrates when sound waves strike it.

\textbf{The middle ear.} An air-filled cavity in the skull containing the three \cle{ossicles}, the smallest bones in the body: the \cle{malleus} (hammer, attached to the eardrum), the \cle{incus} (anvil) and the \cle{stapes} (stirrup, whose footplate fits into the \cle{oval window} of the cochlea). The \cle{Eustachian tube} connects the middle ear to the pharynx and equalises air pressure on both sides of the eardrum — this is why swallowing relieves the "popping" of the ears when a bus climbs the steep road up Mount Cameroon.

\textbf{The inner ear.} A fluid-filled labyrinth in the temporal bone comprising:
\begin{itemize}
\item the \cle{cochlea} — a spiral, snail-shaped tube containing the sensory cells of hearing;
\item the three \cle{semicircular canals} — arranged at right angles to one another, detecting rotational movement;
\item the \cle{utricle} and \cle{saccule} — detecting gravity and linear acceleration.
\end{itemize}

\begin{popfigure}
\begin{tikzpicture}[scale=0.9, font=\small]
% Outer ear
\draw[thick, popdark, fill=popblueL] (0,2.2) .. controls (-0.9,3.4) and (-0.9,1.0) .. (0,1.8);
\node[align=center] at (-1.5,2.2) {pinna};
\draw[thick, popdark] (0,2.35) -- (2.6,2.35);
\draw[thick, popdark] (0,1.65) -- (2.6,1.65);
\node[align=center] at (1.3,2.75) {auditory canal};
% Eardrum
\draw[very thick, poporange] (2.6,2.5) -- (2.6,1.5);
\node[align=center] at (2.6,1.15) {eardrum};
% Ossicles
\draw[thick, popdark] (2.6,2.0) -- (3.3,2.35) -- (4.0,2.0) -- (4.6,2.3);
\fill[popdark] (3.3,2.35) circle (0.07);
\fill[popdark] (4.0,2.0) circle (0.07);
\fill[popdark] (4.6,2.3) circle (0.07);
\node[align=center] at (3.6,3.0) {ossicles: malleus,\\ incus, stapes};
% Eustachian tube
\draw[thick, popmuted] (3.2,1.5) -- (3.2,0.5) -- (4.4,0.5);
\node[align=center] at (3.5,0.15) {Eustachian tube};
% Cochlea (spiral)
\draw[thick, popblue, fill=popblueL] (6.2,2.0) circle (0.95);
\draw[thick, popblue] (6.2,2.0) circle (0.55);
\draw[thick, popblue] (6.2,2.0) circle (0.2);
\node[align=center] at (6.2,0.7) {cochlea\\ (organ of Corti inside)};
% Oval window
\draw[very thick, poporange] (5.25,2.3) -- (5.25,1.9);
\node[align=center] at (4.9,2.75) {oval\\ window};
% Round window
\draw[very thick, poporange] (5.25,1.7) -- (5.25,1.3);
\node[align=center] at (4.9,1.0) {round\\ window};
% Semicircular canals
\draw[thick, popgreen, fill=popgreenL] (5.6,3.6) arc (200:-20:0.7) -- cycle;
\draw[thick, popgreen, fill=popgreenL] (7.0,3.7) arc (160:380:0.55);
\node[align=center] at (6.6,4.6) {semicircular canals};
% Utricle/saccule
\draw[thick, poppurple, fill=poppink!30] (5.9,3.0) circle (0.28);
\draw[thick, poppurple, fill=poppink!30] (6.5,3.05) circle (0.22);
\node[align=center] at (7.9,3.1) {utricle \&\\ saccule};
% Auditory nerve
\draw[very thick, popgold] (7.15,2.0) -- (8.6,2.0);
\node[align=center] at (8.0,1.6) {auditory nerve};
% Vestibular nerve
\draw[very thick, popgold] (7.15,3.0) -- (8.6,3.0);
\node[align=center] at (8.0,3.4) {vestibular nerve};
\node[font=\bfseries] at (1.3,3.6) {outer ear};
\node[font=\bfseries] at (3.6,3.6) {middle ear};
\node[font=\bfseries] at (6.4,5.2) {inner ear};
\end{tikzpicture}
\legende{Section through the human ear showing the outer, middle and inner ear.}
\end{popfigure}

\courssub{The Mechanism of Hearing}

Hearing is a chain of energy conversions: sound waves (mechanical energy in air) $\rightarrow$ mechanical vibrations of bones $\rightarrow$ pressure waves in fluid $\rightarrow$ electrical impulses in the auditory nerve. Follow the sequence carefully:

\begin{methode}[From sound wave to nerve impulse]
\begin{enumerate}
\item Sound waves are collected by the \textbf{pinna} and channelled along the \textbf{auditory canal} to the \textbf{eardrum}, which vibrates at the same frequency as the sound.
\item The vibrations pass across the middle ear through the \textbf{ossicles} (malleus $\rightarrow$ incus $\rightarrow$ stapes). The ossicles act as levers and the stapes footplate is much smaller than the eardrum (area ratio $\approx 17:1$), so the force is concentrated: vibrations are \textbf{amplified} (about twenty-fold in pressure) to move the dense inner-ear fluid efficiently.
\item The stapes pushes on the \textbf{oval window}, setting up pressure waves in the \textbf{perilymph}, the fluid filling the scala vestibuli and scala tympani of the cochlea.
\item The waves travel along the cochlea and make the \textbf{basilar membrane} vibrate. Different frequencies vibrate different regions of the membrane (high pitch near the base, low pitch near the apex) — this \cle{place theory} (tonotopic organisation) is how pitch is discriminated.
\item Sitting on the basilar membrane is the \textbf{organ of Corti}, whose \textbf{hair cells} have sensory hairs (stereocilia) embedded near the overlying tectorial membrane. Vibration of the basilar membrane causes a shearing force that bends the stereocilia.
\item Bending opens mechanically-gated K$^+$ channels in the hair cells (endolymph is high in K$^+$), generating a \textbf{receptor potential} that triggers \textbf{nerve impulses} in fibres of the \textbf{auditory nerve}, which carries them to the auditory cortex of the brain, where they are interpreted as sound.
\end{enumerate}
\end{methode}

\begin{definition}[The organ of Corti]
The \cle{organ of Corti} is the sensory structure of hearing, located on the basilar membrane inside the cochlear duct (scala media). It consists of \cle{inner hair cells} (true sensory receptors) and \cle{outer hair cells} (amplifiers), whose stereocilia project towards the tectorial membrane, together with supporting cells. Bending of the stereocilia by vibrations of the basilar membrane generates receptor potentials that excite the auditory nerve.
\end{definition}

\begin{attention}[The ossicles do not detect sound]
A frequent mistake is to say that the ossicles "receive" or "detect" sound. They do not: the ossicles merely \textbf{transmit and amplify} vibrations from the eardrum to the oval window. Detection (transduction into electrical signals) occurs only in the \textbf{hair cells of the organ of Corti}. Similarly, the eardrum vibrates but detects nothing.
\end{attention}

\courssub{Balance: The Vestibular Apparatus}

Balance uses the same transduction principle as hearing — bending of hair cell stereocilia — but the stimulus is movement of the head rather than sound.

\textbf{Rotational movement: the semicircular canals.} The three semicircular canals lie in three mutually perpendicular planes (horizontal, anterior vertical, posterior vertical). At the base of each canal is a swelling, the \cle{ampulla}, containing a patch of hair cells (the \cle{crista ampullaris}) whose stereocilia are embedded in a jelly-like cup, the \cle{cupula}. When the head rotates, the fluid (endolymph) inside the canal lags behind because of inertia, pushing the cupula and bending the stereocilia. Because the three canals are at right angles, rotation in \emph{any} plane is detected. This is why spinning round and then stopping makes you dizzy: the fluid keeps moving briefly, bending the cupula even though your head has stopped.

\textbf{Gravity and linear acceleration: the utricle and saccule.} These two chambers contain hair cells (the \cle{maculae}) overlain by a jelly layer loaded with tiny calcium carbonate crystals, the \cle{otoliths} (otoconia). When the head tilts or accelerates in a straight line (in a lift, or in a \emph{bendskin} pulling away), gravity or inertia shifts the heavy otoliths, bending the stereocilia. The pattern of bending tells the brain the orientation of the head relative to gravity and any linear acceleration.

Impulses from both systems travel along the \cle{vestibular nerve} (part of the vestibulocochlear nerve, CN VIII) to the \cle{cerebellum} and vestibular nuclei, which coordinate posture and eye movements (vestibulo-ocular reflex) to keep us balanced.

\begin{exoresolu}[Explaining sensitivity and acuity in the retina]
\textbf{Statement.} (a) Explain why rods allow us to see in dim light. (b) Explain why cones, but not rods, give high visual acuity. (c) A patient cannot distinguish red from green in bright light but sees normally at night. Suggest an explanation.
\tcblower
\textbf{Solution.}
\begin{enumerate}
\item \textbf{(a)} Rods contain rhodopsin, which is bleached even by very weak light, so a single photon can contribute to a response. In addition, many rods converge onto a single bipolar cell (\cle{retinal convergence}); their small generator potentials summate, so a threshold is reached even in dim light. Hence rods are highly sensitive.
\item \textbf{(b)} At the fovea each cone has its \emph{own} bipolar cell and ganglion cell (no convergence). Stimulation of two neighbouring cones therefore sends two \emph{separate} signals to the brain, which can distinguish two close points as distinct — this is high acuity. In rods, convergence pools signals from many cells, so the brain cannot tell which rod was stimulated: acuity is poor.
\item \textbf{(c)} Night vision depends on rods, which are normal; colour vision depends on cones. The patient most likely lacks functional red-sensitive or green-sensitive cones (red–green colour blindness), a defect of cone pigments that leaves rod function intact.
\end{enumerate}
\end{exoresolu}

\begin{exoresolu}[Mechanism of hearing and balance]
\textbf{Statement.} (a) Describe how the ossicles amplify sound vibrations. (b) Explain how the organ of Corti converts vibrations into nerve impulses. (c) A person spins rapidly and then stops suddenly. Explain why they feel dizzy.
\tcblower
\textbf{Solution.}
\begin{enumerate}
\item \textbf{(a)} The ossicles act as a lever system: the malleus (attached to the eardrum) moves a greater distance than the stapes (attached to the oval window). Combined with the area ratio of eardrum to stapes footplate ($\approx 17:1$), this amplifies the pressure about 20-fold, overcoming the impedance mismatch between air and cochlear fluid.
\item \textbf{(b)} Vibrations of the stapes create pressure waves in perilymph, which travel through the cochlea and displace the basilar membrane. This shears the stereocilia of inner hair cells against the tectorial membrane. Bending opens K$^+$ channels (endolymph is K$^+$-rich), causing K$^+$ influx, depolarisation (receptor potential), Ca$^{2+}$ entry, and glutamate release onto auditory nerve fibres, generating action potentials.
\item \textbf{(c)} During spinning, endolymph in the horizontal semicircular canal accelerates due to friction with the canal walls. When rotation stops, the endolymph continues moving (inertia), deflecting the cupula in the opposite direction. This signals continued rotation to the brain, conflicting with visual and proprioceptive input, causing dizziness (vertigo).
\end{enumerate}
\end{exoresolu}

\begin{exercice}[Check your understanding]
\begin{enumerate}
\item List, in order, the structures through which a vibration passes from the air to the auditory nerve.
\item Explain why the ears "pop" when altitude changes rapidly, and name the structure responsible.
\item State one structural difference and one functional difference between the semicircular canals and the utricle.
\item Explain why a driver at night sees a pedestrian at the roadside more easily in peripheral vision than by looking directly at them.
\item Describe the role of the tectorial membrane in hearing.
\item What is the function of the round window in the cochlea?
\end{enumerate}
\end{exercice}

\begin{corrige}[Exercise 1]
\begin{enumerate}
\item Air (auditory canal) $\rightarrow$ eardrum $\rightarrow$ malleus $\rightarrow$ incus $\rightarrow$ stapes $\rightarrow$ oval window $\rightarrow$ perilymph of scala vestibuli $\rightarrow$ basilar membrane $\rightarrow$ hair cells of the organ of Corti $\rightarrow$ auditory nerve.
\item A rapid altitude change makes the air pressure outside the eardrum differ from that in the middle ear, bulging the eardrum. Swallowing opens the \cle{Eustachian tube}, equalising the pressure — the "pop".
\item Structural: semicircular canals are three looped tubes with ampullae containing a cupula; the utricle is a chamber whose hair cells bear otoliths. Functional: semicircular canals detect \emph{rotational} movements of the head; the utricle detects \emph{gravity and linear acceleration}.
\item Looking directly at the pedestrian focuses the image on the fovea, which has only cones; cones are insensitive in dim light. Peripheral vision uses rod-rich retina, and rods are highly sensitive to dim light, so the pedestrian is detected more easily.
\item The tectorial membrane overlies the hair cells; as the basilar membrane vibrates, the hair cell stereocilia shear against the tectorial membrane, bending them and opening transduction channels.
\item The round window (covered by the secondary tympanic membrane) bulges out as the stapes pushes in at the oval window, allowing the incompressible perilymph to move and enabling basilar membrane vibration.
\end{enumerate}
\end{corrige}

\begin{savaistu}[Hearing loss in Cameroon]
Prolonged exposure to very loud sound — for example in unprotected work near grinding machines or with powerful sound systems at \emph{makossa} night-clubs — progressively destroys the hair cells of the organ of Corti. Unlike skin cells, human hair cells do not regenerate: the resulting deafness is permanent. Protecting your ears is protecting an irreplaceable sense organ. The World Health Organization estimates that over 1 billion young people worldwide are at risk of hearing loss due to unsafe listening practices.
\end{savaistu}

\begin{aretenir}[Key points of this section]
\begin{itemize}
\item Photoreception is the transduction of light into electrical signals by rods and cones; light bleaches rhodopsin (retinal + opsin), closing Na$^+$ channels, causing hyperpolarisation — the generator potential.
\item Rods: high sensitivity, low acuity, no colour, peripheral retina. Cones: low sensitivity, high acuity, colour (three types — trichromatic theory), concentrated at the fovea.
\item The ear has three parts: outer (pinna, canal, eardrum), middle (ossicles + Eustachian tube), inner (cochlea, semicircular canals, utricle, saccule).
\item Hearing: ossicles amplify (lever + area ratio); only the hair cells of the organ of Corti transduce vibrations into impulses via K$^+$ influx from endolymph.
\item Balance: semicircular canals detect rotation (cupula deflection by endolymph inertia); utricle and saccule detect gravity and linear acceleration (otolith displacement).
\end{itemize}
\end{aretenir}

\coursec{The Nervous System and Neurones}

In the previous section we saw how the eye converts light into electrical signals and how the ear converts sound waves and head movements into impulses. But detection is only the beginning of the story. Once a receptor in the retina or the cochlea has been stimulated, the message must travel rapidly to a coordinating centre, be compared with other incoming information, and an appropriate command must be sent back to muscles or glands. That entire task — reception, transport, processing and command — belongs to the \cle{nervous system}. In this section we study how the nervous system is organised, the structure of its cellular unit, the \cle{neurone}, and how the central nervous system \emph{integrates} information.

\courssub{Organisation of the Nervous System}

Imagine the administration of Cameroon: ministries in Yaoundé make decisions, but they are connected to every division and village by roads and telephone lines. The nervous system is organised in exactly the same way, with a central command and a vast communication network reaching every part of the body.

\begin{definition}[Divisions of the nervous system]
The nervous system is divided into two parts:
\begin{itemize}
\item The \cle{central nervous system} (CNS), made of the \textbf{brain} (protected inside the skull) and the \textbf{spinal cord} (protected inside the vertebral column). It is the coordinating and decision-making centre.
\item The \cle{peripheral nervous system} (PNS), made of all the \textbf{nerves} connecting the CNS to the rest of the body: \textbf{cranial nerves} (arising from the brain, e.g.\ the optic nerve from the eye, the auditory nerve from the ear) and \textbf{spinal nerves} (arising from the spinal cord, one pair per segment, serving the trunk and limbs).
\end{itemize}
\end{definition}

The PNS itself contains two functional categories of fibres: \textbf{afferent (sensory) fibres}, which carry impulses \emph{from receptors to the CNS}, and \textbf{efferent (motor) fibres}, which carry commands \emph{from the CNS to effectors} (muscles and glands). A useful memory trick: \textbf{A}fferent = \textbf{A}rriving at the CNS; \textbf{E}fferent = \textbf{E}xiting the CNS.

\courssub{The Neurone: Unit of the Nervous System}

\begin{definition}[Neurone and nerve]
A \cle{neurone} (nerve cell) is a specialised cell adapted for the rapid reception and transmission of electrical signals called \textbf{nerve impulses}. A \cle{nerve} is a bundle of many neurone fibres (axons) held together by connective tissue, like many electric wires inside one cable.
\end{definition}

\begin{popfigure}
\begin{center}
\begin{tikzpicture}[scale=1.0, every node/.style={font=\small}]
% Cell body
\draw[fill=popblueL, thick, popblue] (0,0) circle (0.9);
\draw[fill=poporange, popdark] (0,0) circle (0.35);
\node at (0,0) {\textbf{N}};
\node[above, popdark] at (0,0.95) {Cell body (soma)};
\node[below, popmuted] at (0,-1.0) {contains nucleus};
% Dendrites
\draw[thick, popblue] (-0.6,0.6) -- (-1.6,1.4);
\draw[thick, popblue] (-1.6,1.4) -- (-2.4,1.6);
\draw[thick, popblue] (-1.6,1.4) -- (-2.1,2.1);
\draw[thick, popblue] (-0.85,0.2) -- (-2.0,0.3);
\draw[thick, popblue] (-2.0,0.3) -- (-2.7,0.7);
\draw[thick, popblue] (-2.0,0.3) -- (-2.8,0.0);
\node[left, popdark, align=center] at (-2.8,1.1) {Dendrites};
% Axon with myelin
\draw[thick, popblue] (0.9,0) -- (1.6,0);
\foreach \x in {1.6, 3.0, 4.4, 5.8} {
  \draw[thick, poppurple, fill=poppurple!25] (\x,-0.28) rectangle (\x+1.1,0.28);
}
\draw[thick, popblue] (1.6,0) -- (6.9,0);
\node[above, popdark] at (3.3,0.35) {Myelin sheath (Schwann cells)};
% Nodes of Ranvier
\foreach \x in {2.7, 4.1, 5.5} {
  \fill[popink] (\x,0) circle (0.05);
}
\draw[<-, popgreen, thick] (4.1,-0.05) -- (4.1,-0.8);
\node[below, popgreen] at (4.1,-0.8) {Node of Ranvier};
% Axon terminals
\draw[thick, popblue] (6.9,0) -- (7.6,0.5);
\draw[thick, popblue] (6.9,0) -- (7.6,-0.5);
\draw[fill=popgold, popdark] (7.7,0.55) circle (0.12);
\draw[fill=popgold, popdark] (7.7,-0.55) circle (0.12);
\node[right, popdark, align=left] at (7.9,0) {Axon terminals\\(synaptic knobs)};
% Direction arrow
\draw[->, very thick, poporange] (-2.2,-1.6) -- (7.2,-1.6);
\node[below, poporange] at (2.5,-1.6) {Direction of impulse};
\node[popdark] at (1.25,-0.55) {Axon};
\end{tikzpicture}
\end{center}
\legende{Structure of a motor neurone. Impulses enter through the dendrites, pass through the cell body and travel along the myelinated axon to the axon terminals.}
\end{popfigure}

\textbf{Parts of a motor neurone and their roles:}
\begin{itemize}
\item \textbf{Dendrites}: highly branched extensions that receive impulses from receptors or from other neurones and conduct them \emph{towards} the cell body.
\item \textbf{Cell body (soma)}: contains the \textbf{nucleus} and abundant cytoplasm with mitochondria, rough endoplasmic reticulum and Golgi apparatus; it controls the cell's metabolism and protein synthesis (including neurotransmitter synthesis).
\item \textbf{Axon}: a single long fibre that conducts impulses \emph{away} from the cell body towards the effector or the next neurone. Motor axons can be remarkably long — the axon controlling the toes runs the whole length of the leg.
\item \textbf{Myelin sheath}: a fatty insulating layer produced by \textbf{Schwann cells} in the PNS and by \textbf{oligodendrocytes} in the CNS, wrapped around the axon.
\item \textbf{Nodes of Ranvier}: short gaps (about \SI{1}{\micro m}) between successive myelinating cells where the axon membrane is exposed and voltage-gated ion channels are concentrated.
\item \textbf{Axon terminals (synaptic knobs)}: branched endings that pass the signal to the next cell across a synapse; they contain mitochondria and synaptic vesicles filled with neurotransmitter.
\end{itemize}

\begin{attention}[Direction of conduction]
A very common examination error is to reverse the direction of travel. Remember: \textbf{dendrites carry impulses \emph{towards} the cell body; the axon carries them \emph{away} from the cell body}. A neurone is a one-way street: dendrite $\rightarrow$ cell body $\rightarrow$ axon $\rightarrow$ terminals.
\end{attention}

\courssub{Types of Neurones}

Three functional types of neurone cooperate in every nervous pathway.

\begin{definition}[The three types of neurones]
\begin{itemize}
\item \textbf{Sensory (afferent) neurone}: carries impulses from a \textbf{receptor} to the CNS. It is typically \textbf{pseudo-unipolar}: a single process emerges from the cell body and divides into a \emph{peripheral branch} (to the receptor) and a \emph{central branch} (to the CNS). Both branches are myelinated. The cell body lies in a \textbf{dorsal root ganglion}, just outside the spinal cord, on a side branch of the fibre.
\item \textbf{Relay (intermediate / association) neurone}: lies entirely \textbf{within the CNS}; it connects sensory neurones to motor neurones and to other relay neurones. It is short, with many dendrites and a short axon; usually unmyelinated.
\item \textbf{Motor (efferent) neurone}: carries impulses from the CNS to an \textbf{effector} (muscle or gland). Its cell body lies \textbf{inside} the CNS (in the grey matter of the spinal cord or brain), and it has a long myelinated axon (in the PNS, myelinated by Schwann cells).
\end{itemize}
\end{definition}

\begin{popfigure}
\begin{center}
\begin{tikzpicture}[scale=0.92, every node/.style={font=\small}]
% ---- Sensory neurone (pseudo-unipolar) ----
\node[popdark, font=\small\bfseries] at (2.0,3.6) {Sensory neurone};
\draw[thick, popblue] (0,2.6) -- (1.2,2.6); % peripheral branch
\draw[fill=popblueL, thick, popblue] (1.6,2.6) circle (0.35); % cell body
\draw[thick, popblue] (1.2,2.6) -- (1.6,2.25); % connection to cell body
\draw[thick, popblue] (1.95,2.6) -- (3.4,2.6); % central branch
\foreach \x in {0.1,2.1} {\draw[thick, poppurple, fill=poppurple!25] (\x,2.45) rectangle (\x+0.9,2.75);} % myelin on both branches
\node[left, popmuted] at (0,2.6) {from receptor};
\node[right, popmuted] at (3.4,2.6) {to CNS};
\node[below, popmuted] at (1.6,2.15) {cell body in ganglion};
% ---- Relay neurone ----
\node[popdark, font=\small\bfseries] at (1.6,1.4) {Relay neurone};
\draw[fill=popgreenL, thick, popgreen] (1.6,0.7) circle (0.3);
\draw[thick, popgreen] (1.35,0.9) -- (0.7,1.2);
\draw[thick, popgreen] (1.35,0.5) -- (0.7,0.2);
\draw[thick, popgreen] (1.9,0.7) -- (2.8,0.7);
\draw[thick, popgreen] (2.8,0.7) -- (3.3,1.0);
\draw[thick, popgreen] (2.8,0.7) -- (3.3,0.4);
\node[below, popmuted] at (1.6,0.3) {short, entirely in CNS};
% ---- Motor neurone ----
\node[popdark, font=\small\bfseries] at (7.6,3.4) {Motor neurone};
\draw[fill=popblueL, thick, popblue] (5.6,2.6) circle (0.35);
\draw[thick, popblue] (5.3,2.8) -- (4.7,3.1);
\draw[thick, popblue] (5.95,2.6) -- (9.6,2.6);
\foreach \x in {6.1,7.3,8.5} {\draw[thick, poppurple, fill=poppurple!25] (\x,2.45) rectangle (\x+1.0,2.75);}
\draw[thick, popblue] (9.6,2.6) -- (10.1,2.9);
\draw[thick, popblue] (9.6,2.6) -- (10.1,2.3);
\draw[fill=popgold, popdark] (10.2,2.95) circle (0.09);
\draw[fill=popgold, popdark] (10.2,2.25) circle (0.09);
\node[left, popmuted] at (5.2,2.6) {cell body in CNS};
\node[right, popmuted, align=left] at (10.35,2.6) {to\\effector};
\node[below, popmuted] at (7.8,2.15) {long myelinated axon};
\end{tikzpicture}
\end{center}
\legende{Comparison of the three neurone types. Note the position of the cell body: in a ganglion outside the CNS (sensory), inside the CNS with a short fibre (relay), and inside the CNS with a long axon (motor). The sensory neurone is pseudo-unipolar: one process splits into peripheral and central myelinated branches.}
\end{popfigure}

\begin{methode}[Identifying a neurone type from a diagram]
When the GCE paper shows you an unlabelled neurone, proceed in order:
\begin{enumerate}
\item \textbf{Locate the cell body.} On a side branch outside the CNS $\Rightarrow$ \textbf{sensory} neurone. At one end of the cell, inside the CNS $\Rightarrow$ \textbf{motor} neurone.
\item \textbf{Look at the length and branching.} A short cell with many dendrites and no long axon, entirely within the CNS $\Rightarrow$ \textbf{relay} neurone.
\item \textbf{Follow the connections.} Attached to a receptor $\Rightarrow$ sensory; attached to a muscle or gland $\Rightarrow$ motor; between the two $\Rightarrow$ relay.
\item \textbf{Check the direction arrows} if given: towards the CNS = sensory; away from the CNS = motor.
\end{enumerate}
\end{methode}

\courssub{The Myelin Sheath and Saltatory Conduction}

Why wrap an axon in fat? For the same reason that electric wires are wrapped in plastic: insulation prevents the signal from leaking away. But myelin does more than insulate — it makes conduction dramatically \emph{faster}.

\begin{definition}[Myelin sheath]
The \cle{myelin sheath} is a fatty, electrically insulating layer formed by \textbf{Schwann cells} (in the PNS) or \textbf{oligodendrocytes} (in the CNS) wound many times around the axon of many neurones. It is interrupted at regular intervals by the \textbf{nodes of Ranvier}, where the axon membrane is exposed and voltage-gated $\mathrm{Na}^+$ and $\mathrm{K}^+$ channels are highly concentrated.
\end{definition}

\begin{propriete}[Saltatory conduction]
In a myelinated axon, the local circuits of current that transmit the impulse cannot cross the insulated regions; they act only at the nodes of Ranvier. The impulse therefore \emph{jumps} from node to node — this is \cle{saltatory conduction} (from Latin \emph{saltare}, to jump). Consequences:
\begin{itemize}
\item Transmission is much \textbf{faster} than in a non-myelinated axon of the same diameter (up to about $100\ \mathrm{m\,s^{-1}}$ in large myelinated fibres, against less than $1\ \mathrm{m\,s^{-1}}$ in small unmyelinated fibres).
\item Less energy is needed, because membrane ion pumps work only at the nodes to restore ionic gradients after each impulse.
\end{itemize}
(The mechanism of the impulse itself is treated in the next section; the jumping property is examinable here.)
\end{propriete}

\begin{savaistu}[When myelin fails]
In the disease \textbf{multiple sclerosis}, the patient's own immune system destroys the myelin sheath in the CNS. Saltatory conduction then fails, impulses slow down or stop, and the person progressively loses control of movement, vision and balance. This illustrates how essential myelin is for rapid, reliable coordination.
\end{savaistu}

\courssub{Nerves, Grey Matter and White Matter}

A single neurone is microscopic, yet you can see a nerve with the naked eye. That is because a \textbf{nerve} is a \emph{bundle} of hundreds or thousands of axons wrapped together by connective tissue (epineurium, perineurium, endoneurium) — exactly like a telephone cable containing many individual wires. Most nerves are \textbf{mixed nerves}, containing both sensory and motor fibres.

Inside the CNS, two kinds of tissue are visible in a cross-section of the spinal cord or brain:

\begin{center}
\begin{tabularx}{\linewidth}{|l|X|X|}
\hline
\rowcolor{popblueL} \textbf{Feature} & \textbf{Grey matter} & \textbf{White matter} \\
\hline
Main content & Cell bodies, dendrites, synapses, unmyelinated axons & Myelinated axons (nerve fibres) \\
\hline
Colour origin & Unmyelinated tissue (greyish) & White fatty myelin \\
\hline
Position in spinal cord & Central (butterfly / H-shape) & Surrounding the grey matter \\
\hline
Position in brain & Outer layer (cortex) and deep nuclei & Beneath the cortex \\
\hline
Function & Integration, processing, synaptic contacts & Rapid transmission between regions \\
\hline
\end{tabularx}
\end{center}

\courssub{Integration: The CNS as a Decision Centre}

At every instant, thousands of receptors — in your eyes, ears, skin, muscles and internal organs — are sending impulses into the CNS. If every message triggered an automatic response, behaviour would be chaotic. The CNS must therefore \emph{select, compare, combine and prioritise} information before commanding effectors.

\begin{definition}[Integration]
\cle{Integration} is the process by which the central nervous system receives impulses from many receptors, processes and coordinates this information (including comparison with stored memory), and sends out appropriate commands to effectors so that the response is adapted to the whole situation.
\end{definition}

\begin{exemplebox}[Integration in everyday life]
You are walking along a road in Bamenda and you hear a motorcycle horn behind you (auditory receptors), you see it approaching in your mirror (photoreceptors), and you feel the weight of your school bag (proprioceptors). Your CNS combines all three messages and commands your leg muscles to step aside \emph{while keeping your balance}. No single reflex could achieve this coordinated, situation-dependent response — it requires integration.
\end{exemplebox}

\textbf{Synapses as decision points.} Integration is physically carried out at \textbf{synapses}, the junctions between neurones. A single relay neurone in the spinal cord may receive synapses from hundreds of other neurones, some \emph{excitatory} (producing EPSPs, pushing it towards firing) and some \emph{inhibitory} (producing IPSPs, holding it back). Whether the relay neurone fires depends on the \emph{sum} of all these influences — \textbf{spatial summation} (simultaneous inputs from different synapses) and \textbf{temporal summation} (rapid successive inputs at the same synapse). Each synapse is therefore a microscopic decision point, and the billions of synapses of the brain together constitute the integrating machinery. We will study synaptic transmission in detail in the next section.

\begin{exoresolu}[Identifying neurones and tracing a pathway]
The diagram of a spinal reflex shows three neurones, A, B and C. Neurone A has its cell body in a swelling (ganglion) on the dorsal root and its ending near the skin. Neurone B is short and lies entirely within the grey matter of the spinal cord. Neurone C has its cell body in the grey matter and its long myelinated axon ends on a biceps muscle. (a) Identify each neurone. (b) State the direction of the impulse through the pathway. (c) Explain why the response of the muscle is rapid.
\tcblower
\textbf{(a)} A is the \textbf{sensory neurone} — its cell body sits in the dorsal root ganglion outside the CNS and it collects impulses from a receptor in the skin. B is the \textbf{relay (intermediate) neurone} — short, entirely within the CNS. C is the \textbf{motor neurone} — cell body in the grey matter of the CNS, long axon running to an effector (the biceps muscle).

\textbf{(b)} Impulse pathway: receptor (skin) $\rightarrow$ A (sensory neurone) $\rightarrow$ B (relay neurone) $\rightarrow$ C (motor neurone) $\rightarrow$ effector (biceps), which contracts.

\textbf{(c)} The response is rapid because (i) nerve impulses are electrical and travel fast, (ii) the motor axon is \textbf{myelinated}, so conduction is \textbf{saltatory} — the impulse jumps between nodes of Ranvier instead of travelling continuously, and (iii) the pathway is short, with few synapses.
\end{exoresolu}

\begin{attention}[Frequent confusions to avoid]
\begin{itemize}
\item Do not confuse a \textbf{neurone} (one cell) with a \textbf{nerve} (a bundle of many axons). Examiners penalise answers that use them interchangeably.
\item The cell body of the \textbf{sensory} neurone is \emph{outside} the CNS (dorsal root ganglion); that of the \textbf{motor} neurone is \emph{inside} the CNS. Reversing these is a classic error.
\item Grey matter is not ``grey neurones'': it is the region rich in \textbf{cell bodies and synapses}; white matter is the region rich in \textbf{myelinated axons}.
\item Myelin does not ``create'' the impulse; it \textbf{insulates} the axon and \textbf{accelerates} conduction by saltatory jumping.
\item Sensory neurones are \textbf{pseudo-unipolar}, not bipolar: a single process divides into two branches. Do not draw two separate processes emerging from the cell body.
\end{itemize}
\end{attention}

\begin{aretenir}[Key points of this section]
\begin{itemize}
\item The nervous system = \textbf{CNS} (brain + spinal cord) + \textbf{PNS} (cranial and spinal nerves).
\item The \textbf{neurone} is the functional unit: dendrites $\rightarrow$ cell body $\rightarrow$ axon $\rightarrow$ terminals.
\item Three types: \textbf{sensory} (receptor $\rightarrow$ CNS, cell body in ganglion, pseudo-unipolar), \textbf{relay} (within CNS, short, unmyelinated), \textbf{motor} (CNS $\rightarrow$ effector, cell body in CNS, long myelinated axon).
\item The \textbf{myelin sheath} (Schwann cells in PNS, oligodendrocytes in CNS) insulates the axon and allows fast \textbf{saltatory conduction} between nodes of Ranvier.
\item A \textbf{nerve} is a bundle of axons wrapped in connective tissue; \textbf{grey matter} = cell bodies and synapses, \textbf{white matter} = myelinated fibres.
\item \textbf{Integration} = reception, processing and coordination of information by the CNS, achieved largely at \textbf{synapses} via spatial and temporal summation of EPSPs and IPSPs.
\end{itemize}
\end{aretenir}

\begin{exercice}[Check your understanding]
\begin{enumerate}
\item Draw and fully label a motor neurone, and add an arrow showing the direction of the impulse.
\item State two structural differences between a sensory neurone and a motor neurone.
\item Explain why an impulse travels faster in a myelinated axon than in a non-myelinated axon of the same diameter.
\item A cross-section of the spinal cord shows a darker central H-shaped region and a paler outer region. Name each region and state what it contains.
\item Using the example of touching a hot cooking pot, explain what is meant by \emph{integration}.
\item Distinguish between spatial summation and temporal summation at a synapse.
\end{enumerate}
\end{exercice}

\begin{corrige}[Exercise 1]
\textbf{1.} Diagram as in the figure of this section: dendrites, cell body with nucleus, axon, myelin sheath with Schwann cells, nodes of Ranvier, axon terminals; arrow pointing from dendrites towards terminals.

\textbf{2.} Sensory neurone: cell body in a dorsal root ganglion \emph{outside} the CNS, on a side branch; pseudo-unipolar (one process splits into peripheral and central branches); carries impulses from receptor to CNS. Motor neurone: cell body \emph{inside} the CNS at the end of the cell; multipolar (many dendrites, one long axon); carries impulses to an effector.

\textbf{3.} Myelin insulates the axon so that local currents act only at the nodes of Ranvier; the impulse jumps from node to node (saltatory conduction), which is faster than continuous conduction along the whole membrane. Voltage-gated channels are concentrated at nodes.

\textbf{4.} Central H-shaped region = \textbf{grey matter} (cell bodies, dendrites, synapses, unmyelinated axons); outer pale region = \textbf{white matter} (myelinated axons running up and down the cord).

\textbf{5.} Touching the hot pot stimulates pain and temperature receptors; sensory neurones carry impulses to the CNS, which processes the information (identifying danger, comparing with experience) and coordinates a command via motor neurones to the arm muscles to withdraw the hand — and possibly to other muscles to keep balance. This selection and coordination of an appropriate response by the CNS is integration.

\textbf{6.} \textbf{Spatial summation}: simultaneous EPSPs from several different presynaptic neurones add together to reach threshold. \textbf{Temporal summation}: rapid successive EPSPs from the same presynaptic neurone add together because the previous depolarisation has not yet decayed.
\end{corrige}

\coursec{Transmission of Impulses and the Synapse}

In the previous section we examined the structure of neurones and the organisation of the nervous system. We now turn to the fundamental question: \emph{how does a neurone actually transmit a signal?} The answer lies in rapid, temporary changes in the electrical potential across the neuronal membrane — the \cle{action potential} — and in the specialised junction where one neurone communicates with another: the \cle{synapse}. Understanding these mechanisms is essential not only for the GCE examination but also for appreciating how drugs, toxins and diseases affect nervous function.

\courssub{Resting Potential: The Polarised Membrane}

Before a neurone can fire, it must establish a \cle{resting potential} — a voltage difference across its membrane with the inside negative relative to the outside. In most mammalian neurones this value is approximately \SI{-70}{mV}.

\begin{definition}[Resting Potential]
The resting potential is the steady membrane potential of an excitable cell (neurone or muscle fibre) when it is not transmitting an impulse. It is typically around \SI{-70}{mV} (inside negative) and is maintained by the unequal distribution of ions and the selective permeability of the membrane.
\end{definition}

Three factors create and maintain this potential:

\begin{enumerate}
\item \textbf{Concentration gradients:} The extracellular fluid is rich in \ce{Na+} and \ce{Cl-}, while the intracellular fluid is rich in \ce{K+} and large organic anions (\ce{A-}). These gradients are established by the \cle{sodium–potassium pump} (\ce{Na+/K+}-ATPase), which actively transports \SI{3}{\ce{Na+}} out and \SI{2}{\ce{K+}} in per ATP hydrolysed.
\item \textbf{Selective permeability:} At rest, the membrane is far more permeable to \ce{K+} (via \cle{leak channels}) than to \ce{Na+} or \ce{Cl-}. \ce{K+} diffuses out down its concentration gradient, leaving behind unbalanced negative charges (mainly \ce{A-}) inside.
\item \textbf{Electrochemical equilibrium:} \ce{K+} efflux continues until the electrical attraction pulling \ce{K+} back in balances the chemical gradient pushing it out. This equilibrium potential for \ce{K+} (calculated by the Nernst equation) is about \SI{-90}{mV}. The resting potential is slightly less negative because of a small \ce{Na+} leak.
\end{enumerate}

\begin{aretenir}[Key Points on Resting Potential]
\begin{itemize}
\item The \ce{Na+/K+} pump is electrogenic (moves 3 positive charges out, 2 in) and contributes directly a few millivolts to the negativity inside.
\item The membrane is \emph{polarised} at rest. Any stimulus that reduces this polarisation (makes the inside less negative) is a \cle{depolarisation}; an increase in negativity is a \cle{hyperpolarisation}.
\item The resting potential is a \emph{steady state}, not a true equilibrium, because it requires continuous ATP expenditure by the pump.
\end{itemize}
\end{aretenir}

\courssub{Action Potential: The Nerve Impulse}

When a stimulus depolarises the membrane to a critical \cle{threshold} (typically around \SI{-55}{mV}), a massive, transient reversal of the membrane potential occurs: the \cle{action potential}.

\begin{definition}[Action Potential]
An action potential is a rapid, all-or-nothing reversal of the membrane potential from about \SI{-70}{mV} to about \SI{+30}{mV} and back, lasting 1–2 ms in myelinated neurones. It is the electrical signal that propagates along the axon.
\end{definition}

\begin{propriete}[Phases of the Action Potential]
The action potential results from the sequential opening and closing of \cle{voltage-gated ion channels}:
\begin{enumerate}
\item \textbf{Depolarisation (rising phase):} Voltage-gated \ce{Na+} channels open rapidly. \ce{Na+} influx drives the membrane potential towards the \ce{Na+} equilibrium potential ($\approx$ \SI{+60}{mV}). The membrane potential overshoots zero.
\item \textbf{Repolarisation (falling phase):} Voltage-gated \ce{Na+} channels \emph{inactivate} (a distinct conformational state). Voltage-gated \ce{K+} channels open more slowly. \ce{K+} efflux restores the negative interior.
\item \textbf{Hyperpolarisation (undershoot):} \ce{K+} channels remain open slightly longer than needed, driving the potential closer to the \ce{K+} equilibrium potential ($\approx$ \SI{-90}{mV}).
\item \textbf{Return to rest:} \ce{K+} channels close; the \ce{Na+/K+} pump restores the exact ion distributions.
\end{enumerate}
\end{propriete}

\begin{popfigure}
\begin{tikzpicture}
\begin{axis}[
    width=\linewidth,
    height=7cm,
    xlabel={Time (ms)},
    ylabel={Membrane potential (mV)},
    xmin=0, xmax=5,
    ymin=-100, ymax=50,
    xtick={0,0.5,1,1.5,2,2.5,3,3.5,4,4.5,5},
    ytick={-90,-70,-55,0,30},
    grid=major,
    axis lines=middle,
    tick label style={font=\small},
    label style={font=\small},
    clip=false
]
% Resting potential line
\addplot[popblue, thick, domain=0:0.5] {-70};
% Action potential curve (schematic)
\addplot[poporange, very thick, smooth, domain=0.5:4.5, samples=200] {
    -70 + 100/(1+exp(-20*(x-0.7))) - 130/(1+exp(-15*(x-1.2))) + 20/(1+exp(-10*(x-2.5)))
};
% Threshold line
\addplot[popmuted, dashed, domain=0:5] {-55};
% Annotations
\node[popblue, anchor=east] at (axis cs:0.2,-70) {Resting potential (\SI{-70}{mV})};
\node[popmuted, anchor=west] at (axis cs:0.2,-55) {Threshold (\SI{-55}{mV})};
\node[poporange, anchor=south] at (axis cs:1.0,30) {Overshoot (\approx \SI{+30}{mV})};
\node[popgreen, anchor=north] at (axis cs:2.5,-85) {Hyperpolarisation};
\draw[popline, <->] (axis cs:0.5,-70) -- (axis cs:0.5,30) node[midway, right, font=\tiny] {Depolarisation};
\draw[popline, <->] (axis cs:1.2,30) -- (axis cs:1.2,-70) node[midway, right, font=\tiny] {Repolarisation};
\end{axis}
\end{tikzpicture}
\legende{Typical action potential waveform in a myelinated mammalian neurone. The threshold is the voltage at which voltage-gated \ce{Na+} channels open regeneratively.}
\end{popfigure}

\begin{definition}[Refractory Periods]
After an action potential, the membrane cannot immediately fire again.
\begin{itemize}
\item \textbf{Absolute refractory period:} From the start of depolarisation until \ce{Na+} channels recover from inactivation (approx. 1–2 ms). \emph{No stimulus, however strong, can trigger a second action potential.} This enforces one-way propagation.
\item \textbf{Relative refractory period:} During hyperpolarisation. A \emph{stronger-than-threshold} stimulus can elicit a second action potential because the membrane is closer to the \ce{K+} equilibrium potential.
\end{itemize}
\end{definition}

\begin{propriete}[All-or-Nothing Law]
The action potential is an \cle{all-or-nothing} event: if the stimulus reaches threshold, the action potential always has the same amplitude and shape for a given neurone; if the stimulus is subthreshold, no action potential occurs. The nervous system encodes stimulus \emph{intensity} by the \emph{frequency} of action potentials, not their size.
\end{propriete}

\begin{attention}[Common Mistake: Graded Potentials vs. Action Potentials]
Do not confuse \emph{graded potentials} (local, decremental, amplitude varies with stimulus strength — e.g. generator potentials at receptors, EPSPs/IPSPs at synapses) with \emph{action potentials} (all-or-nothing, non-decremental, fixed amplitude). Only action potentials propagate over long distances.
\end{attention}

\courssub{Propagation of the Action Potential}

The action potential at one patch of membrane depolarises the adjacent membrane, triggering a new action potential there. This \cle{local current flow} creates a wave of depolarisation that travels along the axon.

\begin{propriete}[Factors Affecting Conduction Speed]
\begin{enumerate}
\item \textbf{Myelination:} In myelinated axons, voltage-gated channels are concentrated at the \cle{nodes of Ranvier}. The action potential \emph{jumps} from node to node (\cle{saltatory conduction}), which is much faster and energetically cheaper than continuous conduction in unmyelinated fibres.
\item \textbf{Axon diameter:} Larger diameter $\to$ lower axial resistance $\to$ faster local current spread $\to$ higher speed.
\item \textbf{Temperature:} Higher temperature increases ion channel kinetics and diffusion rates, speeding conduction (up to a physiological limit).
\end{enumerate}
\end{propriete}

\begin{savaistu}[Cameroon Context: Myelin and Disease]
Demyelinating conditions (e.g. multiple sclerosis, though rarer in sub-Saharan Africa) drastically slow conduction, causing neurological deficits. In Cameroon, nutritional neuropathies (e.g. vitamin B12 deficiency) can also affect myelination. Understanding saltatory conduction explains why even partial demyelination has severe functional consequences.
\end{savaistu}

\courssub{The Synapse: Structure and Chemical Transmission}

When the action potential reaches the axon terminal, it must cross the \cle{synaptic cleft} to excite (or inhibit) the next cell. At a \cle{cholinergic synapse} (using acetylcholine, ACh), the process is as follows.

\begin{popfigure}
\begin{tikzpicture}[scale=0.9, every node/.style={font=\small}]
% Presynaptic terminal
\draw[popblue, thick, fill=popblueL!30] (0,0) rectangle (6,2.5);
% Mitochondria
\foreach \x/\y in {0.8/0.5, 1.5/1.8, 3.2/0.4, 4.5/1.9, 5.2/0.7}
    \draw[popdark, fill=poporange!50] (\x,\y) ellipse (0.25 and 0.15);
% Vesicles
\foreach \x/\y in {1.2/2.1, 1.8/2.15, 2.4/2.1, 3.0/2.15, 3.6/2.1, 4.2/2.15, 4.8/2.1}
    \draw[poppurple, fill=poppurple!50] (\x,\y) circle (0.12);
% Voltage-gated Ca2+ channels
\foreach \x in {1.5, 2.5, 3.5, 4.5}
    \draw[popgreen, thick] (\x,0) rectangle ++(0.3,0.3) node[midway, above=1pt] {Ca\textsuperscript{2+}};
% Synaptic cleft
\draw[popline, thick] (0,0) -- (6,0);
\draw[popline, thick] (0,-0.5) -- (6,-0.5);
\node[popmuted, align=center] at (3,-0.25) {Synaptic cleft (\approx 20--30 nm)};
% Postsynaptic membrane
\draw[popred, thick, fill=poppink!20] (0,-0.5) rectangle (6,-2.5);
% ACh receptors
\foreach \x in {1.2, 2.2, 3.2, 4.2, 5.2}
    \draw[popred, thick] (\x,-0.5) -- ++(0,-0.4) node[midway, left=2pt] {AChR};
% Acetylcholinesterase
\foreach \x in {1.7, 3.7, 4.7}
    \draw[popgold, thick] (\x,-0.5) -- ++(0,-0.3) node[midway, right=2pt] {AChE};
% Labels
\node[popblue, font=\bfseries] at (3,2.8) {Presynaptic terminal (knob)};
\node[popred, font=\bfseries] at (3,-3.0) {Postsynaptic membrane};
% Arrow for impulse arrival
\draw[popink, thick, ->] (3,3.2) -- (3,2.5) node[midway, right] {Action potential arrives};
% Diffusion arrows
\foreach \x in {1.5, 2.5, 3.5, 4.5}
    \draw[poppurple, ->, thick] (\x,2.0) -- (\x,0.2);
% Ca2+ entry arrows
\foreach \x in {1.65, 2.65, 3.65, 4.65}
    \draw[popgreen, ->, thick] (\x,0.3) -- (\x,-0.1);
\end{tikzpicture}
\legende{Structure of a cholinergic synapse. Voltage-gated Ca\textsuperscript{2+} channels open on depolarisation; synaptic vesicles fuse and release ACh; ACh diffuses across the cleft and binds to receptors (AChR) on the postsynaptic membrane; acetylcholinesterase (AChE) rapidly hydrolyses ACh.}
\end{popfigure}

\begin{methode}[Synaptic Transmission in Six Steps]
Learn this sequence for examinations:
\begin{enumerate}
\item \textbf{Arrival:} An action potential depolarises the presynaptic membrane.
\item \textbf{Ca\textsuperscript{2+} entry:} Voltage-gated \ce{Ca^2+} channels open; \ce{Ca^2+} rushes in down its steep electrochemical gradient.
\item \textbf{Vesicle fusion:} \ce{Ca^2+} binds to synaptotagmin, triggering SNARE-mediated fusion of synaptic vesicles with the presynaptic membrane.
\item \textbf{Release:} Acetylcholine (ACh) is expelled into the synaptic cleft by exocytosis.
\item \textbf{Diffusion \& binding:} ACh diffuses across the cleft ($\approx$ \SI{20}{nm}) and binds to nicotinic ACh receptors (ligand-gated \ce{Na+/K+} channels) on the postsynaptic membrane.
\item \textbf{Postsynaptic response:} Channels open; \ce{Na+} influx (and \ce{K+} efflux) causes a local depolarisation — an \cle{excitatory postsynaptic potential (EPSP)}. If the EPSP reaches threshold, an action potential is initiated in the postsynaptic neurone.
\end{enumerate}
\end{methode}

\begin{definition}[Excitatory and Inhibitory Postsynaptic Potentials]
\begin{itemize}
\item \textbf{EPSP:} A depolarising graded potential (e.g. via \ce{Na+} influx) that brings the membrane closer to threshold.
\item \textbf{IPSP:} A hyperpolarising graded potential (e.g. via \ce{Cl-} influx or \ce{K+} efflux) that moves the membrane away from threshold.
\end{itemize}
\end{definition}

\courssub{Synaptic Integration, Modulation and Pharmacology}

A single EPSP is usually too small to reach threshold. The postsynaptic neurone \emph{integrates} thousands of inputs.

\begin{propriete}[Summation]
\begin{itemize}
\item \textbf{Temporal summation:} Rapid successive EPSPs from the \emph{same} presynaptic neurone add up because the membrane time constant ($\approx$ 10–20 ms) allows them to overlap.
\item \textbf{Spatial summation:} Simultaneous EPSPs from \emph{different} presynaptic neurones at different synapses add up.
\end{itemize}
Inhibitory inputs (IPSPs) subtract from this sum. The axon hillock (initial segment) is the \cle{trigger zone} where the net potential is assessed.
\end{propriete}

\begin{attention}[One-Way Transmission]
Synaptic transmission is \emph{unidirectional} (presynaptic $\to$ postsynaptic) because:
\begin{enumerate}
\item Vesicles and release machinery are only in the presynaptic terminal.
\item Receptors are only on the postsynaptic membrane.
\item Acetylcholinesterase (AChE) in the cleft rapidly hydrolyses ACh, preventing back-diffusion and repeated stimulation.
\end{enumerate}
\end{attention}

\begin{aretenir}[Role of Acetylcholinesterase (AChE)]
AChE hydrolyses ACh into choline and acetate within \SI{<1}{ms}. This:
\begin{itemize}
\item Terminates the signal, allowing high-frequency firing.
\item Recycles choline (taken back up by the presynaptic terminal via a high-affinity transporter) for resynthesis of ACh.
\end{itemize}
Inhibition of AChE (e.g. by organophosphates) causes prolonged depolarisation, depolarisation block, and potentially fatal overstimulation of muscles and glands.
\end{aretenir}

\begin{exemplebox}[Drugs and Toxins at the Cholinergic Synapse — Exam Favourites]
\begin{center}
\begin{tabularx}{\linewidth}{|l|X|X|}\hline
\textbf{Agent} & \textbf{Site/Mechanism} & \textbf{Effect} \\ \hline
\textbf{Curare} & Competitive antagonist at nicotinic ACh receptors (postsynaptic) & Blocks EPSPs $\to$ flaccid paralysis (used historically as arrow poison in South America; principle of neuromuscular blockers in anaesthesia). \\ \hline
\textbf{Organophosphates} (e.g. malathion, sarin) & Irreversible inhibition of AChE & ACh accumulates $\to$ sustained depolarisation $\to$ depolarisation block, convulsions, respiratory failure. \\ \hline
\textbf{Neostigmine, physostigmine} & Reversible AChE inhibitors & Used to treat myasthenia gravis (autoimmune loss of ACh receptors) and glaucoma. \\ \hline
\textbf{Botulinum toxin} (Botox) & Prevents vesicle fusion (cleaves SNARE proteins) & Blocks ACh release $\to$ flaccid paralysis (therapeutic: dystonia, cosmetic). \\ \hline
\textbf{Nicotine} & Agonist at nicotinic receptors & Mimics ACh $\to$ stimulation then desensitisation (addiction via dopamine release in reward pathways). \\ \hline
\end{tabularx}
\end{center}
\end{exemplebox}

\begin{experience}[Investigating the Effect of Temperature on Conduction Velocity (Core Practical)]
\textbf{Aim:} Determine how temperature affects the speed of nerve impulse conduction.
\textbf{Apparatus:} Isolated sciatic nerve–gastrocnemius muscle preparation (frog), stimulating electrodes, recording electrodes, water bath at controlled temperatures (5°C, 15°C, 25°C, 35°C), kymograph or digital oscilloscope.
\textbf{Method:}
\begin{enumerate}
\item Stimulate the nerve at a fixed point; record the time for the muscle twitch to begin (latency).
\item Measure the distance between stimulation and recording points.
\item Calculate conduction velocity = distance / (latency – synaptic delay – excitation–contraction coupling time).
\item Repeat at each temperature; allow equilibration.
\end{enumerate}
\textbf{Expected result:} Conduction velocity increases with temperature (approx. doubling per 10°C rise, $Q_{10} \approx 2$) until thermal damage occurs.
\textbf{Examiner's tip:} Be ready to explain \emph{why} (increased ion channel kinetics, faster diffusion) and to identify sources of error (synaptic delay variability, muscle contraction time).
\end{experience}

\begin{exoresolu}[Worked Example: Refractory Period and Maximum Firing Frequency]
\textbf{Question:} A myelinated motor neurone has an absolute refractory period of \SI{1.2}{ms} and a relative refractory period of \SI{3.8}{ms}. What is the theoretical maximum firing frequency of this neurone? Explain why the actual maximum frequency in vivo is often lower.

\textbf{Solution:}
The absolute refractory period sets the hard limit on how quickly a second action potential can be initiated. The minimum interval between two action potentials is the absolute refractory period (ARP).
\[
f_{\text{max}} = \frac{1}{\text{ARP}} = \frac{1}{1.2 \times 10^{-3}\ \mathrm{s}} \approx 833\ \mathrm{Hz}.
\]
\textbf{Why lower in vivo?}
\begin{itemize}
\item During the relative refractory period, a stronger stimulus is needed; natural synaptic inputs may not provide sufficient depolarisation.
\item Synaptic integration (temporal/spatial summation) takes time.
\item Metabolic constraints (ATP for \ce{Na+/K+} pump) limit sustained high-frequency firing.
\item Many neurones exhibit \emph{spike-frequency adaptation} (slow \ce{K+} currents, \ce{Ca^2+}-activated \ce{K+} channels) that progressively raises the threshold.
\end{itemize}
Typical sustained firing rates for motor neurones are 50–100 Hz.
\end{exoresolu}

\begin{exercice}[Practice Question: Synaptic Transmission]
A student applies a toxin that specifically blocks voltage-gated \ce{Ca^2+} channels in the presynaptic terminal of a cholinergic synapse. The presynaptic neurone is stimulated with action potentials at 50 Hz.
\begin{enumerate}
\item Predict the effect on the postsynaptic membrane potential. Justify your answer.
\item Would the application of exogenous acetylcholine (via a micropipette) onto the postsynaptic membrane restore the response? Why or why not?
\item Name one natural toxin that acts via this mechanism.
\end{enumerate}
\end{exercice}

\begin{corrige}[Exercise: Synaptic Transmission]
\begin{enumerate}
\item \textbf{Effect:} No EPSPs will be generated; the postsynaptic membrane will remain at resting potential. \textbf{Justification:} Voltage-gated \ce{Ca^2+} channels are essential for vesicle fusion and ACh release. Blocking them prevents exocytosis of ACh, so no neurotransmitter reaches the cleft, no binding to receptors occurs, and no ion channels open on the postsynaptic membrane.
\item \textbf{Yes.} Exogenous ACh applied directly to the postsynaptic membrane will bind to the nicotinic receptors (which are unaffected by the presynaptic toxin) and open their ion channels, producing a depolarisation (EPSP). This bypasses the blocked release step.
\item \textbf{$\omega$-Conotoxin GVIA} (from the cone snail \emph{Conus geographus}) specifically blocks N-type voltage-gated \ce{Ca^2+} channels at presynaptic terminals. (Also accept: certain spider toxins, though $\omega$-conotoxin is the classic example.)
\end{enumerate}
\end{corrige}

\begin{aretenir}[Summary Checklist for Section 6]
\begin{itemize}
\item Resting potential ($\approx$ \SI{-70}{mV}) is maintained by the \ce{Na+/K+} pump and high resting \ce{K+} permeability.
\item Action potential: threshold $\to$ \ce{Na+} influx (depolarisation) $\to$ \ce{Na+} channel inactivation + \ce{K+} efflux (repolarisation) $\to$ hyperpolarisation $\to$ return to rest.
\item All-or-nothing law; refractory periods enforce one-way propagation and limit frequency.
\item Conduction speed: myelination (saltatory) $>$ diameter $>$ temperature.
\item Synapse: \ce{Ca^2+} entry $\to$ vesicle fusion $\to$ ACh release $\to$ diffusion $\to$ receptor binding $\to$ EPSP/IPSP.
\item AChE terminates signal; summation integrates inputs; transmission is one-way.
\item Key toxins: curare (receptor block), organophosphates (AChE inhibition), botulinum (release block), $\omega$-conotoxin (\ce{Ca^2+} channel block).
\end{itemize}
\end{aretenir}

This section provides the mechanistic foundation for the next topic: how the brain, spinal cord and autonomic nervous system organise these signals into coordinated responses.

\coursec{Reflex Actions, Brain, Spinal Cord and Autonomic Nervous System}

In the previous section we saw how an impulse travels along a neurone and crosses a synapse. But a single impulse travelling along a single neurone does not yet constitute a coordinated response. Real responses involve \emph{chains} of neurones organised into pathways. The simplest and fastest of these pathways is the \cle{reflex arc}, and the great coordinating centres that receive, integrate and dispatch impulses are the \cle{brain} and the \cle{spinal cord}. Finally, much of the body's internal housekeeping — heartbeat, digestion, glandular secretion — is controlled without our conscious will by the \cle{autonomic nervous system}. These three topics form the heart of this section.

\courssub{Reflex Actions}

\textbf{Observation.} Touch a hot cooking pot in the kitchen and your hand is withdrawn \emph{before} you even feel the pain. Dust blows into your eye and you blink instantly. A doctor taps the tendon below your kneecap and your lower leg kicks forward. In each case the response is so fast that there was no time to ``think''. Such responses are called reflex actions.

\begin{definition}[Reflex action]
A \cle{reflex action} is a \textbf{rapid, automatic, involuntary and stereotyped} response to a stimulus, which does not require conscious thought and is usually \textbf{protective} in function.
\end{definition}

\begin{aretenir}[Characteristics of reflex actions]
\begin{itemize}
\item \textbf{Rapid}: the pathway is short and involves few synapses (typically one or two).
\item \textbf{Automatic and involuntary}: they occur without conscious decision; you cannot easily prevent a blink.
\item \textbf{Stereotyped}: the same stimulus always produces the same response.
\item \textbf{Protective}: they remove the body (or a part of it) from harm, or maintain posture and balance.
\end{itemize}
\end{aretenir}

\textbf{The reflex arc.} The nervous pathway followed by the impulse during a reflex is called the \cle{reflex arc}. It always contains the same five components in the same order.

\begin{definition}[Reflex arc]
A \cle{reflex arc} is the nervous pathway along which an impulse travels from a receptor to an effector during a reflex action. Its components are:
\[
\text{receptor} \longrightarrow \text{sensory neurone} \longrightarrow \text{relay neurone (interneurone)} \longrightarrow \text{motor neurone} \longrightarrow \text{effector}
\]
\end{definition}

In a typical \textbf{spinal reflex} (e.g. the withdrawal reflex), the sequence is:

\begin{enumerate}
\item A \textbf{stimulus} (e.g. heat from a hot object) is detected by a \textbf{receptor} (pain/temperature receptors in the skin of the finger).
\item An impulse is generated and travels along the \textbf{sensory neurone}, whose cell body lies in the \textbf{dorsal root ganglion}, into the spinal cord through the \textbf{dorsal root}.
\item Inside the \textbf{grey matter} of the spinal cord, the sensory neurone synapses with a \textbf{relay neurone} (intermediate neurone / interneurone).
\item The relay neurone synapses with a \textbf{motor neurone}, which leaves the cord through the \textbf{ventral root}.
\item The motor neurone carries the impulse to an \textbf{effector} (the biceps muscle of the arm), which contracts and withdraws the hand.
\end{enumerate}

\begin{popfigure}
\begin{center}
\begin{tikzpicture}[font=\small, >=stealth,
box/.style={draw=popblue, thick, rounded corners, fill=popblueL, text width=2.8cm, align=center, minimum height=1.0cm},
greennode/.style={draw=popgreen, thick, rounded corners, fill=popgreenL, text width=2.8cm, align=center, minimum height=1.0cm}]
\node[box, align=center] (rec) at (0,0){\textbf{Receptor}\\ pain receptors in skin};
\node[box, align=center] (sen) at (4.5,0){\textbf{Sensory neurone}\\ dorsal root ganglion $\rightarrow$ dorsal root};
\node[greennode, align=center] (rel) at (9,0){\textbf{Relay neurone}\\ grey matter of spinal cord};
\node[box, align=center] (mot) at (9,-2.8){\textbf{Motor neurone}\\ ventral root};
\node[greennode, align=center] (eff) at (4.5,-2.8){\textbf{Effector}\\ biceps muscle contracts};
\draw[->, very thick, poporange] (rec) -- (sen);
\draw[->, very thick, poporange] (sen) -- (rel);
\draw[->, very thick, poporange] (rel) -- (mot);
\draw[->, very thick, poporange] (mot) -- (eff);
\draw[->, thick, poppurple, dashed] (rel.north) -- ++(0,1.0) node[above, text=poppurple, text width=3cm, align=center]{impulse also sent up to brain (conscious pain felt later)};
\end{tikzpicture}
\end{center}
\legende{Flow diagram of a simple spinal reflex arc (withdrawal reflex).}
\end{popfigure}

\begin{attention}[The brain is not needed first]
A very common examination error is to write that the impulse ``goes to the brain, which decides to withdraw the hand''. In a simple \textbf{spinal} reflex the impulse does \textbf{not} need to reach the brain before the response occurs: the relay takes place entirely in the spinal cord, which is why the response is so fast. An impulse \emph{is} also sent up to the brain, so you become \emph{aware} of the pain — but only \textbf{after} the hand has already been withdrawn.
\end{attention}

\textbf{Common examples of spinal reflexes.}
\begin{itemize}
\item \textbf{Knee jerk (patellar reflex)}: tapping the tendon below the kneecap stretches the thigh muscle (quadriceps); stretch receptors (muscle spindles) trigger a reflex contraction of the same muscle. This is a \textbf{monosynaptic reflex} (only one synapse between sensory and motor neurone). Used by doctors to test the integrity of the spinal cord segments L2--L4.
\item \textbf{Withdrawal reflex}: pulling the hand away from a hot or sharp object (protective). This is a \textbf{polysynaptic reflex} (involves one or more relay neurones).
\item \textbf{Blinking (corneal reflex)}: protects the eye from foreign bodies and bright light; involves the trigeminal (V) and facial (VII) cranial nerves — a \textbf{cranial reflex}.
\item \textbf{Pupil reflex}: constriction of the pupil in bright light, protecting the retina (studied in Section 3).
\item \textbf{Coughing and sneezing reflexes}: expel irritants from the respiratory tract.
\end{itemize}

\textbf{Conditioned reflexes.} All the reflexes above are \textbf{inborn (unconditioned)}. The Russian physiologist \textbf{Ivan Pavlov} showed that a reflex can be \emph{learned}: he rang a bell each time he fed dogs; after repeated association, the dogs salivated at the sound of the bell alone. A \cle{conditioned reflex} is a learned response in which a previously neutral stimulus (the bell) comes to trigger a reflex response (salivation) after repeated association with a natural stimulus (food). Conditioned reflexes involve the \textbf{cerebrum} (higher centres) and are the basis of much learning and habit formation.

\begin{aretenir}[Unconditioned vs. conditioned reflexes]
\begin{itemize}
\item \textbf{Unconditioned (simple) reflex}: inborn, inherited, does not require learning, involves spinal cord or lower brain centres (e.g. knee jerk, withdrawal).
\item \textbf{Conditioned reflex}: acquired during life, requires learning and memory, involves the cerebrum (e.g. salivating at the sound of a bell, a student opening a book when the bell rings).
\end{itemize}
\end{aretenir}

\courssub{The Brain}

The brain is the anterior, enlarged part of the central nervous system, protected by the skull (cranium) and by three membranes, the \textbf{meninges} (dura mater, arachnoid mater, pia mater), and cushioned by \textbf{cerebrospinal fluid (CSF)}. A sagittal section reveals several distinct regions, each with precise functions.

\begin{popfigure}
\begin{center}
\begin{tikzpicture}[font=\small, scale=1.0]
% Cerebrum
\fill[popblueL, draw=popblue, thick] (0,0.4) .. controls (0,2.6) and (2.2,3.6) .. (4.2,3.2)
 .. controls (6.0,2.8) and (6.6,1.6) .. (6.2,0.8)
 .. controls (5.8,0.2) and (4.6,0.1) .. (3.6,0.2)
 .. controls (2.0,0.3) and (0.6,-0.2) .. (0,0.4) -- cycle;
\node[text=popdark] at (3.1,2.1) {\textbf{Cerebrum}};
\node[text=popmuted, text width=3.4cm, align=center] at (3.1,1.5) {(cerebral hemispheres)};
% Cerebellum
\fill[popgreenL, draw=popgreen, thick] (5.0,-0.6) .. controls (5.2,-1.9) and (7.4,-1.9) .. (7.3,-0.5)
 .. controls (6.6,0.0) and (5.6,0.0) .. (5.0,-0.6) -- cycle;
\foreach \x in {5.4,5.8,6.2,6.6,7.0} \draw[popgreen] (\x,-1.55) -- (\x,-0.35);
\node[text=popdark] at (6.2,-2.3) {\textbf{Cerebellum}};
\draw[popgreen, thick] (6.2,-2.05) -- (6.2,-1.75);
% Medulla
\fill[popblueL, draw=popblue, thick] (3.6,-0.4) -- (4.6,-0.4) -- (4.4,-2.2) -- (3.8,-2.2) -- cycle;
\node[text=popdark, text width=2.6cm, align=center] at (1.6,-2.0) {\textbf{Medulla}\\ \textbf{oblongata}};
\draw[popblue, thick] (2.8,-1.9) -- (3.85,-1.5);
% Hypothalamus and pituitary
\fill[popgold, draw=poporange, thick] (4.0,0.35) circle (0.28);
\node[text=popdark] at (2.2,0.9) {\textbf{Hypothalamus}};
\draw[poporange, thick] (3.05,0.85) -- (3.85,0.5);
\fill[poporange] (4.05,-0.05) circle (0.14);
\node[text=popdark] at (5.6,0.55) {\textbf{Pituitary gland}};
\draw[poporange, thick] (5.0,0.45) -- (4.2,0.0);
% Thalamus
\fill[poppurple!30, draw=poppurple, thick] (4.9,0.9) circle (0.32);
\node[text=popdark] at (6.9,1.6) {\textbf{Thalamus}};
\draw[poppurple, thick] (6.35,1.5) -- (5.2,1.05);
% Pons
\fill[poppink!30, draw=poppink, thick] (3.8,-0.2) -- (4.6,-0.2) -- (4.5,-0.6) -- (3.9,-0.6) -- cycle;
\node[text=popdark] at (2.2,-0.4) {\textbf{Pons}};
\draw[poppink, thick] (3.0,-0.35) -- (3.85,-0.3);
% Spinal cord
\draw[popblue, thick] (3.8,-2.2) -- (3.8,-2.9);
\draw[popblue, thick] (4.4,-2.2) -- (4.4,-2.9);
\node[text=popdark] at (5.6,-2.7) {\textbf{Spinal cord}};
\end{tikzpicture}
\end{center}
\legende{Sagittal section of the human brain showing the main regions.}
\end{popfigure}

\textbf{Functions of the main regions.}

\begin{center}
\begin{tabularx}{\linewidth}{|l|X|}
\hline
\rowcolor{popblueL} \textbf{Region} & \textbf{Function(s)} \\
\hline
\textbf{Cerebrum} (cerebral hemispheres) & Centre of \textbf{intelligence, memory, reasoning, learning, imagination and will}; receives and interprets \textbf{sensory information} (sight, hearing, touch, taste, smell); initiates all \textbf{voluntary movements}. Its folded outer layer is the cerebral cortex (grey matter). \\
\hline
\textbf{Cerebellum} & Coordinates \textbf{balance, posture and fine muscular movement}; ensures movements are smooth and precise (e.g. walking, writing, riding a bicycle). Receives input from proprioceptors, eyes and vestibular apparatus. \\
\hline
\textbf{Medulla oblongata} & Controls \textbf{involuntary vital activities}: heart rate (cardiac centre), breathing rate (respiratory centre), blood pressure (vasomotor centre), swallowing, coughing, vomiting and peristalsis. \\
\hline
\textbf{Pons} & Relays signals between cerebrum, cerebellum and medulla; contains the \textbf{pneumotaxic centre} which helps regulate breathing rate. \\
\hline
\textbf{Hypothalamus} & Regulates \textbf{body temperature} (thermoregulation), \textbf{water balance (osmoregulation)}, hunger, thirst and sleep; controls the \textbf{pituitary gland}, linking the nervous and endocrine systems. \\
\hline
\textbf{Thalamus} & A \textbf{relay and sorting station}: directs sensory impulses (except smell) to the correct areas of the cerebrum. \\
\hline
\textbf{Pituitary gland} & The ``master gland'': secretes hormones controlling other endocrine glands (detailed in Sections 8 and 9). \\
\hline
\end{tabularx}
\end{center}

\begin{attention}[Exam tip]
Examiners frequently ask: ``State \textbf{one} function of each named part of the brain.'' Learn one \emph{precise} function per region — e.g. cerebellum = balance and coordination of movement; medulla = control of breathing and heart rate; hypothalamus = temperature and water balance regulation. Vague answers such as ``controls the body'' earn no marks.
\end{attention}

\begin{savaistu}[The blood-brain barrier]
The brain is protected not only by bone and meninges but also by the \textbf{blood-brain barrier}: endothelial cells of brain capillaries are tightly joined, and astrocytes wrap around them, preventing many substances (including most antibiotics and neurotransmitters) from entering the brain tissue. This maintains a stable chemical environment for neurones.
\end{savaistu}

\courssub{The Spinal Cord}

The spinal cord is a cylinder of nervous tissue running from the medulla oblongata down the vertebral column (to about L1/L2 in adults), protected by the vertebrae, the meninges and cerebrospinal fluid. In transverse section it shows:

\begin{itemize}
\item An inner, butterfly-shaped (or H-shaped) region of \textbf{grey matter}, containing the \textbf{cell bodies} of motor and relay neurones, dendrites, and the synapses of reflex arcs. The grey matter has \textbf{dorsal horns} (sensory) and \textbf{ventral horns} (motor).
\item An outer region of \textbf{white matter}, containing \textbf{myelinated axons} running up and down the cord in tracts (ascending sensory tracts, descending motor tracts).
\item A central canal containing cerebrospinal fluid.
\item \textbf{Dorsal roots} (carrying sensory impulses \emph{into} the cord, with a swelling — the \textbf{dorsal root ganglion} — containing sensory neurone cell bodies) and \textbf{ventral roots} (carrying motor impulses \emph{out} of the cord). Dorsal and ventral roots unite to form a \textbf{spinal nerve} (mixed nerve).
\end{itemize}

\begin{popfigure}
\begin{center}
\begin{tikzpicture}[font=\small, scale=1.2]
% White matter outline
\draw[popblue, thick, fill=popblueL] (0,0) ellipse (2.5 and 1.8);
% Grey matter (butterfly shape) - LEFT
\draw[popdark, thick, fill=gray!20] 
(-1.2,0.2) .. controls (-1.5,0.5) and (-1.5,1.2) .. (-0.8,1.4)
.. controls (-0.2,1.5) and (0.2,1.2) .. (0.0,0.8)
.. controls (-0.2,0.4) and (-0.5,0.0) .. (-0.6,-0.3)
.. controls (-0.8,-0.6) and (-1.0,-0.8) .. (-1.2,-0.5)
.. controls (-1.4,-0.2) and (-1.4,0.0) .. (-1.2,0.2) -- cycle;
% Grey matter (butterfly shape) - RIGHT
\draw[popdark, thick, fill=gray!20] 
(1.2,0.2) .. controls (1.5,0.5) and (1.5,1.2) .. (0.8,1.4)
.. controls (0.2,1.5) and (-0.2,1.2) .. (0.0,0.8)
.. controls (0.2,0.4) and (0.5,0.0) .. (0.6,-0.3)
.. controls (0.8,-0.6) and (1.0,-0.8) .. (1.2,-0.5)
.. controls (1.4,-0.2) and (1.4,0.0) .. (1.2,0.2) -- cycle;
% Central canal
\fill[popblueL] (0,0) circle (0.1);
% Dorsal root ganglion and roots
\draw[poporange, very thick] (-2.5,1.2) -- (-1.5,0.8);
\fill[poporange] (-2.7,1.3) circle (0.15);
\node[poporange, text width=2cm, align=left] at (-3.5,1.3) {Dorsal root ganglion (sensory cell bodies)};
\draw[->, poporange, thick] (-2.5,1.2) -- (-1.5,0.8) node[midway, above left, font=\tiny] {sensory in};
% Ventral root
\draw[popgreen, very thick] (1.5,-0.8) -- (2.5,-1.2);
\draw[->, popgreen, thick] (1.5,-0.8) -- (2.5,-1.2) node[midway, below right, font=\tiny] {motor out};
\node[popgreen, text width=2cm, align=right] at (3.5,-1.2) {Ventral root (motor axons)};
% Labels
\node[text=popdark] at (0,1.6) {\textbf{Dorsal}};
\node[text=popdark] at (0,-1.6) {\textbf{Ventral}};
\node[text=popdark] at (-2.2,0) {\textbf{Grey matter}};
\node[text=popdark] at (2.2,0) {\textbf{Grey matter}};
\node[text=popdark] at (0,0) {\textbf{Central canal}};
\node[text=popdark] at (-2.5,-1.0) {\textbf{White matter}};
% Reflex arc path
\draw[->, very thick, poppurple, dashed] (-2.7,1.3) .. controls (-2.0,1.0) and (-1.0,0.5) .. (-0.6,0.2) node[midway, left, font=\tiny, text=poppurple] {sensory};
\draw[->, very thick, poppurple, dashed] (-0.6,0.2) -- (0,0.2) node[midway, above, font=\tiny, text=poppurple] {relay};
\draw[->, very thick, poppurple, dashed] (0,0.2) .. controls (0.6,0.2) and (1.0,-0.2) .. (1.5,-0.8) node[midway, right, font=\tiny, text=poppurple] {motor};
\end{tikzpicture}
\end{center}
\legende{Transverse section of the spinal cord showing grey matter, white matter, central canal, dorsal/ventral roots, and the path of a reflex arc (dashed purple).}
\end{popfigure}

\begin{aretenir}[Functions of the spinal cord]
\begin{enumerate}
\item It is the \textbf{reflex centre} for spinal reflexes (the relay neurones of reflex arcs synapse in its grey matter).
\item It is a \textbf{pathway} conducting sensory impulses \emph{up} to the brain (ascending tracts) and motor impulses \emph{down} from the brain to the effectors (descending tracts).
\end{enumerate}
\end{aretenir}

\begin{attention}[Spinal cord injury level matters]
Damage to the spinal cord in the neck (cervical region) affects all four limbs (tetraplegia/quadriplegia) because both ascending and descending tracts for the arms and legs are interrupted. Damage lower down (thoracic/lumbar) affects only the legs (paraplegia). Local spinal reflexes \emph{below} the injury may persist (e.g. knee jerk) because the reflex arc remains intact, but they are no longer under brain control.
\end{attention}

\courssub{Voluntary and Involuntary Actions; the Autonomic Nervous System}

A \textbf{voluntary action} is initiated consciously by the cerebrum (e.g. picking up a pen, walking to school, speaking). An \textbf{involuntary action} occurs without conscious control. Involuntary actions include reflexes and the continuous automatic regulation of internal organs — heartbeat, gut movements, pupil size, glandular secretion. This internal regulation is the work of the \cle{autonomic nervous system} (ANS).

The ANS consists of motor neurones running from the brain and spinal cord to internal effectors (smooth muscle, cardiac muscle and glands). It has two divisions with \textbf{antagonistic (opposing)} effects.

\begin{propriete}[Antagonism of the two autonomic divisions]
The \cle{sympathetic} and \cle{parasympathetic} divisions of the autonomic nervous system act \textbf{antagonistically} on the same organs: the sympathetic division prepares the body for \textbf{emergency and activity} (``fight or flight''), while the parasympathetic division promotes \textbf{rest, digestion and energy conservation} (``rest and digest''). The balance between the two keeps internal conditions stable.
\end{propriete}

\begin{center}
\begin{tabularx}{\linewidth}{|l|X|X|}
\hline
\rowcolor{popblueL} \textbf{Organ / function} & \textbf{Sympathetic} (emergency) & \textbf{Parasympathetic} (rest) \\
\hline
Heart rate & Increases (positive chronotropic) & Decreases (negative chronotropic) \\
\hline
Force of heart contraction & Increases & Decreases slightly \\
\hline
Pupil of the eye & Dilates (mydriasis) & Constricts (miosis) \\
\hline
Gut (peristalsis, digestion) & Slowed / inhibited & Stimulated \\
\hline
Salivary glands & Secretion reduced (thick, mucus-rich; dry mouth) & Secretion increased (watery) \\
\hline
Bronchi of lungs & Dilate (bronchodilation; more air) & Constrict (bronchoconstriction) \\
\hline
Blood vessels of skin/gut & Constrict (blood diverted to muscles) & Dilate (or no direct effect) \\
\hline
Blood vessels of skeletal muscle & Dilate (via adrenaline on $\beta_2$ receptors) & No direct effect \\
\hline
Adrenal medulla & Stimulated to release adrenaline (epinephrine) & No effect \\
\hline
Bladder / rectum & Emptying inhibited (sphincters contract) & Emptying promoted (detrusor contracts, sphincters relax) \\
\hline
Liver & Glycogenolysis stimulated (glucose release) & Glycogenesis promoted \\
\hline
\end{tabularx}
\end{center}
\legende{Comparison of the antagonistic effects of the sympathetic and parasympathetic divisions.}

\begin{savaistu}[Why your mouth goes dry before an oral exam]
When you are frightened — for example just before speaking in front of the class — the sympathetic division dominates: your heart beats faster, your pupils dilate, blood is diverted to your muscles and saliva production falls, giving you a dry mouth. This is the same ``fight or flight'' preparation that helped our ancestors escape danger on the savanna. Afterwards, the parasympathetic division restores calm: the heart slows and digestion resumes.
\end{savaistu}

\begin{aretenir}[Anatomical differences between the two divisions]
\begin{itemize}
\item \textbf{Sympathetic}: preganglionic neurones have cell bodies in the \textbf{thoracic and lumbar} regions of the spinal cord (thoracolumbar outflow); short preganglionic axons, long postganglionic axons; ganglia form a chain alongside the vertebral column.
\item \textbf{Parasympathetic}: preganglionic neurones have cell bodies in the \textbf{brainstem} (cranial nerves III, VII, IX, X) and \textbf{sacral} spinal cord (craniosacral outflow); long preganglionic axons, short postganglionic axons; ganglia lie near or within the target organs.
\item \textbf{Neurotransmitters}: both divisions release \textbf{acetylcholine (ACh)} at the preganglionic synapse (nicotinic receptors). Parasympathetic postganglionic fibres release \textbf{ACh} (muscarinic receptors). Sympathetic postganglionic fibres mostly release \textbf{noradrenaline (norepinephrine)} (adrenergic receptors), except those to sweat glands and some blood vessels in muscle which release ACh.
\end{itemize}
\end{aretenir}

\textbf{The protective role of reflexes in daily life.} Reflexes protect us constantly and without effort: blinking shields the cornea from dust; coughing and sneezing expel irritants from the airways; the withdrawal reflex removes limbs from hot or sharp objects; postural reflexes (coordinated by the cerebellum and spinal cord) keep us upright; the pupil reflex prevents damage to the retina by bright light. Because they bypass conscious decision-making, they save the precious fractions of a second that often make the difference between injury and safety.

\begin{exoresolu}[Worked exercise: analysing a reflex]
A student accidentally touches a thorn while picking hibiscus flowers and immediately pulls her hand away. (a) Name this type of response and give two of its characteristics. (b) List, in order, the five components of the nervous pathway involved. (c) Explain why the hand is withdrawn \emph{before} the student feels the pain.
\tcblower
\textbf{(a)} This is a \textbf{reflex action} (a spinal withdrawal reflex). Any two characteristics: it is \textbf{rapid}, \textbf{automatic/involuntary} (no conscious thought), \textbf{stereotyped}, and \textbf{protective}.

\textbf{(b)} The reflex arc:
\[
\text{receptor (pain receptor in skin)} \rightarrow \text{sensory neurone} \rightarrow \text{relay neurone (grey matter of spinal cord)} \rightarrow \text{motor neurone} \rightarrow \text{effector (arm muscle)}
\]

\textbf{(c)} The impulse passes directly from the sensory neurone to the relay neurone and then to the motor neurone \emph{within the spinal cord}; this pathway is short and involves few synapses, so the muscle contracts very quickly. The impulse that travels \emph{up the spinal cord to the brain} — producing the conscious sensation of pain — follows a longer route and arrives later. Hence the hand is withdrawn before the pain is consciously felt.
\end{exoresolu}

\begin{exoresolu}[Worked exercise: autonomic control]
A man is suddenly confronted by an aggressive dog. Describe the changes in his body brought about by the sympathetic division and explain how they prepare him for ``fight or flight''.
\tcblower
\textbf{Changes (sympathetic dominance):}
\begin{itemize}
\item Heart rate and force of contraction increase $\rightarrow$ more oxygen and glucose delivered to muscles.
\item Bronchi dilate $\rightarrow$ increased air flow and gas exchange.
\item Blood vessels in skin and gut constrict; vessels in skeletal muscle dilate $\rightarrow$ blood redirected to muscles.
\item Liver glycogen broken down to glucose $\rightarrow$ blood glucose rises for energy.
\item Pupils dilate $\rightarrow$ improved vision in low light / wider field of view.
\item Gut peristalsis and digestion inhibited $\rightarrow$ energy diverted from non-essential processes.
\item Adrenal medulla releases adrenaline $\rightarrow$ prolongs and amplifies sympathetic effects.
\item Sweating increases $\rightarrow$ cooling during anticipated exertion.
\end{itemize}
\textbf{Preparation:} These changes maximise the supply of oxygen and glucose to the brain and muscles, enhance alertness and vision, and suppress non-urgent functions, enabling the man to either fight the dog or run away.
\end{exoresolu}

\begin{exercice}[Practice questions]
\begin{enumerate}
\item Define the term \emph{reflex arc} and distinguish between a simple (unconditioned) reflex and a conditioned reflex, giving one example of each.
\item State one function of each of the following: cerebellum; medulla oblongata; hypothalamus; cerebrum; pons.
\item Draw a labelled transverse section of the spinal cord and indicate the route taken by an impulse during a spinal reflex.
\item Copy and complete: during sympathetic stimulation, the heart rate \ldots, the pupil \ldots, and peristalsis is \ldots; state the opposite effects produced by the parasympathetic division.
\item Explain why damage to the spinal cord in the neck region can abolish both sensation and movement in the legs.
\item A person eats a large meal. Which division of the autonomic nervous system is dominant? List three effects on the digestive system.
\item Compare the anatomical arrangement (outflow, ganglion position, axon lengths) of the sympathetic and parasympathetic divisions.
\end{enumerate}
\end{exercice}

\begin{corrige}[Exercise 1]
\textbf{1.} A reflex arc is the nervous pathway along which an impulse travels from receptor to effector during a reflex: receptor $\rightarrow$ sensory neurone $\rightarrow$ relay neurone $\rightarrow$ motor neurone $\rightarrow$ effector. A simple reflex is inborn and automatic (e.g. knee jerk); a conditioned reflex is learned by association (e.g. Pavlov's dogs salivating at a bell).

\textbf{2.} Cerebellum: coordination of balance and muscular movement. Medulla oblongata: control of breathing rate and heart rate. Hypothalamus: regulation of body temperature and water balance (and control of the pituitary). Cerebrum: intelligence, memory, voluntary movement and interpretation of sensation. Pons: relay between brain regions and contains pneumotaxic centre for breathing.

\textbf{3.} The drawing should show: outer white matter, inner butterfly-shaped grey matter (dorsal and ventral horns), central canal, dorsal root with dorsal root ganglion (sensory neurone entering), ventral root (motor neurone leaving). Route: receptor $\rightarrow$ sensory neurone $\rightarrow$ dorsal root $\rightarrow$ relay neurone in grey matter $\rightarrow$ motor neurone $\rightarrow$ ventral root $\rightarrow$ effector.

\textbf{4.} Heart rate \emph{increases}, the pupil \emph{dilates}, peristalsis is \emph{inhibited}. Parasympathetic: heart rate decreases, pupil constricts, peristalsis is stimulated.

\textbf{5.} The spinal cord is the pathway between the brain and the body below the neck. If it is severed, sensory impulses from the legs can no longer ascend to the brain (loss of sensation) and motor impulses from the brain can no longer descend to the leg muscles (paralysis), although local spinal reflexes below the injury may persist.

\textbf{6.} Parasympathetic division dominates (``rest and digest''). Effects: increased salivation, increased gastric secretion, increased peristalsis, relaxation of sphincters, increased blood flow to gut.

\textbf{7.} Sympathetic: thoracolumbar outflow; ganglia in chain alongside vertebral column; short preganglionic, long postganglionic axons; postganglionic neurotransmitter mostly noradrenaline. Parasympathetic: craniosacral outflow; ganglia near/within target organs; long preganglionic, short postganglionic axons; postganglionic neurotransmitter acetylcholine.
\end{corrige}

\begin{aretenir}[Key points of this section]
\begin{itemize}
\item A reflex action is rapid, automatic, involuntary, stereotyped and protective; its pathway is the reflex arc (receptor $\rightarrow$ sensory $\rightarrow$ relay $\rightarrow$ motor $\rightarrow$ effector).
\item In a spinal reflex, integration occurs in the spinal cord — the brain is informed afterwards.
\item The brain's regions have distinct functions: cerebrum (thought, voluntary action, sensation), cerebellum (balance and coordination), medulla (vital involuntary centres), pons (relay, breathing), hypothalamus (homeostasis and pituitary control), thalamus (sensory relay).
\item The spinal cord is both a reflex centre and a two-way conduction pathway (ascending sensory, descending motor tracts).
\item The autonomic nervous system controls internal organs through two antagonistic divisions: sympathetic (thoracolumbar, ``fight or flight'', noradrenaline) and parasympathetic (craniosacral, ``rest and digest'', acetylcholine).
\end{itemize}
\end{aretenir}

\coursec{The Endocrine System and Mechanism of Hormone Action}

In the previous section we explored how the nervous system provides rapid, short-lived coordination through electrical impulses and neurotransmitters. However, not all responses need to be instantaneous. Many vital processes — growth, metabolism, reproduction, water balance — require slower, longer-lasting, and more widespread regulation. This is the domain of the \cle{endocrine system}: a network of ductless glands that secrete chemical messengers called \cle{hormones} directly into the bloodstream. Together, the nervous and endocrine systems form the two pillars of coordination in animals, often working hand in hand.

\courssub{Endocrine Glands, Hormones and Target Organs}

\begin{definition}[Hormone and Target Organ]
A \textbf{hormone} is a chemical substance, produced in very small quantities by an \textbf{endocrine gland} (ductless gland), secreted directly into the blood plasma, transported by the bloodstream to a specific \textbf{target organ} (or target cells) where it exerts its effect, often at very low concentrations. The target cells possess specific \cle{receptors} for that hormone.
\end{definition}

\begin{propriete}[General Properties of Hormones]
\begin{itemize}
    \item \textbf{Secreted in minute amounts:} Concentrations in blood are typically $10^{-12}$ to $10^{-9}\ \mathrm{mol\,dm^{-3}}$.
    \item \textbf{Transported by blood:} They travel from gland to target via the circulatory system.
    \item \textbf{Specificity:} Each hormone acts only on cells bearing its specific receptor (lock-and-key principle).
    \item \textbf{Action at a distance:} Unlike paracrine/autocrine signals, endocrine signals act on distant targets.
    \item \textbf{Degradation and excretion:} Hormones are inactivated mainly in the liver and kidneys; their half-lives range from minutes (adrenaline) to days (thyroxine).
\end{itemize}
\end{propriete}

\begin{attention}[Endocrine vs Exocrine Glands — Do Not Confuse!]
\begin{itemize}
    \item \textbf{Endocrine glands} are ductless; they secrete hormones directly into the blood capillaries penetrating the gland. Examples: pituitary, thyroid, adrenal glands.
    \item \textbf{Exocrine glands} secrete their products (enzymes, mucus, sweat, saliva) into \cle{ducts} that carry them to a body surface or into a lumen (e.g. digestive tract). Examples: salivary glands, sweat glands, pancreas (acinar cells).
    \item \textbf{The pancreas} is a \cle{mixed gland}: its \cle{acinar cells} (exocrine) secrete digestive enzymes into the pancreatic duct, while its \cle{islets of Langerhans} (endocrine) secrete insulin and glucagon directly into blood. In examination questions, always specify which part you are discussing.
\end{itemize}
\end{attention}

\courssub{Location of the Main Endocrine Glands}

The major endocrine glands in the human body are positioned strategically. Figure~\ref{fig:endocrine_glands} summarises their locations.

\begin{popfigure}
\begin{tikzpicture}[scale=0.9, font=\small]
% Body outline
\draw[popdark, line width=1.5pt] (0,0) ellipse (2.5 and 6);
\draw[popdark, line width=1.5pt] (0,6) ellipse (1.2 and 1.2); % head
% Pituitary (hypophysis)
\fill[popblue] (0,6.8) circle (0.15);
\node[popblue, text width=1.5cm, align=center] at (0,7.2) {Pituitary};
% Thyroid
\draw[poporange, line width=1.2pt] (-0.6,4.5) -- (-0.6,3.5) -- (0.6,3.5) -- (0.6,4.5) -- cycle;
\draw[poporange, line width=1.2pt] (-0.4,4.5) -- (-0.4,3.7) -- (0.4,3.7) -- (0.4,4.5) -- cycle;
\node[poporange, text width=1.5cm, align=center] at (0,5.2) {Thyroid};
% Parathyroids (on posterior of thyroid, shown as small dots)
\foreach \x in {-0.5,0.5} {
    \fill[poppurple] (\x,3.8) circle (0.06);
}
\node[poppurple, text width=1.5cm, align=center] at (0,3.3) {Parathyroids};
% Thymus (in children, upper thorax)
\draw[popgreen, line width=1pt] (-0.8,3.5) .. controls (-1,3) and (-1,2) .. (-0.8,1.5);
\draw[popgreen, line width=1pt] (0.8,3.5) .. controls (1,3) and (1,2) .. (0.8,1.5);
\node[popgreen, text width=1.5cm, align=center] at (0,2.8) {Thymus};
% Adrenals (on kidneys)
\fill[poppink] (-2.0,0.5) ellipse (0.25 and 0.4);
\fill[poppink] (2.0,0.5) ellipse (0.25 and 0.4);
\node[poppink, text width=1.5cm, align=center] at (-2.0,1.2) {Adrenal};
\node[poppink, text width=1.5cm, align=center] at (2.0,1.2) {Adrenal};
% Pancreas (islets)
\draw[popgold, line width=1.2pt] (-1.5,0) .. controls (-1.5,-0.5) and (-0.5,-1) .. (0,-1) .. controls (0.5,-1) and (1.5,-0.5) .. (1.5,0);
\node[popgold, text width=1.5cm, align=center] at (0,-0.3) {Pancreas (Islets)};
% Ovaries / Testes
\fill[popblueL] (-1.2,-2.5) circle (0.2);
\fill[popblueL] (1.2,-2.5) circle (0.2);
\node[popblueL, text width=1.5cm, align=center] at (0,-3.2) {Ovaries / Testes};
% Pineal (in brain)
\fill[popmuted] (0,7.0) circle (0.08);
\node[popmuted, font=\tiny, text width=1cm, align=center] at (0.6,7.0) {Pineal};
% Labels on left side
\node[align=left, font=\small, popdark] at (-4.5,6.5) {\textbf{Brain:}\\Pineal, Pituitary};
\node[align=left, font=\small, popdark] at (-4.5,4.5) {\textbf{Neck:}\\Thyroid, Parathyroids};
\node[align=left, font=\small, popdark] at (-4.5,2.5) {\textbf{Thorax:}\\Thymus};
\node[align=left, font=\small, popdark] at (-4.5,0.5) {\textbf{Abdomen:}\\Adrenals, Pancreas};
\node[align=left, font=\small, popdark] at (-4.5,-2.5) {\textbf{Pelvis:}\\Ovaries / Testes};
\end{tikzpicture}
\legende{Location of the major human endocrine glands. The pancreas is shown for its endocrine portion (islets of Langerhans).}
\label{fig:endocrine_glands}
\end{popfigure}

\begin{aretenir}[Key Endocrine Glands and Their Main Hormones]
\begin{center}
\begin{tabularx}{\linewidth}{|l|X|X|}\hline
\rowcolor{popblueL}
\textbf{Gland} & \textbf{Main Hormones} & \textbf{Primary Functions} \\ \hline
Pituitary (Hypophysis) & GH, TSH, ACTH, LH, FSH, Prolactin, ADH, Oxytocin & Master gland; regulates other glands, growth, water balance, reproduction \\ \hline
Thyroid & Thyroxine (T4), Triiodothyronine (T3), Calcitonin & Basal metabolic rate, growth, calcium homeostasis \\ \hline
Parathyroids & Parathyroid Hormone (PTH) & Raises blood $\mathrm{Ca^{2+}}$ \\ \hline
Adrenal Cortex & Cortisol, Aldosterone, Androgens & Stress response, $\mathrm{Na^+}$/$\mathrm{K^+}$ balance, secondary sex characteristics \\ \hline
Adrenal Medulla & Adrenaline (Epinephrine), Noradrenaline & Fight-or-flight, rapid metabolic changes \\ \hline
Pancreas (Islets) & Insulin ($\beta$-cells), Glucagon ($\alpha$-cells) & Blood glucose homeostasis \\ \hline
Ovaries & Oestrogen, Progesterone & Female reproductive cycle, secondary sexual characteristics \\ \hline
Testes & Testosterone & Male reproductive function, secondary sexual characteristics \\ \hline
\end{tabularx}
\end{center}
\end{aretenir}

\courssub{Chemical Nature of Hormones}

Hormones are classified by their chemical structure, which determines their solubility, transport in blood, receptor location, and mechanism of action.

\begin{propriete}[Classification by Chemical Nature]
\begin{enumerate}
    \item \textbf{Peptide / Protein hormones} (hydrophilic): Insulin, glucagon, ADH, GH, TSH, LH, FSH. Synthesised as preprohormones $\rightarrow$ prohormones $\rightarrow$ active hormones; stored in secretory vesicles; released by exocytosis.
    \item \textbf{Steroid hormones} (lipophilic): Cortisol, aldosterone, testosterone, oestrogen, progesterone. Derived from cholesterol; not stored in vesicles (synthesised on demand); diffuse through plasma membrane; transported in blood bound to carrier proteins.
    \item \textbf{Amine hormones}: Derived from tyrosine. Thyroid hormones (T3, T4) behave like steroids (lipophilic, nuclear receptors). Catecholamines (adrenaline, noradrenaline) behave like peptides (hydrophilic, membrane receptors).
\end{enumerate}
\end{propriete}

\courssub{Mechanism of Hormone Action}

The site of the receptor (cell membrane vs intracellular) dictates the signalling cascade.

\courssub{Protein Hormones: Membrane Receptors and Second Messengers}

Protein hormones cannot cross the phospholipid bilayer. They bind to specific \cle{transmembrane receptors} on the target cell surface. This binding triggers a conformational change that activates an intracellular signalling cascade, often involving a \cle{second messenger}.

\begin{methode}[Mechanism of Protein Hormone Action (cAMP Pathway)]
\begin{enumerate}
    \item \textbf{First messenger:} Hormone (e.g. adrenaline, glucagon) binds to its specific \cle{G-protein-coupled receptor (GPCR)} on the outer surface of the plasma membrane.
    \item \textbf{Receptor activation:} The receptor changes shape, activating the associated \cle{G-protein} ($G_s$) by exchanging GDP for GTP on its $\alpha$-subunit.
    \item \textbf{Effector activation:} The $G_\alpha$-GTP subunit dissociates and activates \cle{adenylyl cyclase} (membrane enzyme).
    \item \textbf{Second messenger formation:} Adenylyl cyclase converts ATP $\rightarrow$ \cle{cyclic AMP (cAMP)} + PP$_i$.
    \item \textbf{Amplification:} One hormone-receptor complex can generate many cAMP molecules.
    \item \textbf{Protein kinase activation:} cAMP activates \cle{Protein Kinase A (PKA)} by binding its regulatory subunits, releasing catalytic subunits.
    \item \textbf{Phosphorylation cascade:} PKA phosphorylates specific target proteins (enzymes, ion channels, transcription factors), altering their activity.
    \item \textbf{Termination:} \cle{Phosphodiesterase} hydrolyses cAMP $\rightarrow$ AMP; phosphatases remove phosphate groups; hormone dissociates.
\end{enumerate}
\end{methode}

\begin{exemplebox}[Adrenaline and Glycogenolysis in Liver]
Adrenaline binds to $\beta$-adrenergic receptors on hepatocytes $\rightarrow$ cAMP $\uparrow$ $\rightarrow$ PKA activated $\rightarrow$ phosphorylates (activates) \cle{phosphorylase kinase} $\rightarrow$ phosphorylates (activates) \cle{glycogen phosphorylase} $\rightarrow$ glycogen breakdown $\rightarrow$ glucose released into blood. Simultaneously, glycogen synthase is phosphorylated (inactivated). This \cle{amplification cascade} ensures a rapid, massive glucose output from a tiny hormone signal.
\end{exemplebox}

\begin{attention}[Other Second Messenger Pathways]
Besides cAMP, two other major pathways exist:
\begin{itemize}
    \item \textbf{IP$_3$/DAG pathway:} Hormone $\rightarrow$ GPCR ($G_q$) $\rightarrow$ Phospholipase C $\rightarrow$ PIP$_2$ $\rightarrow$ IP$_3$ + DAG. IP$_3$ releases $\mathrm{Ca^{2+}}$ from ER; DAG activates Protein Kinase C (PKC). Example: ADH (V1 receptor), Oxytocin.
    \item \textbf{Receptor Tyrosine Kinase (RTK) pathway:} Hormone (e.g. Insulin) binds $\rightarrow$ receptor dimerisation $\rightarrow$ autophosphorylation $\rightarrow$ recruits adaptor proteins (IRS) $\rightarrow$ PI3K/Akt and MAPK cascades. No classic small-molecule second messenger.
\end{itemize}
\end{attention}

\courssub{Steroid Hormones: Intracellular Receptors and Gene Transcription}

Steroid hormones are lipophilic and diffuse freely through the plasma membrane. Their receptors are located in the \cle{cytoplasm} or \cle{nucleus}.

\begin{methode}[Mechanism of Steroid Hormone Action]
\begin{enumerate}
    \item \textbf{Diffusion:} Hormone (e.g. cortisol, testosterone, oestrogen) diffuses through the plasma membrane.
    \item \textbf{Receptor binding:} In cytoplasm, receptor is often bound to \cle{heat-shock proteins (HSP90)} keeping it inactive. Hormone binding causes HSP dissociation.
    \item \textbf{Nuclear translocation:} Hormone-receptor complex dimerises and translocates into the nucleus (if not already nuclear).
    \item \textbf{DNA binding:} Complex binds to specific \cle{Hormone Response Elements (HREs)} — short DNA sequences in promoter/enhancer regions of target genes.
    \item \textbf{Transcription regulation:} Recruits co-activators / co-repressors and RNA polymerase $\rightarrow$ increased or decreased \cle{mRNA synthesis}.
    \item \textbf{Translation:} New mRNA $\rightarrow$ synthesis of specific proteins (enzymes, structural proteins, receptors) $\rightarrow$ long-term cellular response.
\end{enumerate}
\end{methode}

\begin{attention}[Time Course Difference]
Protein hormone actions (via second messengers) take \textbf{seconds to minutes} (modifying existing proteins). Steroid hormone actions (via gene transcription) take \textbf{hours to days} (synthesising new proteins). This is a classic examination distinction.
\end{attention}

\begin{popfigure}
\begin{tikzpicture}[scale=0.85, font=\small]
% LEFT: Protein hormone (membrane receptor)
\begin{scope}[xshift=-5cm]
% Cell membrane
\draw[popdark, line width=2pt] (-2.5,-0.2) -- (2.5,-0.2);
\draw[popdark, line width=2pt] (-2.5,0.2) -- (2.5,0.2);
\foreach \x in {-2.4,-2.2,...,2.4} {
    \fill[popmuted] (\x,0) circle (0.05);
}
\node[popdark, text width=2cm, align=center] at (0,0.6) {\textbf{Plasma Membrane}};
% Receptor
\draw[popblue, line width=1.5pt] (0,-0.2) -- (0,-1.2);
\draw[popblue, fill=popblue!20] (0,-1.2) circle (0.25);
\node[popblue, text width=1.5cm, align=center] at (0,-1.6) {GPCR};
% G-protein
\draw[poporange, line width=1.2pt] (-0.4,-1.2) -- (-0.4,-1.8);
\draw[poporange, fill=poporange!20] (-0.4,-1.8) circle (0.15);
\node[poporange, text width=1.5cm, align=center] at (-0.4,-2.2) {$G_s$};
% Adenylyl cyclase
\draw[popgreen, line width=1.2pt] (0.6,-1.2) -- (0.6,-1.8);
\draw[popgreen, fill=popgreen!20] (0.6,-1.8) circle (0.15);
\node[popgreen, text width=1.5cm, align=center] at (0.6,-2.2) {Adenylyl Cyclase};
% Hormone
\draw[poppink, ->, line width=1pt] (0,1.5) -- (0,0.3);
\node[poppink, text width=1.5cm, align=center] at (0,1.8) {Hormone (1st messenger)};
% cAMP
\draw[poppurple, ->, line width=1pt] (0.6,-1.8) -- (1.5,-1.8);
\node[poppurple, text width=1.5cm, align=center] at (1.5,-1.5) {cAMP (2nd messenger)};
% PKA
\draw[poppurple, ->, line width=1pt] (1.5,-1.8) -- (1.5,-2.8);
\draw[poppurple, fill=poppurple!20] (1.5,-2.8) circle (0.15);
\node[poppurple, text width=1.5cm, align=center] at (1.5,-3.2) {PKA};
% Phosphorylation
\draw[poppurple, ->, line width=1pt] (1.5,-2.8) -- (2.5,-2.8);
\node[poppurple, text width=2cm, align=center] at (2.8,-2.5) {Phosphorylation of target proteins};
% Cellular response
\node[align=center, popdark, text width=2cm] at (0,-4) {Rapid response (seconds--minutes)};
\node[align=center, font=\bfseries, popdark, text width=2.5cm] at (0,2.2) {Protein Hormone Pathway};
\end{scope}

% RIGHT: Steroid hormone (intracellular receptor)
\begin{scope}[xshift=5cm]
% Cell membrane
\draw[popdark, line width=2pt] (-2.5,-0.2) -- (2.5,-0.2);
\draw[popdark, line width=2pt] (-2.5,0.2) -- (2.5,0.2);
\foreach \x in {-2.4,-2.2,...,2.4} {
    \fill[popmuted] (\x,0) circle (0.05);
}
% Nuclear membrane
\draw[popdark, line width=1.5pt, dashed] (-2.5,2.5) -- (2.5,2.5);
\node[popdark, text width=2cm, align=center] at (0,2.8) {\textbf{Nuclear Envelope}};
% Hormone diffusion
\draw[poppink, ->, line width=1pt] (0,1.5) -- (0,0.3);
\node[poppink, text width=1.5cm, align=center] at (0,1.8) {Steroid Hormone};
\draw[poppink, ->, line width=1pt] (0,-0.2) -- (0,-1.0);
% Cytoplasmic receptor + HSP
\draw[popblue, fill=popblue!20] (-1.0,-1.5) rectangle (1.0,-1.0);
\node[popblue, text width=2cm, align=center] at (0,-1.25) {Receptor--HSP Complex};
% Hormone binds, HSP leaves
\draw[poppink, ->, line width=1pt] (0,-1.0) -- (0,-1.25);
\draw[popmuted, ->, line width=0.8pt] (-1.0,-1.25) -- (-1.5,-1.25);
\node[popmuted, font=\tiny, text width=1cm, align=center] at (-1.5,-1.0) {HSP release};
% Hormone-receptor dimer
\draw[popblue, fill=popblue!20] (-0.5,-1.5) rectangle (0.5,-1.0);
\node[popblue, text width=2cm, align=center] at (0,-1.25) {Hormone--Receptor};
% Nuclear translocation
\draw[popblue, ->, line width=1pt] (0,-1.0) -- (0,0.5);
\draw[popblue, ->, line width=1pt] (0,0.5) -- (0,2.0);
% DNA binding
\draw[popdark, line width=1pt] (-1.5,3.0) -- (1.5,3.0); % DNA
\draw[popdark, line width=1pt] (-1.5,3.2) -- (1.5,3.2);
\draw[popblue, fill=popblue!20] (-0.3,3.0) rectangle (0.3,3.4);
\node[popblue, font=\tiny, text width=1cm, align=center] at (0,3.5) {HRE};
% Transcription
\draw[popgreen, ->, line width=1pt] (0,3.4) -- (0,4.0);
\node[popgreen, text width=1.5cm, align=center] at (0,4.3) {mRNA Synthesis};
% Translation
\draw[popgreen, ->, line width=1pt] (1.5,3.4) -- (2.5,3.4);
\node[popgreen, text width=1.5cm, align=center] at (2.8,3.7) {Protein Synthesis};
% Cellular response
\node[align=center, popdark, text width=2cm] at (0,-4) {Slow response (hours--days)};
\node[align=center, font=\bfseries, popdark, text width=2.5cm] at (0,2.2) {Steroid Hormone Pathway};
\end{scope}
\end{tikzpicture}
\legende{Comparison of protein hormone (left, cAMP second messenger) and steroid hormone (right, nuclear receptor / gene transcription) mechanisms of action.}
\label{fig:hormone_mechanisms}
\end{popfigure}

\courssub{Negative Feedback Control of Hormone Secretion}

Endocrine systems are predominantly regulated by \cle{negative feedback loops}: the output (hormone or its effect) inhibits further secretion. This maintains homeostasis.

\begin{methode}[Reading and Interpreting a Negative Feedback Loop]
\begin{enumerate}
    \item \textbf{Identify the stimulus:} What variable changes? (e.g. blood glucose $\uparrow$, blood $\mathrm{Ca^{2+}}$ $\downarrow$, metabolic rate $\downarrow$).
    \item \textbf{Identify the sensor:} Which gland/cells detect the change? (e.g. pancreatic $\beta$-cells, parathyroid chief cells, pituitary thyrotrophs).
    \item \textbf{Identify the hormone secreted:} What is released in response? (e.g. insulin, PTH, TSH).
    \item \textbf{Identify the target and effect:} Where does the hormone act and what does it do? (e.g. liver $\rightarrow$ glycogenesis; bone/kidney $\rightarrow$ $\mathrm{Ca^{2+}}$ release; thyroid $\rightarrow$ thyroxine release).
    \item \textbf{Trace the feedback:} Does the effect \emph{reverse} the original stimulus? (e.g. blood glucose $\downarrow$, blood $\mathrm{Ca^{2+}}$ $\uparrow$, metabolic rate $\uparrow$).
    \item \textbf{State the outcome:} Homeostasis restored; hormone secretion decreases.
\end{enumerate}
\end{methode}

\begin{experience}[Negative Feedback: Thyroxine (T3/T4) Regulation]
\textbf{Scenario:} Cold exposure $\rightarrow$ metabolic rate falls $\rightarrow$ hypothalamus releases TRH $\rightarrow$ anterior pituitary releases TSH $\rightarrow$ thyroid releases T3/T4 $\rightarrow$ basal metabolic rate rises $\rightarrow$ body temperature restored $\rightarrow$ T3/T4 inhibit TRH and TSH secretion (long-loop and short-loop feedback).
\begin{center}
\begin{tikzpicture}[scale=0.8, font=\small, node distance=1.5cm]
% Nodes
\node[draw, rounded corners, fill=popblue!10, popblue, text width=2cm, align=center] (hypo) {Hypothalamus\\TRH};
\node[draw, rounded corners, fill=poporange!10, poporange, right=2.5cm of hypo, text width=2cm, align=center] (pit) {Anterior Pituitary\\TSH};
\node[draw, rounded corners, fill=popgreen!10, popgreen, right=2.5cm of pit, text width=2cm, align=center] (thy) {Thyroid Gland\\T3 / T4};
\node[draw, rounded corners, fill=poppurple!10, poppurple, below=1.5cm of thy, text width=2cm, align=center] (target) {Target Tissues\\\uparrow Metabolic Rate};
% Arrows
\draw[->, thick, popdark] (hypo) -- node[above] {stimulates} (pit);
\draw[->, thick, popdark] (pit) -- node[above] {stimulates} (thy);
\draw[->, thick, popdark] (thy) -- node[right] {stimulates} (target);
% Negative feedback (long loop)
\draw[->, thick, popred, bend left=40] (target) to node[left, popred, text width=1.5cm, align=center] {inhibits (long loop)} (hypo);
% Negative feedback (short loop)
\draw[->, thick, popred, bend right=20] (thy) to node[below, popred, text width=1.5cm, align=center] {inhibits (short loop)} (pit);
% Stimulus
\node[draw, dashed, fill=poppink!10, poppink, left=2cm of hypo, text width=1.8cm, align=center] (stim) {Cold / Low Metabolism};
\draw[->, thick, poppink] (stim) -- (hypo);
\end{tikzpicture}
\end{center}
\legende{Negative feedback control of thyroxine secretion. Long-loop feedback: T3/T4 inhibit hypothalamus. Short-loop feedback: T3/T4 inhibit anterior pituitary.}
\label{fig:thyroxine_feedback}
\end{experience}

\begin{exoresolu}[Blood Glucose Negative Feedback]
\textbf{Statement:} Describe the negative feedback loop that restores blood glucose concentration after a carbohydrate-rich meal.

\textbf{Detailed solution:}
\begin{enumerate}
    \item \textbf{Stimulus:} Blood glucose concentration rises above set point ($\sim 5\ \mathrm{mmol\,dm^{-3}}$).
    \item \textbf{Sensor:} $\beta$-cells of pancreatic islets of Langerhans detect high glucose via GLUT2 transporters and glucokinase.
    \item \textbf{Hormone secreted:} Insulin (protein hormone) released by exocytosis.
    \item \textbf{Target organs and effects:}
    \begin{itemize}
        \item Liver: $\uparrow$ glycogenesis, $\uparrow$ glycolysis, $\downarrow$ gluconeogenesis, $\downarrow$ glycogenolysis.
        \item Muscle/Adipose: $\uparrow$ GLUT4 translocation $\rightarrow$ $\uparrow$ glucose uptake.
    \end{itemize}
    \item \textbf{Result:} Blood glucose concentration falls.
    \item \textbf{Feedback:} Falling glucose removes stimulus $\rightarrow$ insulin secretion decreases. Homeostasis restored.
\end{enumerate}
\textbf{Examiner's tip:} Always mention the \emph{sensor}, the \emph{hormone}, the \emph{target}, the \emph{effect}, and explicitly state that the effect \emph{opposes the original change}.
\end{exoresolu}

\courssub{Comparison of Nervous and Hormonal Coordination}

Both systems coordinate body functions, but their modes of operation differ fundamentally.

\begin{propriete}[Nervous vs Hormonal Coordination]
\begin{center}
\begin{tabularx}{\linewidth}{|l|X|X|}\hline
\rowcolor{popblueL}
\textbf{Feature} & \textbf{Nervous System} & \textbf{Endocrine System} \\ \hline
\textbf{Nature of signal} & Electrical (action potentials) + chemical (neurotransmitters at synapses) & Chemical only (hormones) \\ \hline
\textbf{Transmission pathway} & Along specific neurons (wired) & Via bloodstream (wireless) \\ \hline
\textbf{Speed of transmission} & Very fast (up to $120\ \mathrm{m\,s^{-1}}$) & Slow (seconds to minutes for blood transport) \\ \hline
\textbf{Speed of response} & Milliseconds & Seconds (protein hormones) to hours/days (steroids) \\ \hline
\textbf{Duration of action} & Brief (ms--s); stops when impulses stop & Prolonged (min--days); hormone persists in blood \\ \hline
\textbf{Specificity} & High (specific pathway to specific effector) & High (specific receptor on target cells) but widespread distribution \\ \hline
\textbf{Adaptation} & Rapid adaptation (habituation) & Little adaptation; sustained response \\ \hline
\textbf{Integration with behaviour} & Direct control of voluntary muscle, rapid reflexes & Modulates metabolism, growth, reproduction, mood \\ \hline
\end{tabularx}
\end{center}
\end{propriete}

\begin{savaistu}[Neuroendocrine Integration in Cameroon's Context]
The \cle{hypothalamus} is the master link: it is neural tissue that secretes hormones (releasing/inhibiting factors into the hypophyseal portal system, ADH and oxytocin into posterior pituitary). In stressful situations — like a student facing the GCE Advanced Level exams — the \cle{hypothalamo-pituitary-adrenal (HPA) axis} is activated: neural stress signals $\rightarrow$ hypothalamus (CRH) $\rightarrow$ pituitary (ACTH) $\rightarrow$ adrenal cortex (cortisol). Cortisol then feeds back to the brain to modulate mood and cognition. This neuroendocrine crosstalk explains why chronic stress affects both mental performance and physical health (immune suppression, glucose dysregulation). In Cameroon, rising prevalence of type 2 diabetes in urban areas (Yaoundé, Douala) illustrates the endocrine dimension of lifestyle changes: chronic high-calorie diets $\rightarrow$ sustained insulin demand $\rightarrow$ $\beta$-cell exhaustion $\rightarrow$ hyperglycaemia.
\end{savaistu}

\courssub{Consolidation Exercises}

\begin{exercice}[Exercise 1: Hormone Classification and Mechanism]
Classify each hormone below as \emph{protein/peptide}, \emph{steroid}, or \emph{amine}. For each, state: (i) receptor location (membrane / cytoplasmic / nuclear), (ii) second messenger if applicable, (iii) typical time scale of action.
\begin{enumerate}
    \item Insulin
    \item Cortisol
    \item Adrenaline
    \item Thyroxine (T4)
    \item Antidiuretic hormone (ADH)
    \item Testosterone
\end{enumerate}
\end{exercice}

\begin{exercice}[Exercise 2: Negative Feedback Interpretation]
The diagram below shows a simplified feedback loop for blood calcium regulation.
\begin{center}
\begin{tikzpicture}[scale=0.7, font=\small]
\node[draw, rounded corners, fill=popblue!10, text width=1.8cm, align=center] (lowCa) {Low Blood $\mathrm{Ca^{2+}}$};
\node[draw, rounded corners, fill=poporange!10, right=2.5cm of lowCa, text width=1.8cm, align=center] (parathyroid) {Parathyroid Gland};
\node[draw, rounded corners, fill=popgreen!10, right=2.5cm of parathyroid, text width=1.8cm, align=center] (PTH) {PTH Secreted};
\node[draw, rounded corners, fill=poppurple!10, right=2.5cm of PTH, text width=1.8cm, align=center] (targets) {Bone / Kidney / Gut};
\node[draw, rounded corners, fill=poppink!10, below=1.5cm of targets, text width=1.8cm, align=center] (highCa) {Blood $\mathrm{Ca^{2+}}$ Rises};
\draw[->, thick] (lowCa) -- (parathyroid);
\draw[->, thick] (parathyroid) -- (PTH);
\draw[->, thick] (PTH) -- (targets);
\draw[->, thick] (targets) -- (highCa);
\draw[->, thick, popred, bend left=40] (highCa) to node[left, popred, text width=1.5cm, align=center] {Inhibits} (parathyroid);
\end{tikzpicture}
\end{center}
\begin{enumerate}
    \item Identify the stimulus, sensor, hormone, effectors, and response.
    \item Explain why this is a \emph{negative} feedback loop.
    \item Predict the effect of a tumour in the parathyroid gland that secretes PTH autonomously (ignoring feedback).
\end{enumerate}
\end{exercice}

\begin{corrige}[Exercise 1]
\begin{enumerate}
    \item \textbf{Insulin:} Protein. (i) Membrane receptor (RTK). (ii) No classic second messenger; autophosphorylation cascade (PI3K/Akt). (iii) Seconds--minutes.
    \item \textbf{Cortisol:} Steroid. (i) Cytoplasmic $\rightarrow$ nuclear. (ii) None (direct gene regulation). (iii) Hours--days.
    \item \textbf{Adrenaline:} Amine (catecholamine). (i) Membrane receptor (GPCR). (ii) cAMP (or IP$_3$/DAG via $\alpha$-receptors). (iii) Seconds.
    \item \textbf{Thyroxine (T4):} Amine (iodinated tyrosine) — behaves as steroid. (i) Nuclear receptor. (ii) None. (iii) Hours--days.
    \item \textbf{ADH:} Peptide. (i) Membrane receptor (GPCR, V2 receptor in kidney). (ii) cAMP. (iii) Minutes.
    \item \textbf{Testosterone:} Steroid. (i) Cytoplasmic/nuclear. (ii) None. (iii) Hours--days.
\end{enumerate}
\end{corrige}

\begin{corrige}[Exercise 2]
\begin{enumerate}
    \item \textbf{Stimulus:} Low blood $\mathrm{Ca^{2+}}$. \textbf{Sensor:} Parathyroid chief cells (CaSR receptors). \textbf{Hormone:} Parathyroid Hormone (PTH). \textbf{Effectors:} Bone (osteoclasts $\uparrow$ resorption), Kidney ($\uparrow$ $\mathrm{Ca^{2+}}$ reabsorption, $\uparrow$ $1\alpha$-hydroxylase $\rightarrow$ active vitamin D $\rightarrow$ $\uparrow$ gut $\mathrm{Ca^{2+}}$ absorption). \textbf{Response:} Blood $\mathrm{Ca^{2+}}$ rises.
    \item It is negative feedback because the \emph{response} (rise in blood $\mathrm{Ca^{2+}}$) \emph{inhibits} the \emph{sensor/gland}, reducing further PTH secretion. The output opposes the initial stimulus.
    \item Autonomous PTH secretion $\rightarrow$ persistent high PTH $\rightarrow$ chronic hypercalcaemia (bone resorption $\rightarrow$ osteoporosis, kidney stones, abdominal pain, confusion). Feedback inhibition is lost; the gland no longer responds to high $\mathrm{Ca^{2+}}$. This is \cle{primary hyperparathyroidism}.
\end{enumerate}
\end{corrige}

\begin{aretenir}[Section 8 — Key Points for Revision]
\begin{itemize}
    \item Endocrine glands are ductless; hormones travel in blood to specific target cells with receptors.
    \item Pancreas = mixed gland (exocrine acini + endocrine islets). Do not confuse roles.
    \item Protein hormones $\rightarrow$ membrane receptors $\rightarrow$ second messengers (cAMP, IP$_3$/DAG, Ca$^{2+}$) $\rightarrow$ rapid action via phosphorylation.
    \item Steroid hormones $\rightarrow$ intracellular receptors $\rightarrow$ gene transcription $\rightarrow$ slow, long-lasting action via new protein synthesis.
    \item Negative feedback: hormone (or its effect) inhibits its own secretion $\rightarrow$ homeostasis.
    \item Nervous = fast, wired, brief. Endocrine = slow, wireless, sustained. They interact (neuroendocrine system).
\end{itemize}
\end{aretenir}

\coursec{Functions of the Endocrine Glands}

In the previous section we saw that hormones are chemical messengers released into the blood by ductless glands, and that they act only on target cells bearing the correct receptors. We now examine each major endocrine gland in turn: the hormones it secretes, the target organs, the effects produced, and the disorders that arise when secretion is too high or too low. As you read, notice a recurring theme: most endocrine secretions are controlled by \cle{negative feedback}, so that the internal environment is kept remarkably constant.

\courssub{The Pituitary Gland: the ``Master Gland''}

The \cle{pituitary gland} is a small pea-sized structure hanging from the hypothalamus at the base of the brain. It is called the \cle{master gland} because many of its hormones control other endocrine glands. It has two lobes with different origins and modes of action.

\textbf{The anterior lobe} synthesises and releases several hormones under the control of releasing factors from the hypothalamus:

\begin{itemize}
\item \cle{Growth hormone (GH)}: stimulates growth of bones and soft tissues and promotes protein synthesis. \textbf{Over-secretion in childhood} causes \cle{gigantism} (abnormal height); \textbf{under-secretion in childhood} causes \cle{dwarfism} (stunted growth but normal proportions and intelligence). Over-secretion in adults causes thickening of bones of the hands, feet and jaw (\cle{acromegaly}).
\item \cle{Thyroid-stimulating hormone (TSH)}: stimulates the thyroid gland to secrete thyroxine.
\item \cle{Adrenocorticotrophic hormone (ACTH)}: stimulates the adrenal cortex to release corticosteroids.
\item \cle{Follicle-stimulating hormone (FSH)}: in females stimulates growth of ovarian follicles; in males stimulates sperm production.
\item \cle{Luteinising hormone (LH)}: in females triggers ovulation and formation of the corpus luteum; in males stimulates testosterone secretion.
\item \cle{Prolactin}: stimulates milk production by the mammary glands after childbirth.
\end{itemize}

\textbf{The posterior lobe} does not make hormones; it stores and releases \cle{antidiuretic hormone (ADH)} and oxytocin, both produced by the hypothalamus. ADH increases water reabsorption in the kidney collecting ducts, concentrating the urine; deficiency causes \cle{diabetes insipidus} (production of large volumes of dilute urine).

\begin{attention}[Anterior versus posterior pituitary]
Do not write that the posterior pituitary ``secretes'' ADH in the sense of manufacturing it. ADH is \emph{synthesised} in the hypothalamus and only \emph{stored and released} by the posterior lobe. Examiners penalise this confusion.
\end{attention}

\courssub{The Thyroid Gland}

The \cle{thyroid gland} lies in the neck, on either side of the trachea just below the larynx. It secretes \cle{thyroxine}, an iodine-containing hormone that:

\begin{itemize}
\item controls the \cle{metabolic rate} (rate of cellular respiration and heat production);
\item is essential for normal \cle{growth and development}, especially of the skeleton and nervous system in children.
\end{itemize}

\textbf{Disorders.} Iodine is needed to make thyroxine. Where the diet lacks iodine, the thyroid enlarges in an attempt to capture more iodine, producing a visible swelling of the neck called a \cle{goitre} (simple or endemic goitre). Severe thyroxine deficiency in childhood causes \cle{cretinism} (stunted growth and mental retardation); in adults it causes \cle{myxoedema} (low metabolic rate, weight gain, lethargy). Over-secretion (hyperthyroidism) raises metabolic rate, causing weight loss, rapid heartbeat and nervousness.

\begin{savaistu}[Iodine deficiency in Cameroon]
Goitre due to iodine deficiency has historically been reported in some hilly and mountainous regions of Cameroon, where soils (and therefore crops) are poor in iodine. The use of \textbf{iodised table salt} is a simple public-health measure that prevents this disorder --- a good example of biology applied to everyday life.
\end{savaistu}

\courssub{The Adrenal Glands}

The two \cle{adrenal glands} sit on top of the kidneys. Each has an inner \textbf{medulla} and an outer \textbf{cortex}.

\textbf{Adrenal medulla --- adrenaline.} \cle{Adrenaline} is released in response to stress, fear or danger and produces the \cle{fight-or-flight} response:

\begin{itemize}
\item increases \textbf{heart rate} and the force of contraction, raising blood supply to muscles;
\item stimulates \cle{glycogenolysis} (breakdown of glycogen to glucose) in the liver, raising blood glucose for respiration;
\item dilates the \textbf{pupils} and the bronchioles;
\item diverts blood from the gut and skin to skeletal muscles.
\end{itemize}

\textbf{Adrenal cortex --- corticosteroids.} These include glucocorticoids (e.g.\ cortisol), which raise blood glucose during prolonged stress and suppress inflammation, and mineralocorticoids (e.g.\ aldosterone), which regulate sodium and potassium balance and hence water reabsorption in the kidney.

\begin{attention}[Adrenaline: hormone or nervous effect?]
Adrenaline \emph{is a hormone}: it is secreted into the blood by the adrenal medulla. However, its effects (faster heartbeat, pupil dilation, reduced gut activity) are essentially the \textbf{same as those of sympathetic nervous stimulation}. The distinction to state clearly is this: the sympathetic nervous system acts \emph{rapidly and briefly} through nerves, while adrenaline acts \emph{more slowly but for longer} through the bloodstream. The two systems reinforce each other in emergencies.
\end{attention}

\courssub{The Pancreas and the Control of Blood Glucose}

The pancreas is both an exocrine gland (digestive enzymes) and an endocrine gland. Its endocrine tissue, the \cle{islets of Langerhans}, contains two cell types:

\begin{itemize}
\item \textbf{$\beta$ cells} secrete \cle{insulin} when blood glucose is \emph{high};
\item \textbf{$\alpha$ cells} secrete \cle{glucagon} when blood glucose is \emph{low}.
\end{itemize}

\textbf{Insulin lowers blood glucose} by: (i) stimulating the liver and muscles to convert glucose to glycogen (\cle{glycogenesis}); (ii) increasing \textbf{glucose uptake} and respiration by cells; (iii) promoting conversion of excess glucose to fat.

\textbf{Glucagon raises blood glucose} by stimulating the liver to break down glycogen into glucose (\cle{glycogenolysis}) and to form glucose from amino acids (\cle{gluconeogenesis}).

\begin{propriete}[Antagonistic action of insulin and glucagon]
Insulin and glucagon act \cle{antagonistically}: insulin lowers blood glucose while glucagon raises it. Together, through negative feedback, they keep the blood glucose concentration approximately constant at about $90\ \mathrm{mg}$ per $100\ \mathrm{cm^{3}}$ of blood ($\approx 5\ \mathrm{mmol\,dm^{-3}}$). This is a core example of \cle{homeostasis}.
\end{propriete}

\begin{definition}[Negative feedback in blood glucose homeostasis]
\cle{Negative feedback} is a control mechanism in which a change in a regulated variable triggers a response that \emph{reverses} that change, returning the variable towards its set point. Applied to blood glucose: a \textbf{rise} in blood glucose after a meal is detected by the pancreas, which secretes \textbf{insulin}; insulin lowers glucose, so the stimulus for insulin secretion disappears. A \textbf{fall} in blood glucose triggers \textbf{glucagon} secretion, which raises glucose again. The output thus switches off its own stimulus.
\end{definition}

\begin{methode}[Describing a negative feedback loop]
To explain any negative feedback loop (e.g. blood glucose, thyroxine, ADH, calcium) in an exam answer, follow these four steps:
\begin{enumerate}
\item \textbf{Stimulus:} State the change in the internal environment (e.g. blood glucose rises after a meal).
\item \textbf{Receptor:} Identify the sensor (e.g. $\beta$ cells of the islets of Langerhans detect high glucose).
\item \textbf{Effector response:} Name the hormone secreted and its action (e.g. insulin $\to$ glycogenesis, glucose uptake).
\item \textbf{Return to set point:} Explain how the response reverses the stimulus and switches off further hormone release.
\end{enumerate}
Always use the phrase ``negative feedback'' and mention the \cle{set point}.
\end{methode}

\begin{popfigure}
\begin{center}
\begin{tikzpicture}
\begin{axis}[
  width=12cm, height=7cm,
  xlabel={Time after a meal / min},
  ylabel={Blood glucose / mg per $100\ \mathrm{cm^{3}}$},
  xmin=0, xmax=180, ymin=60, ymax=210,
  xtick={0,30,60,90,120,150,180},
  ytick={60,90,120,150,180,210},
  grid=both, grid style={popline, very thin},
  legend style={at={(0.97,0.97)}, anchor=north east, font=\small},
  axis line style={popink},
]
\addplot[very thick, popblue, smooth] coordinates {
 (0,90) (20,105) (40,135) (60,120) (80,100) (100,92) (120,88) (140,86) (160,88) (180,90)
};
\addlegendentry{Normal person}
\addplot[very thick, poporange, smooth, dashed] coordinates {
 (0,150) (30,185) (60,200) (90,190) (120,175) (150,165) (180,158)
};
\addlegendentry{Diabetic person}
\addplot[popmuted, thin, dashed] coordinates {(0,90) (180,90)};
\addlegendentry{Set point ($\approx 90$)}
\end{axis}
\end{tikzpicture}
\end{center}
\legende{Blood glucose regulation after a meal. In a normal person, insulin brings the post-meal rise back to the set point within about two hours; if glucose later falls, glucagon restores it. In an untreated diabetic, glucose starts high, rises further and falls only slowly.}
\end{popfigure}

\textbf{Diabetes mellitus.} \cle{Diabetes mellitus} is a disorder in which blood glucose is persistently too high because insulin is lacking or ineffective.

\begin{itemize}
\item \textbf{Causes:} failure of the $\beta$ cells to secrete insulin (Type 1, usually appearing in youth) or failure of body cells to respond to insulin (Type 2, associated with obesity and age).
\item \textbf{Symptoms:} high blood glucose (hyperglycaemia), \textbf{glucose in the urine}, excessive urine production, persistent thirst, tiredness and weight loss.
\item \textbf{Control:} Type 1 is managed by regular \textbf{insulin injections} (insulin is a protein and would be digested if swallowed); Type 2 is managed by controlling the diet (reducing sugar intake), exercise, and drugs that improve insulin action.
\end{itemize}

\begin{savaistu}[Diabetes in Cameroon]
Diabetes mellitus is a growing public health concern in Cameroon, driven by urbanisation, reduced physical activity and diets richer in refined sugars and fats. The Cameroon Ministry of Public Health runs awareness campaigns on World Diabetes Day (14 November) promoting early screening, balanced nutrition and regular exercise. Traditional foods like \emph{ndol\'e} (bitterleaf) and \emph{eru} are high in fibre and can help regulate blood glucose when prepared with less oil.
\end{savaistu}

\courssub{The Ovaries and the Testes}

\textbf{Ovaries.} The ovaries secrete \cle{oestrogen} and \cle{progesterone}.

\begin{itemize}
\item \textbf{Oestrogen} (from the developing follicles) controls the development of female \cle{secondary sexual characteristics} at puberty (breast development, widening of hips, growth of pubic hair, onset of the menstrual cycle), and repairs and thickens the uterus lining after menstruation.
\item \textbf{Progesterone} (from the corpus luteum) maintains the thickened uterus lining during the second half of the \cle{menstrual cycle} and throughout \textbf{pregnancy}, and inhibits further ovulation during pregnancy.
\end{itemize}

\textbf{Testes.} The testes secrete \cle{testosterone}, which:

\begin{itemize}
\item stimulates the development of male \cle{secondary sexual characteristics} at puberty (deepening of the voice, growth of beard and body hair, muscular development, enlargement of the testes and penis);
\item is necessary for the continuous \textbf{production of sperm} (spermatogenesis).
\end{itemize}

\courssub{The Parathyroid Glands}

The four small \cle{parathyroid glands} are embedded in the back of the thyroid. They secrete \cle{parathormone (PTH)}, which \textbf{raises blood calcium} by stimulating calcium release from bones, increasing calcium reabsorption in the kidneys, and promoting calcium absorption from the gut (via vitamin D). Calcium is vital for nerve impulse transmission, muscle contraction and blood clotting; over-secretion weakens bones, while under-secretion causes muscle spasms (\cle{tetany}).

\courssub{Summary of the Endocrine Glands}

\begin{popfigure}
\begin{center}
\begin{tikzpicture}[font=\small]
\matrix (summary) [
  matrix of nodes,
  nodes={draw=popline, fill=white, rounded corners, text width=3.2cm, align=center, minimum height=1.8cm, anchor=center},
  column sep=0.5cm, row sep=0.5cm,
  row 1/.style={nodes={fill=popblueL, text=popink, font=\bfseries, minimum height=0.7cm}},
  row 2/.style={nodes={minimum height=2.2cm}},
  row 3/.style={nodes={minimum height=2.2cm}},
] {
Gland & Hormone(s) & Target & Main effect \\
\textbf{Pituitary} & GH, TSH, ACTH, FSH, LH, prolactin, ADH & Various endocrine glands, kidney, uterus, breast & Controls other glands; growth; water balance \\
\textbf{Thyroid} & Thyroxine & All cells & Metabolic rate, growth, development \\
\textbf{Adrenals} & Adrenaline; corticosteroids & Heart, liver, pupils; kidney, liver & Fight-or-flight; long-term stress, salt balance \\
\textbf{Pancreas} & Insulin, glucagon & Liver, muscles, adipose tissue & Blood glucose homeostasis \\
\textbf{Ovaries} & Oestrogen, progesterone & Uterus, breasts, bones & Female sexual characteristics, menstrual cycle, pregnancy \\
\textbf{Testes} & Testosterone & Male tissues, testes & Male sexual characteristics, spermatogenesis \\
\textbf{Parathyroids} & Parathormone (PTH) & Bones, kidneys, gut & Blood calcium homeostasis \\
};
\node[draw=popline, fill=popgreenL, rounded corners, text width=14.5cm, align=center, minimum height=1.2cm, anchor=north] at (summary-3-2.south) {\textbf{Key idea:} Negative feedback loops maintain homeostasis for glucose, calcium, water, metabolic rate and stress responses.};
\end{tikzpicture}
\end{center}
\legende{Summary map of the major endocrine glands, their hormones, targets and effects.}
\end{popfigure}

\begin{exoresolu}[Applying knowledge of endocrine disorders]
\textbf{Statement.} A 14-year-old boy in a village in the North-West Region has a visibly swollen neck, feels tired, and is shorter than most boys of his age. Local crops are grown on iodine-poor soils. (a) Name the gland and hormone most likely involved. (b) Name the disorder shown by the swollen neck and explain its cause. (c) Suggest one simple preventive measure for the community. (d) Explain why the boy's growth is affected.
\tcblower
\textbf{Solution.}
\begin{enumerate}
\item The \textbf{thyroid gland} and its hormone \textbf{thyroxine} are involved.
\item The swollen neck is a \textbf{goitre}. Because the diet lacks iodine, the thyroid cannot make enough thyroxine; the gland enlarges as it works harder to trap the little iodine available.
\item Use of \textbf{iodised table salt} (or iodine-rich foods such as sea fish) prevents the deficiency.
\item Thyroxine is needed for normal \textbf{growth and development}; deficiency in childhood slows skeletal growth (and can impair mental development), so the boy is shorter than average.
\end{enumerate}
\end{exoresolu}

\begin{attention}[Exam tip: one disorder per gland]
For the GCE, learn \textbf{at least one disorder per major gland, with its cause and symptoms}: pituitary --- gigantism or dwarfism (abnormal GH in childhood); thyroid --- goitre (iodine deficiency); pancreas --- diabetes mellitus (lack of, or insensitivity to, insulin); parathyroids --- tetany (under-secretion of parathormone). Be ready to link cause $\to$ hormone level $\to$ symptoms in a logical chain.
\end{attention}

\begin{exercice}[Checking your understanding]
\begin{enumerate}
\item Explain why the pituitary is called the master gland, giving two examples of its control over other glands.
\item A patient has a rapid heartbeat, weight loss despite eating well, and bulging eyes. Which gland and hormone are likely involved, and is secretion too high or too low?
\item Describe, step by step, the negative feedback response when blood glucose rises after a meal of \emph{ndol\'e} and \emph{bobolo}.
\item State two differences between the action of adrenaline and the action of insulin.
\end{enumerate}
\end{exercice}

\begin{corrige}[Exercise 1]
\begin{enumerate}
\item The pituitary secretes hormones that control other endocrine glands, e.g.\ \textbf{TSH} stimulates the thyroid to release thyroxine, and \textbf{ACTH} stimulates the adrenal cortex to release corticosteroids; hence ``master gland''.
\item The \textbf{thyroid gland}, hormone \textbf{thyroxine}; secretion is \textbf{too high} (hyperthyroidism), raising metabolic rate and heart rate.
\item Rise in blood glucose $\to$ detected by $\beta$ cells of the islets of Langerhans $\to$ \textbf{insulin} secreted $\to$ liver and muscles convert glucose to glycogen (glycogenesis) and cells take up more glucose $\to$ blood glucose falls to the set point $\to$ insulin secretion switches off.
\item Adrenaline \emph{raises} blood glucose (glycogenolysis) and acts quickly in emergencies; insulin \emph{lowers} blood glucose (glycogenesis, glucose uptake) and acts during normal feeding. Adrenaline also speeds the heart and dilates pupils; insulin does not.
\end{enumerate}
\end{corrige}

\begin{aretenir}[Key points]
\begin{itemize}
\item The pituitary is the master gland; GH disorders cause gigantism or dwarfism; ADH (stored in the posterior lobe) controls water reabsorption.
\item Thyroxine sets the metabolic rate and drives growth; iodine deficiency causes goitre.
\item Adrenaline produces the fight-or-flight response; corticosteroids act in longer-term stress and salt balance.
\item Insulin and glucagon act antagonistically, by negative feedback, to hold blood glucose near its set point; failure of this system gives diabetes mellitus.
\item Oestrogen and progesterone control female sexual development, the menstrual cycle and pregnancy; testosterone controls male sexual development and sperm production; parathormone raises blood calcium.
\end{itemize}
\end{aretenir}

\coursec{Plant Movements, Plant Hormones, Light and Plant Growth}

In the previous section we saw how animals coordinate their activities through hormones secreted by endocrine glands and transported in the blood. Plants have no blood, no glands and no nervous system, yet a maize seedling in a Cameroonian farm still manages to grow upward out of the soil, bend toward sunlight, and a groundnut plant folds its leaves at night. How is this possible? Plants too produce chemical messengers — the \cle{plant hormones} or \cle{phytohormones} — and they respond to stimuli such as light, gravity, water and touch through slow but precise \cle{growth movements}. In this final section we study these plant responses and show how farmers exploit them in agriculture.

\courssub{Plant Movements: Tropisms and Nastic Movements}

\textbf{Observation.} Place a potted bean seedling on a window sill in Buea. After a few days, the shoot bends toward the window. Touch the leaflets of \emph{Mimosa pudica} (the ``sensitive plant'' common along our roadsides) and they fold within seconds. Both are responses to stimuli, but they differ fundamentally: the bending of the shoot is slow, permanent, and directed by the stimulus; the folding of Mimosa is rapid, reversible, and occurs in the same direction whatever the direction of touch.

\begin{definition}[Tropism]
A \cle{tropism} is a directional growth movement of a plant organ in response to a directional (unilateral) external stimulus. The direction of growth is related to the direction of the stimulus. A tropism is \textbf{positive} when the organ grows \emph{toward} the stimulus and \textbf{negative} when it grows \emph{away} from it.
\end{definition}

\begin{definition}[Nastic movement]
A \cle{nastic movement} is a non-directional movement of a plant organ in response to a non-directional (diffuse) stimulus. The direction of the movement is determined by the structure of the organ, not by the direction of the stimulus. Nastic movements are usually rapid, reversible and caused by changes in turgor pressure, not by growth.
\end{definition}

\textbf{The main types of tropisms} are named after the stimulus involved:

\begin{center}
\begin{tabularx}{\linewidth}{|l|l|X|}
\hline
\rowcolor{popblueL} \textbf{Tropism} & \textbf{Stimulus} & \textbf{Example} \\
\hline
Phototropism & Light & Shoots grow toward light (positive); roots grow away (negative) \\
\hline
Geotropism (gravitropism) & Gravity & Roots grow downward (positive); shoots grow upward (negative) \\
\hline
Hydrotropism & Water & Roots grow toward moist soil (positive) \\
\hline
Thigmotropism (haptotropism) & Touch / contact & Tendrils of passion fruit coil around a support (positive) \\
\hline
Chemotropism & Chemicals & Pollen tube grows toward the ovule (positive) \\
\hline
\end{tabularx}
\end{center}

\textbf{Examples of nastic movements} include:
\begin{itemize}
\item \textbf{Thigmonasty:} folding of leaflets of \emph{Mimosa pudica} when touched, caused by rapid loss of turgor in swollen motor cells (the \cle{pulvini}) at the leaflet bases;
\item \textbf{Nyctinasty:} ``sleep movements'' — leaves of many legumes (e.g. groundnut) fold at night and open in the morning in response to changes in light intensity and temperature;
\item \textbf{Photonasty:} opening of flowers (e.g. hibiscus) in response to light, and their closing at dusk.
\end{itemize}

\begin{attention}[Tropism or nastic movement?]
A very common examination error is to confuse the two. Remember: \textbf{tropisms are growth movements} — they are slow, irreversible (permanent), and their direction depends on the direction of the stimulus. \textbf{Nastic movements are reversible}, often rapid turgor movements, and are \emph{not} oriented by the stimulus. Mimosa folds its leaves whether you touch it from above, below or the side: that is a nastic movement, not a tropism. Also note the spelling \emph{geotropism} = \emph{gravitropism}; both are accepted.
\end{attention}

\courssub{Plant Hormones (Phytohormones)}

Plants coordinate growth and responses through small organic molecules produced in one part of the plant and transported (often in the phloem or from cell to cell) to target tissues where they act at very low concentrations.

\begin{definition}[Plant hormone]
A \cle{plant hormone} (phytohormone) is a chemical substance produced naturally in small quantities in one region of a plant and transported to another region, where it regulates growth, development or responses to stimuli.
\end{definition}

The five classical groups of phytohormones and their main effects are summarised below.

\begin{center}
\begin{tabularx}{\linewidth}{|l|X|X|}
\hline
\rowcolor{popblueL} \textbf{Hormone} & \textbf{Main site of production} & \textbf{Main effects} \\
\hline
Auxins (e.g. IAA) & Shoot tips (apical buds), young leaves, developing seeds & Promote cell elongation in shoots; apical dominance; root initiation; fruit development; mediate phototropism and geotropism \\
\hline
Gibberellins & Young leaves, embryos of seeds, root tips & Promote stem elongation (internode growth); break seed dormancy; stimulate germination; induce flowering and fruit growth \\
\hline
Cytokinins & Root tips, developing fruits and seeds & Promote cell division (cytokinesis); delay leaf senescence (ageing); promote lateral bud growth \\
\hline
Abscisic acid (ABA) & Mature leaves, stems, fruits (especially under stress) & Inhibits growth; promotes seed and bud dormancy; stimulates stomatal closure during water stress (``stress hormone'') \\
\hline
Ethene (ethylene) & Ripening fruits, ageing tissues, nodes & Promotes fruit ripening; promotes leaf and fruit abscission (fall); inhibits stem elongation \\
\hline
\end{tabularx}
\end{center}

\begin{definition}[Auxin]
\cle{Auxins} are plant hormones, the commonest natural one being \cle{indoleacetic acid (IAA)}, synthesised mainly in shoot apical meristems and young leaves. They are transported downwards (basipetally) from cell to cell and promote cell elongation in shoots by loosening cell walls and stimulating proton pumps that acidify the wall space.
\end{definition}

\begin{propriete}[Differential sensitivity of shoots and roots to auxin]
At the same concentration, auxin \textbf{promotes cell elongation in shoots but inhibits cell elongation in roots}. Roots are far more sensitive to auxin than shoots: the optimum concentration for root elongation is very low (about $10^{-10}$ to $10^{-8}$ mol/L), while the optimum for shoots is much higher (about $10^{-5}$ to $10^{-4}$ mol/L). Concentrations that stimulate shoot growth therefore inhibit root growth. This property explains the opposite geotropic responses of shoots and roots. (The interpretation of the dose--response curve is examinable.)
\end{propriete}

\begin{popfigure}
\begin{center}
\begin{tikzpicture}
\begin{axis}[
    width=11cm, height=7.5cm,
    xlabel={Auxin concentration (mol/L, log scale)},
    ylabel={Growth (\% of control)},
    xmode=log,
    xmin=1e-11, xmax=1e-2,
    ymin=-100, ymax=100,
    axis lines=left,
    legend style={at={(0.03,0.97)}, anchor=north west, font=\small},
    grid=both, grid style={popline!40},
]
\addplot[very thick, popgreen, smooth, domain=1e-11:1e-2, samples=100]
    {100*exp(-((ln(x)+9.21)/1.6)^2) - 100*exp(-((ln(x)+6.0)/1.2)^2)};
\addlegendentry{Shoots (stems)}
\addplot[very thick, poporange, smooth, domain=1e-11:1e-2, samples=100]
    {100*exp(-((ln(x)+23.0)/1.3)^2) - 100*exp(-((ln(x)+18.4)/1.1)^2)};
\addlegendentry{Roots}
\addplot[popmuted, dashed, domain=1e-11:1e-2] {0};
\draw[popblue, dashed] (axis cs:1e-5,-100) -- (axis cs:1e-5,100);
\node[popblue, font=\small, align=center] at (axis cs:3e-4,80) {concentration that\\stimulates shoots\\inhibits roots};
\end{axis}
\end{tikzpicture}
\end{center}
\legende{Dose--response curves of auxin on shoots and roots. Note the logarithmic concentration axis: roots respond to concentrations about $10^{4}$ times lower than shoots, and are inhibited by concentrations that strongly stimulate shoot elongation.}
\end{popfigure}

\courssub{The Mechanism of Phototropism}

\textbf{Observation.} A coleoptile (the protective sheath of a grass seedling, e.g. maize) illuminated from one side bends toward the light. Classic experiments (Darwin, then Went and others) showed that the \emph{tip} detects the light, but the \emph{bending} occurs in the elongation zone below the tip — evidence that a chemical messenger travels downward.

\begin{methode}[Demonstrating phototropism with coleoptiles]
\begin{enumerate}
\item Grow maize or oat seedlings in darkness until the coleoptiles are about $2$--$3$ cm tall.
\item Divide them into four groups and treat as follows: (A) control, left intact; (B) tip covered with an opaque foil cap; (C) tip removed; (D) base of coleoptile covered with foil, tip exposed.
\item Illuminate all groups from one side only (unilateral light) for several hours.
\item Observe: A bends toward the light; B and C do not bend (the tip is needed to perceive light and to produce auxin); D bends normally (the light receptor is at the tip, not the base).
\item A \cle{clinostat} (a slowly rotating drum) can be used as a control to neutralise unilateral stimuli: seedlings rotated on a clinostat grow straight because every side receives the stimulus equally.
\end{enumerate}
Conclusion: the tip perceives light and produces a growth-promoting substance (auxin) that is redistributed and causes differential growth below the tip.
\end{methode}

\begin{experience}[Went's agar-block experiment (historical proof that the messenger is chemical)]
Went cut off oat coleoptile tips and placed them on agar blocks for some hours; the agar alone, placed asymmetrically on decapitated coleoptiles in darkness, caused curvature away from the side bearing the block — exactly as if a growth-promoting substance had diffused into the agar and then down one side of the shoot. Control blocks of plain agar (never in contact with a tip) caused no curvature. This proved that the tip releases a diffusible chemical — later identified as IAA — and that its \emph{asymmetric distribution}, not light itself, bends the shoot.
\end{experience}

\textbf{Explanation of the bending.} Auxin (IAA) is synthesised at the shoot tip. Under unilateral light, auxin is transported laterally \emph{away from the illuminated side} toward the shaded side. The shaded side therefore receives more auxin, its cells elongate more, and the shoot curves toward the light.

\begin{popfigure}
\begin{center}
\begin{tikzpicture}[scale=1.0]
% Light source
\draw[popgold, very thick] (-4.5,2.2) circle (0.35);
\foreach \a in {0,45,...,315} \draw[popgold, thick] (-4.5,2.2) ++(\a:0.45) -- ++(\a:0.3);
\node[popgold, font=\small] at (-4.5,3.1) {light};
% Light rays
\foreach \y in {1.4,2.2,3.0} \draw[popgold, -latex, thick] (-4.0,\y) -- (-1.6,\y);
% Coleoptile outline
\draw[popgreen, very thick, fill=popgreenL]
  (-0.9,0) -- (-0.9,2.6) .. controls (-0.9,3.2) and (-0.3,3.5) .. (0,3.7)
  .. controls (0.3,3.5) and (0.9,3.2) .. (0.9,2.6) -- (0.9,0) -- cycle;
% Tip zone
\draw[popdark, dashed] (-0.9,2.6) -- (0.9,2.6);
\node[font=\small, align=center] at (2.6,3.3) {tip: auxin produced\\and redistributed};
\draw[popdark, -latex] (1.7,3.2) -- (0.6,3.0);
% Auxin redistribution arrows
\draw[poppurple, very thick, -latex] (-0.35,3.15) -- (0.45,3.15);
\node[poppurple, font=\small] at (0.1,3.5) {IAA};
% Downward arrows: small on lit side, big on shaded side
\draw[poppurple, very thick, -latex] (-0.45,2.4) -- (-0.45,1.7);
\draw[poppurple, very thick, -latex] (0.45,2.4) -- (0.45,0.9);
\node[poppurple, font=\small, align=center] at (-1.9,0.9) {little auxin:\\less elongation\\(lit side)};
\node[poppurple, font=\small, align=center] at (2.7,0.9) {more auxin:\\more elongation\\(shaded side)};
% Result: bent shoot
\begin{scope}[xshift=7cm]
\draw[popgreen, very thick, fill=popgreenL]
  (-0.9,0) -- (-0.9,1.2) .. controls (-0.95,2.4) and (-1.6,2.9) .. (-2.2,3.3)
  .. controls (-1.5,3.4) and (-0.4,3.1) .. (0.1,2.4) .. controls (0.5,1.8) and (0.9,1.0) .. (0.9,0) -- cycle;
\draw[popgold, -latex, thick] (-4.0,2.6) -- (-2.6,2.9);
\node[font=\small, align=center] at (0.2,3.4) {shoot bends\\toward light};
\end{scope}
\draw[-latex, very thick, popink] (2.6,2.0) -- (4.2,2.0);
\node[font=\small] at (3.4,2.35) {growth};
\end{tikzpicture}
\end{center}
\legende{Auxin redistribution in a shoot tip exposed to unilateral light. Auxin migrates to the shaded side; greater elongation there curves the shoot toward the light.}
\end{popfigure}

\courssub{The Mechanism of Geotropism (Gravitropism)}

Place a germinating seedling horizontally. Within hours the shoot curves upward (negative geotropism) and the root curves downward (positive geotropism). Both responses are caused by the \emph{same} event — auxin accumulating on the \textbf{lower side} — but with opposite outcomes, because of the differential sensitivity stated in the property above.

\begin{itemize}
\item \textbf{In the shoot:} auxin accumulates on the lower side; the higher concentration \emph{promotes} elongation there, so the lower side grows faster and the shoot bends \emph{upward} (away from gravity).
\item \textbf{In the root:} auxin also accumulates on the lower side, but roots are so sensitive that this concentration \emph{inhibits} elongation; the upper side (less auxin) elongates more, so the root bends \emph{downward} (toward gravity).
\end{itemize}

\begin{popfigure}
\begin{center}
\begin{tikzpicture}[scale=1.0]
% Horizontal seedling
\draw[popgreen, very thick] (-3.5,0) -- (0.5,0);
\draw[popgreen, very thick, fill=popgreenL] (0.5,0) .. controls (1.2,0.25) and (1.6,0.9) .. (1.8,1.8) .. controls (1.5,1.0) and (1.4,0.5) .. (1.1,0.1) -- cycle;
\draw[poporange, very thick, fill=poporange!20] (-3.5,0) .. controls (-4.2,-0.25) and (-4.6,-0.9) .. (-4.8,-1.8) .. controls (-4.5,-1.0) and (-4.4,-0.5) .. (-4.1,-0.1) -- cycle;
% gravity arrow
\draw[popink, very thick, -latex] (2.8,1.6) -- (2.8,0.4);
\node[popink, font=\small] at (2.8,1.95) {gravity};
% auxin labels
\node[poppurple, font=\small, align=center] at (0.4,-0.75) {auxin on lower side:\\shoot cells elongate MORE};
\node[poppurple, font=\small, align=center] at (-3.6,0.85) {auxin on lower side:\\root cells elongate LESS};
% dashed original axis
\draw[popmuted, dashed] (-5.2,0) -- (2.2,0);
\node[font=\small, align=center] at (1.9,2.4) {shoot: negative\\geotropism};
\node[font=\small, align=center] at (-4.9,-2.4) {root: positive\\geotropism};
\end{tikzpicture}
\end{center}
\legende{Geotropic responses of a horizontally placed seedling. Auxin settles on the lower side of both organs: it stimulates elongation in the shoot but inhibits it in the root.}
\end{popfigure}

\begin{exoresolu}[Interpreting a geotropism experiment]
A maize seedling is pinned horizontally on a clinostat and another is pinned horizontally on a stationary board, both in darkness. After 24 hours, describe and explain the appearance of each seedling.
\tcblower
\textbf{Solution.}
\begin{itemize}
\item \textbf{Stationary seedling:} the shoot has curved upward and the root downward. Gravity causes auxin to accumulate on the lower side of each organ. In the shoot this higher auxin concentration stimulates cell elongation, so the lower side grows faster and the shoot bends up (negative geotropism). In the root the same concentration inhibits elongation, so the upper side grows faster and the root bends down (positive geotropism).
\item \textbf{Clinostat seedling:} both shoot and root continue to grow horizontally, in a straight line. The slow rotation continuously changes the orientation of the seedling relative to gravity, so auxin never accumulates on any one side; distribution remains symmetrical and growth is straight. The clinostat acts as a control showing that gravity, not another factor, caused the curvature.
\end{itemize}
\end{exoresolu}

\courssub{Light and Plant Growth: Photomorphogenesis and Etiolation}

Light does more than bend shoots: it controls the whole pattern of plant development, a phenomenon called \cle{photomorphogenesis}.

\textbf{Observation.} Potato tubers sprouting in a dark cupboard produce long, pale, weak shoots with tiny leaves. The same tubers sprouting in daylight produce short, sturdy, green shoots. The dark-grown plant is said to be \cle{etiolated}.

\begin{definition}[Etiolation]
\cle{Etiolation} is the syndrome of abnormal development shown by plants grown in darkness: excessively elongated internodes, pale yellow-white colour (no chlorophyll formed), small unexpanded leaves and a weak stem. On exposure to light, internode elongation slows, chlorophyll forms (greening) and leaves expand.
\end{definition}

Etiolation is adaptive: a seedling germinating deep in the soil invests all its reserves in rapid stem elongation to reach the light before its food runs out. Once in the light, energy is redirected to leaf expansion and photosynthesis.

\begin{savaistu}[Phytochrome: the plant's light switch (exam-useful extra)]
Plants detect light quality and day length using a blue-green pigment called \cle{phytochrome}. It exists in two interconvertible forms: $P_r$ (absorbs red light, $\approx 660$ nm) and $P_{fr}$ (absorbs far-red light, $\approx 730$ nm). Red light converts $P_r$ into the active $P_{fr}$ form; far-red light (or darkness, slowly) converts it back. $P_{fr}$ triggers many light responses: greening, leaf expansion, inhibition of etiolation, seed germination in light-sensitive seeds, and the measurement of night length in flowering (photoperiodism). You are not required to know the details, but quoting phytochrome as the receptor of photomorphogenesis earns credit in extended answers.
\end{savaistu}

\courssub{Practical Applications in Agriculture}

Knowledge of plant hormones is exploited daily, including in Cameroonian agriculture:

\begin{itemize}
\item \textbf{Rooting powders:} cuttings of cassava, hibiscus or ornamental shrubs dipped in synthetic auxin powder (e.g. IBA, NAA) form adventitious roots faster and more reliably, making vegetative propagation cheap and efficient.
\item \textbf{Selective weedkillers:} synthetic auxins such as 2,4-D, applied at high concentration, cause uncontrolled, lethal growth in broad-leaved weeds (dicots) while leaving cereal crops like maize and rice (monocots) largely unaffected — widely used in maize fields in the West and North-West Regions.
\item \textbf{Fruit ripening:} ethene gas is used commercially to ripen harvested bananas, plantains and mangoes uniformly during transport and storage; traders exploit the natural ethene released by ripe fruits placed with unripe ones.
\item \textbf{Seedless (parthenocarpic) fruits:} spraying unpollinated flowers with auxin or gibberellin induces fruit development without fertilisation, giving seedless tomatoes, cucumbers or grapes.
\item \textbf{Breaking dormancy:} gibberellins are used to promote uniform germination of seeds and to induce sprouting of seed yams and potato tubers.
\item \textbf{Delaying senescence:} cytokinins keep harvested vegetables and cut flowers fresh longer.
\end{itemize}

\begin{attention}[Common mistakes with plant hormones]
\begin{itemize}
\item Do not write that auxin ``moves the shoot'' — auxin causes \emph{differential cell elongation}, and the resulting growth bends the organ. Movement is by growth, not by muscles.
\item Do not say auxin is ``attracted by darkness'' — it is actively transported away from the illuminated side (phototropism) or accumulates on the lower side under gravity (geotropism).
\item Remember that the same auxin concentration has \emph{opposite} effects in shoots and roots; quoting the property precisely is what earns full marks.
\item Ethene is a \emph{gas}; it is the only gaseous plant hormone.
\end{itemize}
\end{attention}

\begin{exercice}[Phototropism and hormones]
\begin{enumerate}
\item Define the terms \emph{tropism} and \emph{nastic movement}, giving one example of each.
\item A coleoptile tip is separated from the rest of the shoot by a thin block of agar, and the shoot is illuminated from one side. Predict, with reasons, whether the shoot will bend toward the light.
\item Explain why a high concentration of synthetic auxin kills dicot weeds but is harmless to maize.
\item State two differences between the etiolated form and the light-grown form of a bean seedling, and name the pigment system responsible for detecting light.
\end{enumerate}
\end{exercice}

\begin{corrige}[Exercise 1]
\begin{enumerate}
\item A tropism is a directional growth movement in response to a directional stimulus (e.g. a shoot bending toward light — positive phototropism). A nastic movement is a non-directional, usually reversible movement whose direction is fixed by the organ's structure (e.g. folding of \emph{Mimosa pudica} leaflets on touch).
\item Yes, it will bend. Auxin produced in the tip diffuses down through the agar block (agar is permeable to IAA); unilateral light still causes lateral redistribution of auxin to the shaded side, which elongates more, curving the shoot toward the light. (If an impermeable mica sheet were inserted on the shaded side instead, bending would be prevented.)
\item Synthetic auxins at high dose cause uncontrolled, disorganised growth (abnormal elongation, disrupted transport tissues) that is lethal to sensitive broad-leaved dicots. Maize, a monocot, is far less sensitive at that concentration, so it survives — the basis of selective weedkillers such as 2,4-D.
\item Etiolated: long internodes, pale yellow (no chlorophyll), small unexpanded leaves, weak stem; light-grown: short internodes, green, expanded leaves, sturdy stem. The light-detecting pigment system is phytochrome ($P_r$/$P_{fr}$).
\end{enumerate}
\end{corrige}

\begin{aretenir}[Key points of the section]
\begin{itemize}
\item Tropisms are slow, permanent, directional \emph{growth} movements (photo-, geo-, hydro-, thigmo-, chemotropism); nastic movements are rapid, reversible, non-directional turgor movements.
\item Plant hormones: auxins (cell elongation, tropisms), gibberellins (stem elongation, germination), cytokinins (cell division, delay of senescence), abscisic acid (dormancy, stomatal closure), ethene (fruit ripening, abscission).
\item Phototropism: auxin made at the tip moves to the shaded side $\Rightarrow$ more elongation there $\Rightarrow$ bending toward light.
\item Geotropism: auxin accumulates on the lower side; shoots elongate more there (bend up), roots are inhibited there (bend down) — because roots are far more sensitive to auxin.
\item Light controls development (photomorphogenesis) via phytochrome; dark-grown plants are etiolated.
\item Applications: rooting powders, 2,4-D weedkillers, ethene ripening of plantain and mango, seedless fruits, breaking dormancy — all relevant to Cameroonian farming.
\end{itemize}
\end{aretenir}

\newpage

\coursec{Integration activity}

\coursec{Response \& Coordination in Organisms}

\courssub{Aims of the Activity}
\begin{aretenir}[Learning Outcomes]
By the end of this integration workshop, you should be able to:
\begin{itemize}
    \item Consolidate understanding of the fundamental sequence \textbf{stimulus $\rightarrow$ reception $\rightarrow$ coordination $\rightarrow$ response} across animal and plant systems.
    \item Apply knowledge of taxis, kinesis, tropisms and nastic movements to design and interpret a choice-chamber experiment.
    \item Analyse clinical endocrine disorders (goitre, diabetes mellitus) by identifying the gland, hormone, feedback mechanism failure and physiological consequences.
    \item Distinguish between nervous and hormonal coordination in animals, and between hormonal and electrical coordination in plants.
    \item Develop scientific communication skills by synthesising findings into a structured scientific poster.
\end{itemize}
\end{aretenir}

\courssub{Situation}
\begin{experience}[The GBHS Bamenda Integration Workshop]
\textbf{Context:} You are Upper Sixth Biology students at \emph{Government Bilingual High School (GBHS) Bamenda}. The North West Region offers diverse micro-habitats -- from the humid forests around Lake Awing to the drier grasslands of the Bamenda Highlands. Your teacher, Mr. Nfor, has organised a \emph{Response \& Coordination Integration Workshop} to prepare you for the GCE Advanced Level practical (Paper 3) and theory (Papers 1 \& 2) examinations.

The workshop comprises three linked stations. You will rotate in groups of four, spending 45 minutes at each station. At the end, each group produces an A1 poster titled \textbf{``Stimulus $\rightarrow$ Reception $\rightarrow$ Coordination $\rightarrow$ Response: Unity in Diversity''} linking the three stations.
\end{experience}

\begin{popfigure}
\begin{tikzpicture}[scale=0.9, every node/.style={font=\small}]
% Choice Chamber Diagram
\draw[popdark, thick] (0,0) rectangle (10,5);
\draw[popdark, thick] (5,0) -- (5,5); % Central divider
\draw[popdark, thick] (0,2.5) -- (10,2.5); % Horizontal divider -> 4 chambers

% Labels for conditions
\node[popblue, align=center] at (1.25, 4.2) {Chamber A\\(Dry, Light)};
\node[poporange, align=center] at (6.25, 4.2) {Chamber B\\(Dry, Dark)};
\node[popgreen, align=center] at (1.25, 1.7) {Chamber C\\(Humid, Light)};
\node[poppink, align=center] at (6.25, 1.7) {Chamber D\\(Humid, Dark)};

% Connecting tubes (central)
\draw[popdark, fill=popmuted] (4.7,2.3) rectangle (5.3,2.7); % Central square
\draw[popdark] (4.7,2.5) -- (0,2.5); % Left
\draw[popdark] (5.3,2.5) -- (10,2.5); % Right
\draw[popdark] (5,2.3) -- (5,0); % Down
\draw[popdark] (5,2.7) -- (5,5); % Up

% Woodlice start
\node[popdark, circle, fill=popgold, inner sep=1.5pt] at (5,2.5) {};
\node[align=center] at (5, -0.5) {Start: 20 woodlice\\released at centre};

% Arrows for gradients
\draw[<->, popblue, thick] (0.5, 4.5) -- (4.5, 4.5) node[midway, above] {Light gradient};
\draw[<->, poporange, thick] (5.5, 4.5) -- (9.5, 4.5) node[midway, above] {Dark};
\draw[<->, popgreen, thick] (0.5, 0.5) -- (4.5, 0.5) node[midway, below] {Light gradient};
\draw[<->, poppink, thick] (5.5, 0.5) -- (9.5, 0.5) node[midway, below] {Dark};

\draw[<->, popblue, thick, rotate=90] (2.5, 4.8) -- (2.5, 2.8) node[midway, left] {Dry};
\draw[<->, popgreen, thick, rotate=90] (2.5, 2.2) -- (2.5, 0.2) node[midway, left] {Humid};
\draw[<->, poporange, thick, rotate=90] (7.5, 4.8) -- (7.5, 2.8) node[midway, right] {Dry};
\draw[<->, poppink, thick, rotate=90] (7.5, 2.2) -- (7.5, 0.2) node[midway, right] {Humid};

% Silica gel / Cotton wool symbols
\node[align=center, popmuted] at (1.25, 3.5) {Silica gel\\(Dry)};
\node[align=center, popmuted] at (6.25, 3.5) {Silica gel\\(Dry)};
\node[align=center, popmuted] at (1.25, 1.0) {Moist cotton\\(Humid)};
\node[align=center, popmuted] at (6.25, 1.0) {Moist cotton\\(Humid)};

% Light/Dark covers (simple hatching instead of pattern library)
\draw[popdark, fill=popink!20] (0,5) rectangle (5,4.7);
\draw[popdark, fill=popink!20] (5,5) rectangle (10,4.7);
\node[popink, rotate=90] at (-0.3, 4.85) {Opaque cover (Dark side)};
\end{tikzpicture}
\legende{Station 1: Choice chamber apparatus for investigating woodlouse (\emph{Porcellio scaber}) responses to light and humidity gradients.}
\end{popfigure}

\vspace{0.5cm}
\noindent\textbf{Station 1 -- The Choice Chamber (Animal Behaviour)}\\
Mr. Nfor provides a choice chamber (see diagram) with four chambers creating combined gradients of light (Light vs Dark) and humidity (Dry vs Humid). Twenty woodlice (\emph{Porcellio scaber}, locally called ``kankpe'' in Pidgin) are collected from the school compost heap and released at the central junction. After 15 minutes, the distribution is recorded.

\vspace{0.3cm}
\noindent\textbf{Station 2 -- The Clinic (Human Endocrine Disorders)}\\
Two case files from the \emph{Bamenda Regional Hospital} are presented:

\begin{center}
\begin{tabularx}{\linewidth}{|l|X|X|}\hline
\rowcolor{popblueL}
\textbf{Parameter} & \textbf{Case A: Mr. T., 42 years, Farmer from Santa} & \textbf{Case B: Miss N., 16 years, Student, Bamenda Town} \\\hline
\textbf{Presenting Complaint} & Swelling in anterior neck (6 months), fatigue, cold intolerance, weight gain, constipation. & Polyuria, polydipsia, polyphagia, weight loss, fatigue, blurred vision (3 weeks). \\\hline
\textbf{Signs} & Diffuse, non-tender goitre (Grade 2), bradycardia (56 bpm), dry skin, delayed ankle reflex. & BMI 17.2 kg/m$^2$, random blood glucose 14.2 mmol/L, glucose++ in urine, ketones+. \\\hline
\textbf{Background} & Diet: mainly cassava/fufu, low seafood/iodised salt use. No family history. & Sudden onset. No family history of diabetes. Recent severe malaria episode 4 weeks ago. \\\hline
\end{tabularx}
\end{center}

\vspace{0.3cm}
\noindent\textbf{Station 3 -- The School Garden (Plant Responses)}\\
Two observations are made in the school garden behind the science block:
\begin{enumerate}
    \item \textbf{Maize (\emph{Zea mays}) seedlings:} Planted in a row parallel to the western wall of the laboratory. After two weeks, all seedlings show a distinct curvature of the coleoptile and young leaves \emph{towards} the open field (East), away from the wall.
    \item \textbf{Mimosa pudica (``Shame-shame plant''):} A potted plant is brought to the lab. When a student touches one leaflet, the entire leaf folds inward and the petiole droops within 1--2 seconds. The response spreads to adjacent leaves. After 15 minutes, the leaves reopen.
\end{enumerate}

\courssub{Tasks}

\textbf{Station 1: Choice Chamber Analysis (25 marks)}
\begin{enumerate}
    \item Identify the \textbf{stimuli} (adequate and inadequate) and the \textbf{receptors} involved in the woodlouse response. Classify the response type (taxis, kinesis, or both) with justification. [4 marks]
    \item Predict the expected distribution of the 20 woodlice after 15 minutes in the four chambers (A, B, C, D). Explain your prediction using the concepts of \emph{orthokinesis} and \emph{klinokinesis}. [5 marks]
    \item The actual results are: A=2, B=3, C=12, D=3. Perform a Chi-squared ($\chi^2$) test to determine if the distribution differs significantly from random expectation (equal distribution). Use $p=0.05$, critical value $\chi^2_{crit}(3)=7.82$. State your conclusion. [6 marks]
    \item Explain the survival value of this behaviour for woodlice in the leaf litter of the Bamenda Highlands during the dry season (November--March). [3 marks]
    \item Propose \textbf{two} modifications to the apparatus to investigate the effect of \emph{temperature} gradient independently of humidity. [2 marks]
\end{enumerate}

\textbf{Station 2: Clinical Case Analysis (30 marks)}
\begin{enumerate}
    \item For \textbf{Case A (Mr. T.)}:
    \begin{enumerate}
        \item Identify the affected endocrine gland and the deficient hormone. [2 marks]
        \item Explain the \textbf{negative feedback mechanism} involving the hypothalamus, anterior pituitary and thyroid gland that has failed. Use the terms TRH, TSH, Thyroxine ($T_4$)/Triiodothyronine ($T_3$). Why does the goitre develop despite low thyroxine? [6 marks]
        \item Link the symptoms (cold intolerance, bradycardia, constipation) to the cellular role of thyroid hormone (basal metabolic rate, Na$^+$/K$^+$-ATPase, mitochondrial uncoupling). [4 marks]
    \end{enumerate}
    \item For \textbf{Case B (Miss N.)}:
    \begin{enumerate}
        \item Identify the likely type of diabetes mellitus and the deficient hormone. Explain why ketones appear in urine. [3 marks]
        \item Describe the normal \textbf{negative feedback loop} regulating blood glucose concentration involving insulin and glucagon. [4 marks]
        \item Explain the mechanism of \textbf{polyuria} and \textbf{polydipsia} in uncontrolled diabetes (renal threshold, osmotic diuresis, ADH). [4 marks]
        \item Suggest \textbf{two} immediate management steps and \textbf{one} long-term monitoring test (name the test and target value). [3 marks]
    \end{enumerate}
\end{enumerate}

\textbf{Station 3: Plant Response Analysis (25 marks)}
\begin{enumerate}
    \item For the \textbf{maize seedlings}:
    \begin{enumerate}
        \item Name the type of response shown. Is it a tropism or a nastic movement? Justify. [3 marks]
        \item Identify the stimulus, receptor region, and the hormone responsible. Describe the \textbf{asymmetric distribution} mechanism of this hormone causing curvature (Cholodny-Went theory). [6 marks]
    \end{enumerate}
    \item For the \textbf{Mimosa pudica}:
    \begin{enumerate}
        \item Name the type of response. Distinguish it clearly from the maize response (directionality, speed, mechanism). [4 marks]
        \item Outline the sequence: \emph{Mechanical stimulus $\rightarrow$ Receptor $\rightarrow$ Electrical signal (Action Potential) $\rightarrow$ Effector (Pulvinus) $\rightarrow$ Response}. Mention the role of K$^+$ and Cl$^-$ ion fluxes and water potential ($\Psi_w$) changes in the motor cells. [6 marks]
    \end{enumerate}
    \item Compare and contrast the \textbf{coordination systems} in animals (nervous + endocrine) and plants (hormonal + electrical) using examples from Stations 1, 2 and 3. [6 marks]
\end{enumerate}

\textbf{Synthesis Task: The Poster (20 marks)}
\begin{enumerate}
    \item Design the layout of your group's A1 poster. It must contain:
    \begin{itemize}
        \item A central flow diagram: \textbf{Stimulus $\rightarrow$ Reception $\rightarrow$ Coordination $\rightarrow$ Effector $\rightarrow$ Response}.
        \item Three branches (Animal Behaviour, Human Endocrine, Plant Responses) feeding into the central flow.
        \item For each branch: Specific example, Receptor type, Coordination type (Nervous/Hormonal/Electrical), Effector, Response type.
        \item A ``Key Concepts'' box defining: Taxis, Kinesis, Tropism, Nastic movement, Negative Feedback, Homeostasis.
    \end{itemize}
    \item Prepare a 3-minute oral defence: ``How does the woodlouse finding a humid crack, the thyroxine-deficient patient developing goitre, and the maize seedling bending to light all illustrate the same fundamental biological principle?'' [Outline your answer in 150 words.]
\end{enumerate}

\courssub{Model Answers}

\textbf{Station 1: Choice Chamber Analysis}

\begin{enumerate}
    \item \textbf{Stimuli \& Receptors:}
    \begin{itemize}
        \item \textbf{Light:} Adequate stimulus for photoreceptors (compound eyes / ocelli). Woodlice are negatively phototactic (move away from light).
        \item \textbf{Humidity:} Adequate stimulus for hygroreceptors (located on antennae, specifically the flagellum). Woodlice are positively hydrokinetic (orthokinesis: speed decreases in humid air; klinokinesis: turning rate increases in dry air).
    \end{itemize}
    \textbf{Response Classification:} \textbf{Both taxis and kinesis.}
    \begin{itemize}
        \item \emph{Negative phototaxis:} Directed movement \emph{away} from the directional light source (gradient).
        \item \emph{Hygrokinesis (Ortho- + Klinokinesis):} Non-directed movement where \emph{rate of movement} (orthokinesis) and \emph{rate of turning} (klinokinesis) depend on humidity intensity. In dry air: high speed, high turning $\rightarrow$ rapid exit. In humid air: low speed, low turning $\rightarrow$ aggregation.
    \end{itemize}

    \item \textbf{Predicted Distribution:}
    Woodlice seek \textbf{Dark + Humid} (Chamber D) ideally. However, the chamber creates a conflict: Chamber C (Humid, Light) vs Chamber D (Humid, Dark).
    \begin{itemize}
        \item \textbf{Chamber D (Humid, Dark):} Highest aggregation. Both needs met. Low speed, low turning $\rightarrow$ trap.
        \item \textbf{Chamber C (Humid, Light):} High humidity retains them (low orthokinesis), but light causes negative phototaxis $\rightarrow$ movement towards dark partition. Some accumulate near partition.
        \item \textbf{Chambers A \& B (Dry):} High orthokinesis (speed) and klinokinesis (turning) $\rightarrow$ rapid escape. Very few remain.
    \end{itemize}
    \emph{Predicted ranking:} D $>$ C $\gg$ B $\approx$ A.

    \item \textbf{Chi-squared Test:}
    \begin{align*}
    H_0: &\text{ Woodlice are distributed randomly (no preference). Expected } (E) = 20/4 = 5 \text{ per chamber.} \\
    H_1: &\text{ Distribution is not random (preference exists).} \\
    \chi^2 &= \sum \frac{(O-E)^2}{E} \\
    &= \frac{(2-5)^2}{5} + \frac{(3-5)^2}{5} + \frac{(12-5)^2}{5} + \frac{(3-5)^2}{5} \\
    &= \frac{9}{5} + \frac{4}{5} + \frac{49}{5} + \frac{4}{5} = 1.8 + 0.8 + 9.8 + 0.8 = \mathbf{13.2}
    \end{align*}
    Degrees of freedom $df = n-1 = 3$. Critical value $\chi^2_{0.05}(3) = 7.82$.
    Since $13.2 > 7.82$, \textbf{reject $H_0$}. The distribution is highly significant ($p < 0.01$). Woodlice show a strong preference for humid conditions (Chamber C+D = 15 vs A+B = 5) and within humid, a preference for dark (D=3 vs C=12 -- see Note below).

    \begin{attention}[Interpreting C vs D]
    The data shows C (12) $>$ D (3). This suggests \textbf{humidity is the dominant drive (kinesis)} overriding the light taxis in this specific apparatus, or the light gradient was weak. A good student discusses this: ``Humidity kinesis traps them in C before phototaxis moves them to D''. This is excellent evaluation.
    \end{attention}

    \item \textbf{Survival Value (Dry Season Bamenda):}
    During the dry season (Harmattan), relative humidity drops below 30\%, desiccation risk is extreme. Woodlice lack a waxy cuticle and lose water rapidly via transpiration and excretion.
    \begin{itemize}
        \item \textbf{Hygrokinesis} ensures they spend maximum time in moist microhabitats (under logs, deep leaf litter, soil cracks) where $\Psi_w$ is high.
        \item \textbf{Negative phototaxis} keeps them under cover (leaf litter, bark) avoiding UV damage and surface heating, which lowers humidity further.
        \item Combined: Efficient location of the ``humid dark'' refuge maximises water conservation and survival.
    \end{itemize}

    \item \textbf{Modifications for Temperature Gradient:}
    \begin{enumerate}
        \item Use a \textbf{Peltier element} or water jackets (circulating water baths) at each end of a linear chamber to create a stable thermal gradient, replacing the light/dark covers.
        \item \textbf{Control humidity uniformly}: Saturate the entire chamber floor with moist cotton wool / agar gel so humidity is constant (100\% RH) and not a confounding variable. Use a data logger with thermistors to verify the gradient.
    \end{enumerate}
\end{enumerate}

\textbf{Station 2: Clinical Case Analysis}

\begin{enumerate}
    \item \textbf{Case A: Mr. T. (Primary Hypothyroidism / Endemic Goitre)}
    \begin{enumerate}
        \item \textbf{Gland:} Thyroid gland (anterior neck). \textbf{Hormone:} Thyroxine ($T_4$) / Triiodothyronine ($T_3$). Deficiency due to dietary iodine lack (Santa highlands, cassava diet -- goitrogens).
        \item \textbf{Feedback Failure:}
        \begin{center}
        \begin{tikzpicture}[node distance=2.0cm, >=stealth, auto, thick, every node/.style={font=\small}]
        \node (hypo) [draw, rounded corners, fill=popblueL] {Hypothalamus};
        \node (pit) [draw, rounded corners, fill=popblueL, right of=hypo] {Ant. Pituitary};
        \node (thy) [draw, rounded corners, fill=popblueL, right of=pit] {Thyroid Gland};
        \node (blood) [draw, rounded corners, fill=popgreenL, below of=pit] {Blood $T_3/T_4$};
        \draw[->] (hypo) -- node {TRH} (pit);
        \draw[->] (pit) -- node {TSH} (thy);
        \draw[->] (thy) -- node {$T_3/T_4$} (blood);
        \draw[->, popred, dashed] (blood) -- node[near start, left] {Negative Feedback \textbf{FAILS}} (hypo);
        \draw[->, popred, dashed] (blood) -- (pit);
        \end{tikzpicture}
        \end{center}
        Low $T_3/T_4$ $\rightarrow$ No inhibition of Hypothalamus/Pituitary $\rightarrow$ \textbf{High TRH \& High TSH}. High TSH causes \textbf{thyroid hyperplasia/hypertrophy} (goitre) attempting to trap more iodide. But without iodine, hormone synthesis fails $\rightarrow$ colloid accumulation $\rightarrow$ diffuse goitre.
        \item \textbf{Symptoms $\leftrightarrow$ Cellular Role:}
        \begin{itemize}
            \item \textbf{Cold intolerance / Low BMR:} $T_3$ upregulates Na$^+$/K$^+$-ATPase (consumes ATP $\rightarrow$ heat) and mitochondrial uncoupling proteins (UCP1) $\rightarrow$ thermogenesis. Deficiency $\rightarrow$ reduced heat production.
            \item \textbf{Bradycardia:} $T_3$ increases $\beta$-adrenergic receptor density and sino-atrial node automaticity. Deficiency $\rightarrow$ reduced heart rate/contractility.
            \item \textbf{Constipation:} $T_3$ maintains gut smooth muscle tone and peristalsis. Deficiency $\rightarrow$ reduced motility.
        \end{itemize}
    \end{enumerate}

    \item \textbf{Case B: Miss N. (Type 1 Diabetes Mellitus -- likely autoimmune triggered by viral infection/malaria stress)}
    \begin{enumerate}
        \item \textbf{Type 1 DM (Insulin Dependent).} \textbf{Deficient hormone: Insulin} (from $\beta$-cells of Islets of Langerhans). \textbf{Ketones:} Insulin lack $\rightarrow$ uncontrolled lipolysis in adipose tissue (Hormone-Sensitive Lipase active) $\rightarrow$ high free fatty acids $\rightarrow$ liver $\beta$-oxidation $\rightarrow$ acetyl-CoA $\rightarrow$ ketogenesis (acetoacetate, $\beta$-hydroxybutyrate, acetone). Ketones are acidic $\rightarrow$ metabolic ketoacidosis (DKA risk).
        \item \textbf{Glucose Homeostasis Feedback Loop:}
        \begin{center}
        \begin{tikzpicture}[node distance=1.8cm, >=stealth, auto, thick, every node/.style={font=\small}]
        \node (high) [draw, rounded corners, fill=poporange] {High Blood Glucose};
        \node (beta) [draw, rounded corners, fill=popblueL, right of=high] {$\beta$-cells $\rightarrow$ Insulin};
        \node (target) [draw, rounded corners, fill=popgreenL, below of=beta] {Liver/Muscle/Adipose};
        \node (low) [draw, rounded corners, fill=poporange, left of=high] {Low Blood Glucose};
        \node (alpha) [draw, rounded corners, fill=popblueL, left of=low] {$\alpha$-cells $\rightarrow$ Glucagon};
        \draw[->] (high) -- (beta);
        \draw[->] (beta) -- node {$\uparrow$ Glucose uptake, $\uparrow$ Glycogenesis, $\downarrow$ Gluconeogenesis} (target);
        \draw[->] (target) -- node {$\downarrow$ Blood Glucose} (high);
        \draw[->] (low) -- (alpha);
        \draw[->] (alpha) -- node {$\uparrow$ Glycogenolysis, $\uparrow$ Gluconeogenesis} (target);
        \draw[->] (target) -- node {$\uparrow$ Blood Glucose} (low);
        \end{tikzpicture}
        \end{center}
        \item \textbf{Polyuria \& Polydipsia Mechanism:}
        \begin{enumerate}
            \item Hyperglycaemia ($> 10$ mmol/L) exceeds \textbf{renal threshold} ($\approx 10$ mmol/L).
            \item Glucose not fully reabsorbed in PCT (SGLT2 saturation) $\rightarrow$ remains in tubular lumen.
            \item \textbf{Osmotic diuresis:} Glucose lowers tubular fluid $\Psi_w$ $\rightarrow$ water retained in lumen $\rightarrow$ large volume urine (polyuria).
            \item Volume depletion $\rightarrow$ increased plasma osmolarity $\rightarrow$ hypothalamic osmoreceptors stimulate \textbf{thirst centre} $\rightarrow$ polydipsia. Also ADH released but overwhelmed by osmotic load.
        \end{enumerate}
        \item \textbf{Management:}
        \begin{itemize}
            \item \textbf{Immediate:} 1. IV Fluids (0.9\% NaCl) for rehydration + IV Insulin infusion (sliding scale) to lower glucose/ketones. 2. Potassium replacement (insulin drives K$^+$ into cells $\rightarrow$ hypokalaemia risk). Monitor blood gases (pH, bicarbonate).
            \item \textbf{Long-term monitoring:} \textbf{HbA$_{1c}$ (Glycated Haemoglobin)}. Target $< 48$ mmol/mol ($< 6.5\%$) for good control; $< 58$ mmol/mol ($7.5\%$) acceptable in adolescents.
        \end{itemize}
    \end{enumerate}
\end{enumerate}

\textbf{Station 3: Plant Response Analysis}

\begin{enumerate}
    \item \textbf{Maize Seedlings (Phototropism)}
    \begin{enumerate}
        \item \textbf{Positive Phototropism.} It is a \textbf{tropism} (directional growth response where direction of response is determined by direction of stimulus -- light from East).
        \item \textbf{Stimulus:} Unilateral blue light (approx. 450 nm). \textbf{Receptor:} \emph{Phototropins} (phot1, phot2) in the \textbf{coleoptile tip} (and young leaves). \textbf{Hormone:} \textbf{Auxin (IAA -- Indole-3-Acetic Acid)}.
        \item \textbf{Cholodny-Went Mechanism:}
        \begin{enumerate}
            \item Light hits lit side $\rightarrow$ Phototropin activation $\rightarrow$ Phosphorylation of PIN proteins (auxin efflux carriers).
            \item PINs relocalise to lateral membranes $\rightarrow$ Auxin transported \emph{laterally} away from light to \textbf{shaded side}.
            \item \textbf{Asymmetric distribution:} Higher [Auxin] on shaded side.
            \item In shoots: High [Auxin] $\rightarrow$ stimulates \textbf{cell elongation} (via H$^+$-ATPase $\rightarrow$ wall acidification $\rightarrow$ expansin activation $\rightarrow$ wall loosening $\rightarrow$ water influx $\rightarrow$ turgor-driven extension).
            \item Shaded side elongates faster $\rightarrow$ curvature \emph{towards} light.
        \end{enumerate}
    \end{enumerate}

    \item \textbf{Mimosa pudica (Thigmonasty / Seismonasty)}
    \begin{enumerate}
        \item \textbf{Thigmonasty (Seismonasty).} \textbf{Non-directional} (nastic) response to touch (direction of response -- folding -- is fixed by plant anatomy, not stimulus direction). \textbf{Speed:} Very fast (1--2 s) via \textbf{action potentials (APs)}, not growth. \textbf{Mechanism:} Turgor loss in pulvinus motor cells (electro-chemical), not differential growth.
        \item \textbf{Sequence:}
        \begin{enumerate}
            \item \textbf{Stimulus:} Mechanical deformation of mechanoreceptors (hair cells) on leaflets.
            \item \textbf{Receptor Potential:} Depolarisation $\rightarrow$ if threshold reached $\rightarrow$ \textbf{Action Potential (AP)} generated.
            \item \textbf{Transmission:} AP propagates via \textbf{phloem (sieve tubes)} at $\approx$ few cm/s (electrical + chemical/H$^+$ waves) to \textbf{pulvini} (swollen leaf bases).
            \item \textbf{Effector (Pulvinus):} AP reaches extensor/flexor motor cells.
                \begin{itemize}
                    \item \textbf{Flexor side (lower):} AP $\rightarrow$ Voltage-gated K$^+$ \& Cl$^-$ channels OPEN $\rightarrow$ massive efflux of K$^+$, Cl$^-$ (and malate$^{2-}$) $\rightarrow$ $\Psi_s$ increases (less negative) $\rightarrow$ $\Psi_w$ increases $\rightarrow$ Water leaves cells (exosmosis) $\rightarrow$ \textbf{Turgor loss} $\rightarrow$ Cells shrink.
                    \item \textbf{Extensor side (upper):} Less affected or gains water.
                \end{itemize}
            \item \textbf{Response:} Flexor collapse $\rightarrow$ Leaf folds, petiole droops.
            \item \textbf{Recovery (15 min):} Active H$^+$-pump / K$^+$ influx $\rightarrow$ $\Psi_w$ drops $\rightarrow$ Water enters $\rightarrow$ Turgor restored.
        \end{enumerate}
    \end{enumerate}

    \item \textbf{Comparison of Coordination Systems}
    \begin{center}
    \begin{tabularx}{\linewidth}{|l|X|X|X|}\hline
    \rowcolor{popblueL}
    \textbf{Feature} & \textbf{Animal Nervous (Station 1)} & \textbf{Animal Endocrine (Station 2)} & \textbf{Plant (Station 3)} \\\hline
    \textbf{Signal Type} & Electrical (Action Potentials) + Chemical (Neurotransmitters) & Chemical (Hormones in blood) & Chemical (Hormones: Auxin) + Electrical (APs in Mimosa) \\\hline
    \textbf{Speed} & Very fast (ms) -- Woodlouse escape & Slow (mins-hours) -- Thyroxine/Insulin effects & Slow (hours) -- Auxin growth; Fast (s) -- Mimosa AP \\\hline
    \textbf{Target Specificity} & High (specific muscle/gland) & Low (all cells with receptor) & Variable (Auxin: specific tissues; AP: whole leaf) \\\hline
    \textbf{Duration} & Short (ms-s) & Long (days-weeks) & Long (growth) / Short (turgor) \\\hline
    \textbf{Example} & Woodlouse hygroreceptor $\rightarrow$ Interneuron $\rightarrow$ Leg muscles & Hypothalamus $\rightarrow$ Pituitary $\rightarrow$ Thyroid $\rightarrow$ Body cells & Coleoptile tip $\rightarrow$ Auxin transport $\rightarrow$ Shaded cells \\\hline
    \textbf{Integration} & CNS (Brain/Spinal cord) & Feedback loops (Hypothalamo-pituitary axis) & No CNS; Apical dominance, vascular transport \\\hline
    \end{tabularx}
    \end{center}
    \textbf{Key Unifying Principle:} All systems detect a \textbf{change (stimulus)} via a \textbf{specific receptor}, transmit information via a \textbf{coordination pathway} (neural, hormonal, electrical), and activate an \textbf{effector} to produce a \textbf{response} that restores homeostasis or enhances fitness.
\end{enumerate}

\textbf{Synthesis Task: The Poster}

\begin{enumerate}
    \item \textbf{Poster Layout Description:}
    \begin{itemize}
        \item \textbf{Title Banner:} ``Stimulus $\rightarrow$ Reception $\rightarrow$ Coordination $\rightarrow$ Effector $\rightarrow$ Response: Unity in Diversity'' (GBHS Bamenda, Upper Sixth Biology).
        \item \textbf{Central Column (Vertical Flow):} Large arrows connecting 5 boxes: \textsc{Stimulus} (Light, Touch, Hormone level, Humidity) $\downarrow$ \textsc{Reception} (Photoreceptor, Mechanoreceptor, Osmoreceptor, Chemoreceptor) $\downarrow$ \textsc{Coordination} (Nervous System, Endocrine System, Phloem/Electrical) $\downarrow$ \textsc{Effector} (Muscle/Gland, Target Cell, Pulvinus/Meristem) $\downarrow$ \textsc{Response} (Taxis/Kinesis, Metabolic change, Growth/Turgor).
        \item \textbf{Left Branch (Animal Behaviour):} Photo of Choice Chamber. Boxes: Stimulus: Light/Humidity $\rightarrow$ Receptor: Eye/Antenna $\rightarrow$ Coordination: Ventral Nerve Cord $\rightarrow$ Effector: Leg Muscles $\rightarrow$ Response: Kinesis/Taxis.
        \item \textbf{Middle Branch (Human Endocrine):} Diagram of H-P-Thyroid axis \& Glucose loop. Boxes: Stimulus: Low $T_3$/High Glucose $\rightarrow$ Receptor: Hypothalamic neurons / $\beta$-cells $\rightarrow$ Coordination: Hormonal (TSH/Insulin) $\rightarrow$ Effector: Thyroid / Liver/Adipose $\rightarrow$ Response: Metabolic Rate / Glucose Uptake. \textbf{Red Alert Box:} ``Feedback Failure = Disease (Goitre, Diabetes)''.
        \item \textbf{Right Branch (Plant Responses):} Split: Top: Maize (Tropism). Stimulus: Unilateral Light $\rightarrow$ Receptor: Phototropin (Tip) $\rightarrow$ Coordination: Auxin Transport $\rightarrow$ Effector: Shaded Cortex $\rightarrow$ Response: Positive Phototropism. Bottom: Mimosa (Nasty). Stimulus: Touch $\rightarrow$ Receptor: Mechanoreceptor $\rightarrow$ Coordination: Action Potential (Phloem) $\rightarrow$ Effector: Pulvinus Motor Cells $\rightarrow$ Response: Thigmonasty.
        \item \textbf{Bottom Key Concepts Box:} Definitions of Taxis, Kinesis, Tropism, Nastic Movement, Negative Feedback, Homeostasis, Receptor, Effector.
    \end{itemize}

    \item \textbf{Oral Defence Outline (150 words):}
    \emph{``All three scenarios illustrate the universal biological principle of \textbf{Information Processing for Survival}. In the choice chamber, the woodlouse detects a desiccation threat (low humidity stimulus) via antennal hygroreceptors; its ventral nerve cord coordinates a kinesis/taxis response, moving legs to carry it to a humid refuge -- immediate behavioural homeostasis. In the clinic, Mr. T.'s thyroid fails to transduce the metabolic 'low energy' signal into thyroxine due to iodine lack; the pituitary over-secretes TSH (failed negative feedback), causing goitre -- a chronic failure of hormonal homeostasis. In the garden, the maize seedling detects directional light via tip phototropins; auxin redistribution coordinates differential growth, bending the shoot toward energy -- developmental homeostasis. Though the signals (nerves, blood hormones, phloem sap/APs) and timescales (seconds, days, hours) differ, the logic is identical: \textbf{Receptor $\rightarrow$ Transducer $\rightarrow$ Coordinator $\rightarrow$ Effector $\rightarrow$ Adaptive Response}. This unity -- from kankpe in the compost to a student in the ward -- is the essence of physiology.''}
\end{enumerate}

\newpage

\coursec{Methods and worked examples}

\coursec{Response and Coordination in Organisms}

\courssub{Skill 1: Identifying and Interpreting a Reflex Arc}

\begin{definition}[Reflex Action]
A \cle{reflex action} is a rapid, automatic, involuntary and stereotyped response to a stimulus, mediated by a \cle{reflex arc} that bypasses conscious brain centres for speed. It is protective and inborn (not learnt).
\end{definition}

\begin{propriete}[Components of the Reflex Arc]
The reflex arc consists of five essential components in sequence:
\begin{enumerate}
\item \textbf{Receptor:} detects the stimulus (e.g. free nerve endings for pain, muscle spindle for stretch).
\item \textbf{Sensory (afferent) neurone:} carries the impulse from receptor to CNS; cell body in \cle{dorsal root ganglion}.
\item \textbf{Relay (intermediate) neurone:} lies entirely within the CNS (spinal cord or brainstem); integrates the signal.
\item \textbf{Motor (efferent) neurone:} carries the impulse from CNS to effector; cell body in \cle{ventral horn} of spinal cord grey matter.
\item \textbf{Effector:} muscle (contracts) or gland (secretes) producing the response.
\end{enumerate}
\end{propriete}

\begin{methode}[Analysing a Reflex Arc Diagram]
\begin{enumerate}
\item \textbf{Locate the receptor.} Identify the structure in contact with the stimulus (skin, muscle spindle, sense organ). The stimulus arrow points \emph{towards} the receptor.
\item \textbf{Trace the afferent pathway.} Follow the sensory neurone from receptor to CNS. Note the dorsal root ganglion (swelling on dorsal root).
\item \textbf{Identify the integration centre.} In the spinal cord, the grey matter (butterfly-shaped) contains the synapse between sensory and relay neurone, and relay to motor neurone.
\item \textbf{Trace the efferent pathway.} Follow the motor neurone from ventral horn, out through ventral root, to the effector.
\item \textbf{Identify the effector.} Muscle or gland.
\item \textbf{State the sequence:} stimulus $\rightarrow$ receptor $\rightarrow$ sensory neurone $\rightarrow$ relay neurone $\rightarrow$ motor neurone $\rightarrow$ effector $\rightarrow$ response.
\end{enumerate}
\textbf{Key anatomical landmarks:} Dorsal root = sensory (afferent), ventral root = motor (efferent). Dorsal root ganglion contains sensory cell bodies. Grey matter = cell bodies and synapses; white matter = myelinated axons.
\end{methode}

\begin{attention}[Common Mistakes]
\begin{itemize}
\item Confusing dorsal (sensory) and ventral (motor) roots.
\item Thinking the relay neurone cell body is in a ganglion (it is in the CNS grey matter).
\item Forgetting that a reflex is \emph{involuntary} -- the brain is informed \emph{after} the response.
\item Calling the knee-jerk a "withdrawal reflex" (it is a \cle{stretch reflex}, monosynaptic; withdrawal is polysynaptic).
\end{itemize}
\end{attention}

\begin{savaistu}[Cameroon Context]
In rural clinics, testing the \cle{knee-jerk reflex} (patellar reflex) and \cle{ankle-jerk reflex} helps assess spinal cord integrity (L2--L4 and S1--S2) after trauma from road accidents on the Douala--Yaound\'e highway. Absent reflexes may indicate spinal shock or peripheral neuropathy (e.g. from poorly controlled diabetes).
\end{savaistu}

\begin{exoresolu}[The Knee-Jerk and the Withdrawal Reflex]
A student accidentally touches a hot object in the laboratory and immediately withdraws her hand before feeling any pain.\par
(a) Name the type of action described. \hfill (1 mark)\par
(b) List, in order, the five components of the nervous pathway involved. \hfill (3 marks)\par
(c) Explain why the hand is withdrawn \emph{before} pain is consciously felt. \hfill (2 marks)\par
(d) State two characteristics of this type of action. \hfill (2 marks)
\tcblower
\textbf{(a)} This is a \textbf{reflex action} (specifically, a spinal \textbf{withdrawal reflex}, polysynaptic).\par
\textbf{(b)} The five components, in order, are:
\begin{enumerate}
\item \textbf{Receptor:} pain/heat receptors (free nerve endings) in the skin of the finger;
\item \textbf{Sensory (afferent) neurone:} carries the impulse from the receptor to the spinal cord via the dorsal root;
\item \textbf{Relay (intermediate) neurone:} in the grey matter of the spinal cord, transmits the impulse across a synapse to the motor neurone;
\item \textbf{Motor (efferent) neurone:} carries the impulse from the spinal cord via the ventral root to the effector;
\item \textbf{Effector:} the flexor muscles (e.g. biceps brachii) of the arm, which contract to withdraw the hand.
\end{enumerate}
\textbf{(c)} The reflex pathway passes only through the \textbf{spinal cord} (three neurones, two synapses), which is short and fast. The impulse travelling \emph{up} to the brain (where pain is consciously perceived) follows a longer ascending pathway with more synapses, so it arrives \emph{after} the withdrawal has already occurred. This delay is protective: the body is removed from danger before the brain even registers the pain.\par
\textbf{(d)} Any two of: it is \textbf{rapid}; \textbf{automatic/involuntary} (does not require conscious thought); \textbf{stereotyped} (the same stimulus always gives the same response); \textbf{protective} in function; the pathway is \textbf{inborn} (not learnt).
\end{exoresolu}

\begin{experience}[Testing the Knee-Jerk Reflex]
\textbf{Aim:} Demonstrate a monosynaptic stretch reflex.\par
\textbf{Method:} Subject sits with legs hanging freely. Examiner taps the patellar tendon sharply with a reflex hammer.\par
\textbf{Observation:} Sudden extension of the leg (quadriceps contraction).\par
\textbf{Explanation:} Tap stretches quadriceps muscle $\rightarrow$ muscle spindles stimulated $\rightarrow$ sensory neurone $\rightarrow$ synapse directly on motor neurone in spinal cord (monosynaptic) $\rightarrow$ motor neurone $\rightarrow$ quadriceps contracts $\rightarrow$ leg extends. No interneurone involved, hence very fast.
\end{experience}

\courssub{Skill 2: Calculating the Speed of Nerve Impulse Conduction}

\begin{propriete}[Factors Affecting Conduction Speed]
\begin{itemize}
\item \textbf{Myelination:} Myelin sheath (Schwann cells in PNS, oligodendrocytes in CNS) insulates the axon, forcing depolarisation to jump between \cle{nodes of Ranvier} (\cle{saltatory conduction}). This dramatically increases speed (up to $100\ \mathrm{m\,s^{-1}}$ in humans).
\item \textbf{Axon diameter:} Larger diameter $\rightarrow$ lower axial resistance $\rightarrow$ faster conduction.
\item \textbf{Temperature:} Higher temperature increases ion diffusion rates and channel kinetics, increasing speed (Q10 effect).
\end{itemize}
\end{propriete}

\begin{methode}[Speed of Transmission Along a Neurone]
\begin{enumerate}
\item \textbf{Write the formula:} speed $v = \dfrac{d}{t}$, where $d$ is the distance travelled by the impulse and $t$ is the time taken.
\item \textbf{Convert all units} to metres (m) and seconds (s) before substituting: $1\ \mathrm{mm} = 10^{-3}\ \mathrm{m}$; $1\ \mathrm{ms} = 10^{-3}\ \mathrm{s}$.
\item \textbf{Substitute and compute.} Express the answer in $\mathrm{m\,s^{-1}}$, in standard form if appropriate.
\item \textbf{Interpret the result:} compare myelinated vs non-myelinated fibres; note that synapses introduce delays (about $0.5$--$1\ \mathrm{ms}$ each).
\item \textbf{If several synapses are involved,} subtract total synaptic delay from the measured time before calculating conduction speed along the axons.
\end{enumerate}
\end{methode}

\begin{attention}[Unit Conversion Traps]
\begin{itemize}
\item Always convert \emph{before} dividing: $90\ \mathrm{mm} = 0.09\ \mathrm{m}$, not $90/1000$ inside the formula.
\item $1\ \mathrm{ms} = 0.001\ \mathrm{s}$, not $0.01\ \mathrm{s}$.
\item Speed unit is $\mathrm{m\,s^{-1}}$ (metres per second), not $\mathrm{mm\,ms^{-1}}$.
\end{itemize}
\end{attention}

\begin{exoresolu}[Conduction Speed in a Reflex Arc]
In an experiment, an impulse was found to travel along a myelinated sensory axon of length $90\ \mathrm{mm}$ in $1.5\ \mathrm{ms}$.\par
(a) Calculate the speed of conduction of the impulse along this axon. \hfill (3 marks)\par
(b) The same impulse takes $4.5\ \mathrm{ms}$ in total to cross the whole reflex arc, which contains two synapses. If each synapse causes a delay of $1.0\ \mathrm{ms}$, calculate the time spent in actual conduction along the neurones, and comment. \hfill (3 marks)\par
(c) State two structural features of a neurone that increase the speed of conduction. \hfill (2 marks)
\tcblower
\textbf{(a)} Convert units first:
\[ d = 90\ \mathrm{mm} = 90 \times 10^{-3}\ \mathrm{m} = 0.09\ \mathrm{m}; \qquad t = 1.5\ \mathrm{ms} = 1.5 \times 10^{-3}\ \mathrm{s}. \]
Apply the formula:
\[ v = \frac{d}{t} = \frac{0.09}{1.5 \times 10^{-3}} = \frac{0.09}{0.0015} = 60\ \mathrm{m\,s^{-1}}. \]
The speed of conduction is $\boxed{60\ \mathrm{m\,s^{-1}}}$, a typical value for a myelinated mammalian neurone.\par
\textbf{(b)} Total synaptic delay $= 2 \times 1.0 = 2.0\ \mathrm{ms}$.\par
Time spent in conduction $= 4.5 - 2.0 = 2.5\ \mathrm{ms}$.\par
\textbf{Comment:} nearly half of the total reflex time ($2.0$ out of $4.5\ \mathrm{ms}$) is lost at the synapses. Synaptic transmission (release, diffusion and action of neurotransmitter) is much slower than conduction along an axon; synapses are therefore the main factor limiting the speed of reflexes.\par
\textbf{(c)} Any two of:
\begin{itemize}
\item presence of a \textbf{myelin sheath} (allows saltatory conduction between nodes of Ranvier);
\item \textbf{large axon diameter} (lower internal resistance to ion flow);
\item presence of regularly spaced \textbf{nodes of Ranvier}.
\end{itemize}
\end{exoresolu}

\courssub{Skill 3: Explaining Events at the Synapse}

\begin{definition}[Synapse]
A \cle{synapse} is a specialised junction between two neurones (or a neurone and an effector) where an electrical impulse triggers the release of a chemical neurotransmitter, which diffuses across the \cle{synaptic cleft} and binds to receptors on the post-synaptic membrane, potentially initiating a new impulse.
\end{definition}

\begin{methode}[Describing Synaptic Transmission in Sequence (Cholinergic Synapse)]
\begin{enumerate}
\item \textbf{Arrival of the impulse:} the action potential reaches the \cle{pre-synaptic knob}; depolarisation opens voltage-gated \cle{calcium channels}.
\item \textbf{Calcium entry:} $\mathrm{Ca^{2+}}$ ions diffuse into the knob down their electrochemical gradient.
\item \textbf{Vesicle migration and fusion:} $\mathrm{Ca^{2+}}$ causes \cle{synaptic vesicles} (containing neurotransmitter, e.g. \cle{acetylcholine}) to move to and fuse with the pre-synaptic membrane.
\item \textbf{Release of transmitter:} the neurotransmitter is released by \emph{exocytosis} into the \cle{synaptic cleft}.
\item \textbf{Diffusion:} the transmitter diffuses across the cleft (this causes the \cle{synaptic delay} of $\approx 0.5$--$1\ \mathrm{ms}$).
\item \textbf{Binding:} transmitter molecules bind to specific \cle{receptor proteins} on the post-synaptic membrane (ligand-gated ion channels).
\item \textbf{New depolarisation (EPSP):} sodium channels open, $\mathrm{Na^{+}}$ enters, causing an \cle{excitatory post-synaptic potential (EPSP)}; if threshold is reached, a new \cle{action potential} is generated in the post-synaptic neurone.
\item \textbf{Removal of transmitter:} the enzyme \cle{acetylcholinesterase} (AChE) rapidly hydrolyses acetylcholine to choline and acetate, terminating the signal and allowing the synapse to reset.
\end{enumerate}
\textbf{One-way transmission:} vesicles and $\mathrm{Ca^{2+}}$ channels are only on the pre-synaptic side; receptors and AChE are only on the post-synaptic side.
\end{methode}

\begin{propriete}[Excitatory vs Inhibitory Synapses]
\begin{itemize}
\item \textbf{Excitatory (e.g. acetylcholine at neuromuscular junction):} opens $\mathrm{Na^{+}}$/$\mathrm{Ca^{2+}}$ channels $\rightarrow$ depolarisation (EPSP).
\item \textbf{Inhibitory (e.g. GABA, glycine in CNS):} opens $\mathrm{Cl^{-}}$ or $\mathrm{K^{+}}$ channels $\rightarrow$ hyperpolarisation (IPSP), making the post-synaptic neurone less likely to fire.
\end{itemize}
\cle{Spatial summation} (multiple simultaneous inputs) and \cle{temporal summation} (rapid successive inputs) determine whether threshold is reached.
\end{propriete}

\begin{attention}[Synapse Exam Traps]
\begin{itemize}
\item Forgetting $\mathrm{Ca^{2+}}$ entry -- it is the essential trigger for vesicle fusion.
\item Saying "impulse jumps across the gap" -- it is chemical transmission, not electrical (except rare gap junctions).
\item Confusing AChE (breaks down ACh) with ACh (the transmitter).
\item Not mentioning that the post-synaptic potential is \emph{graded}, not all-or-none; only the resulting action potential is all-or-none.
\end{itemize}
\end{attention}

\begin{savaistu}[Toxins and Synapses]
Many poisons target synapses: \textbf{curare} (arrow poison) blocks ACh receptors $\rightarrow$ paralysis; \textbf{strychnine} blocks glycine receptors $\rightarrow$ convulsions; \textbf{organophosphates} (insecticides, nerve gases) inhibit AChE $\rightarrow$ continuous stimulation $\rightarrow$ spasm then paralysis. \textbf{Botulinum toxin} prevents vesicle fusion $\rightarrow$ flaccid paralysis (used medically in tiny doses for dystonia and cosmetics).
\end{savaistu}

\begin{exoresolu}[Synaptic Transmission and the Effect of a Toxin]
The diagram of a cholinergic synapse shows structures labelled A (vesicles), B (synaptic cleft) and C (receptor proteins).\par
(a) Describe how an impulse arriving at the pre-synaptic knob leads to the generation of an impulse in the post-synaptic neurone. \hfill (5 marks)\par
(b) Explain why transmission across a synapse occurs in one direction only. \hfill (2 marks)\par
(c) A certain pesticide inhibits the enzyme acetylcholinesterase. Predict and explain its effect on the post-synaptic neurone. \hfill (3 marks)
\tcblower
\textbf{(a)} Award one mark for each correct step in logical order:
\begin{enumerate}
\item The action potential depolarises the pre-synaptic membrane, opening voltage-gated \textbf{calcium channels}, so $\mathrm{Ca^{2+}}$ enters the knob;
\item $\mathrm{Ca^{2+}}$ causes synaptic \textbf{vesicles (A)} containing acetylcholine to fuse with the pre-synaptic membrane;
\item Acetylcholine is released by \textbf{exocytosis} into the synaptic cleft (B) and \textbf{diffuses} across it;
\item Acetylcholine binds to \textbf{specific receptor proteins (C)} on the post-synaptic membrane, opening sodium channels;
\item $\mathrm{Na^{+}}$ floods in, depolarising the post-synaptic membrane; if the threshold is reached, a \textbf{new action potential} is generated.
\end{enumerate}
\textbf{(b)} Synaptic vesicles and transmitter are present \emph{only} in the pre-synaptic knob, and specific receptor proteins are found \emph{only} on the post-synaptic membrane. Transmitter can therefore be released and received in one direction only, so the impulse cannot travel backwards across the synapse.\par
\textbf{(c)} If acetylcholinesterase is inhibited, acetylcholine is \textbf{not broken down} in the cleft and accumulates. It therefore binds \textbf{repeatedly/continuously} to the post-synaptic receptors. The post-synaptic membrane is \textbf{continuously depolarised}, producing prolonged, uncontrolled firing of the post-synaptic neurone (in muscles: continuous contraction/spasm, then paralysis). This is the basis of the toxicity of many insecticides and nerve gases.
\end{exoresolu}

\courssub{Skill 4: The Iris Reflex and Control of Light Intensity}

\begin{definition}[Pupil Reflex (Iris Reflex)]
The \cle{pupil reflex} is a brainstem reflex that adjusts the diameter of the pupil in response to light intensity falling on the retina, regulating the amount of light entering the eye. It is an autonomic (involuntary) reflex.
\end{definition}

\begin{propriete}[Antagonistic Muscles of the Iris]
The iris contains two sets of smooth muscle fibres:
\begin{itemize}
\item \textbf{Circular (sphincter) muscles:} arranged concentrically around the pupil; contraction $\rightarrow$ \cle{pupil constriction} (miosis). Innervated by \cle{parasympathetic} fibres via the oculomotor nerve (CN III).
\item \textbf{Radial (dilator) muscles:} arranged like spokes of a wheel; contraction $\rightarrow$ \cle{pupil dilation} (mydriasis). Innervated by \cle{sympathetic} fibres from the superior cervical ganglion.
\end{itemize}
They act antagonistically.
\end{propriete}

\begin{methode}[Explaining Pupil Size Changes]
\begin{enumerate}
\item \textbf{Identify the stimulus:} change in light intensity detected by photoreceptors (rods/cones) in the retina.
\item \textbf{Bright light:} high intensity $\rightarrow$ strong stimulation $\rightarrow$ parasympathetic activation $\rightarrow$ \textbf{circular muscles contract}, radial muscles relax $\rightarrow$ pupil \textbf{constricts} $\rightarrow$ less light enters (protects retina from bleaching).
\item \textbf{Dim light:} low intensity $\rightarrow$ weak stimulation $\rightarrow$ sympathetic activation $\rightarrow$ \textbf{radial muscles contract}, circular muscles relax $\rightarrow$ pupil \textbf{dilates} $\rightarrow$ more light enters.
\item \textbf{Control:} reflex arc: retina $\rightarrow$ optic nerve $\rightarrow$ pretectal nucleus (midbrain) $\rightarrow$ Edinger-Westphal nucleus $\rightarrow$ oculomotor nerve $\rightarrow$ ciliary ganglion $\rightarrow$ iris muscles. \emph{Consensual reflex:} shining light in one eye constricts both pupils.
\end{enumerate}
\textbf{Do not confuse:} the iris reflex (seconds, neural) with \emph{retinal adaptation} to light/dark (minutes, photochemical -- rhodopsin breakdown/resynthesis).
\end{methode}

\begin{attention}[Common Errors]
\begin{itemize}
\item Saying "circular muscles dilate the pupil" -- they constrict it.
\item Attributing the reflex to the spinal cord -- the centre is in the \emph{midbrain} (pretectal area).
\item Confusing the iris reflex (pupil size) with accommodation (lens shape) -- different effectors, different pathways.
\end{itemize}
\end{attention}

\begin{savaistu}[Clinical Relevance]
Doctors test the \textbf{pupillary light reflex} to assess brainstem function (CN II and III). Fixed, dilated pupils ("blown pupils") indicate severe brainstem compression (e.g. from raised intracranial pressure after head trauma). In Cameroon, this is a vital sign checked in emergency wards in Yaound\'e and Douala.
\end{savaistu}

\begin{exoresolu}[From a Dark Cinema into Bright Sunshine]
A student walks out of a dark cinema hall in Yaound\'e into bright midday sunlight. At first she can hardly see and shades her eyes; after a few minutes her vision becomes comfortable.\par
(a) Describe and explain the immediate change that occurs in her pupils on entering the bright light. \hfill (3 marks)\par
(b) Explain why the pupil reflex alone cannot fully account for her comfortable vision after several minutes. \hfill (3 marks)\par
(c) Name the pigment involved in the slower adjustment and state where it is found. \hfill (2 marks)
\tcblower
\textbf{(a)} In bright light, the pupil \textbf{constricts} (becomes smaller). The \textbf{circular muscles of the iris contract} while the \textbf{radial muscles relax}. This reduces the amount of light entering the eye and protects the retina from damage by excessive light. It is a \textbf{reflex action} controlled by the autonomic nervous system (parasympathetic branch), with the receptor being the retina and the effector the iris muscles.\par
\textbf{(b)} The pupil can only vary light entry over a limited range (its area changes by a factor of about 16 at most), whereas light intensity outdoors can be thousands of times greater than in the dark hall. Full adjustment requires \textbf{retinal (light) adaptation}: in bright light the visual pigment \textbf{rhodopsin is bleached} (broken down into retinal and opsin) faster than it is resynthesised, so the \textbf{rods become less sensitive} and vision shifts to the \textbf{cones}, which work well in bright light. This photochemical adjustment takes several minutes, which matches the gradual improvement described.\par
\textbf{(c)} The pigment is \textbf{rhodopsin} (visual purple), found in the \textbf{rod cells} of the retina. (Cone pigments -- iodopsins -- mediate bright-light and colour vision.)
\end{exoresolu}

\courssub{Skill 5: Accommodation and Defects of Vision}

\begin{definition}[Accommodation]
\cle{Accommodation} is the process by which the eye changes the focal length of its lens to focus images of objects at different distances sharply on the retina.
\end{definition}

\begin{propriete}[Mechanism of Accommodation]
\begin{itemize}
\item \textbf{Near object (e.g. reading at 25 cm):} ciliary muscles \textbf{contract} $\rightarrow$ suspensory ligaments \textbf{slacken} $\rightarrow$ elastic lens becomes \textbf{more convex} (shorter focal length, higher power) $\rightarrow$ diverging rays focused on retina.
\item \textbf{Distant object (parallel rays):} ciliary muscles \textbf{relax} $\rightarrow$ suspensory ligaments become \textbf{taut} $\rightarrow$ lens pulled \textbf{flatter} (longer focal length, lower power) $\rightarrow$ parallel rays focused on retina.
\end{itemize}
The lens is naturally elastic and rounded; the suspensory ligaments keep it flattened when ciliary muscles are relaxed.
\end{propriete}

\begin{definition}[Defects of Vision]
\begin{itemize}
\item \textbf{Myopia (short sight):} eyeball too long or lens/cornea too powerful $\rightarrow$ image forms \textbf{in front of} retina for distant objects. Corrected by \textbf{diverging (concave) lens} which diverges rays before they enter the eye.
\item \textbf{Hypermetropia (long sight):} eyeball too short or lens/cornea too weak $\rightarrow$ image forms \textbf{behind} retina for near objects (and sometimes distant). Corrected by \textbf{converging (convex) lens} which converges rays before they enter the eye.
\item \textbf{Presbyopia:} age-related loss of lens elasticity and ciliary muscle power $\rightarrow$ difficulty with near vision. Corrected by converging lenses (reading glasses).
\item \textbf{Astigmatism:} uneven curvature of cornea/lens $\rightarrow$ rays in different planes focus at different points. Corrected by cylindrical lenses.
\end{itemize}
\end{definition}

\begin{methode}[Explaining Accommodation and Correcting Defects]
\begin{enumerate}
\item \textbf{For accommodation:} state ciliary muscle action $\rightarrow$ ligament tension $\rightarrow$ lens shape $\rightarrow$ focal length $\rightarrow$ image on retina.
\item \textbf{For defects:} always state (i) where the image falls relative to retina, (ii) the cause (eyeball length or lens power), (iii) the correcting lens type and its effect on incident rays.
\item \textbf{Diagram:} draw rays from object $\rightarrow$ correcting lens $\rightarrow$ eye lens $\rightarrow$ retina.
\end{enumerate}
\end{methode}

\begin{attention}[Exam Tips]
\begin{itemize}
\item "Short sight = see near, not far" -- the name describes what you \emph{can} see.
\item Concave lens = diverging = thinner at centre; Convex lens = converging = thicker at centre.
\item In accommodation for near vision, the ciliary muscles \textbf{contract} (many students wrongly say relax).
\end{itemize}
\end{attention}

\begin{exoresolu}[A Student Who Cannot Read the Board]
A Form Five student can read her book clearly but cannot read writing on the classroom board unless she sits in the front row.\par
(a) Name the defect of vision she suffers from. \hfill (1 mark)\par
(b) Explain, in terms of where images are focused, why she cannot see distant objects clearly. \hfill (2 marks)\par
(c) State the type of spectacle lens that would correct this defect and explain how it works. \hfill (2 marks)\par
(d) Describe how her eye accommodates when she reads her book at $25\ \mathrm{cm}$. \hfill (3 marks)
\tcblower
\textbf{(a)} \textbf{Short sightedness (myopia)}: near objects are seen clearly, distant objects are blurred.\par
\textbf{(b)} In her eye, parallel rays from a distant object are brought to a focus \textbf{in front of the retina} (because the eyeball is too long or the lens/cornea system converges the rays too strongly). By the time the rays reach the retina they have diverged again, so the image on the retina is \textbf{blurred}.\par
\textbf{(c)} A \textbf{diverging (concave) lens} is used. It spreads the incoming parallel rays slightly apart \emph{before} they enter the eye, so that the eye's own lens then focuses them exactly \textbf{on the retina} instead of in front of it, giving a sharp image of distant objects.\par
\textbf{(d)} When reading at $25\ \mathrm{cm}$ (near object):
\begin{enumerate}
\item the \textbf{ciliary muscles contract};
\item the \textbf{suspensory ligaments slacken}, releasing tension on the lens;
\item the elastic lens bulges, becoming \textbf{more convex} (shorter focal length), so the diverging rays from the near book are focused sharply on the retina.
\end{enumerate}
\end{exoresolu}

\courssub{Skill 6: Analysing Hormonal Control and Negative Feedback}

\begin{definition}[Hormone]
A \cle{hormone} is a chemical messenger secreted by an \cle{endocrine gland} directly into the bloodstream, which carries it to specific \cle{target cells} bearing complementary receptors, where it exerts a regulatory effect.
\end{definition}

\begin{propriete}[Mechanism of Hormone Action]
\begin{itemize}
\item \textbf{Protein/peptide hormones (e.g. insulin, adrenaline, ADH, FSH, LH):} water-soluble, cannot cross phospholipid bilayer. Bind to \textbf{cell surface receptors} $\rightarrow$ activate \textbf{second messenger} (usually \cle{cyclic AMP (cAMP)} via adenyl cyclase) $\rightarrow$ cascade of enzyme activations $\rightarrow$ rapid response (seconds to minutes).
\item \textbf{Steroid hormones (e.g. oestrogen, testosterone, cortisol, aldosterone):} lipid-soluble, diffuse through membrane. Bind to \textbf{intracellular receptors} (cytoplasm or nucleus) $\rightarrow$ hormone-receptor complex acts as \textbf{transcription factor} on DNA $\rightarrow$ alters gene expression $\rightarrow$ new protein synthesis $\rightarrow$ slower response (hours), longer lasting.
\end{itemize}
\end{propriete}

\begin{propriete}[Negative Feedback]
\cle{Negative feedback} is the principal mechanism of hormonal homeostasis: a rise in the controlled variable stimulates hormone secretion that lowers the variable; as the variable returns to normal, the stimulus for secretion diminishes, switching off further hormone release. This maintains the variable within narrow limits (set point).
\end{propriete}

\begin{methode}[Interpreting a Hormone/Feedback Question]
\begin{enumerate}
\item \textbf{Identify the gland and hormone} from the symptoms or data given (e.g.\ high blood glucose $\rightarrow$ insulin from pancreas $\beta$-cells; goitre $\rightarrow$ thyroid/thyroxine).
\item \textbf{State the stimulus} that triggers secretion (e.g.\ rising blood glucose detected by $\beta$-cells).
\item \textbf{State the target organ(s) and effect} (e.g.\ liver and muscles convert glucose to glycogen; cells increase glucose uptake).
\item \textbf{Describe the feedback loop:} when the controlled factor returns to normal, the change is detected and hormone secretion is \emph{switched off/reduced} -- this is negative feedback.
\item \textbf{Contrast with the antagonistic hormone} where relevant (insulin vs glucagon; thyroxine vs TSH; they act in opposite directions to keep the variable constant).
\item \textbf{Mechanism of action:} specify whether the hormone acts via membrane receptor/second messenger (protein) or intracellular receptor/gene activation (steroid).
\end{enumerate}
\end{methode}

\begin{savaistu}[Cameroon Context]
\textbf{Diabetes mellitus} is a growing public health concern in Cameroon (urbanisation, diet change). Type 1 = autoimmune destruction of $\beta$-cells (no insulin); Type 2 = insulin resistance (receptors defective). \textbf{Goitre} (iodine deficiency) was endemic in the Adamawa and North regions; salt iodisation programmes have reduced prevalence. \textbf{Adrenaline} surge during "fright" (e.g. encountering a snake on Mount Cameroon) illustrates neuroendocrine interaction.
\end{savaistu}

\begin{exoresolu}[Blood Glucose After a Meal of Fufu and Eru]
After a heavy meal rich in carbohydrate, a student's blood glucose concentration rose from $90\ \mathrm{mg}$ per $100\ \mathrm{cm^{3}}$ to $140\ \mathrm{mg}$ per $100\ \mathrm{cm^{3}}$, then fell back to about $90\ \mathrm{mg}$ per $100\ \mathrm{cm^{3}}$ within two hours.\par
(a) Name the hormone responsible for the fall and the gland that secretes it. \hfill (2 marks)\par
(b) Explain how this hormone brings about the fall in blood glucose. \hfill (3 marks)\par
(c) Explain what prevents the blood glucose from falling far \emph{below} the normal level. \hfill (3 marks)\par
(d) The hormone in (a) is a protein. Explain why it cannot enter its target cells and state how its message is transmitted inside the cell. \hfill (2 marks)
\tcblower
\textbf{(a)} The hormone is \textbf{insulin}, secreted by the \textbf{$\beta$-cells of the islets of Langerhans in the pancreas}.\par
\textbf{(b)} Insulin lowers blood glucose by:
\begin{itemize}
\item stimulating the \textbf{liver and muscles} to convert excess glucose into \textbf{glycogen} for storage (glycogenesis);
\item increasing the \textbf{uptake and respiration of glucose by cells} (increased membrane permeability to glucose via GLUT4 translocation);
\item promoting conversion of glucose to \textbf{fat} in adipose tissue (lipogenesis).
\end{itemize}
\textbf{(c)} This is \textbf{negative feedback}: as blood glucose falls towards normal, the stimulus for insulin secretion disappears, so insulin release is \textbf{switched off}. In addition, if glucose falls below normal, the \textbf{$\alpha$-cells of the pancreas secrete glucagon}, which is \textbf{antagonistic} to insulin: it stimulates the liver to convert stored \textbf{glycogen back into glucose} (glycogenolysis) and release it into the blood, and promotes gluconeogenesis. The two hormones together keep blood glucose fluctuating narrowly around the set point.\par
\textbf{(d)} Insulin is a \textbf{large, water-soluble protein molecule} and cannot cross the \textbf{phospholipid bilayer} of the cell surface membrane. It therefore binds to a \textbf{specific receptor on the membrane} (a receptor tyrosine kinase); this activates the enzyme adenyl cyclase inside the cell, which produces the \textbf{second messenger cyclic AMP (cAMP)}, which then triggers the internal response (e.g.\ activation of enzymes for glycogen synthesis).
\end{exoresolu}

\courssub{Skill 7: Designing and Interpreting Tropism Experiments}

\begin{definition}[Tropism]
A \cle{tropism} is a directional growth response of a plant to a unidirectional external stimulus. \textbf{Positive tropism} = growth towards the stimulus; \textbf{negative tropism} = growth away. Main types: \cle{phototropism} (light), \cle{geotropism} (gravity), \cle{hydrotropism} (water), \cle{thigmotropism} (touch).
\end{definition}

\begin{definition}[Nastic Movement]
A \cle{nastic movement} is a non-directional response to a stimulus (e.g. \cle{photonasty} -- flower opening/closing with light intensity; \cle{seismonasty} -- \emph{Mimosa} leaf folding on touch). The response direction is determined by plant structure, not stimulus direction.
\end{definition}

\begin{propriete}[Role of Auxin (IAA) in Phototropism and Geotropism]
\begin{itemize}
\item \textbf{Shoots (positive phototropism):} Auxin produced in tip $\rightarrow$ redistributed to \textbf{shaded side} $\rightarrow$ \textbf{promotes cell elongation} on shaded side $\rightarrow$ shoot bends \textbf{towards light}.
\item \textbf{Roots (positive geotropism):} Auxin redistributed to \textbf{lower side} $\rightarrow$ root cells are \textbf{more sensitive} -- auxin \textbf{inhibits elongation} on lower side $\rightarrow$ upper side elongates more $\rightarrow$ root bends \textbf{downwards}.
\end{itemize}
\cle{Acid growth hypothesis:} Auxin activates proton pumps $\rightarrow$ wall acidification $\rightarrow$ expansins loosen cellulose microfibrils $\rightarrow$ turgor-driven elongation.
\end{propriete}

\begin{methode}[Investigating Phototropism / Geotropism]
\begin{enumerate}
\item \textbf{Define variables:} independent = direction of light (or gravity); dependent = angle/direction of growth; controls = temperature, water, seedling type, age, uniform light for control.
\item \textbf{Always include a control} (seedling in uniform light/darkness or vertical) to show bending is a response to \emph{direction} of stimulus.
\item \textbf{Use the auxin explanation:} production in tip, lateral redistribution, differential elongation.
\item \textbf{Classic experiments to quote:}
\begin{itemize}
\item \textbf{Darwin:} tip perceives light; covering tip prevents bending.
\item \textbf{Boysen-Jensen:} tip separated by mica (impermeable) blocks bending; gelatin (permeable) allows bending $\rightarrow$ chemical messenger.
\item \textbf{Went:} tip on agar block $\rightarrow$ auxin diffuses into agar $\rightarrow$ agar placed asymmetrically on decapitated stump causes bending without light $\rightarrow$ auxin identified as the messenger.
\end{itemize}
\end{enumerate}
\end{methode}

\begin{attention}[Tropism Experiment Pitfalls]
\begin{itemize}
\item Forgetting the control (uniform light).
\item Saying auxin "moves away from light" -- it is \emph{redistributed} to the shaded side (active transport).
\item Confusing shoot and root responses to auxin (shoots: promotes elongation; roots: inhibits at same concentration).
\item Not mentioning that the tip is the site of both perception and auxin synthesis.
\end{itemize}
\end{attention}

\begin{savaistu}[Mount Cameroon Forest Floor]
On the shaded slopes of Mount Cameroon, seedlings of \emph{Maesopsis eminii} or \emph{Terminalia superba} exhibit strong positive phototropism. Auxin-mediated bending towards fleeting sunflecks maximises photosynthetic gain in the understorey, a key survival strategy in the competitive rainforest environment.
\end{savaistu}

\begin{exoresolu}[Phototropism in Maize Seedlings]
Three groups of maize coleoptiles were treated as follows and left for 48 hours with light shining from one side only:\par
Group P: tips left intact. \quad Group Q: tips cut off. \quad Group R: tips cut off and replaced with a thin agar block between tip and stump.\par
Results: P bent towards the light; Q did not bend and grew little; R bent towards the light.\par
(a) State the aim of this experiment. \hfill (1 mark)\par
(b) Explain the result obtained in group P. \hfill (3 marks)\par
(c) Explain the results of groups Q and R, and state the conclusion that can be drawn about the role of the tip. \hfill (4 marks)\par
(d) State the biological advantage of this response to a plant growing on a shaded forest floor on Mount Cameroon. \hfill (2 marks)
\tcblower
\textbf{(a)} To investigate the role of the \textbf{shoot tip (coleoptile tip)} in the \textbf{phototropic response} (bending of a shoot towards unilateral light).\par
\textbf{(b)} In group P, light from one side causes \textbf{auxin (IAA)} produced in the tip to be \textbf{redistributed to the shaded side} of the coleoptile. Auxin \textbf{stimulates cell elongation} in shoots, so cells on the shaded side elongate \textbf{more} than those on the lit side. This \textbf{unequal growth} causes the coleoptile to bend \textbf{towards the light} (positive phototropism).\par
\textbf{(c)} In group Q, removing the tip removes the \textbf{source of auxin}; without auxin there is little elongation growth and no differential growth, so the coleoptile \textbf{neither grows nor bends}. In group R, the agar block is \textbf{permeable to auxin}: auxin produced in the replaced tip \textbf{diffuses down through the agar} into the stump, is redistributed by the unilateral light, and bending towards the light is restored.\par
\textbf{Conclusion:} the tip is the site of \textbf{auxin production} (and light perception); auxin is a \textbf{chemical messenger} that moves downwards from the tip and controls the growth response. The response depends on a diffusible substance, not on the tip cells themselves being present at the bending zone.\par
\textbf{(d)} Positive phototropism makes the shoot grow \textbf{towards the light source}, positioning the leaves in brighter light. This maximises \textbf{light absorption for photosynthesis}, giving the seedling a competitive advantage in the dim, shaded conditions of the forest floor, where light is the main limiting factor.
\end{exoresolu}

\courssub{Skill 8: Comparing Nervous and Hormonal Coordination}

\begin{propriete}[Comparison of Nervous and Hormonal Coordination]
\begin{center}
\begin{tabularx}{\linewidth}{|l|X|X|}
\hline
\rowcolor{popblueL} \textbf{Criterion} & \textbf{Nervous Coordination} & \textbf{Hormonal Coordination} \\
\hline
Nature of message & Electrical impulses (action potentials) along axons; chemical (neurotransmitter) at synapses & Chemical (hormones) \\
\hline
Pathway & Along specific neurones (fixed, wired pathways) & Through the bloodstream (circulation) \\
\hline
Speed of transmission & Very fast (up to $\approx 100\ \mathrm{m\,s^{-1}}$ in myelinated fibres) & Slow (limited by blood flow, seconds to minutes) \\
\hline
Duration of response & Short-lived, brief (milliseconds to seconds) & Longer-lasting (minutes to days) \\
\hline
Spread of response & Localised (specific effector organ) & Widespread (all target organs with receptors reached by blood) \\
\hline
Target specificity & High (specific neurone-to-effector connection) & High (only cells with specific receptors respond) \\
\hline
Example & Withdrawal reflex, knee-jerk & Blood glucose regulation (insulin/glucagon), metamorphosis \\
\hline
\end{tabularx}
\end{center}
\end{propriete}

\begin{methode}[Building a Comparison Answer]
\begin{enumerate}
\item \textbf{Choose clear criteria} and compare \emph{both} systems under each one (never describe one system fully and then the other).
\item \textbf{Use these criteria:} nature of message, pathway, speed, duration, spread, target specificity.
\item \textbf{Give one example of each} (reflex action vs blood glucose control).
\item \textbf{Note the overlap:} some responses use both (e.g.\ the adrenal medulla is stimulated by a nerve but releases adrenaline into the blood -- neuroendocrine link; the hypothalamus links the two systems).
\end{enumerate}
\end{methode}

\begin{attention}[Comparison Traps]
\begin{itemize}
\item Saying hormones are "slow" without qualifying (blood flow limited).
\item Saying nervous system is "only electrical" -- synapses are chemical.
\item Forgetting that both systems can be specific (hormones only affect target cells with receptors).
\item Not mentioning the neuroendocrine integration (hypothalamus-pituitary, adrenal medulla).
\end{itemize}
\end{attention}

\begin{exoresolu}[Nervous versus Hormonal Coordination]
(a) Construct a table comparing nervous and hormonal coordination under four headings. \hfill (4 marks)\par
(b) A man touches a thorn and withdraws his hand; later the same day, a fright causes his heart to race for several minutes. For each event, state the type of coordination involved and justify your answer. \hfill (4 marks)
\tcblower
\textbf{(a)}
\begin{center}
\begin{tabularx}{\linewidth}{|l|X|X|}
\hline
\rowcolor{popblueL} \textbf{Criterion} & \textbf{Nervous coordination} & \textbf{Hormonal coordination} \\
\hline
Nature of message & Electrical impulses (and chemical transmitter at synapses) & Chemical (hormones) \\
\hline
Pathway & Along neurones (fixed pathways) & Through the bloodstream \\
\hline
Speed of transmission & Very fast (up to about $100\ \mathrm{m\,s^{-1}}$) & Slow (depends on blood flow) \\
\hline
Duration of response & Short-lived, brief & Longer-lasting \\
\hline
Spread of response & Localised (specific effector) & Widespread (all target organs reached by blood) \\
\hline
\end{tabularx}
\end{center}
\textbf{(b)} \emph{Withdrawal from the thorn:} \textbf{nervous coordination} (a spinal reflex). Justification: the response is \textbf{very rapid, brief and localised} to the arm muscles, and travels along a fixed pathway of neurones -- all characteristic of nervous control.\par
\emph{Racing heart after fright:} mainly \textbf{hormonal coordination}. The fright causes the \textbf{adrenal medulla} to release \textbf{adrenaline} into the blood; adrenaline reaches the heart via the circulation and increases heart rate. The effect is \textbf{slower in onset but persists for several minutes} and affects many organs (heart, liver, muscles) -- characteristic of hormonal control. (Note: the initial stimulation of the adrenal medulla is nervous, showing that the two systems interact.)
\end{exoresolu}

\courssub{Skill 9: The Ear -- Hearing and Balance}

\begin{definition}[Ear Structure and Function]
The ear is a sensory organ for \cle{hearing} (detection of sound waves) and \cle{balance} (detection of head position and movement). It has three regions: outer, middle, inner.
\end{definition}

\begin{propriete}[Hearing Mechanism]
\begin{enumerate}
\item \textbf{Outer ear:} pinna collects sound waves $\rightarrow$ auditory canal $\rightarrow$ \cle{tympanic membrane} (eardrum) vibrates.
\item \textbf{Middle ear:} three \cle{ossicles} (malleus, incus, stapes) amplify and transmit vibrations $\rightarrow$ \cle{oval window} (membrane on cochlea). \cle{Eustachian tube} equalises pressure.
\item \textbf{Inner ear (cochlea):} stapes vibrates oval window $\rightarrow$ pressure waves in \cle{perilymph} $\rightarrow$ \cle{basilar membrane} vibrates $\rightarrow$ \cle{hair cells} of \cle{organ of Corti} bent against \cle{tectorial membrane} $\rightarrow$ generator potentials $\rightarrow$ action potentials in \cle{auditory nerve} (CN VIII) $\rightarrow$ auditory cortex.
\item \textbf{Frequency discrimination:} high frequencies peak near base (oval window), low frequencies near apex (place theory).
\end{enumerate}
\end{propriete}

\begin{propriete}[Balance Mechanism]
\begin{itemize}
\item \textbf{Static balance (position):} \cle{utricle} and \cle{saccule} (vestibule) contain \cle{maculae} with hair cells covered by \cle{otoliths} (calcium carbonate crystals). Gravity pulls otoliths $\rightarrow$ bends hair cells $\rightarrow$ signals head position.
\item \textbf{Dynamic balance (movement):} three \cle{semicircular canals} (mutually perpendicular) contain \cle{ampullae} with \cle{cristae} (hair cells in gelatinous cupula). Rotational acceleration moves endolymph $\rightarrow$ deflects cupula $\rightarrow$ bends hair cells $\rightarrow$ signals rotation.
\end{itemize}
\end{propriete}

\begin{attention}[Ear Exam Points]
\begin{itemize}
\item Ossicles amplify force $\approx 20\times$ (lever action) and match impedance (air to fluid).
\item Cochlea unrolled is a tube with basilar membrane; hair cells are the receptors.
\item Deafness types: \textbf{conductive} (outer/middle ear blockage, ossicle fusion) vs \textbf{sensorineural} (hair cell/nerve damage).
\item Balance and hearing both use hair cells (mechanoreceptors) but different structures.
\end{itemize}
\end{attention}

\begin{exoresolu}[Hearing and Balance]
(a) Describe how sound waves are converted into nerve impulses in the cochlea. \hfill (4 marks)\par
(b) Explain how the semicircular canals detect rotational movement of the head. \hfill (3 marks)\par
(c) A person suffers from dizziness and nausea after a severe cold. Suggest a possible explanation. \hfill (2 marks)
\tcblower
\textbf{(a)} Sound waves cause the tympanic membrane to vibrate. Ossicles amplify and transmit vibrations to the oval window. Pressure waves in the perilymph cause the basilar membrane to vibrate. Hair cells of the organ of Corti are bent against the tectorial membrane, generating generator potentials. If threshold is reached, action potentials are generated in the auditory nerve fibres and transmitted to the brain.\par
\textbf{(b)} The three semicircular canals are oriented in three perpendicular planes. Each has an ampulla containing a crista (hair cells embedded in a gelatinous cupula). During rotational acceleration, the endolymph lags behind due to inertia, pushing the cupula and bending the hair cells. This generates nerve impulses proportional to the rate of rotation, informing the brain of the direction and speed of head rotation.\par
\textbf{(c)} The cold may have caused inflammation blocking the \textbf{Eustachian tube}, preventing pressure equalisation in the middle ear, or spread to the \textbf{inner ear (labyrinthitis)} affecting the vestibular apparatus (semicircular canals/utricle/saccule), disrupting balance signals and causing vertigo/nausea.
\end{exoresolu}

\courssub{Skill 10: The Brain, Spinal Cord and Autonomic Nervous System}

\begin{definition}[Central Nervous System (CNS)]
The CNS consists of the \cle{brain} and \cle{spinal cord}. It is the site of integration, processing, and command. Protected by meninges (dura mater, arachnoid, pia mater) and cerebrospinal fluid (CSF).
\end{definition}

\begin{propriete}[Major Brain Regions and Functions]
\begin{itemize}
\item \textbf{Cerebrum (cerebral hemispheres):} higher functions -- consciousness, memory, reasoning, voluntary movement, sensory perception. Grey matter (cortex) outside, white matter inside.
\item \textbf{Cerebellum:} coordination of movement, balance, posture, motor learning. "Little brain" posterior to brainstem.
\item \textbf{Medulla oblongata:} vital reflex centres -- cardiac, respiratory, vasomotor, vomiting, coughing.
\item \textbf{Hypothalamus:} homeostasis -- temperature, water balance (ADH), hunger, endocrine control via pituitary.
\item \textbf{Pituitary gland:} "master gland" -- anterior (trophic hormones: TSH, ACTH, FSH, LH, GH, prolactin) and posterior (ADH, oxytocin).
\item \textbf{Midbrain:} reflex centres for vision (pupil reflex) and hearing.
\end{itemize}
\end{propriete}

\begin{propriete}[Spinal Cord]
\begin{itemize}
\item Runs through vertebral column; segmental organisation (31 pairs of spinal nerves).
\item \textbf{Grey matter (central):} butterfly-shaped; cell bodies of motor neurones (ventral horn), interneurones, sensory fibre terminals (dorsal horn).
\item \textbf{White matter (peripheral):} myelinated axons forming ascending (sensory) and descending (motor) tracts.
\item \textbf{Function:} reflex centre (spinal reflexes); conduction pathway to/from brain.
\end{itemize}
\end{propriete}

\begin{propriete}[Autonomic Nervous System (ANS)]
Involuntary control of cardiac muscle, smooth muscle, glands. Two antagonistic divisions:
\begin{center}
\begin{tabularx}{\linewidth}{|l|X|X|}
\hline
\rowcolor{popblueL} \textbf{Feature} & \textbf{Sympathetic} & \textbf{Parasympathetic} \\
\hline
Origin & Thoracolumbar (T1--L2) & Craniosacral (CN III, VII, IX, X; S2--S4) \\
\hline
Pre-ganglionic fibre & Short & Long \\
\hline
Post-ganglionic fibre & Long & Short \\
\hline
Ganglia & Paravertebral chain, prevertebral & Near/on effector organs \\
\hline
Neurotransmitter (pre) & Acetylcholine & Acetylcholine \\
\hline
Neurotransmitter (post) & Noradrenaline (mostly) & Acetylcholine \\
\hline
General effect & "Fight or flight" -- arousal, energy mobilisation & "Rest and digest" -- conservation, restoration \\
\hline
Heart rate & Increases & Decreases \\
\hline
Pupil & Dilates (radial muscle) & Constricts (circular muscle) \\
\hline
Gut motility & Decreases & Increases \\
\hline
\end{tabularx}
\end{center}
\end{propriete}

\begin{savaistu}[Neuroendocrine Integration]
The \textbf{hypothalamus} is the key link: it receives neural inputs (e.g. stress, temperature, osmotic pressure) and outputs hormonal signals to the pituitary. The \textbf{adrenal medulla} is a modified sympathetic ganglion -- preganglionic fibres stimulate chromaffin cells to secrete adrenaline/noradrenaline into blood (neuroendocrine transducer).
\end{savaistu}

\courssub{Skill 11: Plant Hormones -- Auxins, Gibberellins, Cytokinins, Abscisic Acid, Ethylene}

\begin{definition}[Plant Hormones (Phytohormones)]
Chemical messengers produced in one part of the plant, transported (often via phloem or diffusion) to target tissues where they regulate growth and development at very low concentrations. Major classes:
\end{definition}

\begin{propriete}[Major Plant Hormones and Functions]
\begin{center}
\begin{tabularx}{\linewidth}{|l|X|X|}
\hline
\rowcolor{popblueL} \textbf{Hormone} & \textbf{Main Sites of Synthesis} & \textbf{Key Functions} \\
\hline
\textbf{Auxins (IAA)} & Shoot apex, young leaves, developing seeds & Cell elongation (phototropism, geotropism), apical dominance, root initiation (cuttings), fruit development, prevents abscission \\
\hline
\textbf{Gibberellins (GA)} & Young leaves, shoot tips, seeds, roots & Stem elongation (bolting), seed germination (mobilises reserves), breaks dormancy, fruit growth, flowering \\
\hline
\textbf{Cytokinins} & Root tips, developing seeds, fruits & Cell division (cytokinesis), delays senescence (leaf aging), promotes lateral bud growth (releases apical dominance), chloroplast development \\
\hline
\textbf{Abscisic Acid (ABA)} & Leaves, roots, seeds & Stomatal closure (water stress), seed dormancy, inhibits growth, promotes leaf abscission \\
\hline
\textbf{Ethene (Ethylene)} & Ripening fruits, senescing leaves, stressed tissues & Fruit ripening, leaf/flower/fruit abscission, senescence, stress response, flower sex determination \\
\hline
\end{tabularx}
\end{center}
\end{propriete}

\begin{methode}[Commercial Applications of Plant Hormones]
\begin{enumerate}
\item \textbf{Auxins:} rooting powders for cuttings (IBA, NAA); selective herbicides (2,4-D, MCPA -- broadleaf weeds); parthenocarpic fruit set (seedless tomatoes).
\item \textbf{Gibberellins:} increase grape bunch size and berry size; break dormancy in barley for malting (beer); induce flowering in long-day plants.
\item \textbf{Cytokinins:} tissue culture (shoot proliferation); delay leaf senescence in cut flowers/vegetables.
\item \textbf{Abscisic Acid:} potential antitranspirant (spray to close stomata during drought).
\item \textbf{Ethylene:} controlled fruit ripening (bananas, mangoes in Cameroon markets); degreening citrus; synchronising flowering in pineapples.
\end{enumerate}
\end{methode}

\begin{attention}[Hormone Interaction]
Plant responses usually result from \textbf{hormone ratios}, not single hormones. E.g. apical dominance: high auxin:cytokinin ratio suppresses lateral buds; removing apex lowers auxin, allowing cytokinin to promote branching. Root:shoot ratio in tissue culture controlled by auxin:cytokinin balance.
\end{attention}

\begin{savaistu}[Cameroon Agriculture]
In Cameroon, \textbf{ethylene} is used to ripen plantains and bananas uniformly for urban markets (Douala, Yaound\'e). \textbf{Gibberellic acid} improves germination of uniform barley for the brewing industry (SOCABO). \textbf{Auxin-based herbicides} (2,4-D) control broadleaf weeds in maize and rice fields. \textbf{Cytokinins} in tissue culture labs (IRAD, universities) mass-propagate disease-free plantain and cassava planting material.
\end{savaistu}

\courssub{Skill 12: Light and Plant Growth -- Photoperiodism and Phytochrome}

\begin{definition}[Photoperiodism]
\cle{Photoperiodism} is the response of plants to the relative lengths of day and night (photoperiod), particularly controlling \textbf{flowering}. The critical factor is the length of the \textbf{uninterrupted dark period}.
\end{definition}

\begin{propriete}[Photoperiodic Classification]
\begin{itemize}
\item \textbf{Short-day plants (SDPs):} flower when night length \textbf{exceeds} a critical duration (e.g. chrysanthemum, poinsettia, rice, soybean, cotton). Actually "long-night plants".
\item \textbf{Long-day plants (LDPs):} flower when night length is \textbf{shorter} than a critical duration (e.g. wheat, barley, spinach, clover).
\item \textbf{Day-neutral plants:} flower regardless of photoperiod (e.g. tomato, cucumber, maize).
\end{itemize}
\textbf{Night-break experiment:} a brief flash of light during the long night prevents flowering in SDPs but promotes it in LDPs $\rightarrow$ proves dark period is critical.
\end{propriete}

\begin{propriete}[Phytochrome -- The Photoreceptor]
\cle{Phytochrome} is a blue-green pigment (protein + chromophore) existing in two interconvertible forms:
\begin{itemize}
\item \textbf{Pr (red-absorbing, $\lambda_{\max} \approx 660\ \mathrm{nm}$):} inactive form.
\item \textbf{Pfr (far-red absorbing, $\lambda_{\max} \approx 730\ \mathrm{nm}$):} active form.
\end{itemize}
\textbf{Conversions:} Red light ($660\ \mathrm{nm}$) $\rightarrow$ Pr $\xrightarrow{\text{absorbs}}$ Pfr (fast). Far-red light ($730\ \mathrm{nm}$) $\rightarrow$ Pfr $\xrightarrow{\text{absorbs}}$ Pr (fast). In darkness, Pfr slowly reverts to Pr (or is degraded).
\end{propriete}

\begin{methode}[Explaining Photoperiodic Response via Phytochrome]
\begin{enumerate}
\item \textbf{Daylight:} rich in red $\rightarrow$ mostly Pfr at sunset.
\item \textbf{Night:} Pfr slowly reverts to Pr.
\item \textbf{SDPs:} need \textbf{low Pfr} (long night) to flower. If night is interrupted by red light $\rightarrow$ Pfr rises $\rightarrow$ flowering inhibited. Far-red after red cancels the effect.
\item \textbf{LDPs:} need \textbf{high Pfr} (short night) to flower.
\item \textbf{Flowering hormone (florigen):} produced in leaves under appropriate photoperiod, transported via phloem to shoot apex $\rightarrow$ induces flowering.
\end{enumerate}
\end{methode}

\begin{attention}[Phytochrome Details]
\begin{itemize}
\item Pr $\leftrightarrow$ Pfr interconversion is rapid (seconds). Dark reversion is slow (hours).
\item The \textbf{last light quality} received determines the response (red $\rightarrow$ promotes LDP, inhibits SDP; far-red $\rightarrow$ opposite).
\item Phytochrome also controls seed germination (lettuce needs red light), de-etiolation, shade avoidance.
\end{itemize}
\end{attention}

\begin{savaistu}[Cameroon Crops and Photoperiod]
\textbf{Rice} (short-day) flowers as days shorten after the June solstice in the Sudano-Sahelian zone. \textbf{Soybean} varieties are selected for specific critical night lengths to match planting dates in different agro-ecological zones. \textbf{Chrysanthemum} growers in the western highlands use blackout curtains to induce flowering year-round for the cut-flower market.
\end{savaistu}

\begin{exoresolu}[Photoperiodism and Phytochrome]
(a) Explain why a short-day plant fails to flower if its long night is interrupted by a brief flash of red light, but flowers if the red flash is followed by a far-red flash. \hfill (4 marks)\par
(b) A farmer wants to delay flowering in a long-day crop (e.g. spinach) to extend the vegetative harvest. Suggest a light treatment and explain its effect. \hfill (3 marks)\par
(c) Name two other plant responses controlled by phytochrome besides flowering. \hfill (2 marks)
\tcblower
\textbf{(a)} In a short-day plant, flowering requires a long uninterrupted night during which \textbf{Pfr reverts to Pr} to a low level. A \textbf{red light flash} converts remaining Pr to \textbf{Pfr}, raising the Pfr level $\rightarrow$ the plant "perceives" a short night $\rightarrow$ flowering \textbf{inhibited}. A subsequent \textbf{far-red flash} converts Pfr back to Pr, restoring the low Pfr level $\rightarrow$ the plant "perceives" a long night $\rightarrow$ flowering \textbf{restored}.\par
\textbf{(b)} To delay flowering in a long-day plant (which needs short nights/high Pfr), give \textbf{long nights} or \textbf{end the day with far-red light} (or give a far-red pulse during the night). This ensures \textbf{low Pfr} at the end of the night, so the plant does not receive the inductive signal for flowering. Vegetative growth continues.\par
\textbf{(c)} Any two of: \textbf{seed germination} (e.g. lettuce requires red light), \textbf{de-etiolation} (greening, inhibition of hypocotyl elongation in seedlings), \textbf{shade avoidance} (stem elongation in response to low red:far-red ratio under canopy), \textbf{leaf abscission}, \textbf{dormancy}.
\end{exoresolu}

\courssub{Skill 13: Photoreception -- Rods and Cones}

\begin{definition}[Photoreception]
\cle{Photoreception} is the process by which light energy is converted into nerve impulses in the retina. The retina contains two types of photoreceptor cells: \cle{rods} and \cle{cones}.
\end{definition}

\begin{propriete}[Comparison of Rods and Cones]
\begin{center}
\begin{tabularx}{\linewidth}{|l|X|X|}
\hline
\rowcolor{popblueL} \textbf{Feature} & \textbf{Rods} & \textbf{Cones} \\
\hline
Visual pigment & \textbf{Rhodopsin} (visual purple) & \textbf{Iodopsins} (photopsins) -- three types (S, M, L) \\
\hline
Sensitivity & Very high (dim light, scotopic vision) & Low (bright light, photopic vision) \\
\hline
Colour vision & None (monochromatic) & Yes (trichromatic -- blue, green, red) \\
\hline
Acuity (detail) & Low (many rods converge on one bipolar cell) & High (one-to-one connection in fovea) \\
\hline
Distribution & Peripheral retina (absent from fovea) & Concentrated in \textbf{fovea} (yellow spot) \\
\hline
Adaptation & Dark adaptation (slow, 20--30 min) & Light adaptation (fast, seconds) \\
\hline
\end{tabularx}
\end{center}
\end{propriete}

\begin{propriete}[Mechanism of Phototransduction (Rods)]
\begin{enumerate}
\item \textbf{Dark (resting):} $\mathrm{Na^{+}}$ channels open (cGMP-gated) $\rightarrow$ "dark current" inward $\rightarrow$ cell depolarised ($\approx -40\ \mathrm{mV}$) $\rightarrow$ continuous glutamate release (inhibitory) $\rightarrow$ bipolar cells inhibited.
\item \textbf{Light hits rhodopsin:} retinal (11-cis) $\rightarrow$ all-trans retinal $\rightarrow$ conformational change $\rightarrow$ activates \cle{transducin} (G-protein) $\rightarrow$ activates \cle{phosphodiesterase (PDE)} $\rightarrow$ hydrolyses cGMP $\rightarrow$ cGMP levels fall $\rightarrow$ $\mathrm{Na^{+}}$ channels close $\rightarrow$ cell \textbf{hyperpolarises} $\rightarrow$ glutamate release stops $\rightarrow$ bipolar cells excited $\rightarrow$ signal to ganglion cells $\rightarrow$ optic nerve.
\end{enumerate}
\textbf{Key point:} Light causes \textbf{hyperpolarisation} (not depolarisation) of photoreceptors.
\end{propriete}

\begin{attention}[Photoreception Traps]
\begin{itemize}
\item Saying rods/cones "fire action potentials" -- they produce \emph{graded potentials}; only ganglion cells fire action potentials.
\item Confusing rhodopsin (rods) with iodopsins (cones).
\item Thinking the fovea has rods -- it is rod-free, cone-only, for maximum acuity in bright light.
\item Forgetting that photoreceptors \emph{hyperpolarise} in response to light (the opposite of most sensory receptors).
\end{itemize}
\end{attention}

\begin{methode}[Answering a Rods vs Cones Question]
\begin{enumerate}
\item \textbf{State the stimulus conditions} (dim vs bright light) and match to the receptor type.
\item \textbf{Link structure to function:} pigment type $\rightarrow$ sensitivity; neural wiring (convergence vs one-to-one) $\rightarrow$ acuity; retinal distribution $\rightarrow$ where each works best.
\item \textbf{For colour vision:} three cone types with different iodopsins; the brain compares their relative stimulation (trichromatic theory).
\item \textbf{For adaptation:} dark adaptation = rhodopsin resynthesis (slow); light adaptation = rhodopsin bleaching (fast).
\end{enumerate}
\end{methode}

\begin{savaistu}[Night Vision and Vitamin A]
The retinal part of rhodopsin is derived from \textbf{vitamin A}. Vitamin A deficiency -- still observed where diets lack green leafy vegetables, palm oil or liver -- causes \textbf{night blindness} (nyctalopia): rhodopsin cannot be resynthesised fast enough, so rod vision fails in dim light. Red palm oil, abundant in Cameroonian cooking, is an excellent source of provitamin A carotenoids.
\end{savaistu}

\begin{exoresolu}[Rods, Cones and Night Vision]
(a) Explain why, when looking at a faint star at night, it is easier to see it by looking slightly \emph{to one side} of it rather than directly at it. \hfill (3 marks)\par
(b) Explain why colours cannot be distinguished in very dim light. \hfill (2 marks)\par
(c) Describe the events in a rod cell between the absorption of a photon and the generation of a signal to the brain. \hfill (4 marks)
\tcblower
\textbf{(a)} The \textbf{fovea} (centre of the retina, where the image falls when we look directly at an object) contains \textbf{only cones}, which have \textbf{low sensitivity} to dim light. The \textbf{rods}, which are highly sensitive in dim light, are found in the \textbf{peripheral retina}. Looking slightly to one side makes the star's image fall on a rod-rich region of the retina, so the faint light is detected.\par
\textbf{(b)} Colour vision depends on the \textbf{three types of cones} (with different iodopsins), but cones are \textbf{insensitive in dim light} and do not respond. Only \textbf{rods} function, and rods contain a single pigment (rhodopsin) and cannot distinguish wavelengths; hence vision in dim light is \textbf{monochromatic} (shades of grey).\par
\textbf{(c)} Award one mark for each correct step:
\begin{enumerate}
\item A photon is absorbed by \textbf{rhodopsin}; 11-cis retinal is converted to \textbf{all-trans retinal}, changing the shape of the opsin (bleaching);
\item Activated rhodopsin activates the G-protein \textbf{transducin}, which activates \textbf{phosphodiesterase};
\item Phosphodiesterase hydrolyses \textbf{cGMP}, so cGMP-gated \textbf{sodium channels close} and the dark current stops;
\item The rod cell \textbf{hyperpolarises}, \textbf{glutamate release decreases}, which excites the connected bipolar cell; the signal passes to \textbf{ganglion cells}, which generate \textbf{action potentials} transmitted along the optic nerve to the brain.
\end{enumerate}
\end{exoresolu}

\begin{aretenir}[Key Points of the Chapter]
\begin{itemize}
\item A reflex arc has five components: receptor $\rightarrow$ sensory neurone $\rightarrow$ relay neurone $\rightarrow$ motor neurone $\rightarrow$ effector; dorsal root = sensory, ventral root = motor.
\item Conduction speed $v = d/t$; convert mm and ms before calculating; synapses cause delays of about $0.5$--$1\ \mathrm{ms}$ each.
\item Synaptic transmission: impulse $\rightarrow$ $\mathrm{Ca^{2+}}$ entry $\rightarrow$ vesicle fusion $\rightarrow$ transmitter release $\rightarrow$ diffusion $\rightarrow$ receptor binding $\rightarrow$ EPSP $\rightarrow$ AChE removes transmitter. Transmission is one-way.
\item The iris reflex uses antagonistic circular (constrict, parasympathetic) and radial (dilate, sympathetic) muscles; accommodation uses ciliary muscles and suspensory ligaments to change lens shape.
\item Myopia (image in front of retina) is corrected by a concave lens; hypermetropia (image behind retina) by a convex lens.
\item Hormones act via membrane receptors and second messengers (proteins) or intracellular receptors and gene activation (steroids); negative feedback keeps controlled variables near a set point.
\item Tropisms are directional growth responses; auxin is redistributed and causes differential elongation (promotes in shoots, inhibits in roots).
\item Nervous coordination is fast, brief and localised; hormonal coordination is slower, longer-lasting and widespread; the hypothalamus and adrenal medulla link the two systems.
\item The ear detects sound (cochlea, organ of Corti) and balance (utricle, saccule, semicircular canals) using hair cells.
\item Plant hormones (auxins, gibberellins, cytokinins, ABA, ethylene) control growth and have major agricultural applications in Cameroon.
\item Photoperiodism is controlled by night length via the phytochrome system (Pr $\leftrightarrow$ Pfr); the last light quality received determines the flowering response.
\item Rods (rhodopsin, dim light, peripheral retina, no colour) and cones (iodopsins, bright light, fovea, colour, high acuity); light \emph{hyperpolarises} photoreceptors.
\end{itemize}
\end{aretenir}

\newpage

\coursec{Exercises}

\courssub{I check my knowledge}

\begin{exercice}[Multiple Choice Questions]
For each question, choose the correct answer (A, B, C or D).

\begin{enumerate}
\item Which of the following structures in the eye is responsible for detecting light intensity and initiating the pupillary reflex?
\begin{itemize}
\item[A] Rods and cones in the retina
\item[B] Ganglion cells in the retina
\item[C] Bipolar cells in the retina
\item[D] Pigmented epithelium
\end{itemize}

\item During the depolarisation phase of an action potential in a myelinated axon, which ion movement occurs?
\begin{itemize}
\item[A] $\mathrm{Na^+}$ ions move out of the axon
\item[B] $\mathrm{K^+}$ ions move into the axon
\item[C] $\mathrm{Na^+}$ ions move into the axon
\item[D] $\mathrm{Ca^{2+}}$ ions move out of the axon
\end{itemize}

\item The hormone that stimulates the conversion of glycogen to glucose in the liver is:
\begin{itemize}
\item[A] Insulin
\item[B] Glucagon
\item[C] Adrenaline
\item[D] Thyroxine
\end{itemize}

\item In a cholinergic synapse, the enzyme that hydrolyses acetylcholine in the synaptic cleft is:
\begin{itemize}
\item[A] Acetylcholine synthetase
\item[B] Acetylcholinesterase
\item[C] Choline acetyltransferase
\item[D] Monoamine oxidase
\end{itemize}
\end{enumerate}
\end{exercice}

\begin{exercice}[True / False with Justification]
State whether each statement is true or false. Justify your answer with a brief biological explanation.

\begin{enumerate}
\item The fovea contains only rods and is responsible for vision in dim light.
\item The refractory period of a neurone ensures that action potentials travel in one direction only.
\item Auxins promote root elongation at high concentrations.
\item The semicircular canals detect linear acceleration and gravity.
\end{enumerate}
\end{exercice}

\begin{exercice}[Definitions]
Define the following terms precisely, as required in the GCE Advanced Level examination.

\begin{enumerate}
\item \cle{Stimulus}
\item \cle{Receptor}
\item \cle{Action potential}
\item \cle{Synaptic cleft}
\item \cle{Phototropism}
\item \cle{Second messenger} (in hormone action)
\end{enumerate}
\end{exercice}

\begin{exercice}[Direct Calculations]
\begin{enumerate}
\item The resting membrane potential of a neurone is $-70\ \mathrm{mV}$. During an action potential, the membrane potential rises to $+30\ \mathrm{mV}$. Calculate the magnitude of the change in membrane potential during depolarisation.
\item A student measures the conduction velocity of an action potential in a myelinated axon of length $1.2\ \mathrm{m}$. The time taken for the impulse to travel this distance is $6\ \mathrm{ms}$. Calculate the conduction velocity in $\mathrm{m\,s^{-1}}$.
\item In an experiment on phototropism, a coleoptile bends $15^\circ$ towards a unilateral light source after $2\ \mathrm{hours}$. If the rate of bending is constant, what angle will it have bent after $5\ \mathrm{hours}$?
\end{enumerate}
\end{exercice}

\courssub{I practise}

\begin{exercice}[Eye Structure, Image Formation and Light Adaptation]
The diagram below shows a longitudinal section of the human eye. (In the examination, a diagram would be provided; here you should draw and label your own simplified diagram.)

\begin{enumerate}
\item Draw a large, labelled diagram of the human eye showing the following structures: cornea, aqueous humour, pupil, lens, vitreous humour, retina, fovea, optic nerve, iris, ciliary muscle, suspensory ligaments, sclera, choroid.
\item Explain how the eye focuses on a near object when a person moves from a dimly lit room into bright sunlight. Your answer should include:
\begin{itemize}
\item Changes in the lens shape and the role of the ciliary muscles and suspensory ligaments (accommodation).
\item Changes in pupil diameter and the role of the iris muscles (pupillary reflex).
\item The sequence of events from light hitting the retina to the brain interpreting the image.
\end{itemize}
\end{enumerate}
\end{exercice}

\begin{exercice}[Analysis of an Action Potential]
The graph below shows the change in membrane potential of a motor neurone during an action potential. (Sketch a typical action potential graph: resting potential at $-70\ \mathrm{mV}$, rapid rise to $+30\ \mathrm{mV}$, fall below resting potential, return to resting potential.)

\begin{enumerate}
\item On the graph, label and define: resting potential, depolarisation, repolarisation, hyperpolarisation (after-potential), refractory period.
\item Explain the ion movements responsible for each phase (depolarisation, repolarisation, hyperpolarisation). Mention the specific voltage-gated channels involved.
\item Why does the refractory period prevent the action potential from travelling backwards along the axon?
\end{enumerate}
\end{exercice}

\begin{exercice}[Comparison of Rods and Cones]
\begin{enumerate}
\item Complete the table below to compare rods and cones in the human retina.

\begin{center}
\begin{tabularx}{\linewidth}{|l|X|X|}\hline
\textbf{Feature} & \textbf{Rods} & \textbf{Cones} \\ \hline
Visual pigment & & \\ \hline
Distribution in retina & & \\ \hline
Sensitivity to light intensity & & \\ \hline
Visual acuity (detail) & & \\ \hline
Colour vision & & \\ \hline
Convergence (neuronal pathways) & & \\ \hline
\end{tabularx}
\end{center}

\item Explain why the fovea provides the sharpest vision, referring to the distribution of photoreceptors and their connections to bipolar and ganglion cells.
\item Why is it difficult to see a faint star at night when looking directly at it, but easier when looking slightly to the side?
\end{enumerate}
\end{exercice}

\begin{exercice}[Cholinergic Synapse and Drug Action]
\begin{enumerate}
\item Describe the sequence of events at a cholinergic synapse from the arrival of an action potential at the presynaptic knob to the generation of a postsynaptic potential. Include the roles of voltage-gated $\mathrm{Ca^{2+}}$ channels, vesicles, acetylcholine (ACh), receptors, and acetylcholinesterase.
\item Explain why synaptic transmission is unidirectional.
\item A drug inhibits acetylcholinesterase. Predict and explain the effect of this drug on:
\begin{itemize}
\item The concentration of ACh in the synaptic cleft.
\item The frequency of action potentials in the postsynaptic neurone.
\item The duration of the postsynaptic response.
\end{itemize}
\item Give one medical use and one toxic effect of an acetylcholinesterase inhibitor.
\end{enumerate}
\end{exercice}

\courssub{I prepare for the GCE}

\begin{exercice}[{Essay: Nervous vs Hormonal Coordination] [15 marks}]
Compare and contrast nervous coordination and hormonal (endocrine) coordination in mammals. Your answer should include:
\begin{itemize}
\item Nature of the signal (electrical/chemical) and speed of transmission.
\item Duration of effect and reversibility.
\item Pathway of the message (neurones vs bloodstream).
\item Specificity of target cells.
\item Two named examples of each type of coordination in humans, describing the stimulus, receptor, coordinator, effector and response.
\end{itemize}
\end{exercice}

\begin{exercice}[{Data Analysis: Blood Glucose Regulation] [12 marks}]
The graph shows blood glucose concentration over time for two individuals, A (normal) and B (diabetic), after each consumed a glucose-rich meal at time $t=0$. (Sketch typical curves: A rises to a peak at ~30 min, returns to normal by 120 min; B rises higher, stays elevated longer.)

\begin{enumerate}
\item Describe the changes in blood glucose concentration for individual A between $0$ and $120$ minutes. [3]
\item Explain the hormonal mechanisms responsible for the return of blood glucose to normal in individual A. Name the hormones, their source, and their target tissues/actions. [5]
\item Suggest why individual B fails to regulate blood glucose effectively. Distinguish between Type 1 and Type 2 diabetes in your explanation. [4]
\end{enumerate}
\end{exercice}

\begin{exercice}[{Experimental Design: Phototropism in Coleoptiles] [13 marks}]
Design a controlled experiment to demonstrate that unilateral light causes phototropic curvature in oat coleoptiles (or maize seedlings) and that the tip is the site of perception/production of the hormone involved.

Your answer must include:
\begin{itemize}
\item A clear hypothesis.
\item Description of at least four experimental groups (including controls) with diagrams/labels.
\item The independent variable, dependent variable, and at least three controlled variables.
\item The expected results for each group.
\item A conclusion linking the results to the role of auxin (IAA) distribution.
\end{itemize}
\end{exercice}

\newpage

\coursec{Solutions to the exercises}

\begin{corrige}[Exercise 1 — Multiple Choice Questions]
\begin{enumerate}
\item \textbf{Answer: B — Ganglion cells in the retina.}\par
A sub-population of retinal ganglion cells contains the photopigment \cle{melanopsin} and is directly sensitive to light intensity. These cells send signals via the optic nerve to the pretectal nucleus of the midbrain, which initiates the \cle{pupillary reflex} (constriction of the pupil in bright light). Rods and cones (A) detect light for image formation, not primarily for the pupillary reflex; bipolar cells (C) are intermediate relay neurones; the pigmented epithelium (D) absorbs stray light and regenerates visual pigments.

\item \textbf{Answer: C — $\mathrm{Na^+}$ ions move into the axon.}\par
Depolarisation is caused by the opening of voltage-gated $\mathrm{Na^+}$ channels. Because the electrochemical gradient favours $\mathrm{Na^+}$ entry (higher external concentration and negative interior), $\mathrm{Na^+}$ rushes \emph{into} the axon, making the inside positive (up to about $+30\ \mathrm{mV}$). $\mathrm{K^+}$ moving \emph{out} (not in) causes repolarisation.

\item \textbf{Answer: B — Glucagon.}\par
\cle{Glucagon}, secreted by the $\alpha$-cells of the islets of Langerhans in the pancreas when blood glucose falls, binds to liver (hepatocyte) receptors and stimulates \cle{glycogenolysis} (glycogen $\rightarrow$ glucose) and gluconeogenesis. Insulin (A) does the opposite (glucose $\rightarrow$ glycogen). Adrenaline (C) can also promote glycogen breakdown, but it is primarily a stress hormone; the specific blood-glucose-regulating hormone asked for here is glucagon. Thyroxine (D) controls metabolic rate.

\item \textbf{Answer: B — Acetylcholinesterase.}\par
\cle{Acetylcholinesterase} (AChE), located in the synaptic cleft, hydrolyses acetylcholine into choline and ethanoic (acetic) acid, terminating the signal. Choline is reabsorbed into the presynaptic knob for resynthesis of ACh by choline acetyltransferase (C). Monoamine oxidase (D) acts on monoamine transmitters such as noradrenaline, not ACh.
\end{enumerate}
\end{corrige}

\begin{corrige}[Exercise 2 — True / False with Justification]
\begin{enumerate}
\item \textbf{False.} The fovea (fovea centralis) contains \emph{only cones}, not rods — rods are completely absent from it. Cones give high-acuity colour vision in bright light. Vision in dim light is the function of \emph{rods}, which are concentrated in the peripheral retina and contain the pigment rhodopsin, which is very sensitive to low light intensities.

\item \textbf{True.} During the \cle{absolute refractory period}, the voltage-gated $\mathrm{Na^+}$ channels are inactivated and cannot reopen, so no new action potential can be generated in the region just behind the travelling impulse. The impulse can therefore only propagate forwards into the region where channels are still in their resting (closed but activatable) state. This ensures \emph{unidirectional} transmission along the axon and keeps successive impulses separate.

\item \textbf{False.} Auxins (e.g.\ IAA) \emph{inhibit} root elongation at high concentrations; roots are far more sensitive to auxin than shoots. Root elongation is promoted only by very \emph{low} auxin concentrations. In shoots, by contrast, higher auxin concentrations promote cell elongation — this differential sensitivity explains the opposite curvatures in gravitropism of roots and shoots.

\item \textbf{False.} The semicircular canals detect \emph{rotational (angular) acceleration} of the head: movement of endolymph deflects the cupula in each ampulla. \emph{Linear} acceleration and gravity are detected by the \cle{utricle and saccule} (the otolith organs), where otoliths embedded in a gelatinous membrane press on hair cells.
\end{enumerate}
\end{corrige}

\begin{corrige}[Exercise 3 — Definitions]
\begin{enumerate}
\item \textbf{Stimulus:} A change in the internal or external environment of an organism which is detected by a receptor and which brings about a response.

\item \textbf{Receptor:} A specialised cell, or structure, that detects a specific stimulus and converts (transduces) the energy of that stimulus into an electrical signal — a generator potential / nerve impulse.

\item \textbf{Action potential:} A rapid, transient, self-propagating reversal of the potential difference across the membrane of a neurone (from about $-70\ \mathrm{mV}$ to about $+30\ \mathrm{mV}$ and back), caused by the sequential opening of voltage-gated $\mathrm{Na^+}$ and $\mathrm{K^+}$ channels; it is the electrical signal by which information travels along a neurone.

\item \textbf{Synaptic cleft:} The narrow gap (about $20\ \mathrm{nm}$) separating the presynaptic membrane of one neurone from the postsynaptic membrane of the next cell (neurone or effector), across which the signal is carried by diffusion of a chemical neurotransmitter.

\item \textbf{Phototropism:} A directional growth movement (curvature) of a plant organ in response to the direction of unilateral light; shoots are positively phototropic (grow towards light) because auxin accumulates on the shaded side and stimulates greater elongation there.

\item \textbf{Second messenger:} An intracellular molecule (e.g.\ cyclic AMP, cAMP) produced inside a target cell when a water-soluble hormone (the \emph{first} messenger, e.g.\ adrenaline or glucagon) binds to a receptor on the cell surface membrane; the second messenger relays and amplifies the signal by activating enzymes inside the cell, producing the cellular response.
\end{enumerate}
\end{corrige}

\begin{corrige}[Exercise 4 — Direct Calculations]
\begin{enumerate}
\item The change in membrane potential is:
\[
\Delta V = V_{\text{peak}} - V_{\text{rest}} = (+30) - (-70) = +100\ \mathrm{mV}.
\]
The \textbf{magnitude of the change is $100\ \mathrm{mV}$} (an increase of $100\ \mathrm{mV}$ during depolarisation).

\item Conduction velocity $v = \dfrac{d}{t}$ with $d = 1.2\ \mathrm{m}$ and $t = 6\ \mathrm{ms} = 6 \times 10^{-3}\ \mathrm{s}$:
\[
v = \frac{1.2}{6 \times 10^{-3}} = 200\ \mathrm{m\,s^{-1}}.
\]
\textbf{Conduction velocity $= 200\ \mathrm{m\,s^{-1}}$.} (This high value is typical of saltatory conduction in a myelinated axon.)

\item Rate of bending $= \dfrac{15^\circ}{2\ \mathrm{h}} = 7.5^\circ\ \mathrm{h^{-1}}$. After $5$ hours:
\[
\theta = 7.5 \times 5 = 37.5^\circ.
\]
\textbf{The coleoptile will have bent $37.5^\circ$ after 5 hours} (assuming the rate remains constant).
\end{enumerate}
\end{corrige}

\begin{corrige}[Exercise 5 — Eye Structure, Image Formation and Light Adaptation]
\textbf{1. Labelled diagram of the human eye.}\par
The diagram should show the three layers (sclera, choroid, retina), the optical pathway (cornea $\rightarrow$ aqueous humour $\rightarrow$ pupil $\rightarrow$ lens $\rightarrow$ vitreous humour $\rightarrow$ retina) and the accessory structures.

\begin{popfigure}
\begin{tikzpicture}[scale=1.0]
% Sclera (outer)
\draw[very thick, popink] (0,0) ellipse (3.4 and 2.6);
% Choroid
\draw[thick, poppurple] (0,0) ellipse (3.15 and 2.35);
% Retina
\draw[thick, popgreen] (0,0) ellipse (2.9 and 2.1);
% Cornea bulge
\draw[very thick, popblue] (2.9,-1.0) .. controls (4.1,-0.7) and (4.1,0.7) .. (2.9,1.0);
% Iris
\draw[very thick, poporange] (2.35,0.45) -- (2.35,1.05);
\draw[very thick, poporange] (2.35,-0.45) -- (2.35,-1.05);
% Lens
\draw[thick, popgold, fill=popgold!20] (1.85,0) ellipse (0.28 and 0.95);
% Ciliary muscles
\draw[thick, popdark] (1.85,0.95) -- (2.35,1.35);
\draw[thick, popdark] (1.85,-0.95) -- (2.35,-1.35);
% Optic nerve
\draw[very thick, poppink] (-3.3,0.35) -- (-4.2,0.35);
\draw[very thick, poppink] (-3.3,-0.35) -- (-4.2,-0.35);
\draw[very thick, poppink] (-3.3,0.35) .. controls (-3.5,0) .. (-3.3,-0.35);
% Fovea
\fill[popgreen] (-2.88,0.55) circle (0.06);
% Labels
\node[font=\small, align=center] at (0,-1.7) {vitreous\\humour};
\node[font=\small] at (3.6,0) {cornea};
\node[font=\small] at (2.75,-0.75) {pupil};
\node[font=\small] at (1.85,0) {\scriptsize lens};
\node[font=\small] at (3.3,1.5) {iris};
\draw[->, thin] (3.3,1.35) -- (2.4,1.0);
\node[font=\small] at (1.0,1.7) {ciliary muscle};
\draw[->, thin] (1.35,1.6) -- (2.1,1.25);
\node[font=\small] at (0.6,2.55) {suspensory ligaments};
\draw[->, thin] (1.1,2.4) -- (1.75,1.0);
\node[font=\small] at (-1.3,2.6) {sclera};
\draw[->, thin] (-1.3,2.45) -- (-1.3,2.35);
\node[font=\small] at (-2.9,2.1) {choroid};
\draw[->, thin] (-2.6,2.0) -- (-2.2,2.05);
\node[font=\small] at (-2.3,1.35) {retina};
\draw[->, thin] (-2.3,1.25) -- (-2.5,1.75);
\node[font=\small] at (-4.35,0) {optic nerve};
\node[font=\small] at (-3.9,0.9) {fovea};
\draw[->, thin] (-3.6,0.8) -- (-2.95,0.6);
\node[font=\small] at (2.7,0.55) {\scriptsize aqueous};
\end{tikzpicture}
\legende{Longitudinal section through the human eye (schematic).}
\end{popfigure}

\textbf{2. Response to moving into bright sunlight and focusing on a near object.}\par
\emph{(a) Accommodation (focusing on the near object):}
\begin{itemize}
\item The \cle{ciliary muscles contract}, releasing tension on the suspensory ligaments.
\item The \cle{suspensory ligaments slacken}, so the elastic \cle{lens} becomes \emph{more convex} (shorter, fatter), increasing its refractive power.
\item Light rays from the near object (which are diverging strongly) are now refracted enough to be focused sharply on the retina.
\end{itemize}
\emph{(b) Pupillary reflex (response to bright light):}
\begin{itemize}
\item High light intensity is detected by retinal ganglion cells; impulses pass via the optic nerve to the midbrain and back along parasympathetic motor neurones.
\item The \cle{circular muscles of the iris contract} and the radial muscles relax.
\item The \cle{pupil constricts}, reducing the amount of light entering the eye and protecting the retina (especially the rods, whose rhodopsin is bleached).
\end{itemize}
\emph{(c) From retina to perception:}
\begin{enumerate}
\item Light passes through the cornea, aqueous humour, pupil, lens and vitreous humour and is focused on the retina, forming a real, inverted, diminished image.
\item Photoreceptors (rods and cones) absorb light; visual pigments (rhodopsin/iodopsin) are bleached, generating a generator potential.
\item Impulses pass from photoreceptors to bipolar neurones, then to ganglion cells.
\item Ganglion cell axons form the optic nerve, which carries impulses to the visual cortex of the cerebrum (occipital lobe).
\item The brain interprets the impulses, re-inverting the image so that the object is perceived upright and in detail.
\end{enumerate}
\end{corrige}

\begin{corrige}[Exercise 6 — Analysis of an Action Potential]
\textbf{1. Labelled graph and definitions.}

\begin{popfigure}
\begin{tikzpicture}[scale=1.0]
\draw[->] (0,0) -- (10.5,0) node[below left, font=\small]{time / ms};
\draw[->] (0,-2.2) -- (0,2.6) node[above, font=\small]{$V$ / mV};
\draw[dashed, gray] (0,1.5) node[left, font=\small]{$+30$} -- (9.5,1.5);
\draw[dashed, gray] (0,-1.4) node[left, font=\small]{$-70$} -- (9.5,-1.4);
\draw[very thick, popblue] (0,-1.4) -- (2.5,-1.4)
 .. controls (2.9,-1.4) and (3.0,1.5) .. (3.5,1.5)
 .. controls (4.0,1.5) and (4.3,-1.9) .. (5.2,-1.9)
 .. controls (5.9,-1.9) and (6.1,-1.4) .. (7.0,-1.4) -- (9.8,-1.4);
\node[font=\small, popink] at (1.2,-1.0) {resting potential};
\node[font=\small, popink, rotate=62] at (2.75,0.35) {depolarisation};
\node[font=\small, popink, rotate=-62] at (4.35,0.3) {repolarisation};
\node[font=\small, popink] at (5.6,-2.05) {hyperpolarisation};
\draw[<->, thick, poporange] (3.0,-2.6) -- (5.6,-2.6) node[midway, below, font=\small]{refractory period};
\end{tikzpicture}
\legende{Membrane potential during an action potential.}
\end{popfigure}

\begin{itemize}
\item \textbf{Resting potential:} the stable potential difference across the membrane of an unstimulated neurone (about $-70\ \mathrm{mV}$, inside negative), maintained by the $\mathrm{Na^+}$/$\mathrm{K^+}$ pump and the greater permeability to $\mathrm{K^+}$.
\item \textbf{Depolarisation:} the rapid reversal of the membrane potential from $-70\ \mathrm{mV}$ to about $+30\ \mathrm{mV}$.
\item \textbf{Repolarisation:} the rapid return of the membrane potential towards the resting value.
\item \textbf{Hyperpolarisation (after-potential):} a brief period when the potential falls below $-70\ \mathrm{mV}$ (about $-90\ \mathrm{mV}$) before returning to rest.
\item \textbf{Refractory period:} the interval during and just after an action potential when a new action potential cannot (absolute) or can only with difficulty (relative) be generated.
\end{itemize}

\textbf{2. Ionic basis of each phase.}
\begin{itemize}
\item \emph{Depolarisation:} the stimulus raises the potential above the threshold (about $-55\ \mathrm{mV}$); \cle{voltage-gated $\mathrm{Na^+}$ channels} open, $\mathrm{Na^+}$ floods \emph{into} the axon down its electrochemical gradient, making the inside positive ($+30\ \mathrm{mV}$) — positive feedback opens more $\mathrm{Na^+}$ channels.
\item \emph{Repolarisation:} at the peak, $\mathrm{Na^+}$ channels \emph{inactivate} and \cle{voltage-gated $\mathrm{K^+}$ channels} open; $\mathrm{K^+}$ diffuses \emph{out} of the axon, making the inside negative again.
\item \emph{Hyperpolarisation:} the $\mathrm{K^+}$ channels are slow to close, so excessive $\mathrm{K^+}$ leaves the cell and the potential briefly falls below $-70\ \mathrm{mV}$. The channels then close and the $\mathrm{Na^+}$/$\mathrm{K^+}$ pump (3 $\mathrm{Na^+}$ out : 2 $\mathrm{K^+}$ in, using ATP) restores the resting distribution of ions.
\end{itemize}

\textbf{3. Why the impulse cannot travel backwards.}\par
Immediately behind the propagating action potential, the voltage-gated $\mathrm{Na^+}$ channels are \emph{inactivated} (absolute refractory period) and cannot reopen whatever the stimulus. The local currents from the active region can therefore only depolarise the membrane \emph{ahead} of the impulse, where the channels are closed but activatable. Hence propagation is strictly unidirectional, and successive impulses remain discrete (they cannot merge).
\end{corrige}

\begin{corrige}[Exercise 7 — Comparison of Rods and Cones]
\textbf{1. Completed table.}

\begin{center}
\begin{tabularx}{\linewidth}{|l|X|X|}\hline
\rowcolor{popblueL}\textbf{Feature} & \textbf{Rods} & \textbf{Cones} \\ \hline
Visual pigment & Rhodopsin (one type only) & Iodopsin (three types, sensitive to red, green or blue light) \\ \hline
Distribution in retina & Peripheral retina; absent from the fovea & Concentrated in the fovea; few elsewhere \\ \hline
Sensitivity to light intensity & High — function in dim light (scotopic vision) & Low — require bright light (photopic vision) \\ \hline
Visual acuity (detail) & Low acuity, poor detail & High acuity, sharp detail \\ \hline
Colour vision & None — black/white/grey only & Responsible for colour vision (trichromatic) \\ \hline
Convergence (neuronal pathways) & High: many rods converge onto one bipolar/ganglion cell (summation) & Low: at the fovea each cone connects to its own bipolar and ganglion cell \\ \hline
\end{tabularx}
\end{center}

\textbf{2. Why the fovea gives the sharpest vision.}
\begin{itemize}
\item The fovea contains \emph{only cones}, packed very densely, and cones give high acuity and colour vision.
\item Each foveal cone synapses with its \emph{own} bipolar cell, which connects to its \emph{own} ganglion cell — a 1:1:1 pathway with \emph{no convergence}. The brain can therefore distinguish signals from two adjacent cones separately, resolving fine detail.
\item Overlying neurones and blood vessels are displaced to the sides of the fovea, so light strikes the cones directly without scattering.
\end{itemize}

\textbf{3. Why a faint star is seen better when looking slightly to the side.}\par
Looking directly at the star focuses its image on the \emph{fovea}, which contains only cones — and cones are insensitive to dim light, so the star is not detected. Looking slightly to the side moves the image onto the \emph{peripheral retina}, which is rich in \emph{rods}. Rods are highly sensitive to low light intensities (and their convergence allows \cle{summation} of weak signals), so the faint star becomes visible — although with poor detail and no colour.
\end{corrige}

\begin{corrige}[Exercise 8 — Cholinergic Synapse and Drug Action]
\textbf{1. Sequence of events at a cholinergic synapse.}
\begin{enumerate}
\item An action potential arrives at the presynaptic knob and depolarises the presynaptic membrane.
\item \cle{Voltage-gated $\mathrm{Ca^{2+}}$ channels} open; $\mathrm{Ca^{2+}}$ ions diffuse \emph{into} the knob down their electrochemical gradient.
\item The rise in $\mathrm{Ca^{2+}}$ causes synaptic \cle{vesicles} containing acetylcholine (ACh) to move to and fuse with the presynaptic membrane, releasing ACh into the synaptic cleft by \emph{exocytosis}.
\item ACh diffuses across the cleft ($\approx 20\ \mathrm{nm}$) and binds to specific \cle{ACh receptors} (ligand-gated $\mathrm{Na^+}$ channels) on the postsynaptic membrane.
\item The channels open; $\mathrm{Na^+}$ enters the postsynaptic neurone, causing depolarisation — an \emph{excitatory postsynaptic potential} (EPSP). If the EPSP reaches threshold, an action potential is generated in the postsynaptic neurone.
\item \cle{Acetylcholinesterase} in the cleft rapidly hydrolyses ACh into choline and ethanoic acid, terminating the signal; choline is reabsorbed into the presynaptic knob and used to resynthesise ACh.
\end{enumerate}

\textbf{2. Why transmission is unidirectional.}\par
Vesicles of neurotransmitter are present \emph{only} in the presynaptic knob, and the specific receptors are located \emph{only} on the postsynaptic membrane. Transmitter can therefore be released and detected in one direction only, so the impulse cannot travel backwards across the synapse.

\textbf{3. Effect of an acetylcholinesterase inhibitor.}
\begin{itemize}
\item \emph{Concentration of ACh in the cleft:} it \textbf{increases}, because ACh is released normally but is no longer broken down, so it accumulates.
\item \emph{Frequency of postsynaptic action potentials:} it \textbf{increases}. Persisting ACh repeatedly opens the receptor channels, so the postsynaptic membrane is re-depolarised again and again, firing repeated action potentials.
\item \emph{Duration of the postsynaptic response:} it is \textbf{prolonged}, because each release of ACh continues to stimulate the receptors until the ACh finally diffuses away or is degraded by other routes.
\end{itemize}

\textbf{4. Medical use and toxic effect.}
\begin{itemize}
\item \emph{Medical use:} treatment of \textbf{myasthenia gravis} (an autoimmune disease in which ACh receptors are destroyed); AChE inhibitors such as neostigmine raise ACh levels in the cleft so that the remaining receptors are adequately stimulated, restoring muscle contraction.
\item \emph{Toxic effect:} \textbf{organophosphate insecticides and nerve gases} irreversibly inhibit AChE; ACh accumulates, causing continuous stimulation of muscles — twitching, paralysis of respiratory muscles, excessive glandular secretion — which can be fatal.
\end{itemize}
\end{corrige}

\begin{corrige}[Exercise 9 — Essay: Nervous vs Hormonal Coordination (15 marks)]
\textbf{Indicative mark scheme and model content.}

\emph{Comparison table (up to 7 marks — one mark per correct contrasted pair).}

\begin{center}
\begin{tabularx}{\linewidth}{|l|X|X|}\hline
\rowcolor{popblueL}\textbf{Feature} & \textbf{Nervous coordination} & \textbf{Hormonal coordination} \\ \hline
Nature of signal & Electrical impulses (plus chemical transmitter at synapses) & Chemical only (hormones) \\ \hline
Speed & Very fast (up to $\sim 120\ \mathrm{m\,s^{-1}}$) & Slow (speed of blood circulation) \\ \hline
Duration of effect & Short-lived, brief & Longer-lasting, sustained \\ \hline
Reversibility & Rapidly reversed when stimulus stops & Effects persist while hormone circulates; slower to reverse \\ \hline
Pathway & Along specific neurones to a precise effector & Through the bloodstream to the whole body \\ \hline
Specificity of target & Highly localised (one muscle/gland) & Widespread; only cells with the correct receptors respond \\ \hline
\end{tabularx}
\end{center}

\emph{Similarities (2 marks):} both involve chemical messengers (neurotransmitters / hormones); both link a stimulus to a response via receptors and effectors; both maintain homeostasis.

\emph{Two named nervous examples (3 marks):}
\begin{itemize}
\item \textbf{Withdrawal reflex:} stimulus — a pin prick on the finger; receptor — pain receptors in the skin; coordinator — spinal cord (relay neurone); effector — biceps muscle of the arm; response — rapid withdrawal of the hand.
\item \textbf{Pupillary reflex:} stimulus — bright light; receptor — retinal ganglion cells; coordinator — midbrain; effector — circular muscles of the iris; response — pupil constricts.
\end{itemize}

\emph{Two named hormonal examples (3 marks):}
\begin{itemize}
\item \textbf{Blood glucose regulation:} stimulus — rise in blood glucose after a meal; receptor — $\beta$-cells of the islets of Langerhans; coordinator — pancreas; effector — liver and muscle cells; response — insulin secreted, glucose converted to glycogen, blood glucose falls.
\item \textbf{``Fight or flight'':} stimulus — sudden danger; receptor/coordinator — hypothalamus stimulates the adrenal medulla; effector — heart, liver, bronchioles; response — adrenaline released, heart rate and blood glucose rise, preparing the body for action.
\end{itemize}

\textbf{Examiner expectations:} a clear table or paired statements (not two separate unlinked accounts); accurate terminology (neurone, effector, target cells, receptors); correct stimulus--receptor--coordinator--effector--response chains for each example; at least one similarity stated. Award full marks only if all four bullet points of the question are addressed.
\end{corrige}

\begin{corrige}[Exercise 10 — Data Analysis: Blood Glucose Regulation (12 marks)]
\textbf{1. Description for individual A [3 marks].}
\begin{itemize}
\item From $t = 0$ blood glucose rises steeply as glucose is absorbed from the gut, reaching a peak at about 30 minutes. [1]
\item It then falls progressively as glucose is removed from the blood. [1]
\item By about 120 minutes it has returned to (and is stabilised at) the normal resting level (about $90\ \mathrm{mg}$ per $100\ \mathrm{cm^3}$). [1]
\end{itemize}

\textbf{2. Hormonal mechanism of the return to normal [5 marks].}
\begin{itemize}
\item The rise in blood glucose is detected by the \cle{$\beta$-cells of the islets of Langerhans} in the pancreas. [1]
\item They secrete \cle{insulin} into the blood. [1]
\item Insulin binds to receptors on \emph{liver cells (hepatocytes)} and \emph{muscle cells}, stimulating the conversion of glucose to \cle{glycogen} (glycogenesis) for storage. [1]
\item It also increases the uptake of glucose by cells and its oxidation in respiration, and inhibits glycogen breakdown. [1]
\item As blood glucose falls, insulin secretion falls (negative feedback); if glucose drops below normal, \cle{glucagon} from the $\alpha$-cells stimulates the liver to convert glycogen back to glucose, restoring the level. [1]
\end{itemize}

\textbf{3. Why individual B fails to regulate [4 marks].}
\begin{itemize}
\item B's glucose rises higher and remains elevated because the insulin mechanism is defective. [1]
\item \textbf{Type 1 diabetes:} the $\beta$-cells are destroyed (autoimmune reaction), so \emph{no insulin is produced}; blood glucose cannot be lowered despite the stimulus. Managed by insulin injections. [1.5]
\item \textbf{Type 2 diabetes:} insulin is produced, but the \emph{target cells' receptors are insensitive/resistant} to insulin (often associated with obesity and age), so glucose uptake and glycogenesis are inadequate. Managed by diet, exercise and drugs. [1.5]
\end{itemize}
\textbf{Examiner expectations:} precise description quoting times and trend; both hormones named with sources and targets; clear distinction between insulin deficiency (Type 1) and insulin resistance (Type 2).
\end{corrige}

\begin{corrige}[Exercise 11 — Experimental Design: Phototropism in Coleoptiles (13 marks)]
\textbf{Hypothesis [1]:} Unilateral light causes a coleoptile to bend towards the light because the \emph{tip} perceives the light and produces a growth hormone (auxin/IAA) which is redistributed to the shaded side, where it stimulates greater cell elongation.

\textbf{Experimental groups [4 — one per group with correct treatment].} Use uniform oat/maize coleoptiles of the same age and length, all exposed to unilateral light (except Group E):

\begin{popfigure}
\begin{tikzpicture}[scale=0.85]
% Group A
\draw[thick, popgreen] (0,0) -- (0,1.6) .. controls (0,1.9) .. (0.0,2.0);
\draw[thick, popgreen] (0.25,0) -- (0.25,1.6) .. controls (0.25,1.9) .. (0.0,2.0);
\node[font=\small] at (0.12,-0.35) {A: intact};
% Group B
\draw[thick, popgreen] (1.6,0) -- (1.6,1.3);
\draw[thick, popgreen] (1.85,0) -- (1.85,1.3);
\draw[thick, popink] (1.5,1.3) -- (1.95,1.3);
\node[font=\small] at (1.72,-0.35) {B: tip cut off};
% Group C
\draw[thick, popgreen] (3.2,0) -- (3.2,1.6) .. controls (3.2,1.9) .. (3.2,2.0);
\draw[thick, popgreen] (3.45,0) -- (3.45,1.6) .. controls (3.45,1.9) .. (3.2,2.0);
\draw[thick, popink, fill=popink!40] (3.13,1.55) rectangle (3.52,2.05);
\node[font=\small] at (3.32,-0.35) {C: tip covered};
% Group D
\draw[thick, popgreen] (4.8,0) -- (4.8,1.6) .. controls (4.8,1.9) .. (4.8,2.0);
\draw[thick, popgreen] (5.05,0) -- (5.05,1.6) .. controls (5.05,1.9) .. (4.8,2.0);
\draw[thick, popink, fill=popink!40] (4.73,0.6) rectangle (5.12,1.3);
\node[font=\small] at (4.92,-0.35) {D: base covered};
% Group E
\draw[thick, popgreen] (6.4,0) -- (6.4,1.6) .. controls (6.4,1.9) .. (6.4,2.0);
\draw[thick, popgreen] (6.65,0) -- (6.65,1.6) .. controls (6.65,1.9) .. (6.4,2.0);
\node[font=\small] at (6.52,-0.35) {E: intact, dark};
% Light arrow
\draw[->, very thick, popgold] (-1.0,1.6) -- (-0.25,1.6) node[midway, above, font=\small]{light};
\end{tikzpicture}
\legende{Five groups of coleoptiles under unilateral light (E in darkness).}
\end{popfigure}

\begin{itemize}
\item \textbf{Group A (control):} intact coleoptile, unilateral light.
\item \textbf{Group B:} tip removed, unilateral light.
\item \textbf{Group C:} tip covered with an opaque cap, unilateral light.
\item \textbf{Group D:} tip exposed but the \emph{base} covered with an opaque sleeve, unilateral light.
\item \textbf{Group E (control):} intact coleoptile kept in darkness (or uniform light).
\end{itemize}

\textbf{Variables [3]:}
\begin{itemize}
\item \emph{Independent variable:} the treatment of the tip (present / removed / covered) and the direction of light. [1]
\item \emph{Dependent variable:} the angle/direction of curvature (measured with a protractor) and the growth in length. [1]
\item \emph{Controlled variables (any three):} same species and age of seedlings; same initial coleoptile length; same light intensity, duration and wavelength; same temperature; same humidity/water supply; same number of seedlings per group (replicates). [1]
\end{itemize}

\textbf{Expected results [3]:}
\begin{itemize}
\item \textbf{A:} grows and bends \emph{towards} the light (positive phototropism).
\item \textbf{B:} little or no growth and \emph{no bending} — the tip is needed both for perception and as the source of auxin.
\item \textbf{C:} grows straight up, \emph{no bending} — the tip is present but cannot perceive the light direction, so auxin is distributed evenly.
\item \textbf{D:} bends \emph{towards} the light like A — covering the base does not matter, since perception occurs at the tip and the response (elongation) still occurs below.
\item \textbf{E:} grows straight up (etiolated, tall and pale) with no curvature.
\end{itemize}

\textbf{Conclusion [2]:} Only coleoptiles with an \emph{illuminated tip} bend towards unilateral light. The tip is therefore the site of light perception and of auxin (IAA) production. Unilateral light causes auxin to be transported to the \emph{shaded side}, where the higher auxin concentration stimulates greater cell elongation; this differential growth curves the coleoptile towards the light. [2]

\textbf{Examiner expectations:} a testable hypothesis; at least four clearly labelled groups including a positive and a negative control; correct identification of all three variable types; results consistent with the auxin-redistribution theory; a conclusion explicitly linking tip, light and unequal auxin distribution to differential elongation.
\end{corrige}

\newpage

\coursec{Summary sheet}

\courssub{Summary Sheet}

\begin{popfigure}
\begin{tikzpicture}[
  concept/.style={circle, draw, very thick, fill=#1!20, text width=2.8cm, align=center, font=\small\bfseries, inner sep=4pt},
  branch/.style={very thick, draw=#1},
  level 1/.style={sibling distance=4cm, level distance=3.5cm}
]
\node[concept=popdark] (center) {Response \& Coordination}
  child[grow=30, branch=popblue] { node[concept=popblue] {Stimuli \& Reception} }
  child[grow=90, branch=popgreen] { node[concept=popgreen] {Eye \& Vision} }
  child[grow=150, branch=poporange] { node[concept=poporange] {Ear: Hearing \& Balance} }
  child[grow=210, branch=poppurple] { node[concept=poppurple] {Nervous System} }
  child[grow=270, branch=poppink] { node[concept=poppink] {Endocrine System} }
  child[grow=330, branch=popgold] { node[concept=popgold] {Plant Coordination} };
\end{tikzpicture}
\legende{Mind map of Chapter BIO02: Response \& Coordination in Organisms}
\end{popfigure}

\begin{bilan}

\textbf{Key Definitions}
\begin{itemize}
  \item \cle{Stimulus}: A detectable change in the internal or external environment.
  \item \cle{Response}: The activity of an organism or part of an organism resulting from a stimulus.
  \item \cle{Receptor}: A specialised cell or organ that converts a stimulus into a nerve impulse (transduction).
  \item \cle{Effector}: A muscle or gland that brings about a response.
  \item \cle{Neurone}: The structural and functional unit of the nervous system; transmits electrical impulses.
  \item \cle{Synapse}: The junction between two neurones or a neurone and an effector; transmission is chemical (neurotransmitters).
  \item \cle{Hormone}: A chemical messenger secreted by an endocrine gland into the blood, acting on specific target cells.
  \item \cle{Tropism}: A directional growth response of a plant to a unilateral stimulus (e.g. phototropism, geotropism).
  \item \cle{Nastic response}: A non-directional response to a stimulus (e.g. Mimosa leaf folding).
\end{itemize}

\textbf{Core Principles \& Formulas}
\begin{itemize}
  \item \textbf{Reflex Arc}: Receptor $\to$ Sensory neurone $\to$ Interneurone (spinal cord) $\to$ Motor neurone $\to$ Effector. \emph{Fast, involuntary, protective}.
  \item \textbf{All-or-None Law}: An action potential is generated only if the stimulus reaches threshold; amplitude is constant.
  \item \textbf{Refractory Periods}: Absolute (no new AP possible) and Relative (stronger stimulus needed); ensures unidirectional propagation.
  \item \textbf{Accommodation (Eye)}: Ciliary muscles contract $\to$ suspensory ligaments slacken $\to$ lens becomes more convex $\to$ increased refractive power for near objects.
  \item \textbf{Pupillary Light Reflex}: Circular muscles (parasympathetic) constrict pupil in bright light; radial muscles (sympathetic) dilate in dim light.
  \item \textbf{Phototransduction}: Light $\to$ Rhodopsin bleaching $\to$ Transducin activation $\to$ PDE activation $\to$ cGMP hydrolysis $\to$ Na$^+$ channels close $\to$ Hyperpolarisation.
  \item \textbf{Hormone Mechanisms}:
    \begin{itemize}
      \item Peptide hormones: Bind surface receptor $\to$ Second messenger (cAMP/IP$_3$) $\to$ Protein kinase cascade $\to$ Rapid response.
      \item Steroid hormones: Diffuse through membrane $\to$ Bind intracellular receptor $\to$ Hormone-receptor complex enters nucleus $\to$ Gene transcription $\to$ Slow, long-lasting.
    \end{itemize}
  \item \textbf{Plant Hormones}: Auxins (IAA) promote cell elongation, apical dominance, root initiation; Gibberellins break dormancy, stem elongation; Cytokinins promote cell division, delay senescence; Abscisic acid (ABA) closes stomata, induces dormancy; Ethene promotes fruit ripening, leaf abscission.
\end{itemize}

\textbf{Key Methods (Exam Skills)}
\begin{enumerate}
  \item \textbf{Diagram Annotation}: Draw and label the eye (retina, fovea, blind spot, iris, lens, ciliary body), ear (cochlea, semi-circular canals, ossicles), neurone (myelin sheath, nodes of Ranvier, synaptic knob), reflex arc.
  \item \textbf{Action Potential Graph}: Label axes (mV vs ms), phases: Resting potential ($-70$ mV), Depolarisation (Na$^+$ influx), Overshoot ($+30$ mV), Repolarisation (K$^+$ efflux), Hyperpolarisation, Return to rest.
  \item \textbf{Synaptic Transmission Steps}: 1. AP arrives at knob $\to$ 2. Voltage-gated Ca$^{2+}$ channels open $\to$ 3. Vesicles fuse $\to$ 4. Acetylcholine released $\to$ 5. Binds Na$^+$ channel receptors $\to$ 6. Depolarisation (EPSP) $\to$ 7. Acetylcholinesterase hydrolyses ACh.
  \item \textbf{Comparing Nervous vs Endocrine}: Speed (ms vs mins/hrs), Duration (short vs long), Transmission (neurones vs blood), Specificity (precise vs widespread), Adaptation (rapid vs slow).
  \item \textbf{Plant Tropism Experiments}: Know the classic Darwin/Coleoptile tip experiments: Tip perceives light $\to$ Auxin migrates to shaded side $\to$ Asymmetric elongation $\to$ Bending towards light.
\end{enumerate}

\textbf{Common Mistakes (Traps)}
\begin{itemize}
  \item Confusing \cle{rods} (dim light, low acuity, no colour, rhodopsin) and \cle{cones} (bright light, high acuity, colour vision, iodopsin).
  \item Stating the lens changes shape by muscle contraction \emph{directly}; it is the \emph{release of tension} on suspensory ligaments.
  \item Thinking the \cle{cochlea} detects balance; it detects sound. Balance is via \cle{semi-circular canals} (dynamic) and \cle{utricle/saccule} (static).
  \item Writing ``neurones secrete hormones'' at synapses; they secrete \cle{neurotransmitters}.
  \item Forgetting the \cle{myelin sheath} is formed by Schwann cells (PNS) / Oligodendrocytes (CNS) and enables \cle{saltatory conduction}.
  \item Mixing up \cle{sympathetic} (fight/flight: dilates pupil, inhibits digestion, increases heart rate) and \cle{parasympathetic} (rest/digest: constricts pupil, stimulates digestion, decreases heart rate).
  \item Claiming auxin \emph{inhibits} root growth at high concentration but \emph{promotes} shoot growth; the concentration optimum differs.
  \item Omitting \cle{acetylcholinesterase} in synaptic transmission, leading to continuous stimulation.
\end{itemize}

\textbf{GCE Advanced Level Tips}
\begin{itemize}
  \item \textbf{Diagrams}: Large ($>10$ cm), clear lines, labels horizontal, label lines touching structure (no arrows), title + magnification if relevant.
  \item \textbf{Definitions}: Quote precisely: ``A hormone is a chemical substance produced by an endocrine gland and transported by the blood to a target organ where it exerts its effect.''
  \item \textbf{Comparison Tables}: Use a ruler, clear columns (Feature / Nervous / Endocrine). Fill every cell.
  \item \textbf{Calculations}: Magnification $= \frac{\text{Image size}}{\text{Actual size}}$ (same units). Scale bars: $\text{Actual size} = \frac{\text{Measured length}}{\text{Scale factor}}$.
  \item \textbf{Experimental Design}: State \cle{independent}, \cle{dependent}, \cle{controlled} variables. Mention \cle{replicates} and \cle{statistical test} (e.g. t-test, chi-squared) if asked.
  \item \textbf{Contextual Knowledge}: Link to Cameroon: Sickle cell trait (malaria resistance), Onchocerciasis (river blindness -- eye), Goitre (iodine deficiency -- thyroid), Agricultural use of auxins (rooting powders, weedkillers like 2,4-D).
\end{itemize}

\end{bilan}

\findecours

\end{document}
